% !TeX program = xelatex
% Standalone English volume. Use XeLaTeX until the references stabilize.
\documentclass[a4paper,11pt]{article}
\usepackage[T1]{fontenc}
\usepackage{lmodern}
\usepackage[english]{babel}
\usepackage[left=26mm,right=26mm,top=25mm,bottom=25mm,headheight=15pt,headsep=7mm,footskip=11mm]{geometry}
\usepackage{amsmath,amssymb,mathtools,amsthm}
\newtheorem{proposition}{Proposition}[section]
\newtheorem{lemma}[proposition]{Lemma}
\newtheorem{theorem}[proposition]{Theorem}
\newtheorem{corollary}[proposition]{Corollary}
\newtheorem{maintheorem}[proposition]{Main Theorem}
\theoremstyle{definition}
\newtheorem{definition}[proposition]{Definition}
\theoremstyle{remark}
\newtheorem{remark}[proposition]{Remark}
\makeatletter
\renewcommand*\l@section{\@dottedtocline{1}{0em}{2.4em}}
\renewcommand*\l@subsection{\@dottedtocline{2}{1.45em}{3.45em}}
\renewcommand*\l@subsubsection{\@dottedtocline{3}{3.8em}{4.55em}}
\makeatother
\usepackage{xcolor,fancyhdr,titlesec,needspace,xurl}
\usepackage{longtable,multicol}
\usepackage{hyperref}
\definecolor{ink}{HTML}{1A3542}
\definecolor{muted}{HTML}{4F606A}
\hypersetup{unicode=true,colorlinks=true,linkcolor=ink,filecolor=ink,urlcolor=ink,citecolor=ink,bookmarksnumbered=true,
 pdfdisplaydoctitle=true,pdflang={en},
 pdftitle={The generalized Riemann hypothesis for Dirichlet L-functions. Edition 930, Volume IV},
 pdfsubject={A natural-language reconstruction of the proposed proof},
 pdfauthor={A mathematical reconstruction}}
\titleformat{\section}{\Large\bfseries\color{ink}}{\thesection}{.65em}{}
\titlespacing*{\section}{0pt}{.8em}{.65em}
\setlength{\parindent}{1.2em}
\linespread{1.045}
\setlength{\parskip}{.22em plus .05em minus .03em}
\setlength{\emergencystretch}{2em}
\setcounter{topnumber}{2}
\setcounter{tocdepth}{3}\setcounter{secnumdepth}{3}
\widowpenalty=4500
\clubpenalty=4500
\pagestyle{fancy}\fancyhf{}
\fancyhead[L]{\small\color{muted}Dirichlet GRH \textperiodcentered\ Volume I}
\fancyhead[R]{\small\color{muted}Dirichlet zeros and power intervals}
\fancyfoot[C]{\small\thepage}
\renewcommand{\headrulewidth}{.3pt}
\newcommand{\lead}[1]{\par\smallskip\noindent{\bfseries\color{ink}#1.}\enspace}
\newcommand{\sour}[2]{\unskip\penalty150\hspace{.35em}{\small\hyperlink{ref#1}{[#2]}}}
\titleformat{\subsection}{\large\bfseries}{\thesubsection}{.6em}{}
\titlespacing*{\subsection}{0pt}{1.15em plus .15em}{.4em}
\setlength{\abovedisplayskip}{9pt plus 2pt minus 3pt}
\setlength{\belowdisplayskip}{9pt plus 2pt minus 3pt}
\newcommand{\RR}{\mathbb R}\newcommand{\CC}{\mathbb C}\newcommand{\QQ}{\mathbb Q}
\newcommand{\ZZ}{\mathbb Z}\newcommand{\ord}{\operatorname{ord}}
\newcommand{\Frob}{\operatorname{Frob}}\newcommand{\Gal}{\operatorname{Gal}}
\newcommand{\len}{\operatorname{length}}\newcommand{\Res}{\operatorname*{Res}}
\newcommand{\wmean}{\operatorname{av}}\newcommand{\Ker}{\operatorname{ker}}
\newcommand{\Lh}{\mathcal L}\newcommand{\Gg}{\mathcal G}
\newcommand{\real}{\operatorname{Re}}\newcommand{\imag}{\operatorname{Im}}
\newcommand{\abs}[1]{\left|#1\right|}\newcommand{\eps}{\varepsilon}
\newcommand{\anchoritem}[2]{\par\Needspace{2.8\baselineskip}\hypertarget{ref#1}{}\noindent{\bfseries\color{ink}#2.}\enspace}
\newcommand{\srcloc}[2]{{\small\nolinkurl{#1.lean}}, line~#2}
\numberwithin{equation}{section}
\raggedbottom

\newcommand{\Sel}{\mathcal S}
\newcommand{\Cl}{\operatorname{Cl}}
\newcommand{\Aut}{\operatorname{Aut}}
\newcommand{\GL}{\operatorname{GL}}
\newcommand{\OO}{\mathcal O}
\hypersetup{bookmarksdepth=3}


\addto\captionsenglish{\renewcommand{\partname}{Volume}}


\makeatletter
\expandafter\def\csname r@cor:alg97\endcsname{{87.3}{34}{}{proposition.87.3}{Dirichlet-GRH-v930-Volume-XVII.pdf}}
\expandafter\def\csname r@cor:criticaledges333\endcsname{{34.25}{89}{}{proposition.34.25}{Dirichlet-GRH-v930-Volume-VI.pdf}}
\expandafter\def\csname r@cor:exp139\endcsname{{25.8}{30}{}{proposition.25.8}{Dirichlet-GRH-v930-Volume-V.pdf}}
\expandafter\def\csname r@cor:family143\endcsname{{25.13}{35}{}{proposition.25.13}{Dirichlet-GRH-v930-Volume-V.pdf}}
\expandafter\def\csname r@cor:gram233\endcsname{{94.19}{53}{}{proposition.94.19}{Dirichlet-GRH-v930-Volume-XVIII.pdf}}
\expandafter\def\csname r@cor:levels113\endcsname{{11.23}{21}{}{proposition.11.23}{Dirichlet-GRH-v930-Volume-II.pdf}}
\expandafter\def\csname r@cor:normbudget340\endcsname{{28.17}{102}{}{proposition.28.17}{Dirichlet-GRH-v930-Volume-V.pdf}}
\expandafter\def\csname r@cor:original200\endcsname{{83.7}{57}{}{proposition.83.7}{Dirichlet-GRH-v930-Volume-XVI.pdf}}
\expandafter\def\csname r@cor:periodfieldweights644\endcsname{{57.52}{107}{}{proposition.57.52}{Dirichlet-GRH-v930-Volume-XI.pdf}}
\expandafter\def\csname r@cor:polarsplit593\endcsname{{61.19}{63}{}{proposition.61.19}{Dirichlet-GRH-v930-Volume-XII.pdf}}
\expandafter\def\csname r@cor:quartic231\endcsname{{25.25}{44}{}{proposition.25.25}{Dirichlet-GRH-v930-Volume-V.pdf}}
\expandafter\def\csname r@cor:shared201\endcsname{{25.21}{41}{}{proposition.25.21}{Dirichlet-GRH-v930-Volume-V.pdf}}
\expandafter\def\csname r@cor:stopped101\endcsname{{24.21}{18}{}{proposition.24.21}{Dirichlet-GRH-v930-Volume-V.pdf}}
\expandafter\def\csname r@cor:sum79\endcsname{{24.17}{15}{}{proposition.24.17}{Dirichlet-GRH-v930-Volume-V.pdf}}
\expandafter\def\csname r@cor:transfer100\endcsname{{11.18}{17}{}{proposition.11.18}{Dirichlet-GRH-v930-Volume-II.pdf}}
\expandafter\def\csname r@eq:V141\endcsname{{25.8}{31}{}{equation.25.8}{Dirichlet-GRH-v930-Volume-V.pdf}}
\expandafter\def\csname r@eq:count201\endcsname{{25.17}{41}{}{equation.25.17}{Dirichlet-GRH-v930-Volume-V.pdf}}
\expandafter\def\csname r@eq:floor115\endcsname{{24.19}{21}{}{equation.24.19}{Dirichlet-GRH-v930-Volume-V.pdf}}
\expandafter\def\csname r@eq:gap\endcsname{{6.3}{26}{}{equation.6.3}{Dirichlet-GRH-v930-Volume-I.pdf}}
\expandafter\def\csname r@eq:int201\endcsname{{25.15}{40}{}{equation.25.15}{Dirichlet-GRH-v930-Volume-V.pdf}}
\expandafter\def\csname r@eq:last182\endcsname{{104.10}{22}{}{equation.104.10}{Dirichlet-GRH-v930-Volume-XX.pdf}}
\expandafter\def\csname r@eq:length12\endcsname{{{L}}{30}{}{AMS.3}{Dirichlet-GRH-v930-Volume-I.pdf}}
\expandafter\def\csname r@eq:level146\endcsname{{25.13}{38}{}{equation.25.13}{Dirichlet-GRH-v930-Volume-V.pdf}}
\expandafter\def\csname r@eq:low136\endcsname{{87.5}{40}{}{equation.87.5}{Dirichlet-GRH-v930-Volume-XVII.pdf}}
\expandafter\def\csname r@eq:minors33\endcsname{{6.6}{46}{}{equation.6.6}{Dirichlet-GRH-v930-Volume-I.pdf}}
\expandafter\def\csname r@eq:moment141\endcsname{{25.7}{31}{}{equation.25.7}{Dirichlet-GRH-v930-Volume-V.pdf}}
\expandafter\def\csname r@eq:norm201\endcsname{{25.16}{40}{}{equation.25.16}{Dirichlet-GRH-v930-Volume-V.pdf}}
\expandafter\def\csname r@eq:orbit143\endcsname{{25.10}{35}{}{equation.25.10}{Dirichlet-GRH-v930-Volume-V.pdf}}
\expandafter\def\csname r@eq:window\endcsname{{{W}}{2}{}{AMS.2}{Dirichlet-GRH-v930-Volume-I.pdf}}
\expandafter\def\csname r@lem:coarse17\endcsname{{11.11}{11}{}{proposition.11.11}{Dirichlet-GRH-v930-Volume-II.pdf}}
\expandafter\def\csname r@lem:collisionclass457\endcsname{{55.33}{49}{}{proposition.55.33}{Dirichlet-GRH-v930-Volume-XI.pdf}}
\expandafter\def\csname r@lem:coord139\endcsname{{25.6}{28}{}{proposition.25.6}{Dirichlet-GRH-v930-Volume-V.pdf}}
\expandafter\def\csname r@lem:des28\endcsname{{24.8}{8}{}{proposition.24.8}{Dirichlet-GRH-v930-Volume-V.pdf}}
\expandafter\def\csname r@lem:dual241\endcsname{{36.2}{2}{}{proposition.36.2}{Dirichlet-GRH-v930-Volume-VII.pdf}}
\expandafter\def\csname r@lem:fit13\endcsname{{11.2}{3}{}{proposition.11.2}{Dirichlet-GRH-v930-Volume-II.pdf}}
\expandafter\def\csname r@lem:fourier231\endcsname{{25.22}{42}{}{proposition.25.22}{Dirichlet-GRH-v930-Volume-V.pdf}}
\expandafter\def\csname r@lem:fullinteger517\endcsname{{9.31}{132}{}{proposition.9.31}{Dirichlet-GRH-v930-Volume-I.pdf}}
\expandafter\def\csname r@lem:good175\endcsname{{72.18}{60}{}{proposition.72.18}{Dirichlet-GRH-v930-Volume-XIV.pdf}}
\expandafter\def\csname r@lem:int13\endcsname{{11.3}{4}{}{proposition.11.3}{Dirichlet-GRH-v930-Volume-II.pdf}}
\expandafter\def\csname r@lem:loss143\endcsname{{25.11}{33}{}{proposition.25.11}{Dirichlet-GRH-v930-Volume-V.pdf}}
\expandafter\def\csname r@lem:markmultiplicity425\endcsname{{54.2}{4}{}{proposition.54.2}{Dirichlet-GRH-v930-Volume-XI.pdf}}
\expandafter\def\csname r@lem:period30\endcsname{{24.11}{11}{}{proposition.24.11}{Dirichlet-GRH-v930-Volume-V.pdf}}
\expandafter\def\csname r@lem:ray88\endcsname{{7.10}{96}{}{proposition.7.10}{Dirichlet-GRH-v930-Volume-I.pdf}}
\expandafter\def\csname r@lem:shared201\endcsname{{25.18}{39}{}{proposition.25.18}{Dirichlet-GRH-v930-Volume-V.pdf}}
\expandafter\def\csname r@lem:simuldet340\endcsname{{28.15}{100}{}{proposition.28.15}{Dirichlet-GRH-v930-Volume-V.pdf}}
\expandafter\def\csname r@lem:size97\endcsname{{87.1}{33}{}{proposition.87.1}{Dirichlet-GRH-v930-Volume-XVII.pdf}}
\expandafter\def\csname r@nav:I-28\endcsname{{6.6}{32}{}{subsection.6.6}{Dirichlet-GRH-v930-Volume-I.pdf}}
\expandafter\def\csname r@nav:II-04\endcsname{{7.1}{88}{}{subsection.7.1}{Dirichlet-GRH-v930-Volume-I.pdf}}
\expandafter\def\csname r@nav:II-10\endcsname{{11.3}{9}{}{subsection.11.3}{Dirichlet-GRH-v930-Volume-II.pdf}}
\expandafter\def\csname r@nav:II-21\endcsname{{11.5}{14}{}{subsection.11.5}{Dirichlet-GRH-v930-Volume-II.pdf}}
\expandafter\def\csname r@nav:III-02\endcsname{{24.1}{3}{}{subsection.24.1}{Dirichlet-GRH-v930-Volume-V.pdf}}
\expandafter\def\csname r@nav:III-04\endcsname{{24.3}{6}{}{subsection.24.3}{Dirichlet-GRH-v930-Volume-V.pdf}}
\expandafter\def\csname r@nav:III-14\endcsname{{24.7}{16}{}{subsection.24.7}{Dirichlet-GRH-v930-Volume-V.pdf}}
\expandafter\def\csname r@nav:III-19\endcsname{{25.3}{28}{}{subsection.25.3}{Dirichlet-GRH-v930-Volume-V.pdf}}
\expandafter\def\csname r@prop:alloriginalpoints424\endcsname{{44.42}{34}{}{proposition.44.42}{Dirichlet-GRH-v930-Volume-IX.pdf}}
\expandafter\def\csname r@prop:curve71\endcsname{{34.3}{66}{}{proposition.34.3}{Dirichlet-GRH-v930-Volume-VI.pdf}}
\expandafter\def\csname r@prop:digits231\endcsname{{25.24}{44}{}{proposition.25.24}{Dirichlet-GRH-v930-Volume-V.pdf}}
\expandafter\def\csname r@prop:example30\endcsname{{24.15}{14}{}{proposition.24.15}{Dirichlet-GRH-v930-Volume-V.pdf}}
\expandafter\def\csname r@prop:field12\endcsname{{6.4}{27}{}{proposition.6.4}{Dirichlet-GRH-v930-Volume-I.pdf}}
\expandafter\def\csname r@prop:filteredalgebra427\endcsname{{54.11}{11}{}{proposition.54.11}{Dirichlet-GRH-v930-Volume-XI.pdf}}
\expandafter\def\csname r@prop:finite15\endcsname{{11.4}{6}{}{proposition.11.4}{Dirichlet-GRH-v930-Volume-II.pdf}}
\expandafter\def\csname r@prop:kummer30\endcsname{{24.12}{12}{}{proposition.24.12}{Dirichlet-GRH-v930-Volume-V.pdf}}
\expandafter\def\csname r@prop:level241\endcsname{{36.1}{2}{}{proposition.36.1}{Dirichlet-GRH-v930-Volume-VII.pdf}}
\expandafter\def\csname r@prop:localmult399\endcsname{{41.4}{32}{}{proposition.41.4}{Dirichlet-GRH-v930-Volume-VIII.pdf}}
\expandafter\def\csname r@prop:moment267\endcsname{{9.5}{110}{}{proposition.9.5}{Dirichlet-GRH-v930-Volume-I.pdf}}
\expandafter\def\csname r@prop:normexample340\endcsname{{28.18}{102}{}{proposition.28.18}{Dirichlet-GRH-v930-Volume-V.pdf}}
\expandafter\def\csname r@prop:planelabels379\endcsname{{40.11}{8}{}{proposition.40.11}{Dirichlet-GRH-v930-Volume-VIII.pdf}}
\expandafter\def\csname r@prop:quadratic345\endcsname{{36.3}{3}{}{proposition.36.3}{Dirichlet-GRH-v930-Volume-VII.pdf}}
\expandafter\def\csname r@prop:realclass100\endcsname{{11.19}{18}{}{proposition.11.19}{Dirichlet-GRH-v930-Volume-II.pdf}}
\expandafter\def\csname r@prop:res23\endcsname{{6.11}{32}{}{proposition.6.11}{Dirichlet-GRH-v930-Volume-I.pdf}}
\expandafter\def\csname r@prop:res80\endcsname{{7.7}{94}{}{proposition.7.7}{Dirichlet-GRH-v930-Volume-I.pdf}}
\expandafter\def\csname r@prop:short107\endcsname{{9.9}{114}{}{proposition.9.9}{Dirichlet-GRH-v930-Volume-I.pdf}}
\expandafter\def\csname r@sec:coho12\endcsname{{6.4}{28}{}{subsection.6.4}{Dirichlet-GRH-v930-Volume-I.pdf}}
\expandafter\def\csname r@sec:cyclicray657\endcsname{{26.10}{62}{}{subsection.26.10}{Dirichlet-GRH-v930-Volume-V.pdf}}
\expandafter\def\csname r@sec:missing\endcsname{{6.6}{32}{}{subsection.6.6}{Dirichlet-GRH-v930-Volume-I.pdf}}
\expandafter\def\csname r@sec:sourcelevels563\endcsname{{59.3}{7}{}{subsection.59.3}{Dirichlet-GRH-v930-Volume-XII.pdf}}
\expandafter\def\csname r@seriesguide\endcsname{{1.1}{4}{}{subsection.1.1}{Dirichlet-GRH-v930-Volume-I.pdf}}
\expandafter\def\csname r@thm:adjacentweights427\endcsname{{54.13}{13}{}{proposition.54.13}{Dirichlet-GRH-v930-Volume-XI.pdf}}
\expandafter\def\csname r@thm:adjoint240\endcsname{{34.16}{78}{}{proposition.34.16}{Dirichlet-GRH-v930-Volume-VI.pdf}}
\expandafter\def\csname r@thm:all24\endcsname{{24.6}{6}{}{proposition.24.6}{Dirichlet-GRH-v930-Volume-V.pdf}}
\expandafter\def\csname r@thm:alltails648\endcsname{{57.59}{113}{}{proposition.57.59}{Dirichlet-GRH-v930-Volume-XI.pdf}}
\expandafter\def\csname r@thm:alltrunc329\endcsname{{9.16}{120}{}{proposition.9.16}{Dirichlet-GRH-v930-Volume-I.pdf}}
\expandafter\def\csname r@thm:auxdescent361\endcsname{{36.18}{18}{}{proposition.36.18}{Dirichlet-GRH-v930-Volume-VII.pdf}}
\expandafter\def\csname r@thm:average269\endcsname{{8.3}{102}{}{proposition.8.3}{Dirichlet-GRH-v930-Volume-I.pdf}}
\expandafter\def\csname r@thm:char\endcsname{{24.1}{3}{}{proposition.24.1}{Dirichlet-GRH-v930-Volume-V.pdf}}
\expandafter\def\csname r@thm:circindex267\endcsname{{27.1}{75}{}{proposition.27.1}{Dirichlet-GRH-v930-Volume-V.pdf}}
\expandafter\def\csname r@thm:common111\endcsname{{9.13}{118}{}{proposition.9.13}{Dirichlet-GRH-v930-Volume-I.pdf}}
\expandafter\def\csname r@thm:common16\endcsname{{11.9}{10}{}{proposition.11.9}{Dirichlet-GRH-v930-Volume-II.pdf}}
\expandafter\def\csname r@thm:common178\endcsname{{72.30}{71}{}{proposition.72.30}{Dirichlet-GRH-v930-Volume-XIV.pdf}}
\expandafter\def\csname r@thm:compare30\endcsname{{24.13}{13}{}{proposition.24.13}{Dirichlet-GRH-v930-Volume-V.pdf}}
\expandafter\def\csname r@thm:copies137\endcsname{{87.12}{41}{}{proposition.87.12}{Dirichlet-GRH-v930-Volume-XVII.pdf}}
\expandafter\def\csname r@thm:critical13\endcsname{{11.1}{3}{}{proposition.11.1}{Dirichlet-GRH-v930-Volume-II.pdf}}
\expandafter\def\csname r@thm:criticalall310\endcsname{{28.6}{90}{}{proposition.28.6}{Dirichlet-GRH-v930-Volume-V.pdf}}
\expandafter\def\csname r@thm:effect146\endcsname{{25.17}{38}{}{proposition.25.17}{Dirichlet-GRH-v930-Volume-V.pdf}}
\expandafter\def\csname r@thm:evenbudget306\endcsname{{34.21}{83}{}{proposition.34.21}{Dirichlet-GRH-v930-Volume-VI.pdf}}
\expandafter\def\csname r@thm:exact241\endcsname{{8.1}{98}{}{proposition.8.1}{Dirichlet-GRH-v930-Volume-I.pdf}}
\expandafter\def\csname r@thm:exception15\endcsname{{11.7}{8}{}{proposition.11.7}{Dirichlet-GRH-v930-Volume-II.pdf}}
\expandafter\def\csname r@thm:forced310\endcsname{{28.9}{94}{}{proposition.28.9}{Dirichlet-GRH-v930-Volume-V.pdf}}
\expandafter\def\csname r@thm:graphtransport337\endcsname{{71.21}{38}{}{proposition.71.21}{Dirichlet-GRH-v930-Volume-XIV.pdf}}
\expandafter\def\csname r@thm:grh16\endcsname{{1.1}{1}{}{proposition.1.1}{Dirichlet-GRH-v930-Volume-I.pdf}}
\expandafter\def\csname r@thm:integer\endcsname{{30.11}{10}{}{proposition.30.11}{Dirichlet-GRH-v930-Volume-VI.pdf}}
\expandafter\def\csname r@thm:integer201\endcsname{{25.19}{40}{}{proposition.25.19}{Dirichlet-GRH-v930-Volume-V.pdf}}
\expandafter\def\csname r@thm:inter141\endcsname{{25.10}{32}{}{proposition.25.10}{Dirichlet-GRH-v930-Volume-V.pdf}}
\expandafter\def\csname r@thm:last182\endcsname{{104.26}{21}{}{proposition.104.26}{Dirichlet-GRH-v930-Volume-XX.pdf}}
\expandafter\def\csname r@thm:layers308\endcsname{{28.4}{88}{}{proposition.28.4}{Dirichlet-GRH-v930-Volume-V.pdf}}
\expandafter\def\csname r@thm:levels88\endcsname{{7.11}{97}{}{proposition.7.11}{Dirichlet-GRH-v930-Volume-I.pdf}}
\expandafter\def\csname r@thm:localvalue309\endcsname{{28.8}{92}{}{proposition.28.8}{Dirichlet-GRH-v930-Volume-V.pdf}}
\expandafter\def\csname r@thm:logunits319\endcsname{{28.11}{96}{}{proposition.28.11}{Dirichlet-GRH-v930-Volume-V.pdf}}
\expandafter\def\csname r@thm:long232\endcsname{{5.6}{23}{}{proposition.5.6}{Dirichlet-GRH-v930-Volume-I.pdf}}
\expandafter\def\csname r@thm:low136\endcsname{{87.10}{40}{}{proposition.87.10}{Dirichlet-GRH-v930-Volume-XVII.pdf}}
\expandafter\def\csname r@thm:markedglobal339\endcsname{{71.25}{41}{}{proposition.71.25}{Dirichlet-GRH-v930-Volume-XIV.pdf}}
\expandafter\def\csname r@thm:moment141\endcsname{{25.9}{30}{}{proposition.25.9}{Dirichlet-GRH-v930-Volume-V.pdf}}
\expandafter\def\csname r@thm:norm233\endcsname{{94.15}{49}{}{proposition.94.15}{Dirichlet-GRH-v930-Volume-XVIII.pdf}}
\expandafter\def\csname r@thm:normsample340\endcsname{{28.16}{101}{}{proposition.28.16}{Dirichlet-GRH-v930-Volume-V.pdf}}
\expandafter\def\csname r@thm:openindex269\endcsname{{27.3}{77}{}{proposition.27.3}{Dirichlet-GRH-v930-Volume-V.pdf}}
\expandafter\def\csname r@thm:originalorder424\endcsname{{44.44}{35}{}{proposition.44.44}{Dirichlet-GRH-v930-Volume-IX.pdf}}
\expandafter\def\csname r@thm:params238\endcsname{{34.15}{77}{}{proposition.34.15}{Dirichlet-GRH-v930-Volume-VI.pdf}}
\expandafter\def\csname r@thm:positive306\endcsname{{28.1}{84}{}{proposition.28.1}{Dirichlet-GRH-v930-Volume-V.pdf}}
\expandafter\def\csname r@thm:prime210\endcsname{{93.17}{34}{}{proposition.93.17}{Dirichlet-GRH-v930-Volume-XVIII.pdf}}
\expandafter\def\csname r@thm:quartic115\endcsname{{24.25}{21}{}{proposition.24.25}{Dirichlet-GRH-v930-Volume-V.pdf}}
\expandafter\def\csname r@thm:ray101\endcsname{{24.20}{18}{}{proposition.24.20}{Dirichlet-GRH-v930-Volume-V.pdf}}
\expandafter\def\csname r@thm:res100\endcsname{{11.17}{17}{}{proposition.11.17}{Dirichlet-GRH-v930-Volume-II.pdf}}
\expandafter\def\csname r@thm:ring39\endcsname{{32.1}{23}{}{proposition.32.1}{Dirichlet-GRH-v930-Volume-VI.pdf}}
\expandafter\def\csname r@thm:scalar231\endcsname{{25.23}{43}{}{proposition.25.23}{Dirichlet-GRH-v930-Volume-V.pdf}}
\expandafter\def\csname r@thm:scale232\endcsname{{110.18}{53}{}{proposition.110.18}{Dirichlet-GRH-v930-Volume-XXI.pdf}}
\expandafter\def\csname r@thm:shared201\endcsname{{25.20}{41}{}{proposition.25.20}{Dirichlet-GRH-v930-Volume-V.pdf}}
\expandafter\def\csname r@thm:simultaneous307\endcsname{{34.22}{85}{}{proposition.34.22}{Dirichlet-GRH-v930-Volume-VI.pdf}}
\expandafter\def\csname r@thm:strict113\endcsname{{11.22}{21}{}{proposition.11.22}{Dirichlet-GRH-v930-Volume-II.pdf}}
\expandafter\def\csname r@thm:strictcyclic657\endcsname{{26.23}{64}{}{proposition.26.23}{Dirichlet-GRH-v930-Volume-V.pdf}}
\expandafter\def\csname r@thm:strictlarge425\endcsname{{54.3}{5}{}{proposition.54.3}{Dirichlet-GRH-v930-Volume-XI.pdf}}
\expandafter\def\csname r@thm:ststrict616\endcsname{{32.23}{47}{}{proposition.32.23}{Dirichlet-GRH-v930-Volume-VI.pdf}}
\expandafter\def\csname r@thm:traces\endcsname{{30.8}{8}{}{proposition.30.8}{Dirichlet-GRH-v930-Volume-VI.pdf}}
\expandafter\def\csname r@thm:trunc15\endcsname{{11.5}{6}{}{proposition.11.5}{Dirichlet-GRH-v930-Volume-II.pdf}}
\expandafter\def\csname r@thm:twists\endcsname{{5.1}{19}{}{proposition.5.1}{Dirichlet-GRH-v930-Volume-I.pdf}}
\expandafter\def\csname r@thm:two17\endcsname{{11.12}{13}{}{proposition.11.12}{Dirichlet-GRH-v930-Volume-II.pdf}}
\expandafter\def\csname r@thm:units180\endcsname{{104.25}{20}{}{proposition.104.25}{Dirichlet-GRH-v930-Volume-XX.pdf}}
\expandafter\def\csname r@thm:unramdet330\endcsname{{28.13}{98}{}{proposition.28.13}{Dirichlet-GRH-v930-Volume-V.pdf}}
\expandafter\def\csname r@thm:wild98\endcsname{{25.1}{24}{}{proposition.25.1}{Dirichlet-GRH-v930-Volume-V.pdf}}
\makeatother
\begin{document}
\setcounter{part}{3}
\fancyhead[L]{\small\color{muted}Dirichlet GRH \textperiodcentered\ Volume IV}
\fancyhead[R]{\small\color{muted}Cyclotomic descent and norms}
\thispagestyle{plain}\begin{center}
{\LARGE\bfseries The generalized Riemann hypothesis\\[2pt]for Dirichlet \(L\)-functions}\\[6pt]
{\small 27 November 2026 \quad\textperiodcentered\quad Edition 930}\end{center}
\part{Cyclotomic descent and relative strict orders}\label{part:IV}
\begin{abstract}\noindent We descend in the real cyclotomic tower, determine the ranges of coefficient primes and track exceptional twists. Norm maps compare the strict orders at consecutive layers.\end{abstract}
\setcounter{section}{18}
\noindent The \hyperref[seriesguide]{guide in Volume~I} lists all twenty-seven volumes.
\begingroup\small\setlength{\parskip}{0pt}\tableofcontents\endgroup
\clearpage

\section{Descent in the real cyclotomic tower}\label{nav:g-II-1}\label{nav:volume-IV}

\subsection{A bound at all primes for a field of given degree}\label{nav:II-13}
For a number field \(L\), a rational prime \(p\) and an integer \(r\ge1\), set
\[
 A_{p,r}=\QQ_p/\ZZ_p(-r),\qquad
 \Sel_{p,r}(L)=\ker\!\left(H^1(G_{L,S_p},A_{p,r})
                  \longrightarrow\prod_{w\mid p}H^1(L_w,A_{p,r})\right).
\]
Here \(G_{L,S_p}\) allows ramification at the finite places over \(p\) only. The twist means that the cyclotomic character \(\kappa_p\) acts to the power \(-r\). All cohomology is continuous. We write \(D_L\) for the absolute discriminant and \(d_L=[L:\QQ]\).

\begin{theorem}\label{thm:bound}
There is an absolute constant \(C\) with the following property. Let \(L\subseteq\QQ(\zeta_M)\) be totally imaginary and \(r\ge1\). Then
\begin{equation}
 \sum_p(\log p)\len_{\ZZ_p}\Sel_{p,r}(L)
 \ \le\ \exp\!\big(C[d_L(r+2)]^{10}\log(2M d_L(r+2))\big).
 \label{eq:all}
\end{equation}
The sum includes both conjugation eigenspaces at every odd prime and the full exponent of each group.
\end{theorem}

We prove an exponent bound and a bound on the number of generators separately. It matters that a single finite group accounts for the sum at all sufficiently large primes.

\begin{lemma}\label{lem:top}
Set \(i=2r,\ n=2i+8,\ d=d_L,\ W=3^{dn}\). There are positive integers \(J,T\) such that
\[
 \log(2JT)\le\exp\!\big(C(nd)^9\log(2Mnd)\big).
\]
The integer \(p^{i+4}JT\) kills \(\Sel_{p,r}(L)\) for \(p\le W\). All larger primes together satisfy
\[
          \sum_{p>W}\log|\Sel_{p,r}(L)|\le\log T .
\]
\end{lemma}
\begin{proof}
First, the dual of \(\Sel_{p,r}(L)\) embeds in \(K_i(\OO_L)[p^\infty]\). Here and below the dual is the group of continuous characters into \(\QQ_p/\ZZ_p\). Indeed, Poitou--Tate duality places this dual in \(H^2_{\mathrm{\acute et}}(\OO_L[1/p],\ZZ_p(r+1))\) as the kernel of localization. The archimedean terms vanish since all archimedean places are complex. The Quillen--Lichtenbaum comparison identifies this \(H^2\) with the \(p\)-adic completion of \(K_{2r}(\OO_L[1/p])\). The other possible term in even degree is \(H^0(\ZZ_p(r))=0\) because the cyclotomic image is open and \(r>0\). In the localization sequence the finite residue fields at \(p\) have zero \(K_{2r}\) and \(K_{2r-1}\) of order prime to \(p\). Thus one can use \(\OO_L\), and \(K_{2r}(\OO_L)\) is finite by finite generation and Borel's rank formula. These are the comparisons used in Proposition \ref{prop:finite15}.~\cite{IV-II-II-duality13,IV-II-III-weibel,IV-II-III-kahn}

Let \(\Gamma=\GL_n(\OO_L)\), and let \(\Gamma_3\) be the kernel of its action on \(\OO_L^n/3\OO_L^n\). Its index \(J\) divides \(|\GL_{dn}(\mathbf F_3)|\le3^{(dn)^2}\). The principal congruence kernel is torsion free. Indeed, an element of prime order, viewed in \(\GL_{dn}(\ZZ)\), would have the form \(I+3^aB\), \(a\ge1\), \(B\not\equiv0\bmod3\). Taking its \(p\)-th power gives a nonzero leading term \(p3^aB\) modulo \(3^{a+1}\) for \(p\ne3\), and \(3^{a+1}B\) modulo \(3^{a+2}\) for \(p=3\). This excludes every element of prime order, hence every element of finite order other than the identity.

Put \(T=|H_i(\Gamma_3,\ZZ)_{\rm tors}|\). We track the dependence on \(d,n\) in the cell count of Lemma \ref{lem:coarse17}. Ash's well-rounded retract for the free lattice \(\OO_L^n\) is contractible and has finite cell stabilizers. Its cells are indexed by sets of shortest vectors of normalized metrics. Corollary 5.2 of Bayer-Fluckiger, Emery and Houriet gives a finite set \(\Phi\) into which an element of \(\Gamma\) moves all shortest vectors of any one such metric. Its bound, valid since \(L\ne\QQ\) and \(n\ge5\), is
\[
 |\Phi|\le n^{3n^2d^2}\,d^{5n^3d^4}\,D_L^{9n^3d^3}.
\]
In particular \(\log(2+|\Phi|)\le C(nd)^4\log(2ndD_L)\).~\cite[Sec.~2 and Cor.~5.2]{IV-II-III-beh}

For completeness, the real vector space of metrics before positivity and normalization has dimension at most \(D=dn^2\). The equation that a vector has squared trace length one is affine linear there. A consistent set of these equations is implied by at most \(D+1\) of its members. Thus a cell can be encoded by at most \(D+1\) of its shortest vectors. Indeed, if two choices of metrics give the same encoding, the full equations from either cell hold on both metrics, so the two sets of shortest vectors coincide. Along a chain of incident cells the shortest-vector sets are nested in reverse order. Move the largest such set into \(\Phi\). This encodes an entire flag by at most \(D+1\) lists, each of length at most \(D+1\). Allowing a padding symbol gives at most \((2+|\Phi|)^{(D+2)^2}\) orbits of simplices in the barycentric subdivision. The quotient by \(\Gamma_3\) has at most \(J\) times that many simplices and computes its group homology. Call this last upper bound \(H\). Each column of its integral boundary in degree \(i+1\) has sum of absolute values at most \(i+2\), even if faces become identified. Smith normal form and Hadamard's inequality give \(\log T\le H\log(i+2)\), since homology embeds in the cokernel of this boundary map. The conductor--discriminant formula gives \(D_L\le M^d\). Substitution, using \(D+2\ll(nd)^2\), proves the displayed bound for \(\log(2JT)\).

Stability in the range \(n\ge2i+1\) identifies \(H_i(\Gamma,\ZZ)\) with the homology of \(B\GL(\OO_L)^+\). Transfer from \(\Gamma\) to \(\Gamma_3\), followed by inclusion, is multiplication by \(J\). It follows that \(JT\) kills the torsion in \(H_i(\Gamma,\ZZ)\). By Quillen's plus construction the ordinary Hurewicz map on this space has source \(K_i(\OO_L)\). Its kernel at \(p\) is killed by \(p^{i+4}\). To see this last assertion, the adjunction from the suspension spectrum of the basepoint component of \(\Omega^\infty\mathbf K(\OO_L)\) to the connective \(K\)-spectrum takes the ordinary Hurewicz image of a sphere to its stable Hurewicz image. Thus Mathew's upper bound for the kernel of the stable map applies to the ordinary kernel. His Theorem 1.3(a,b) gives exponents at most \(i+4\) at all primes. This proves the annihilator assertion.~\cite[Sec.~3.1]{IV-II-III-beh}\cite[Thm.~1.3]{IV-II-III-mathew}

Every prime divisor of \(J\) is at most \(3^{dn}\), since the factors of \(|\GL_{dn}(\mathbf F_3)|\) are powers of three and factors \(3^h-1\), \(h\le dn\). For \(p>W\) the ordinary Hurewicz map is injective on \(p\)-torsion by Arlettaz's theorem, which applies already for \(p>(i+1)/2\). Transfer is injective on this primary part as well, since \(p\nmid J\). The direct sum of the large-primary parts of \(K_i(\OO_L)\) therefore embeds in the single finite torsion group of order \(T\). Apply the injection at the beginning of the proof and take logarithms.~\cite[Sec.~3.1]{IV-II-III-beh}
\end{proof}

\begin{lemma}\leavevmode\label{lem:rank}
For every prime \(p\), the minimal number \(b_{p,r}\) of cyclic summands of \(\Sel_{p,r}(L)\) satisfies
\[
 b_{p,r}\log p\le C d_L p\log(2Mp).
\]
\end{lemma}
\begin{proof}
The group is finite by Lemma~\ref{lem:top}, so \(b_{p,r}\) is its dimension of \(p\)-torsion. The multiplication-by-\(p\) exact sequence on \(A_{p,r}\) bounds that dimension by \(\dim_{\mathbf F_p} H^1(G_{L,S_p},\mathbf F_p(-r))\). Put \(K=L(\mu_p)\), \(a=[K:\QQ]\le d_Lp\). Restriction to \(K\) is injective by corestriction since \([K:L]\mid p-1\). This extension is unramified outside \(p\). Over \(K\) the coefficient becomes isomorphic to \(\mu_p\). For \(R=\OO_K[1/p]\), Kummer theory gives an extension of \(\Cl(R)[p]\) by \(R^\times/(R^\times)^p\) as this degree-one group. The unit term has dimension at most \(2a\) by the \(S\)-unit theorem. Each ideal class in \(K\) has a representative of norm at most \(B=2^aD_K^{1/2}\) by Minkowski's bound. There are at most \(B^2\zeta_K(2)\le B^2\zeta(2)^a\) integral ideals of such norm, by the Dirichlet series and the Euler factors. Since \(\Cl(R)\) is a quotient of \(\Cl(\OO_K)\) and \(D_K\le(Mp)^a\), the class group contributes at most \(Ca\log(2Mp)/\log p\) to the dimension. Adding the unit term proves the lemma.
\end{proof}

\begin{proof}[Proof of Theorem~\ref{thm:bound}]
For \(p\le W\), a cyclic summand has order at most \(p^{i+4}JT\). Its logarithmic order is therefore at most \((i+4)\log W+\log(2JT)\). Sum over the generator bound of Lemma~\ref{lem:rank}, using \(\log p\ge\log 2\) and \(\sum_{p\le W}p\le W^2\). The sum at these primes is bounded by
\[
 C d W^2\log(2MW)\big((i+4)\log W+\log(2JT)\big).
\]
Add the single \(\log T\) from Lemma~\ref{lem:top}. Since \(\log W=dn\log 3\) and \(n=4r+8\), increasing an absolute constant absorbs all factors in \eqref{eq:all}.
\end{proof}

\subsection{Descent in the real cyclotomic tower}\label{nav:II-14}
Fix a primitive Dirichlet character \(\psi\) of conductor \(F\) and order \(o\), with fixed field \(K_\psi\) cut out by \(\ker\psi\). Set \(0<\eps<1/52\), \(\eta=\eps/128\), \(C_j=\QQ(\zeta_{2^{j+2}})^+\) for \(j\ge0\), and let \(l\) be the least integer with \(2^l\ge x^{2+8\eps}\). The field in the prime count and a subfield of it are
\[
 \begin{split}
  E_x&=K_\psi C_l(\sqrt{-q},\sqrt{-d}),\\
  L_x&=K_\psi C_s(\sqrt{-q},\sqrt{-d}),\qquad
          s=\lfloor\eta\log_2 x\rfloor .
 \end{split}
\]
The auxiliary primes \(q,d\) are those of Theorem \ref{thm:twists}. On its range they can be held fixed with \(\log(qd)\ll_{\psi,\eps}x^{\eps/5}(\log(2x))^2\), and satisfy \(q,d\equiv3\bmod4\). In particular \(M_x=2^{l+2}Fqd\) is a cyclotomic modulus for both fields. We may use the same primes on the shorter intervals below.

\begin{lemma}[Descent of strict classes]\label{lem:desc}
Let \(E/L\) be a finite Galois extension of degree prime to \(p\), with group \(H\). Restriction identifies \(\Sel_{p,r}(L)\) with \(\Sel_{p,r}(E)^H\).
\end{lemma}
\begin{proof}
Work first with the absolute Galois groups so that ramification of \(E/L\) away from \(p\) is allowed. In inflation--restriction the higher cohomology of \(H\) on \(A_{p,r}^{G_E}\) vanishes. Multiplication by \(|H|\) is invertible there, and averaging gives exactness of invariants. Hence restriction is an isomorphism onto invariants on global \(H^1\). For \(v\nmid p\) the twist is unramified, so a cocycle restricted to inertia is a homomorphism and a coboundary restricts to zero. If it vanishes on inertia over \(E\), it factors through an inertia quotient of order dividing \(|H|\). Such a homomorphism into a \(p\)-primary group vanishes. Thus the descended class is unramified away from \(p\), which is precisely the condition for inflation from \(G_{L,S_p}\). At \(v\mid p\), restriction followed by corestriction is multiplication by the local degree, a \(p\)-adic unit. Vanishing after restriction implies vanishing over \(L_v\). The two local conditions are preserved in the forward direction as well.
\end{proof}

\begin{theorem}[Both signs on a growing subfield]\label{thm:growing}
Write \(H_x=\Gal(E_x/L_x)\). On the common intervals of Theorem \ref{thm:two17}, uniformly for all sufficiently large \(x\) in the interval,
\begin{equation}
 \sum_{1\le r\le x^\eta}\ \sum_{p>2}(\log p)
              \len_{\ZZ_p}\Sel_{p,r}(E_x)^{H_x}
       \le \exp(x^{2\eps/3}).
 \label{eq:grow}
\end{equation}
The inner sum is over primes. The invariants are a direct summand over \(\ZZ_p\); the estimate also holds for either complex-conjugation eigenspace separately.
\end{theorem}
\begin{proof}
The extension \(E_x/L_x\) has degree a power of two. For odd \(p\) the averaging idempotent \(|H_x|^{-1}\sum_{\sigma\in H_x}\sigma\) projects onto its invariants. It commutes with conjugation since \(E_x/\QQ\) is abelian. Apply Lemma~\ref{lem:desc}. The field \(L_x\) contains \(\QQ(\sqrt{-q})\) and has degree between \(2^s\) and \(4o2^s\). Moreover \(\log(2M_x d_{L_x}(r+2))\ll_{\psi,\eps}x^{\eps/4}\) on the indicated range. The exponent on the right of \eqref{eq:all} is then
\[
 \ll_{\psi,\eps} x^{20\eta+\eps/4}=x^{13\eps/32}.
\]
Summing at most \(x^\eta\) copies of its exponential proves \eqref{eq:grow}. In particular the bound is \(o(\exp(x^\eps/2))\), where the exponent in the latter expression is \(\tfrac12x^\eps\).
\end{proof}

\subsection{Fixed odd primes through the full range of twists}\label{nav:II-15}
Fix an integer \(a\ge0\) and a finite set \(\Pi\) of odd primes. The field \(B_a=K_\psi C_a\) embeds in every sufficiently high \(E_x\). Write \(H_{a,x}=\Gal(E_x/B_a)\), a group of order a power of two. The parameters \(a,\Pi\) can be prescribed before choosing an interval.

\begin{theorem}\label{thm:fixed}
There are arbitrarily large \(V\) and a single pair \(q,d\) for each interval \([V,V^{3/2}]\) such that at every \(x\) there, with \(t(x)=\lfloor x^{401}\rfloor\),
\begin{equation}
 \sum_{p\in\Pi}(\log p)\sum_{1\le r<t(x)}
       \len_{\ZZ_p}\Sel_{p,r}(E_x)^{H_{a,x}}
       \le \exp(x^{2\eps/3}).
 \label{eq:fixed}
\end{equation}
These same fields and intervals satisfy \eqref{eq:grow}, the full dyadic bound of Theorem \ref{thm:two17}, and the odd-prime bound for the sign \((-1)^{r+1}\) in Theorem~\ref{thm:critical13}.
\end{theorem}

\begin{lemma}\label{lem:sparse}
Let \(B\) be a fixed totally imaginary abelian number field and \(p\) a fixed odd prime. There is an integer \(D=D(B,p)\ge0\) such that, for every \(T_0\ge2\), all but at most \(D\) of the integers \(1\le r\le T_0\) satisfy
\[
       \len_{\ZZ_p}\Sel_{p,r}(B)\ll_{B,p}\log(2T_0).
\]
\end{lemma}
\begin{proof}
Set \(B_\infty=B(\mu_{p^\infty})\), \(\Gamma=\Gal(B_\infty/B)\). Over \(B(\mu_p)\) this is the cyclotomic \(\ZZ_p\)-extension. Its unramified abelian pro-\(p\) class module is finitely generated over \(\ZZ_p\). Indeed, Iwasawa's theorem realizes it as a finitely generated torsion module over the one-variable Iwasawa algebra, and Ferrero--Washington gives \(\mu=0\) because \(B(\mu_p)\) is abelian over \(\QQ\). The structure theorem then gives finite generation over \(\ZZ_p\), including its finite submodule. Let \(X'\) be the quotient imposing complete splitting at the primes above \(p\). It is also finitely generated over \(\ZZ_p\).~\cite{IV-II-II-iwasawa15,IV-II-III-fw}

To relate \(X'\) to degree one, restrict a cocycle to the kernel of the cyclotomic character in \(G_{B,S_p}\). The action on \(A_{p,r}\) becomes trivial and the restriction is a homomorphism. A strict class restricts to zero on all decomposition groups over \(p\), so its homomorphism factors through \(X'\). This uses only the maximal unramified quotient, since the initial Galois group already forbids ramification outside \(p\). This description also follows by taking the maximal pro-\(p\) quotient of the kernel and killing its decomposition groups over \(p\); the resulting extension is everywhere unramified and split over \(p\), so it is a quotient of the unramified class module. Conjugation makes the homomorphism \(\Gamma\)-equivariant with scalar action \(\kappa_p^{-r}\) on its values.

Restriction is injective here. Write \(\Gamma_1=\Gamma\cap(1+p\ZZ_p)\) under \(\kappa_p\), choose its topological generator \(\gamma\), and put \(u=\kappa_p(\gamma)\). The procyclic resolution computes \(H^1(\Gamma_1,A_{p,r})=A_{p,r}/(u^{-r}-1)A_{p,r}=0\), since \(u^{-r}-1\ne0\) and the coefficient group is divisible. The index \([\Gamma:\Gamma_1]\) is prime to \(p\), so \(H^1(\Gamma,A_{p,r})=0\) by inflation--restriction and averaging. Thus \(\Sel_{p,r}(B)\) injects into the dual of \(X'/(\gamma-u^{-r})X'\). An injection suffices for the estimate.

Take a minimal generating surjection \(\pi:\ZZ_p^D\to X'\) and lift the images under \(\gamma\) of its \(D\) generators to a matrix \(A\) satisfying \(\pi A=\gamma\pi\). Reduction modulo \(p\) makes \(\pi\) an isomorphism on generator spaces, so \(A\) is invertible. If \(D=0\), all groups vanish. Otherwise right exactness and Smith normal form give
\[
 \len_{\ZZ_p}\Sel_{p,r}(B)\le v_p(\det(A-u^{-r}I))
\]
whenever the determinant is nonzero. All valuations are normalized by \(v_p(p)=1\). Put \(e=v_p(u-1)\ge1\), \(J_0=\lceil\log_p(T_0+1)\rceil\). The binomial formula at odd \(p\) gives
\[
 v_p(u^{-r}-u^{-s})=e+v_p(r-s)\le e+J_0\qquad(1\le r\ne s\le T_0).
\]
For the equality, raising a principal unit to a power prime to \(p\) preserves the valuation of its difference from one, and taking its \(p\)-th power increases it by exactly one. If the determinant valuation exceeds \(D(e+J_0)\), including when it is infinite, some eigenvalue \(\alpha\) of \(A\) satisfies \(v_p(\alpha-u^{-r})>e+J_0\). Two distinct indices cannot satisfy this for the same eigenvalue by the displayed spacing and the ultrametric inequality. At most \(D\) indices are exceptional. The other indices satisfy the asserted bound.
\end{proof}

\begin{proof}[Proof of Theorem~\ref{thm:fixed}]
Let \(B=B_a(i)\). It is totally imaginary and abelian. Lemma~\ref{lem:desc} gives an injection from the strict groups of \(B_a\) into those of \(B\) at every odd prime, since their relative degree divides two. It also identifies the groups in \eqref{eq:fixed} with those over \(B_a\). For \(p\in\Pi\), take \(D_p\) from Lemma~\ref{lem:sparse} for \(B\). Let \(R_2(\psi)\) be the uniform bound on the number of exceptional twists at two in Theorem \ref{thm:exception15}, and set
\[
 R=R_2(\psi)+\sum_{p\in\Pi}D_p,\qquad b=401/\eta,\qquad K=2(2b)^R .
\]
These constants are fixed before the auxiliary primes. For arbitrarily large \(Z\), choose \(q,d\) by Theorem \ref{thm:twists} at the initial scale \(Z\) with upper endpoint \(Z^{K+1}\). More explicitly, that theorem supplies fixed \(q,d\) with the character estimates throughout \([Z,Z^{K+1}]\) and \(\log(qd)\ll_{\psi,\eps,K}Z^{\eps/5}(\log Z)^2\). Apply Lemma~\ref{lem:sparse} with \(T_0=\lfloor Z^{401K}\rfloor\) and take the union of its exceptional sets over \(\Pi\). Adjoin the exceptional set at two for the whole master interval \([Z,Z^K]\). Theorem \ref{thm:common16} constructs the latter by using this same \(T_0\) in its dyadic determinant estimate. Its cardinality is at most \(R_2(\psi)\), independently of \(K,Z,q,d\). The union has at most \(R\) elements.

Start with \(U=K\log Z\) and associate to each exceptional \(r\) the point \(a_r=(\log r)/401\). If any point is in \((U/(2b),U]\), discard it and replace \(U\) by \(U/(2b)\). Discarded points are above the new range. After at most \(R\) divisions the band is empty. Put \(V=\exp(U/2)\). Then \(V\ge Z,\ V^2\le Z^K\). If \(x\in[V,V^2]\) and an exceptional \(r\le x^{401}\), then \(a_r\le U\), so \(a_r\le U/(2b)\) and \(r\le x^\eta\). Restrict to \([V,V^{3/2}]\). The field and weight constructions in Theorem \ref{thm:common16} now apply. The proof of Theorem~\ref{thm:two17} works with any smaller positive \(\eta\) than \(\eps/64\), since this only shortens the exceptional twist range. This gives the asserted simultaneous dyadic bound. The bounds for the critical sign and \eqref{eq:grow} held already before shortening.

For a nonexceptional twist the contribution from \(\Pi\) in \eqref{eq:fixed} is \(O_{\psi,\eps,a,\Pi}(\log(2x))\), since \(\log(2T_0)\ll_K\log(2x)\). The \(O(R)\) exceptional twists with \(r<t(x)\) satisfy \(r\le x^\eta\). For them apply \eqref{eq:all} to the fixed field \(B\) and a fixed modulus for \(B\). Its exponent is \(O_{\psi,a,\eps}(x^{10\eta}\log(2x))\). Consequently the entire sum is at most
\[
 C_{\psi,\eps,a,\Pi}\,t(x)\log(2x)
       +R\exp\!\big(C_{\psi,\eps,a} x^{10\eta}\log(2x)\big).
\]
For sufficiently large \(Z\) this is at most \(\exp(x^{2\eps/3})\), as claimed.
\end{proof}

\begin{corollary}\label{cor:diag}
There is a sequence of disjoint intervals with \(V_j\to\infty\) on which \eqref{eq:fixed} holds with \(a=j\) and \(\Pi\) the first \(j\) odd primes.
\end{corollary}
\begin{proof}
At step \(j\), fix these two parameters and use Theorem~\ref{thm:fixed} with a lower endpoint larger than the preceding right endpoint and all thresholds for the new parameters.
\end{proof}

This choice by successive intervals gives no prescribed rate of growth for the number of primes in terms of \(V_j\).

\subsubsection{Stabilization over the full real cyclotomic tower}\label{nav:II-16}
\label{sec:stable27}

Fix an odd prime \(p\). The fields \(C_j=\QQ(\zeta_{2^{j+2}})^+\) are the layers of the cyclotomic \(\ZZ_2\)-extension. We first describe exactly which degree-one classes detect the \(p\)-torsion of a strict group. This description will be independent of the positive twist.

\begin{lemma}[Characters that are locally cyclotomic]\label{lem:local27}
Let \(L\) be a number field containing \(\mu_p\). Write \(Y(L)\) for the space of continuous additive characters \(G_L\to\mathbf F_p\) unramified outside \(p\) whose restriction at each \(w\mid p\) factors through \(\kappa_p(G_{L_w})\). For every \(r\ge1\), identify \(A_{p,r}[p]\) with \(\mathbf F_p\) using the coordinate \(a/p\mapsto a\bmod p\). There is an exact sequence
\begin{equation}
 0\longrightarrow \mathbf F_p\lambda_L\longrightarrow Y(L)
      \longrightarrow\Sel_{p,r}(L)[p]\longrightarrow0,        \label{eq:line27}
\end{equation}
where \(\lambda_L\) is a nonzero additive character of the global cyclotomic image of \(G_L\), pulled back to \(G_L\). The line it spans is independent of \(r\). This sequence is equivariant under automorphisms over any base containing \(\mu_p\).
\end{lemma}
\begin{proof}
The finite coefficient action is trivial over \(L\), so finite-coefficient cocycles are homomorphisms. From multiplication by \(p\) on \(A=A_{p,r}\), every \(p\)-torsion cohomology class lifts to one of these homomorphisms. If its divisible-coefficient class is strict, then at every inertia group away from \(p\) it is a coboundary. The cyclotomic character is trivial there, so the homomorphism is zero there. On \(G_{L_w}\) for \(w\mid p\) it has the form \(g\mapsto(\kappa_p(g)^{-r}-1)z_w\) in \(A\). Hence it factors through the local cyclotomic image.

Conversely that image is an open subgroup of \(1+p\ZZ_p\), since \(L_w/\QQ_p\) is finite and contains \(\mu_p\). For a topological generator \(u\), multiplication by \(u^{-r}-1\ne0\) is surjective on \(A\). A continuous cocycle on this procyclic group with values in \(A\) is therefore a coboundary: match its value at the generator and then use continuity and density of the integral powers. Thus every character of \(Y(L)\) maps to a strict class.

The global image is likewise \(1+p^a\ZZ_p\) for some \(a\ge1\). The logarithm identifies its power \(-r\) with multiplication by \(-r\) on \(p^a\ZZ_p\). Hence the ideal generated by all \(\kappa_p(g)^{-r}-1\) is \(p^s\ZZ_p\), with \(s=a+v_p(r)\). The global invariants of \(A\) are \(p^{-s}\ZZ_p/\ZZ_p\). The kernel of the finite-to-divisible map is their quotient by multiplication by \(p\), of dimension one. The connecting map sends a generator to the homomorphism \(g\mapsto(\kappa_p(g)^{-r}-1)p^{-s-1}\bmod\ZZ_p\). This factors through the global cyclotomic image and is nonzero on a generator. The space of additive characters of that image into \(\mathbf F_p\) has dimension one, giving the stated independence. An automorphism fixing \(\mu_p\) acts trivially on the finite coefficient coordinate and preserves \(\kappa_p\) under conjugation. These observations give equivariance.
\end{proof}

\begin{theorem}[Simultaneous stabilization in every twist]\label{thm:stable27}
Let \(B/\QQ\) be a fixed finite abelian field and \(B_j=BC_j\). There is an integer \(a=a(B,p)\), independent of \(r\), such that restriction is an isomorphism
\[
       \Sel_{p,r}(B_a)\ \longrightarrow\ \Sel_{p,r}(B_j)
       \qquad(j\ge a,\ r\ge1).
\]
The assertion concerns the full groups, including all exponents and both conjugation signs.
\end{theorem}
\begin{proof}
Suppose first that \(\mu_p\subset B\). Washington's class number theorem for the cyclotomic \(\ZZ_2\)-extension of an abelian number field states that the \(p\)-parts of the class numbers of its layers eventually have constant order. The fields \(B_j\) give these layers after any repetitions at the bottom. Extension of ideals on \(p\)-primary ideal class groups is injective, since extension followed by norm is multiplication by the relative degree, a power of two. After a fixed layer it is therefore bijective. In particular \(H=\Gal(B_j/B_a)\) acts trivially on the \(p\)-primary ideal class group for all \(j\ge a\) once \(a\) is sufficiently large.~\cite{IV-II-III-wash27}

The primes above \(p\) also stop splitting after a fixed layer. In the real cyclotomic tower over \(\QQ\), Frobenius at \(p\) has image \(p\) in \(\ZZ_2^\times/\{\pm1\}\simeq\ZZ_2\). Its closed cyclic subgroup is open: \(p\ne \pm1\) in \(\ZZ_2\), and every nonzero closed subgroup of \(\ZZ_2\) is open. Passing to the absolute Galois group of a finite local extension at \(p\) reduces this image by at most finite index. Thus the decomposition group of every prime of \(B\) over \(p\) is open in the tower over \(B\). Enlarge \(a\) to ensure that \(H\) fixes every prime of \(B_j\) over \(p\).

Consider the restriction map from \(Y(B_j)\) to the direct sum over \(w\mid p\) of the additive character spaces of \(\kappa_p(G_{B_{j,w}})\). Its kernel consists of characters unramified everywhere and trivial at \(w\mid p\). By the Hilbert class field theorem it is a subspace of the dual modulo \(p\) of the ideal class group. Here \(B_j\) is totally imaginary because it contains \(\mu_p\). The preceding stabilization makes the action of \(H\) on this kernel trivial. The action on the direct sum is trivial as well. Each place is fixed, and on its cyclotomic image conjugation by a lift of an element of \(H\) is trivial since \(\kappa_p\) is abelian. These statements concern actual Galois actions under the class field correspondence.

Since \(|H|\) is a unit in \(\mathbf F_p\), averaging over \(H\) is a projection. Its complementary projection kills the direct sum and the kernel just described. On \(Y(B_j)\) its image is therefore contained in the kernel, where another application kills it. Idempotence then gives zero. By Lemma~\ref{lem:local27}, \(H\) also acts trivially on \(\Sel_{p,r}(B_j)[p]\) for all \(r\). Apply the same complementary projection to the full strict group, now over \(\ZZ_p\). A nonzero image would contain an element of order \(p\), since the group is \(p\)-primary. The preceding conclusion excludes this. Lemma~\ref{lem:desc} identifies restriction with the invariant subgroup, proving the result.

For general \(B\) put \(F_j=B_j(\mu_p)\). The first case gives a layer \(s\) above which \(\Gal(F_j/F_s)\) acts trivially on every strict group. Choose \(a\ge s\) with \(F_s\cap\bigcup_j B_j\subset B_a\). Such an \(a\) exists since the intersection is a finite extension. For \(j\ge a\), every element of \(\Gal(B_j/B_a)\) fixes \(B_j\cap F_s\), so it extends to an automorphism of \(F_j=B_jF_s\) fixing \(F_s\). This follows from the restriction isomorphism for a compositum of finite Galois fields. Restriction from \(B_j\) to \(F_j\) is injective by Lemma~\ref{lem:desc}, because \([F_j:B_j]\mid p-1\). Equivariance of this restriction now gives triviality of \(\Gal(B_j/B_a)\) on the strict group over \(B_j\). Apply the lemma once more.
\end{proof}

\begin{corollary}[The tower before the auxiliary quadratic extensions]\label{cor:tower27}
For a prescribed finite set \(\Pi\) of odd primes, write \(B(x)=K_\psi C_l\) and \(H_x^0=\Gal(E_x/B(x))\). There are arbitrarily distant common intervals of Theorem~\ref{thm:fixed} on which
\begin{equation}
 \sum_{p\in\Pi}(\log p)\sum_{1\le r<t(x)}
        \len_{\ZZ_p}\Sel_{p,r}(E_x)^{H_x^0}
       \le \exp(x^{2\eps/3}).                                 \label{eq:tower27}
\end{equation}
One may in particular include every odd prime dividing \(\operatorname{ord}(\psi)\). The bound holds throughout each interval with the same \(q,d\) as the dyadic and critical-sign estimates.
\end{corollary}
\begin{proof}
Apply Theorem~\ref{thm:stable27} to \(B=K_\psi\) for each \(p\in\Pi\) and take the largest of the finitely many stabilization indices. In Theorem~\ref{thm:fixed} use this fixed layer, increased if necessary to meet any previously prescribed lower bound on \(a\). For large \(x\) the index \(l\) exceeds it. At every twist the strict group over \(B(x)\) is precisely the restriction of the one at the fixed layer. Lemma~\ref{lem:desc} identifies it with the invariants on the left of \eqref{eq:tower27}, since \([E_x:B(x)]=4\). Substitution in \eqref{eq:fixed} proves the claim. Every map commutes with complex conjugation, so both signs are included.
\end{proof}

Only the three summands on which at least one of the two auxiliary quadratic involutions acts by \(-1\) can contribute at the primes in \(\Pi\) outside \eqref{eq:tower27}. This follows from the four averaging idempotents of \(H_x^0\simeq(\ZZ/2)^2\) at an odd prime. For a fixed pair \(q,d\), Theorem~\ref{thm:stable27} also applies with \(B=K_\psi(\sqrt{-q},\sqrt{-d})\). Its stabilization index can depend on that pair. As the pair varies with the interval, this last application supplies no uniform upper bound for the three remaining summands.

\subsection{The conductor and the range of full strict orders}\label{nav:II-19}
\label{sec:smaller197}

Let \(0<\varepsilon<1/52\), fix the primitive character \(\psi\)
of conductor \(F\) and order \(o\), and let
\(C_j=\QQ(\zeta_{2^{j+2}})^+\). On a selected interval let
\(q,d_0\) be the prescribed auxiliary primes. For each \(x\) put
\[
 L_j=K_\psi C_j(\sqrt{-q},\sqrt{-d_0}),\qquad
 B_j=K_\psi C_j,\qquad
 E=L_J,\qquad t=\lfloor x^{401}\rfloor,\qquad b=\exp(x^\varepsilon).
\]
Here \(J\) is the least index for which
\([K_\psi C_J:\QQ]\ge x^{2+8\varepsilon}\).
The integral theta construction in Volume~I uses this field \(E\).
The absolute logarithmic size of the two auxiliary primes can be
reduced by placing the final interval farther into the range
where the same two primes give the character estimates.

\begin{lemma}[Choosing an interval with smaller auxiliary conductor]
\label{lem:delay197}
For every fixed \(A\ge1\) there are arbitrarily large \(Y\), a
single pair \(q,d_0\) for each \(Y\), and intervals
\([V,V^{3/2}]\) with \(V\ge Y^A\) on which the prime count,
the uniform dyadic estimate, and the growing odd-prime estimate
of Theorem~\ref{thm:log147} all hold. The period field and its
finite-level and real estimates for the terminal periods may be taken
on these intervals. Uniformly for \(x\) in them one has
\begin{equation}
 \log(2^{J+2}F qd_0)
     \ll_{\psi,\varepsilon,A}
        \log(2x)+x^{\varepsilon/(5A)}(\log(2x))^2.
 \label{eq:delay197}
\end{equation}
\end{lemma}
\begin{proof}
Set \(\eta=\varepsilon/80\), \(b_1=401/\eta\), and let
\(R_2=R_2(\psi)\) bound the number of exceptional dyadic twists
in Theorem \ref{thm:exception15}. Choose real constants
\[
       K_2=2A(2b_1)^{R_2},\qquad W>K_2+2 .
\]
These are fixed before any field or prime is chosen. Begin the
proof of Theorem~\ref{thm:log147} at \(Z\), using this \(W\).
For clarity, \(P_0=\lceil\log\log Z\rceil\),
\(R=\lceil C_\psi P_0^5\rceil\) and \(K=W(Wb_1)^R\). By
Lemma~\ref{lem:budget147} the exceptional twists of the fields
\(B_{a_p}\) for odd primes \(p\le P_0\) occupy at most \(R\)
indices up to \(\lceil Z^{401K}\rceil+1\). These fields depend
on \(\psi,p\) alone. The interval argument of Theorem~\ref{thm:log147}
therefore supplies \(Y\ge Z\) with \([Y,Y^W]\subset[Z,Z^K]\).
Throughout that interval any of the listed exceptions below
\(x^{401}\) is at most \(x^\eta\).

Now choose the two auxiliary primes by Theorem \ref{thm:twists}
at the initial scale \(Y\), with fixed upper power \(K_2+1\).
The twisted character estimates then hold on \([Y,Y^{K_2+1}]\)
and \(\log(qd_0)\ll Y^{\varepsilon/5}(\log(2Y))^2\).
Use the dyadic determinant of Theorem \ref{thm:exception15} for all
twists up to \(\lfloor Y^{401K_2}\rfloor\). Its count of
exceptional indices is at most \(R_2\), independently of
\(Y,q,d_0\). Start with \(U=K_2\log Y\). If an exceptional
index \(r\) has \((\log r)/401\) in \((U/(2b_1),U]\), replace
\(U\) by \(U/(2b_1)\). Such an index is then above every
later endpoint. After at most \(R_2\) replacements the
half-open interval contains no exceptional index. Set
\(V=\exp(U/2)\). Then
\[
 V\ge\exp\!\left(\frac{K_2\log Y}{2(2b_1)^{R_2}}\right)=Y^A,
 \qquad V^2\le Y^{K_2}.
\]
For \(x\in[V,V^2]\), each dyadic exception with \(r\le x^{401}\)
satisfies \((\log r)/401\le U/(2b_1)\le(\log x)/b_1\), so that
\(r\le x^\eta\). All nonexceptional dyadic lengths and the fixed
field for the exceptional ones are exactly those in Theorem \ref{thm:common16}. Its integral rounding keeps the required
arguments of the prime sums in \([Y,Y^{K_2+1}]\). Its Theorem
11.2 applies since \(\eta<\varepsilon/64\). The earlier
odd-prime argument also applies since \([V,V^2]\subset[Y,Y^W]\).
In particular its conversion of \(P_0\) into a constant times
\((\log\log x)^{1/5}\) still follows from
\(\log\log x\le\log\log Z+\log K\). Constants may now depend
on the fixed \(A\).

Take \([V,V^{3/2}]\), as in the prime-count and theta
constructions. It is a subinterval of the range just verified,
with the same primes and the same field \(E\) at each \(x\).
The shortened cyclotomic index satisfies \(J=O_\varepsilon(\log(2x))\)
by minimality and \([C_j:\QQ]=2^j\). Since \(Y\le x^{1/A}\)
the conductor estimate follows.
\end{proof}

For an odd prime \(p\) write \(S_{p,r}=\Sel_{p,r}(E)\). If
\(L\subset E\) with \([E:L]\) a power of two, write
\(e_L=[E:L]^{-1}\sum_{g\in\Gal(E/L)}g\). Its image on
\(S_{p,r}\) is exactly \(\Sel_{p,r}(L)\) by restriction.
Indeed averaging gives the isomorphism on global degree-one
cohomology. Corestriction makes restriction injective locally
at \(p\), and at other finite primes a homomorphism on inertia
vanishing upstairs factors through a group of two-power order.
This proves descent with all strict local conditions and at
every primary exponent.

\begin{theorem}[A bound with a separate field degree for each twist]
\label{thm:region197}
Fix \(0<\delta<1\). For each of the following alternatives, use
the indicated rule for \(j(r)\) and set \(e(r)=e_{L_{j(r)}}\).
Choose the largest \(0\le j\le J\) satisfying the rule. When
there is no such \(j\), set \(e(r)=0\).
\begin{align}
 &[2^j(r+2)]^5\le x^{(1-\delta)\varepsilon},
                  &&\text{using the published lattice bound},\notag\\
 &2^j(r+2)^2\le x^{(1-\delta)\varepsilon},
                  &&\text{using the cell count of de la Torre Dur\'an}.
                                                        \label{eq:region197}
\end{align}
There are arbitrarily distant common intervals on which each
alternative satisfies
\begin{equation}
   \sum_{1\le r<t}\ \sum_{p\ {\rm odd\ prime}}
            \log|e(r) S_{p,r}|
     \le \exp\!\big(C_{\psi,\varepsilon,\delta}
                          x^{(1-\delta/2)\varepsilon}\big).
                                                         \label{eq:block197}
\end{equation}
In particular the sum is \(o(b^a)\) for every fixed \(a>0\).
Both signs of complex conjugation and the full orders at all
odd primes are included.
\end{theorem}
\begin{proof}
Choose \(A\) so that \(1/(5A)<\delta/2\), then select the
intervals of Lemma~\ref{lem:delay197}. Each \(L_j\) is totally
imaginary, contains \(\sqrt{-q}\), and has degree
\(d_j\le4o\,2^j\). Its cyclotomic modulus divides
\(M=2^{J+2}Fqd_0\). Since \(r<t\) and \(d_j\le[E:\QQ]\) are
polynomially bounded in \(x\), equation~\eqref{eq:delay197} gives
\[
          \log(2M d_j(r+2))\ll_{\psi,\varepsilon,\delta}
                           x^{\delta\varepsilon/2}.
\]
The strict inequality in the choice of \(A\) absorbs the
logarithmic factors. For one admitted twist apply
Theorem~\ref{thm:cells103} to \(L_j\) with this modulus. Its
first bound uses \([d_j(r+2)]^5\) and its second uses
\(d_j(r+2)^2\) as the multiplier of this logarithm in the
exponent. The respective hypotheses of \eqref{eq:region197}
bound either multiplier by a constant times
\(x^{(1-\delta)\varepsilon}\). Thus the sum at this twist,
over all primes, is at most the exponential in
\eqref{eq:block197} with a fixed constant. Descent to the actual
invariant group and summation over the at most \(x^{401}\) twists
prove the claim after increasing the constant. The two cell
estimates are those established in Section~\ref{sec:cells103}
from the published bound of Bayer-Fluckiger, Emery and Houriet
and from the preprint of de la Torre Dur\'an, respectively.
~\cite{IV-II-III-beh,IV-II-III-torre103}
\end{proof}

\begin{corollary}[Vanishing above the large coefficient cutoff]
\label{cor:cut197}
Let \(P=\exp(x^\varepsilon/12)\). There is a positive constant
\(c\) depending on \(\psi,\varepsilon\) such that, on every
sufficiently distant admissible interval, for all odd primes
\(p>P\), all \(0\le j\le J\) and all \(1\le r<t\),
\[
 2^j(r+2)\le c\left(\frac{\log x}{\log\log x}\right)^{1/5}
        \quad\Longrightarrow\quad e_{B_j}S_{p,r}=0 .
\]
Using the second cell bound, the same conclusion follows from
\(2^j(r+2)^2\le c(\log x)/(\log\log x)\).
\end{corollary}
\begin{proof}
Restriction first identifies \(e_{B_j}S_{p,r}\) with
\(\Sel_{p,r}(B_j)\) since \([E:B_j]\) is a power of two.
Restriction from there to \(D_j=B_j(i)\) is injective by the
same prime-to-\(p\) argument. The latter is a totally imaginary
field of degree at most \(2o\,2^j\) and has a modulus at most
\(2^{j+2}F\). These bounds involve neither auxiliary quadratic
prime. Under either displayed size condition with \(c\le1\),
\(\log(2\operatorname{cond}(D_j)[D_j:\QQ](r+2))
\le C_\psi\log\log x\) eventually. In the first case the
exponent in \eqref{eq:pub103} is at most \(C_\psi c^5\log x\).
In the second case the exponent in \eqref{eq:vor103} is at
most \(C_\psi c\log x\). Take \(c\) small enough that both
are at most \((\varepsilon/2)\log x\). The sum of logarithmic
orders over all primes for \(D_j\) is then at most
\(x^{\varepsilon/2}\). One nonzero \(p\)-primary summand would
already contribute at least \(\log p>x^\varepsilon/12\), which
exceeds the entire bound for sufficiently large \(x\).
This proves the asserted vanishing on \(E\).
\end{proof}

\subsection{Polynomial conductor on the common interval}\label{nav:II-20}
\begin{theorem}[The same period fields with a polynomial conductor]
\label{thm:conductor232}
Fix \(\psi,\varepsilon\) as in Lemma~\ref{lem:delay197}. There
are arbitrarily distant intervals \([V,V^{3/2}]\), each with
one pair of auxiliary primes \(q,d_0\), on which one may use
\[
  E_x=K_\psi C_J(\sqrt{-q},\sqrt{-d_0}),\qquad
  N_x=2^{J+2}Fqd_0,\qquad
  x^{2+8\varepsilon}\le[K_\psi C_J:\QQ]<2x^{2+8\varepsilon}
\]
at every \(x\), with the same minimal index \(J\) as above.
There is a fixed \(D=D(\psi,\varepsilon)\) for which
\(N_x\le x^D\) everywhere on the interval. The prime-sum bounds
for the auxiliary characters, the dyadic strict-order bound and
the growing odd-prime estimate of Theorem~\ref{thm:log147} on the
fields \(B_J\) hold there. The integral weights and terminal-period
data can be chosen on these same \(E_x\).

There is also a choice retaining both original prime-size
conditions at the initial scale of Theorem~\ref{thm:twists} in Volume~I.
For this choice the conclusions hold with
\(\log N_x\ll_{\psi,\varepsilon}
 \log(2x)(1+\log\log(3x))^2\).
\end{theorem}
\begin{proof}
Put \(\eta=\varepsilon/80\), \(b_1=401/\eta\), and let
\(R_2=R_2(\psi)\) be the bound from Theorem
\ref{thm:exception15}. Fix
\(K_0=2(2b_1)^{R_2}\) and \(W>K_0+2\).
Carry out the first, odd-prime selection in
Lemma~\ref{lem:delay197} with this \(W\). At arbitrarily large
\(H\) it gives a range \([H,H^W]\) on which each of the finitely
many listed odd-prime exceptions below \(x^{401}\) is at most
\(x^\eta\). The count in Lemma~\ref{lem:budget147} is uniform
in the cutoff for the twists. The fields on which these lists are
made are the \(K_\psi C_{a_p}\), independent of both auxiliary
primes. This explains the order of selection.

Let \(c,C\) be the constants of Theorem~\ref{thm:long232}.
Set \(z=2W\log H/c\) and \(Y=z^{5/\varepsilon}\).
Then \(Y<H\) for large \(H\) and \(e^{cz}=H^{2W}\).
Use that theorem at \(Y\), first with \(U_2=e^{Cz}\).
It supplies a single pair with
\[
       \log(qd_0)\le(1+C)z\ll_{\psi,\varepsilon}\log H.
\]
Both auxiliary characters and their products with every permitted
base character have the required estimates throughout
\([H,H^W+1]\). The congruences at two distinguish the two
quadratic layers as in Theorem~\ref{thm:exception15}.
In particular the number of dyadic exceptional indices is at most
\(R_2\) through \(\lfloor H^{401K_0}\rfloor\) for this fixed
pair. The ramified primes are distinct, avoid \(2F\), and exceed
\(H^{401W}\), since \(\log q>W\log H/c\) and \(1/c>401\).
Every fixed odd prime needed to split in
\(\QQ(\sqrt{-d_0})\) is eventually covered by \(d_0\equiv-1\bmod p\).

For clarity, start with \(U=K_0\log H\). If
\((U/(2b_1),U]\) contains \((\log r)/401\) for a dyadic
exception, divide \(U\) by \(2b_1\) and repeat. Each division
puts a different exceptional index above the new upper endpoint.
After at most \(R_2\) divisions the range is empty of such indices.
Set \(V=e^{U/2}\). Then \(H\le V\), \(V^2\le H^{K_0}\), and at
every \(V\le x\le V^2\) a dyadic exception with
\(r\le x^{401}\) has \(r\le x^\eta\).
Theorem~\ref{thm:log147} therefore gives its growing odd-prime
estimate by exactly its bounds for the large and small twists,
and the dyadic descent and bounded-degree estimate of Theorems
\ref{thm:common16} and
\ref{thm:two17} bound the full dyadic sum.
Their conductor input is an upper bound for \(\log(2Fqd_0)\),
which here is \(O_{\psi,\varepsilon}(\log H)\).

Restrict to \([V,V^{3/2}]\). The rounded weight arguments
\(\lceil bx\rceil/b\), \(b=\exp(x^\varepsilon)\), belong to
\([H,H^W+1]\). The field construction and character average in
Proposition~\ref{prop:field12}, use the distinct
quadratic ramification and these auxiliary character estimates.
Restriction to the minimal layer \(J\) uses the descent through a two-power extension in Proposition~\ref{prop:short107}. Since \([C_j:\QQ]=2^j\)
and successive degrees \([K_\psi C_j:\QQ]\) grow by at most two,
minimality gives the displayed degree range and
\(2^{J+2}\ll_\psi x^{2+8\varepsilon}\). This proves the
polynomial bound for the actual modulus.
The ray-class estimate of Lemma~\ref{lem:ray88} is uniform in the
quadratic discriminant and modulus. For the split level primes
and the common terminal prime its uses in Theorem~\ref{thm:common111}
require the bound \(\log(2N_x)=o(\log b)\), which follows here.
The weights, ranks and tests in that construction are therefore
available without a further change of field or interval.

Finally repeat the same selection using
\(U_2=\exp(z(\log Y)^2)\). The initial prime-size and congruence
conditions are now exactly those of Theorem~\ref{thm:twists} in Volume~I.
Here \(\log Y=(5/\varepsilon)\log z\) and \(z\asymp\log H\).
This gives the second conductor estimate, and it still supplies
all the conductor bounds just used. Both constructions concern
the primes and intervals obtained by the stated selections.
\end{proof}

\subsection{All odd primes on a growing part of the same tower}\label{nav:II-29}
Fix the primitive character \(\psi\) of order \(o\), its conductor
\(F\), and the parameter \(\eps\) of Section~\ref{sec:missing}.
Throughout a selected interval write \(q,d_0\) for the actual
two auxiliary primes. We use \(d_0\) here to distinguish this
prime from a field degree. Put \(C_j=\QQ(\zeta_{2^{j+2}})^+\),
and recall that \(l\) is least with \(2^l\ge x^{2+8\eps}\).
For fixed \(0<u,v<1\) set
\[
\begin{split}
 E_x&=K_\psi C_l(\sqrt{-q},\sqrt{-d_0}),\qquad
 s=\lfloor u\eps\log_2 x\rfloor,\\
 L_{u,x}&=K_\psi C_s(\sqrt{-q},\sqrt{-d_0}),\qquad
 H_{u,x}=\Gal(E_x/L_{u,x}).
\end{split}
\]
The sum in the next theorem is formed from the very strict groups
on the left side of \eqref{eq:length12}. We write logarithmic
orders, so each cyclic summand of order \(p^a\) contributes
\(a\log p\). Define
\[
 B_{u,v}(x)=
 \sum_{1\le r\le x^{v\eps}}\;\sum_{p\ {\rm odd\ prime}}
       \log\big|\mathcal S_{p,r}(E_x)^{H_{u,x}}\big|.
\]
As elsewhere on these intervals, constants denoted by
\(C_{\psi,\eps}\) may depend on additional parameters fixed
before the choice of the two auxiliary primes.

\begin{theorem}[Uniform bounds through the full primary exponents]
\label{thm:tower104}
On the common intervals of Section~\ref{sec:missing},
\[
 B_{u,v}(x)=o\big(\exp(x^\eps/2)\big)
 \quad\text{if}\quad 5(u+v)+\tfrac15<1 .
\]
Using the cell-orbit bound of de la Torre Dur\'an in
Theorem \ref{thm:cells103}
gives the same conclusion whenever \(u+2v+\tfrac15<1\).
In particular the respective estimates at equal exponents are
\begin{align}
 B_{1/16,1/16}(x)&\le\exp(C_{\psi,\eps}x^{5\eps/6}),
                                                     \label{eq:pub104}\\
 B_{1/4,1/4}(x)&\le\exp(C_{\psi,\eps}x^{23\eps/24}).
                                                     \label{eq:vor104}
\end{align}
Both eigenspaces of complex conjugation are included at each
prime. Each assertion holds uniformly as the left endpoints of
the intervals tend to infinity.
\end{theorem}
\begin{proof}
The index \([E_x:L_{u,x}]\) is a power of two, since adjoining
\(C_l\) completes the extension. For odd \(p\), restriction
identifies \(\mathcal S_{p,r}(L_{u,x})\) with
\(\mathcal S_{p,r}(E_x)^{H_{u,x}}\). For the global degree-one
group this is inflation--restriction, since higher cohomology
of the finite group \(H_{u,x}\) on a \(p\)-primary group vanishes.
At a place above \(p\), corestriction composed with restriction
is multiplication by a local degree prime to \(p\), so local
triviality descends. Away from \(p\) the twist is unramified.
An inertia homomorphism that becomes zero upstairs factors
through a group of order prime to \(p\), so it vanishes.
This also proves descent at places that ramify in the field
extension. The invariant group is a direct summand under the
averaging idempotent, which commutes with complex conjugation.

Write \(d_L=[L_{u,x}:\QQ]\). Since \([C_s:\QQ]=2^s\),
\[
  2^s\le d_L\le 4o\,2^s\le 4o\,x^{u\eps}.
\]
The field contains \(\sqrt{-q}\), is totally imaginary, and is
contained in \(\QQ(\zeta_N)\) for \(N=2^{l+2}Fqd_0\). The common
interval construction in
Sections \ref{nav:II-04} and \ref{nav:II-10}
has \(\log(qd_0)\ll_{\psi,\eps}Z^{\eps/5}(\log(2Z))^2\) at its
initial scale \(Z\le x\). Any fixed extra parameters in that
construction are held fixed here. Since \(l=O_\eps(\log(2x))\),
for \(1\le r\le x^{v\eps}\) this gives
\[
 \log(2Nd_L(r+2))\ll_{\psi,\eps}
           x^{\eps/5}(\log(2x))^2
      \ll_{\psi,\eps} x^{5\eps/24}.
\]
No auxiliary prime or interval has been reselected.

Use first the published lattice bound in Theorem \ref{thm:cells103}.
At one twist the logarithmic order summed over all primes is at
most \(\exp(C[d_L(r+2)]^5\log(2Nd_L(r+2)))\). Inserting the
preceding estimates and summing at most \(x^{v\eps}\) such terms gives
\[
 B_{u,v}(x)\le x^{v\eps}
      \exp\!\big(C x^{[5(u+v)+1/5]\eps}(\log(2x))^2\big).
\]
If \(5(u+v)+1/5<1\), the logarithm of this upper bound is
\(o(x^\eps)\), proving the first assertion. With \(u=v=1/16\),
using \(5\eps/24\) for the logarithmic factor yields
\(5(1/16+1/16)+5/24=5/6\). Absorbing \(x^{v\eps}\) gives
\eqref{eq:pub104}.

For the second assertion use the other bound in that same
theorem. Its exponent is \(C d_L(r+2)^2\log(2Nd_L(r+2))\).
This replaces \(5(u+v)\) by \(u+2v\) in the preceding argument.
At \(u=v=1/4\) the exponent with \(5/24\) is \(3/4+5/24=23/24\).
This proves \eqref{eq:vor104} with all odd primes included. In
each calculation all constants are independent of \(q,d_0,x\)
once the interval construction and the indicated parameters are fixed.
\end{proof}

\begin{corollary}[The contribution outside the bounded summands]
\label{cor:rest104}
Choose either range of Theorem~\ref{thm:tower104}, with its
stated hypotheses. Write \(t=\lfloor x^{401}\rfloor\).
On \(S=\mathcal S_{p,r}(E_x)\) for odd \(p\) put
\[
 e_r=\frac{1+(-1)^r c}{2},\qquad
 \pi_{p,r}=
 \begin{cases}
 |H_{u,x}|^{-1}\displaystyle\sum_{\sigma\in H_{u,x}}\sigma,
                       & r\le x^{v\eps},\\
 0,                    & r>x^{v\eps},
 \end{cases}
\]
where \(c\) is complex conjugation. Then
\begin{equation}
 \mathcal L(x)-
 \sum_{1\le r<t}\sum_{p\ {\rm odd\ prime}}
   \log\big|(1-\pi_{p,r})e_r\mathcal S_{p,r}(E_x)\big|
             =o\big(\exp(x^\eps/2)\big).                 \label{eq:rest104}
\end{equation}
One may also omit from the sum the character components certified
by Theorem~\ref{thm:stable102} at \(p\nmid N[E_x:\QQ]\).
\end{corollary}
\begin{proof}
For odd \(p\) the commuting idempotents give the direct sum
\[
 S=(1-e_r)S\ \oplus\ \pi_{p,r}e_r S\
             \oplus\ (1-\pi_{p,r})e_r S.
\]
The sum of the logarithmic orders of the first terms is bounded
by the polynomial critical-sign estimate in
Theorem \ref{thm:critical13}.
The middle terms together contribute at most \(B_{u,v}(x)\).
The dyadic groups in \(\mathcal L(x)\) contribute
\(o(\exp(x^\eps/2))\) on these same intervals by
Theorem~\ref{thm:res100} and its proof. This gives
\eqref{eq:rest104}. At a prime of the stated kind we may
extend scalars to the unramified character ring of
Theorem~\ref{thm:stable102}. Its character projections commute
with \(\pi_{p,r}\) and \(e_r\). Projecting a certified component
again only decreases its length. Thus the total cost of
omitting these components, with lengths normalized as in that
theorem, is polynomial in \(x\). This proves the final assertion.
\end{proof}

For \(r>x^{v\eps}\) the last sum still contains the whole
complementary eigenspace, before the possible character omissions.
For smaller \(r\) it still contains the part that does not descend
to \(L_{u,x}\). The full field contains \(C_l\) of degree
\(2^l\ge x^{2+8\eps}\), whereas \([C_s:\QQ]\le x^{u\eps}\).
A uniform bound for those remaining orders is open. The height
estimate for the arithmetically normalized theta-period map
in Corollary~\ref{cor:alg97} is also open on these fields.

\subsection{Bounds for deep torsion on the period field}\label{nav:II-30}
Let \(E=\widetilde E=K_\psi C_{l'}(\sqrt{-q},\sqrt{-d_0})\) be
the field carrying the periods in Proposition~\ref{prop:short107}.
Throughout this paragraph \(x\in[V,V^{3/2}]\), with the same
\(V,q,d_0\) as in that proposition. Write \(b=\exp(x^\eps)\),
\(P=b^{1/12}\), \(t=\lfloor x^{401}\rfloor\), \(n=[E:\QQ]\),
and take \(N\) from the common construction, so
\[
 E\subset\QQ(\zeta_N),\qquad n\ll_\psi x^{2+8\eps},\qquad
 \log(2N)\ll_{\psi,\eps}x^{\eps/4}.
\]
We use the original small-\(\eps\) range \(0<\eps<1/52\).
For an odd prime \(p\), write \(S_{p,r}=\mathcal S_{p,r}(E)\).
Its elements are global classes with coefficients
\((\QQ_p/\ZZ_p)(-r)\), unramified away from \(p\) and locally
zero at every prime above \(p\). Put
\[
     D_{p,r}=\frac{p-1}{\gcd(p-1,r)},\qquad
     M_{p,r}=\tfrac12(1+(-1)^rc)S_{p,r}.
\]

\begin{lemma}[Two bounds for the subgroup killed by \(p\)]
\label{lem:rank114}
For every odd prime \(p\) and every \(r\ge1\),
\begin{equation}
    \log |S_{p,r}[p]|
      \ll nD_{p,r}\log(2Np).                              \label{eq:first114}
\end{equation}
For \(p>P\), all sufficiently large \(x\), and the full range of
twists used by the periods, one also has
\begin{equation}
    \sum_{1\le r<t}\log |M_{p,r}[p]|
      \ll np\log(2Np).                                    \label{eq:many114}
\end{equation}
The constants in these two inequalities are absolute.
\end{lemma}
\begin{proof}
Let \(F/E\) be the field cut out by \(\bar\kappa_p^{-r}\).
Then \(F\subset E(\mu_p)\) and \([F:E]\le D_{p,r}\) since the
\(r\)-th powers in \(\mathbf F_p^\times\) form a group of that
order. The estimate at \(j=1\) in
Theorem \ref{thm:strict113}
gives
\[
    \log |S_{p,r}[p]|\le\log h(F)+[F:\QQ]\log p .
\]
Here \(h(F)\) is the ordinary class number. For any subfield
\(H\) of \(\QQ(\zeta_{Np})\) of degree \(m\), the
conductor--discriminant formula gives
\(|\operatorname{disc} H|\le(Np)^m\). Each ideal class has an
integral representative of norm at most
\(Y=2^m|\operatorname{disc}H|^{1/2}\) by Minkowski's theorem.
The number of integral ideals of norm at most \(Y\) is bounded
by \(Y^2\zeta_H(2)\le Y^2\zeta(2)^m\), by comparison of Euler
factors. Thus \(\log h(H)\ll m\log(2Np)\), which proves
\eqref{eq:first114}.

For \(p>P\) one has \(p>t+1\) and \(p\nmid nN\) for all
sufficiently large \(x\), by the displayed bounds on \(n,N\).
Let \(T=(E\QQ(\mu_p))^+\). The groups in \eqref{eq:many114}
use distinct eigenspaces in its ideal class group modulo \(p\).
We recall the argument. Work over an unramified coefficient ring
containing the values of the characters \(\alpha\) of
\(\operatorname{Gal}(E/\QQ)\). Descent by the character
idempotent, as in
Theorem \ref{thm:char} and Corollary \ref{cor:sum79}, gives the strict character over \(\QQ\)
with residual action \(\alpha^{-1}\bar\kappa_p^{-r}\).
The descent uses averaging over the prime-to-\(p\) group, on
global cohomology and on each local restriction. In the
complementary summand \(\alpha(c)=(-1)^r\), so this character
factors through \(\operatorname{Gal}(T/\QQ)\). It is nontrivial
on inertia at \(p\) since \(0<r<p-1\) and \(p\nmid N\). Finite
and divisible coefficients therefore give the same zero local
condition at \(p\) for a class killed by \(p\). Over \(T\) it
gives an unramified homomorphism, hence one from the class group.
Restriction is injective since \([T:\QQ]\) is prime to \(p\).

The cyclotomic factors of conductors \(N,p\) are disjoint.
On their product the second factor determines \(r\) in this
range and the first determines \(\alpha\). Thus the pairs
\((r,\alpha)\) use pairwise distinct idempotents. Sum their
dimensions over the residue field. Extension of the coefficient
field preserves dimension, so after multiplication by \(\log p\)
the sum is at most \(\log h(T)\). Applying the preceding
class-number estimate with \([T:\QQ]\le np\) proves
\eqref{eq:many114}.
\end{proof}

Define two depths, with \(\lfloor y\rfloor=0\) for \(0\le y<1\),
by
\[
 u_{p,r}=\left\lfloor\frac{x^{-420}b^{1/2}}{pD_{p,r}}\right\rfloor,
 \qquad
 j_p=\begin{cases}
   \left\lfloor x^{-4}b^{1/2}/p^2\right\rfloor,&p>P,\\
   0,&p\le P.
 \end{cases}
 \qquad k_{p,r}=\max(u_{p,r},j_p).
\]
We use \(M[p^0]=0\).

\begin{theorem}[Growing depths at all the twists of the period field]
\label{thm:deep114}
Uniformly on the selected intervals,
\begin{equation}
 \sum_{1\le r<t}\ \sum_{p\ {\rm odd\ prime}}
       \log|M_{p,r}[p^{k_{p,r}}]|
      \ll_{\psi,\eps} x^{-1}b^{1/2}.                    \label{eq:deep114}
\end{equation}
For all \(p\le P\), every \(k_{p,r}\) is at least
\(\lfloor x^{-420}b^{1/3}\rfloor\). For
\(P<p\le b^{1/5}\) it is at least
\(\lfloor x^{-4}b^{1/10}\rfloor\). The whole first layer at
every \(P<p\le x^{-2}b^{1/4}\) is included.
\end{theorem}
\begin{proof}
If \(A\) is a finite abelian \(p\)-group with elementary
divisors \(p^{d_i}\), then for every integer \(j\ge0\)
\[
     \log|A[p^j]|=(\log p)\sum_i\min(j,d_i)
                        \le j\log|A[p]|.
\]
This permits estimates at growing \(j\) on the field of the
first layer. For a maximum of two depths we can bound the
logarithmic order by the sum of their two bounds.

Set \(U=x^{-420}b^{1/2}\). The inequality \(u_{p,r}\ge1\)
implies \(p(p-1)\le U\gcd(p-1,r)\le Ut\). Thus
\(p\le Z=1+\sqrt{Ut}\) whenever this depth contributes.
Using \(M_{p,r}[p]\subset S_{p,r}[p]\) and
\eqref{eq:first114}, the sum for the depths \(u_{p,r}\) is
at most
\[
 C t n U\log(2NZ)\sum_{p\le Z}\frac1p
       \ \le\ C t n U\log(2NZ)(1+\log Z).
\]
Here \(\log(2NZ)\) and \(1+\log Z\) are both \(O_{\psi,\eps}(x)\),
and \(n\ll_\psi x^3\). This bound is
\(O_{\psi,\eps}(x^{-14}b^{1/2})\).

Put \(U'=x^{-4}b^{1/2}\). For the depths \(j_p\) sum first
over \(r\) using \eqref{eq:many114}. We obtain
\[
 C nU'\sum_{P<p\le\sqrt{U'}}\frac{\log(2Np)}p
 \ll_{\psi,\eps} nU' x^{2\eps}
 \ll_{\psi,\eps} x^{-2+10\eps}b^{1/2}.
\]
Indeed \(\log(2Np)\ll_{\psi,\eps} x^\eps\) in this range,
and the harmonic sum over all integers up to \(\sqrt{U'}\)
is \(O(x^\eps)\). Since \(10\eps<1\), this proves
\eqref{eq:deep114}. For \(p\le P\) use \(D_{p,r}\le p\)
and \(P^2=b^{1/6}\). The two assertions above \(P\) follow
from \(p^2\le b^{2/5}\) and \(p^2\le x^{-4}b^{1/2}\),
respectively.
\end{proof}

\begin{corollary}[The order still required on the same interval]
\label{cor:tail114}
Use the averaging projector \(\pi_{p,r}\) of
Proposition~\ref{prop:short107} on \(E\), in either range allowed
there and with its corresponding hypotheses. Put
\(Q_{p,r}=(1-\pi_{p,r})M_{p,r}\), and let \(\mathcal L_E(x)\)
denote the left side of \eqref{eq:length12} on \(E\). Then
\begin{equation}
 \mathcal L_E(x)=
    \sum_{1\le r<t}\ \sum_{p\ {\rm odd\ prime}}
          \log|p^{k_{p,r}}Q_{p,r}|+o(b^{1/2}).          \label{eq:tail114}
\end{equation}
Each pair for which the finite orders of \(S_{p,r}[p^{k_{p,r}}]\)
and \(S_{p,r}[p^{k_{p,r}+1}]\) agree contributes zero to this
sum. Both orders are given by
Corollary \ref{cor:levels113}.
\end{corollary}
\begin{proof}
Proposition~\ref{prop:short107} gives \eqref{eq:rest104} on
\(E\) with these projectors. For each \(Q=Q_{p,r}\) and
\(k=k_{p,r}\), multiplication by \(p^k\) gives
\(\log|Q|-\log|p^k Q|=\log|Q[p^k]|\). Its right side is
nonnegative and at most \(\log|M_{p,r}[p^k]|\). Substitution
of Theorem~\ref{thm:deep114} into \eqref{eq:rest104} proves
\eqref{eq:tail114}, with the stated uniform error. Equal finite
orders at the indicated consecutive depths force
\(p^{k_{p,r}}S_{p,r}=0\) by the cyclic decomposition, including
when \(k_{p,r}=0\). This also kills \(Q_{p,r}\).
\end{proof}

For comparison, Proposition~\ref{prop:res23} removes depth
\(\lfloor b^{1/8}\rfloor\) at \(p\le P\), and
Proposition~\ref{prop:res80} removes at most four layers above
\(P\). Theorem~\ref{thm:deep114} controls growing depths on
both sides of \(P\), on the field used by the period itself.
For any fixed \(\delta>0\), a bound by
\((1-\delta)b^{1/2}\) for the sum in \eqref{eq:tail114} on
arbitrarily distant common intervals would give
\eqref{eq:length12}. The sizes of the terms surviving
these depths and a nonzero integral comparison on the projected
theta image are still required.

\subsection{Strict cohomology on the fields of the norm theorem}\label{nav:II-31}
Let \(E=K_\psi C_J(\sqrt{-q},\sqrt{-d_0})\) be the field of
Theorem~\ref{thm:low136}, with \(J=l'\), \(t=\lfloor x^{401}\rfloor\)
and \(A=\operatorname{Gal}(E/C_J)\). We give a full-depth
estimate for a collection of summands in the strict cohomology
sum \eqref{eq:length12}. At \(p\equiv3\bmod8\), \(p>3\), put
\(\lambda(v)=(v/p)\) and take the twists
\[
             r+1\equiv(p-1)/2\pmod{p-1}.
\]
They are even, so the \(A\)-invariant classes on this real
cyclotomic subfield have the complementary sign in
\eqref{eq:tail114}.

Fix a nontrivial character \(\alpha\) of \(C_J\) with primitive
conductor \(f_0\), and fix an embedding of its values in an
unramified coefficient ring \(O\) over \(\ZZ_p\). Choose a
positive \(a_0\equiv1\bmod p\) prime to \(f_0\) with
\(\alpha(a_0)\ne1\). This choice may be kept for the entire
Frobenius orbit of \(\alpha\). For any \(p\)-adic unit \(v\)
write \(\omega(v)\) for its Teichm\"uller lift and
\(\langle v\rangle=v/\omega(v)\). For integers \(s\ge0\) put
\begin{equation}
 A_m(s)=\sum_{\substack{1\le v\le f_0p^m\\(v,f_0p)=1}}
 \left\lfloor\frac{a_0 v}{f_0p^m}\right\rfloor
      \bar\alpha(v)\lambda(v)\langle v\rangle^{-s}
                     \pmod{p^mO}.                         \label{eq:As140}
\end{equation}
When \(s=r+1\) in the indicated congruence this is exactly the
annihilator for the \(p^m\)-torsion in the strict component,
by the floor-sum test in
Sections \ref{nav:III-04} and \ref{nav:III-19}.

\begin{lemma}[Finite differences of the annihilator]
\label{lem:difference140}
The residues in \eqref{eq:As140} define \(A(s)\in O\).
Put \(\bar q(v)=(\langle v\rangle-1)/p\bmod p\) and let \(M\)
be the sum at modulus \(f_0p^2\) with weight
\(\bar\alpha(v)\lambda(v)\bar q(v)\) in the residue field.
For distinct nonnegative integers \(s,s'\) one has
\[
 \frac{A(s')-A(s)}{p(s'-s)}\in O,\qquad
 \frac{A(s')-A(s)}{p(s'-s)}\equiv-M\pmod{pO}.
\]
At \(s=0\) the value is
\((a_0-\alpha(a_0)\lambda(a_0))B_\alpha\), where \(B_\alpha\)
is the algebraic integer of Lemma \ref{lem:coord139}.
\end{lemma}
\begin{proof}
Let \(n=f_0p^i\) and pass from \(n\) to \(kn\) with \(k\) a
power of \(p\). Above a unit \(v\bmod n\) the representatives
are \(v+bn,\ 0\le b<k\). The sum of their floor coefficients is
\[
        \left\lfloor a_0v/n\right\rfloor
                       +(a_0-1)(k-1)/2 .
\]
Indeed the residues of \(a_0 b\bmod k\) run through a complete
system, and summing the integer parts followed by the fractional
parts gives the displayed identity. The extra constant has
zero sum against the reduced weights, which are odd under
\(v\mapsto-v\) since \(\lambda\) is odd and
\(\langle-v\rangle=\langle v\rangle\). This proves
compatibility modulo \(p^i\). At \(s=0\) the weight comes from
modulus \(f_0p\) itself. The same calculation gives exact
equality of the weighted floor sums at every modulus. The
permutation by \(a_0\) at the first modulus gives the stated
value, as in equation \ref{eq:floor115} with the identity for floors
proved there.

Let \(d=v_p(s'-s)\). At modulus \(f_0p^{d+2}\), for each
unit \(v\) the binomial expansion or the \(p\)-adic logarithm
gives
\[
 \langle v\rangle^{-s'}-\langle v\rangle^{-s}
       \equiv -p(s'-s)\bar q(v)\pmod{p^{d+2}} .
\]
For instance \(\log\langle v\rangle/p\equiv\bar q(v)\bmod p\),
and exponentiating \(-(s'-s)\log\langle v\rangle\) preserves
its linear term at this precision since \(p\) is odd.
Also \(\langle v\rangle^{-s}\equiv1\bmod p\). After summing
and division by \(p(s'-s)\), the remaining weights factor
modulo \(f_0p^2\) and are still odd under negation. The floor
identity therefore reduces their sum to \(M\). This proves both
claims about the finite difference.
\end{proof}

\begin{theorem}[At most one exceptional twist per conductor]
\label{thm:except140}
For every sufficiently large \(x\) on the common intervals and
every prime
\[
              x^4<p\le2t+1,\qquad p\equiv3\bmod8 ,
\]
all but at most \(J-2\) of the twists in the stated congruence
with \(1\le r<t\) satisfy
\begin{equation}
                p^{101}\mathcal S_{p,r}(E)^A=0 .
                \label{eq:except140}
\end{equation}
More precisely, at each conductor \(2^{j+2}\), \(3\le j\le J\),
and for every \(m\ge1\), at most one class of \(r+1\bmod p^m\)
can escape the exponent bound \(p^m\) for its character
components. This class, when present, is nonzero modulo \(p\).
The conditions can be tested by \eqref{eq:As140} with precision
\(p^{m+1}\).
\end{theorem}
\begin{proof}
The bound on the period degree in Proposition~\ref{prop:short107}
gives \(2^{J+2}\le x^3\) for all sufficiently large \(x\).
Throughout the prime range
\(2^{J+2}(2+\log(2^{J+2}p))^2<p\), since \(\log p\ll\log x\).
Also \(p>|A|\) since \(|A|\le4[K_\psi:\QQ]\). Prime-to-\(p\)
descent with the local conditions identifies
\(\mathcal S_{p,r}(E)^A\) with \(\mathcal S_{p,r}(C_J)\).

The first four labels vanish by Theorem \ref{thm:quartic115}.
At a conductor \(2^{j+2}\) with \(j\ge3\) there are just two
Frobenius orbits of primitive labels. By Lemma \ref{lem:coord139} and the proof of Corollary~\ref{cor:exp139}, the sum of the
two valuations of \(B_\alpha\) is strictly less than two.
They are nonnegative integers. Thus there can be only one
orbit where the valuation is nonzero, and there it is one.
The factor multiplying \(B_\alpha\) in Lemma~\ref{lem:difference140}
is a unit. The finite differences show \(A(s)\equiv A(0)\bmod pO\)
for every \(s\). A unit at zero therefore kills its strict
component for every relevant twist by the test at \(m=1\).

For the possible other orbit \(v_p(A(0))=1\). If \(M=0\),
the lemma gives \(A(s)\equiv A(0)\bmod p^2O\) for all \(s\),
and \(p\) kills every component in that orbit. If \(M\ne0\),
the same lemma gives exactly
\[
               v_p(A(s')-A(s))=1+v_p(s'-s).
\]
If two values of \(A\) are both zero modulo \(p^{m+1}\),
their arguments must thus coincide modulo \(p^m\).
Conversely that congruence of arguments preserves the residue
of \(A\) modulo \(p^{m+1}\). No such argument is zero modulo
\(p\), as that would give \(A(s)\equiv A(0)\bmod p^2O\).
The Frobenius-conjugate labels in the orbit give conjugate
annihilators at each integer, with the same valuation. If
\(A(s)\ne0\bmod p^{m+1}\) for \(s=r+1\), the test at precision
\(p^{m+1}\) kills the entire strict component by \(p^m\).
Indeed its valuation is at most \(m\), and an element of larger
order would yield one of order \(p^{m+1}\). This proves the
assertion about classes at every conductor.

Finally \(p^{101}>x^{404}>t\). At each of the \(J-2\)
conductors the possible exceptional residue at \(m=101\)
can contain at most one integer \(r+1\) in the required range.
All of the other twists satisfy the bound at every character.
The character idempotents are integral after unramified
coefficient extension since \(p\) is odd. Faithful scalar
extension gives \eqref{eq:except140} for the entire invariant
group.
\end{proof}

\begin{corollary}[The logarithmic order on the common interval]
\label{cor:sum140}
Sum the full logarithmic orders of \(\mathcal S_{p,r}(E)^A\)
over the primes and quadratic residual twists covered by
\eqref{eq:except140}, omitting the possible exceptional twists.
Uniformly on the same intervals used for the period construction
the sum is
\[
                \ll [E:\QQ]\,t^2\log(2Nt)
                 =o\bigl(\exp(x^\eps/2)\bigr).
\]
At every included pair this invariant group is killed by the
depth \(p^{k_{p,r}}\) of Theorem~\ref{thm:deep114}.
\end{corollary}
\begin{proof}
By cyclic decomposition, \eqref{eq:except140} bounds each
logarithmic order by \(101\log|\mathcal S_{p,r}(E)[p]|\).
Lemma~\ref{lem:rank114} bounds this by
\(C[E:\QQ]p\log(2Np)\). For each prime the number of twists is
at most \(1+t/(p-1)\). Summing over the integers from two to
\(2t+1\) bounds \(\sum_p p(1+t/(p-1))\) by \(O(t^2)\).
The degree, \(t\) and \(\log(2Nt)\) are all polynomially
bounded in \(x\), proving the little-oh estimate for fixed
\(\eps>0\). All those coefficient primes satisfy
\(p\le2t+1\le P=\exp(x^\eps/12)\) eventually. The lower bound
for \(k_{p,r}\) in Theorem~\ref{thm:deep114} consequently
exceeds 101.
\end{proof}

These groups lie on the field carrying the automorphic periods.
The estimate leaves the possible exceptional twists, the
components nontrivial on \(A\), and other ranges of primes
and twists in the full cohomological sum. The norm identity
of Theorem~\ref{thm:low136} on these fields still requires
control of the omitted real tail and an integral normalization
for the prescribed theta periods.

\subsection{Strict cohomology in cyclic extensions of prime power degree}\label{nav:IV-44}
Let \(B/\QQ\) be cyclic of degree \(p^a\), for an odd prime \(p\),
with conductor \(b\) prime to \(p\). Fix a primitive character
\(\chi\) of order prime to \(p\) and conductor \(f\) prime to
\(pb\). Put \(O=\ZZ_p[\chi]\), \(J=\operatorname{Frac}O\).
Take \(r\ge1\) with \(\chi(-1)=(-1)^r\), and put
\(A=(J/O)(\chi^{-1}\kappa_p^{-r})\).
Write \(S(B,A)\) for the subgroup of \(H^1(B,A)\) of classes
unramified away from \(p\) and zero at every place over \(p\).
For each \(0\le j\le a\), choose a character \(\lambda_j\) of
\(\Gal(B/\QQ)\) of exact order \(p^j\), with \(\lambda_0=1\).
Write \(d_j=d_{p,r,\chi\lambda_j}(fb)\) for the cyclotomic
index on the open set of Theorem~\ref{thm:openindex269}.

\begin{theorem}[Integral bounds through the cyclic character quotients]\leavevmode
\label{thm:cyclic278}
Suppose that \(\bar\chi^{-1}\omega^{-r}\) is nontrivial on
\(G_{\QQ_p}\), where \(\omega=\kappa_p\bmod p\). Suppose also
that \(\bar\chi(q)\ne1\) at every prime \(q\mid b\). Then
\[
       \max_{0\le j\le a}d_j
          \ \le\ \len_O S(B,A)\ \le\ \sum_{j=0}^a d_j .
\]
The group vanishes exactly when \(d_0=0\).
\end{theorem}
\begin{proof}
Put \(\Sigma=\{v:v\mid pfb\}\) for the finite primes in question.
Every \(q\mid b\) is \(1\bmod p\). Indeed its inertia in the
abelian extension is a nontrivial tame \(p\)-group, hence a
quotient of \(\mathbb F_q^\times\) by local reciprocity.
For \(v\mid f\) the residual finite character is nontrivial on
inertia. At \(q\mid b\) its Frobenius value differs from one and
\(\omega(q)=1\). Thus at every \(v\mid fb\) both the residual
line \(\bar A=\bar\chi^{-1}\omega^{-r}\) and its Tate dual have
zero invariants. The local Euler characteristic formula away
from \(p\), together with local duality in degree two, gives
zero degree-one cohomology as well. The same conclusions apply
after a local extension of \(p\)-power degree, since a nontrivial
character of order prime to \(p\) stays nontrivial there.
Filtering finite coefficients by residual lines and passing to
their union proves local vanishing for the divisible coefficients
below. Every simple constituent of the regular representation of
a cyclic \(p\)-group in characteristic \(p\) is trivial, by
\(X^{p^a}-1=(X-1)^{p^a}\). This justifies the filtration also
for the regular representation and its lattice quotients.
At \(p\) we use only degree-zero vanishing, which follows from
the first hypothesis by the same filtration.

For these coefficients \(M\) over \(\QQ\), write \(S(M)\) for
the kernel of localization at all \(v\in\Sigma\) in cohomology
unramified outside \(\Sigma\). At primes dividing \(fb\) the
local condition is automatic by the preceding paragraph. Set
\(O_j=O[\zeta_{p^j}]\), \(J_j=\operatorname{Frac}(O_j)\) and
\(A_j=(J_j/O_j)((\chi\lambda_j)^{-1}\kappa_p^{-r})\).
In the Poitou--Tate sequence the dual of \(S(A_j)\) identifies
with the kernel of localization on degree two for the dual
lattice. All its local degree-two terms vanish, by local
duality and the degree-zero vanishing just proved. The trace
dual of \(O_j\) is a free fractional \(O_j\)-ideal, so choosing
a generator identifies this lattice with
\(O_j(\chi\lambda_j\kappa_p^{r+1})\).
Theorem~\ref{thm:openindex269} now gives
\(\len_{O_j}S(A_j)=d_j\). Indeed any extra prime factor in its
Euler product divides \(b\), and the reduction of the factor is
\(1-\bar\chi(q)\). The extension \(O_j/O\) is totally ramified,
since \(O/\ZZ_p\) is unramified and the translated cyclotomic
polynomial is Eisenstein. The two lengths agree over \(O\) on
every finite \(O_j\)-module, by filtering it by the common
residue field.

Put \(R_i=O[X]/(X^{p^i}-1)\), with \(X\) the action of a
generator of the cyclic quotient. Shapiro's lemma identifies
\(S(A\otimes_O R_a)\) with \(S(B,A)\), including each local
restriction. Sending \(X\) to the inverse character value chosen
for \(\lambda_i\) gives quotient \(O_i\), for \(i\ge1\).
Multiplication by \(\Phi_{p^i}(X)\) identifies the kernel with
\(R_{i-1}\), by cancellation of this monic factor in \(O[X]\).
Tensoring with \(A\) preserves exactness because the quotient
is \(O\)-free. In a resulting short exact sequence the global
degree-one map from the first term is injective, since the last
term has zero global invariants. Its image on strict groups is
the kernel of the map to the third strict group. Indeed a global
lift from the first term is locally zero at every required prime
by injectivity on local degree one, which follows from local
degree-zero vanishing for the third term. Applying the resulting
order inequality successively proves the upper bound. For the
lower bound, \(O_j\) embeds in \(R_a\) by multiplication by
\((X^{p^a}-1)/\Phi_{p^j}(X)\), using \(\Phi_1=X-1\) at \(j=0\).
Its quotient is again \(O\)-free and has only the indicated
residual line after twisting. The same injection therefore gives
the lower bound for each \(j\).

Finally let \(b_j\) be the conductor of \(\lambda_j\).
It is squarefree, as the higher unit groups at a tame prime
\(q\mid b\) are pro-\(q\). Modulo the maximal ideal of \(O_j\),
\(\lambda_j\) becomes trivial. The corestriction defining the
reduced class may therefore be taken first from level \(fpb_j\)
to level \(fp\), with the tensor basis carried along. Removing
each new \(q\) gives
\[
  N_{\QQ(\zeta_{mq})/\QQ(\zeta_m)}(1-\zeta_{mq})
      =\frac{1-\zeta_m}{1-\zeta_m^{\,q^{-1}\bmod m}}
                \qquad(q\nmid m,\ \zeta_{mq}^{\,q}=\zeta_m).
\]
This follows by taking the product over the \(q\) roots of
\(X^q-\zeta_m\) and omitting the unique root in \(\mu_m\).
Moving inverse Frobenius on the second Kummer class across the
corestriction acts on the tensor basis by \(\chi(q)q^r\).
Thus the reduced class is the reduction of the \(j=0\) class
times \(\prod_{q\mid b_j}(1-\bar\chi(q)q^r)\), a unit.
Both classes are compared on the same open set with the same
residue field. The multiplication exact sequence identifies
the kernel of reduction on the free degree-one lattice with its
maximal-ideal multiple. Hence \(d_j=0\) exactly when \(d_0=0\).
The two bounds prove the last assertion.
\end{proof}

\subsection{The exact length over a cyclic group ring}\label{nav:IV-45}
Let \(p\) be odd, \(r\ge1\), and \(\chi\) a finite character
of order prime to \(p\) with \(\chi(-1)=(-1)^r\).
Write \(O=\ZZ_p[\chi]\), \(k=O/pO\) and
\(T=O(\chi\kappa_p^{r+1})\). Let \(P\) be the Galois group of
a cyclic degree-\(p^a\) extension of \(\QQ\). Choose \(N\)
divisible by its conductor and that of \(\chi\), and put
\(U=\operatorname{Spec}\ZZ[1/(Np)]\). Fix a generator of \(P\).
For \(i=0,\ldots,a\) put \(R_i=O[X]/(X^{p^i}-1)\), with the
chosen generator acting by multiplication by \(X\). Put
\(D_i=O[\zeta_{p^i}]\), with the generator acting by
\(\zeta_{p^i}\), and \(D_0=O\). In every tensor with \(T\) use
the diagonal Galois action. Denote by \(M_i\) and \(H_i\)
the \(O\)-lengths of \(H^2(U,T\otimes R_i)\) and
\(H^2(U,T\otimes D_i)\), respectively.

\begin{theorem}[The length lost on separating two integral lattices]\leavevmode
\label{thm:intersection279}
Fix \(1\le i\le a\), put \(m=p^{i-1}\) and
\(W=k[X]/(X^m-1)\). Map \(R_{i-1}\) to \(W\) by reduction
and map \(D_i\) to \(W\) by \(\zeta_{p^i}\mapsto X\).
Let \(I_0,I_1\) be the respective images of integral
degree-one cohomology in \(H^1(U,T\otimes W)\).
Then both have \(k\)-dimension \(m\), and
\[
       M_i=M_{i-1}+H_i-m+\dim_k(I_0\cap I_1).
\]
In particular \(0\le \sum_{j=0}^aH_j-M_a\le(p^a-1)/(p-1)\).
For a cyclic extension of degree \(p\) the correction is zero
when the two residue lines agree and one when they are distinct.
\end{theorem}
\begin{proof}
Write \(g=X^m-1\), \(h=\Phi_{p^i}(X)\). In \(O[X]/g\) one has
\(h=p\), by summing the first \(p\) powers of \(X^m\).
Consequently \(O[X]/(g,h)=W\). The intersection of the ideals
\((g),(h)\) is \((gh)\). Indeed the two monic polynomials are
coprime over \(\operatorname{Frac}O\), and monic division gives
the asserted divisibility already over \(O\).
Taking the difference of the two quotient maps therefore gives
an exact sequence of Galois modules
\[
 0\longrightarrow R_i\longrightarrow R_{i-1}\oplus D_i
                \longrightarrow W\longrightarrow0 .
\]
For surjectivity onto the pairs with equal images, if polynomial
representatives \(u,v\) differ by \(gA+hB\), use the common lift
\(u-gA=v+hB\). In \(D_i\) reduction means quotient by the
element \(\delta=\zeta_{p^i}^{\,m}-1\). The field extension for
\(D_i/O\) is totally ramified, so \(D_i/\delta D_i\simeq W\)
has dimension \(m\) over \(k\).

After tensoring with the free lattice \(T\), all the integral
groups in degree two are finite. For \(D_i\) this follows from
Theorem~\ref{thm:openindex269}, since the orders of the two
finite characters are coprime and their product generates the
coefficient ring. For \(R_i\) use its decomposition over the
fraction field and finite generation of global cohomology.
There is no cohomology above degree two here. For finite
coefficients this is the odd-\(p\) cohomological dimension
bound with \(p\) inverted. Passage to compact lattices preserves
it since the finite cohomology groups satisfy the
Mittag--Leffler condition. The degree-one group for \(D_i\)
is free of rank one over \(D_i\) by the same theorem.
For \(R_{i-1}\) the rational decomposition gives rank \(m\)
over \(O\). Complex conjugation acts by \(-1\) on the lattice
modulo \(p\), since its image in \(P\) is trivial. Thus the
invariant subgroup is zero modulo \(p\), and the multiplication
sequence shows that degree one is torsion free, hence free over
\(O\).

Reduction of degree one for \(R_{i-1}\) has kernel its multiple
by \(p\). For \(D_i\) the kernel is the multiple by \(\delta\),
by the corresponding multiplication sequence. These prove the
two dimensions in the statement. Put
\(u=\dim_k H^2(U,T\otimes R_{i-1})[p]\).
Reduction by \(p\) and vanishing in degree three give
\[
 \dim_k H^1(U,T\otimes W)=m+u,\qquad
 \dim_k H^2(U,T\otimes W)=u.
\]
For the second equality a finite \(O\)-module has equally many
elements in its \(p\)-torsion and its quotient modulo \(p\).
In the long exact sequence of the displayed pair of lattices,
the cokernel in degree one has dimension
\(m+u-\dim_k(I_0+I_1)\). It injects into degree two for \(R_i\).
The next image has length \(M_{i-1}+H_i-u\), by surjectivity
onto degree two for \(W\). Adding and using
\(\dim_k(I_0+I_1)=2m-\dim_k(I_0\cap I_1)\) gives the formula.
Summing in \(i\) and using \(M_0=H_0\) proves the bounds.
\end{proof}

\subsection{Determining a compact module from finitely many twists}\label{nav:IV-46}
\begin{theorem}[A finite determination of a cyclic cyclotomic module]\leavevmode
\label{thm:specializations280}
Let \(p\) be prime, \(\Lambda=\ZZ_p[[T]]\), and \(Y\) an inverse
limit of finite abelian \(p\)-groups with surjective maps and a
compatible continuous \(\Lambda\)-action. Suppose \(1+T\) acts
on each finite group with \(p\)-power order. Write
\(Y(z)=Y/(T-z)Y\) for \(z\in p\ZZ_p\). Let \(b\ge2\).
Suppose \(|Y(z_*)|=p\) for some \(v_p(z_*)=1\), and
\(|Y(z_j)|=p^b\) for \(p\) elements \(z_j\) representing all
classes of \(p^b\ZZ_p/p^{b+1}\ZZ_p\). Then \(Y\) itself is
cyclic of order \(p^b\), and \(T\) acts by zero.
\end{theorem}
\begin{proof}
The quotient \(Y/(p,T)Y\) equals \(Y(z_*)/pY(z_*)\), hence is
one-dimensional over \(\mathbb F_p\). Lift a generator to \(y\).
Its image generates each finite group as a \(\Lambda\)-module.
Indeed the quotient of that group by the image of
\(\Lambda y\) is equal to its multiple by \((p,T)\).
On this finite quotient \(p\) is nilpotent, and modulo \(p\)
some \(p\)-power of \(T\) is zero by the hypothesis on
\(1+T\). A power of the ideal \((p,T)\) consequently kills
the quotient, forcing it to vanish. The image of \(\Lambda y\)
is compact and surjects onto every finite level, so it equals
\(Y\). Write \(Y=\Lambda/I\).

Evaluation at \(z\in p\ZZ_p\) gives
\(\Lambda/(T-z)=\ZZ_p\), with convergence of every series.
Its kernel description also follows by dividing \(H(T)-H(z)\)
by \(T-z\) term by term. Thus the ideal of values \(I(z)\)
has index \(|Y(z)|\). For \(H\in I\) one has
\(H(0)\equiv H(z_j)\equiv0\bmod p^b\). Choose \(H\in I\)
with \(v_p(H(z_*))=1\), using \(I(z_*)=p\ZZ_p\).
In \(H(z_*)\) the constant and all terms of degree at least
two are zero modulo \(p^2\). Hence the linear coefficient is
a unit. Starting from \(0\bmod p^b\), the unit derivative
gives a unique Hensel lift at each successive precision.
The limit is a root \(a\in p^b\ZZ_p\).

Write \(H(T)=(T-a)U(T)\). The coefficient of \(T^i\) in
\(U\) is \(\sum_{n>i}H_n a^{n-1-i}\), with \(H_n\) the
coefficients of \(H\). It converges in \(\ZZ_p\), and the
constant coefficient reduces to \(H_1\ne0\bmod p\).
Thus \(U\) is a unit and \(T-a\in I\). It follows that
\(Y=\ZZ_p/J\) for an ideal \(J\), with \(T\) acting by \(a\).
Choose \(j\) with \(z_j\equiv a\bmod p^{b+1}\). The quotient
by \(J+(a-z_j)\) has order \(p^b\), so this ideal is
\(p^b\ZZ_p\). Since \(a-z_j\in p^{b+1}\ZZ_p\), the valuation
of ideals in \(\ZZ_p\) forces \(J=p^b\ZZ_p\). Finally
\(a\in p^b\ZZ_p\) acts as zero on \(Y\).
\end{proof}

\section{Uniform ranges of coefficient primes}\label{nav:g-II-2}

\subsection{Cell counts and full primary orders}\label{nav:II-17}
\label{sec:cells103}

Let \(L\subseteq\QQ(\zeta_M)\) be totally imaginary, with degree
\(d\) and absolute discriminant \(D_L\). Fix \(r\ge1\), put
\(i=2r\), \(n=2i+8\), and write \(\Gamma=\GL_n(\OO_L)\).
We use the free lattice \(\Lambda=\OO_L^n\). In this section all
logarithms are natural. We will bound the order of a cohomology
group by counting cells in just two degrees. This allows us to
retain torsion at the small primes when passing from the lattice
estimate to Galois cohomology.

\begin{lemma}[Cells and their incidences]\label{lem:inc103}
Let \(a\) be an upper bound for the number of \(\Gamma\)-orbits of
cells in each of dimensions \(i\) and \(i+1\) of Ash's well-rounded
retract. Put \(J=[\Gamma:\Gamma_3]\), where \(\Gamma_3\) is the
principal subgroup of level three. Write \(D\) for the real
dimension of the space of Hermitian metrics before positivity and
normalization, so \(D=dn^2/2\) for this totally imaginary field.
If \(T=|H_i(\Gamma_3,\ZZ)_{\rm tors}|\), then
\[
 \log T\ \le\ J a\big\{\log(2Ja)+(D+1)\log(1+3^{dn})\big\}.
\]
\end{lemma}
\begin{proof}
A metric \(h\) gives a positive definite real quadratic form on the
rank \(dn\) lattice \(\Lambda\) by taking the archimedean trace.
Write \(h[v]\) for its value and \(S(h)\) for the set of vectors
of squared length one when the minimum is one. For a well-rounded
metric \(S(h)\) spans \(L^n\) over \(L\). Ash's retract is
contractible, has finite cell stabilizers, and has one open cell
for each such set \(S(h)\). Closing a cell allows additional
shortest vectors.~\cite[Sec.~2.3]{IV-II-III-beh}

Here is the geometry of these closures. Write \(P\) for the set of
positive definite metrics with \(h[v]\ge1\) for all
\(v\in\Lambda\setminus\{0\}\). For \(S=S(h_*)\) the closed cell is
\[
                 F_S=\{h\in P:\ h[v]=1\quad(v\in S)\}.
\]
Indeed the open segment from any point of this set to \(h_*\) has
exactly \(S\) as its shortest vectors. To see compactness, take an
\(L\)-basis among \(S\). In the corresponding basis at every
embedding, each diagonal coefficient is bounded above since it
contributes positively to a trace of value one. Cauchy--Schwarz
then bounds every entry. A limit could only leave the positive
cone if the determinant of the real trace form tended to zero in
a fixed lattice basis. To apply Minkowski's first theorem, observe
that the ellipsoid of squared length less than one contains no nonzero
lattice vector, and its volume is a fixed positive constant
divided by the square root of that determinant. Minkowski's bound
on this volume bounds the determinant away from zero. Thus \(F_S\) is
compact inside the open cone. Enclose it in the interior of a
compact polytope in that cone, for instance by taking the convex
hull of finitely many small closed cubes about points of \(F_S\).
On that polytope the smallest eigenvalue of the trace form has a
positive lower bound. Only finitely many vectors can therefore
reach length at most one anywhere on it. The displayed intersection
is a convex polytope. At a relative interior point, the inactive
inequalities are locally strict.

The functional \(h\mapsto h[v]\) is real linear in the metric.
A basis of the span of the functionals for \(v\in S\) determines
the whole face by setting those basis values equal to one and
intersecting with \(P\). In fact every further functional in \(S\)
is a linear combination of this basis, and its coefficients sum
to one by evaluation at \(h_*\). If \(F_S\) is a face of another
cell, that cell has a defining set contained in \(S\). A list of
at most \(D\) members of \(S\) therefore specifies each cell
containing \(F_S\). Open balls of radius \(1/2\) about the shortest
vectors in the trace norm have disjoint interiors and fit inside
the ball of radius \(3/2\) about zero. Comparing volumes gives
\(|S|\le3^{dn}\). Thus a fixed cell lies in at most
\(B=(1+3^{dn})^{D+1}\) closed cells.

The level-three subgroup is torsion free by Lemma~\ref{lem:top},
so its cell stabilizers are trivial. In each of the two degrees the
quotient has at most \(Ja\) cells. Count orbits of pairs consisting
of an \(i\)-cell and a containing \((i+1)\)-cell, beginning with
the \(i\)-cell. There are at most \(JaB\) such pairs up to
\(\Gamma_3\). The distinct facets of a fixed \((i+1)\)-cell give
distinct pair orbits, since an identification between the pairs
would stabilize that cell. Its boundary therefore has sum of
absolute coefficients at most \(JaB\), including after faces are
identified in the quotient. Each incidence in the polytope before
identification contributes one sign. The quotient is a
classifying space for \(\Gamma_3\). The torsion in its degree
\(i\) cellular homology embeds in the torsion of the cokernel of
the degree \(i+1\) boundary. If that boundary has rank \(s>0\),
the latter torsion order is the greatest common divisor of its
nonzero \(s\)-rowed minors, by Smith normal form. Hadamard's
inequality bounds each by \((JaB)^s\), with \(s\le Ja\).
For \(s=0\) the cokernel is free. This proves the estimate.
\end{proof}

\begin{theorem}[Two uniform bounds for full strict orders]\label{thm:cells103}
With \(b=r+2\), write \(U(L,r)=\sum_p\log|\Sel_{p,r}(L)|\).
The sum includes every rational prime and every elementary
divisor at that prime. There is an absolute constant \(C\) with
\begin{equation}
 U(L,r)\ \le\ \exp\!\big(C(db)^5\log(2Mdb)\big).          \label{eq:pub103}
\end{equation}
Using the cell-orbit estimate in Theorem 3.9 of the November 2026
preprint of de la Torre Dur\'an gives the further bound
\begin{equation}
 U(L,r)\ \le\ \exp\!\big(Cdb^2\log(2Mdb)\big).             \label{eq:vor103}
\end{equation}
\end{theorem}
\begin{proof}
For \eqref{eq:pub103}, use the finite set \(\Phi\) in
Corollary 5.2 of Bayer-Fluckiger, Emery and Houriet. Some element
of \(\Gamma\) moves the entire shortest-vector set of any given
well-rounded metric into \(\Phi\). The assumptions \(n\ge5\) and
\(L\ne\QQ\) hold here. The basis of defining functionals in
Lemma~\ref{lem:inc103} then bounds the number of cell orbits in
each degree by \(a=(1+|\Phi|)^{D+1}\). By the explicit bound for
\(\Phi\) in Lemma~\ref{lem:top},
\[
 \begin{split}
 \log(2+|\Phi|)&\ll
 n^2d^2\log(2n)+n^3d^4\log(2d)+n^3d^3\log D_L,\\
 \log(2a)&\ll (nd)^5\log(2Mnd).
 \end{split}
\]
Here \(D_L\le M^d\) by the conductor--discriminant formula.
This proves the orbit bound needed for the first assertion.~\cite{IV-II-III-beh}

For \eqref{eq:vor103} we relate the two required cell degrees
to the dimensions in the preprint. For a metric \(h\) in a
\(j\)-cell take the cone \(V(h)\) generated by \(vv^*\) for
\(v\in S(h)\) in the real space of Hermitian matrices. The trace
pairing identifies these matrices with the defining functionals
above. The cone has linear dimension \(D-j\), since near \(h\)
the defining equalities cut out precisely the cell. It is a face
of a perfect-form cone. To see this choose a vertex \(h_0\) of the
compact polytope \(F_{S(h)}\). If its active functionals failed
to span the dual space, a short segment about \(h_0\) annihilating
their linear parts would lie in that polytope, contradicting
extremality. Thus \(h_0\) is perfect. On \(V(h_0)\) the functional
\(h-h_0\) is nonnegative, and its zero face is exactly \(V(h)\).
Moreover \(\sum_{v\in S(h)} vv^*\) is positive definite at every
embedding because \(S(h)\) contains an \(L\)-basis. A relative
interior point of this cone is a positive combination of its
generating vectors, so the relative interior lies in the positive
cone.

Distinct cells give distinct cones. In fact, suppose \(V(h)=V(h')\).
For \(v\in S(h)\) express \(vv^*=\sum_{w\in S(h')}c_w ww^*\)
with \(c_w\ge0\). Pairing with \(h\) gives
\(1\ge\sum c_w\), and pairing with \(h'\) gives
\(h'[v]=\sum c_w\ge1\). Hence \(v\in S(h')\). Reversing the
roles gives equality of shortest-vector sets. Change of lattice
basis preserves this correspondence, acting on normals by the
inverse adjoint to its action on metrics. It therefore gives an
injection on \(\Gamma\)-orbits.

In the dimension convention of the preprint, \(V(h)\) has
dimension \(k=D-j-1\). For \(j=i,i+1\) one has
\(dn-1\le k\le D-1\), because \(D=dn^2/2\) and \(n=2i+8\).
This verifies the lower dimension cutoff in its Theorem 3.9.
The field \(L\) is abelian and totally imaginary. Its subfield
fixed by complex conjugation is totally real of index two, so the
CM case of the cited theorem counts all the required cone orbits
with an upper bound
\[
 a\ \le\ \exp\!\big(Cdn^2\log(2nd)\big)\,D_L^{\,4n(n+1)} .
\]
Thus we may use \(\log(2a)\ll dn^2\log(2Mnd)\).
Its separate Theorem 4.5 estimates \(K\)-torsion above a
prime cutoff. Here Theorem 3.9 supplies the cell count before
the passage to cohomology.~\cite{IV-II-III-torre103}

We finish the arithmetic argument for either choice of \(a\).
Reduction modulo three gives \(J\le3^{dn^2}\), by counting
matrices over the ring \(\OO_L/3\) of cardinality \(3^d\).
Each prime factor of \(J\) is at most \(W=3^{dn}\), by the
embedding of the reduction group in \(\GL_{dn}(\mathbf F_3)\).
Write \(A=(nd)^5\log(2Mnd)\) for the first bound and
\(A=dn^2\log(2Mnd)\) for the second. Lemma~\ref{lem:inc103}
now gives \(\log(2JT)\le\exp(CA)\) in both cases. Indeed
\(\log(2Ja)\ll A\), while
\((D+1)\log(1+3^{dn})\ll d^2n^3\) appears only as a
multiplicative factor in the bound for \(\log T\).

Recall the arithmetic steps of Lemmas~\ref{lem:top}
and~\ref{lem:rank}. Duality and the comparison in even
\(K\)-theory embed the dual of \(\Sel_{p,r}(L)\) in
\(K_i(\OO_L)\{p\}\). A stable homology torsion class transfers
into the group of order \(T\) for \(\Gamma_3\). Composing this
transfer with inclusion multiplies by \(J\), so \(JT\) kills the class.
The Hurewicz kernel on its \(p\)-part is killed by \(p^{i+4}\).
For \(p>W\), both Hurewicz and transfer are injective on that
part. Hence all \(p>W\) together contribute at most \(\log T\)
because there is a single target torsion group of that order.
For \(p\le W\), each cyclic summand has logarithmic order at
most \((i+4)\log W+\log(2JT)\). The Kummer bound of
Lemma~\ref{lem:rank} gives at most
\(Cd p\log(2Mp)/\log p\) summands at \(p\), independently of
their exponents. Summing and using \(\sum_{p\le W}p\le W^2\)
gives
\[
 U(L,r)\le\log T+
 C dW^2\log(2MW)\big((i+4)\log W+\log(2JT)\big).
\]
Since \(\log W=dn\log3\) and \(n=4(r+2)\), enlarging the
absolute constant gives \eqref{eq:pub103} and \eqref{eq:vor103}.
\end{proof}

\subsection{All twists at a growing set of odd primes}\label{nav:II-18}
\label{sec:log147}

Fix \(\psi\) of conductor \(F\) and retain \(B_j=K_\psi C_j\).
Write \(M_{p,r}(L)=\tfrac12(1+(-1)^r c)\Sel_{p,r}(L)\) at
an odd prime \(p\). The layer \(a_p\) of
Theorem \ref{thm:effect146}
works at every twist \(r\). We use this fact to make the number
of excluded indices independent of the length of the interval
over which the auxiliary primes will be selected.

\begin{lemma}[Exceptional twists on the fixed prime-dependent layers]
\label{lem:budget147}
For every \(P\ge3\) and integer \(T\ge2\) there is a set
\(\mathcal D\subseteq[1,T]\cap\ZZ\) with
\(|\mathcal D|\ll_\psi P^5\) such that, for every \(1\le r\le T\)
outside \(\mathcal D\) and every \(J\ge0\),
\begin{equation}
      \sum_{\substack{2<p\le P\\p\ {\rm prime}}}
           \log|M_{p,r}(B_J)|
          \ll_\psi P^{10}(1+\log(2T)).
      \label{eq:budget147}
\end{equation}
\end{lemma}
\begin{proof}
For each \(p\) it suffices to prove the estimates in
\(B_{a_p}\), by Theorem \ref{thm:effect146}. Let \(n_p\) be its degree
and \(N_p=F2^{a_p+2}\) a cyclotomic modulus. Both are
\(O_\psi((p+1)4^{e_p})\), with \(e_p=v_2(p^D-1)\) and
\(D=2\varphi([K_\psi:\QQ])\). Both are \(O_\psi(p^3)\).
For \(p\) fixed, apply Corollary \ref{cor:family143} to the
characters of the prime-to-\(p\) part of this field's Galois
group, on \([1,T]\) using the upper endpoint \(T+1\) in that
corollary. If its \(p\)-part has order \(d\), a rational
conjugacy orbit \(\mathcal O\) contributes
\(O(d|\mathcal O|p\log(2N_pp)/\log p)\) exceptions for each
residual index \(h\). Sum over at most \(p-1\) indices and use
\(\sum_{\mathcal O}d|\mathcal O|=n_p\). Since
\(\log(2N_pp)/\log p\ll_\psi1\), the resulting number is
\(O_\psi(n_p p^2)\). This argument includes \(p\mid n_p\).

To sum the exceptions, write \(w_p=\max(v_2(p-1),v_2(p+1))\).
Then \(e_p=w_p+v_2(D)\), by equation \ref{eq:level146} and its proof.
For \(P\ge3\) one has
\[
  \sum_{2<p\le P}4^{e_p}\ll_\psi
       \sum_{1\le m\le P+1}4^{v_2(m)}\ll P^2 .
\]
Indeed the map from primes \(p\) to either \(p-1\) or \(p+1\)
at which the maximum is attained has multiplicity at most two.
There are at most \((P+1)/2^a\) integers of valuation \(a\).
Sum the geometric progression for \(2^a\le P+1\). The union
over the primes thus has size
\(O_\psi(\sum_{p\le P}p^3 4^{e_p})=O_\psi(P^5)\).

At every other twist the exponent of \(M_{p,r}(B_{a_p})\),
written as a power of \(p\), has exponent
\(O_\psi(P^4(1+\log(2T)))\). Indeed \eqref{eq:orbit143}
bounds each \(1+V_{\alpha,h}\) by \(O_\psi(n_p p)\).
The interpolation estimate in Corollary~\ref{cor:family143} multiplies this
by \(O(1+\log(2T))\) and adds at most
\(2(1+v_p(n_p)+\log_p(T+1))\). The components exhaust the group.

The logarithmic order of its subgroup killed by \(p\) is
\(O(n_p p\log(2N_pp))\). To recall the first-layer argument,
pass to the field \(L/B_{a_p}\) cut out by the residual
cyclotomic twist, of relative degree at most \(p-1\).
Theorem \ref{thm:strict113} bounds the first-layer logarithm by
\(\log h(L)+[L:\QQ]\log p\), where \(h(L)\) is the class
number. The conductor--discriminant formula for
\(L\subseteq\QQ(\zeta_{N_pp})\) gives
\(|\operatorname{disc}L|\le (N_pp)^{[L:\QQ]}\). Minkowski's
bound for representatives and
\(\zeta_L(2)\le\zeta(2)^{[L:\QQ]}\) give the asserted estimate,
also for totally real fields at an odd prime. For a finite
abelian group killed by \(p^j\), its logarithmic order is at
most \(j\) times the first-layer logarithm, term by term in
its elementary divisors. Multiplying the two bounds above,
then summing over at most \(P\) primes, proves \eqref{eq:budget147}.
\end{proof}

\begin{theorem}[An explicit growing prime range at every twist]
\label{thm:log147}
For fixed \(0<\eps<1/52\) there is \(c_{\psi,\eps}>0\).
On arbitrarily distant intervals \([V,V^{3/2}]\) carrying the
auxiliary fields and weights of Section \ref{nav:I-28}, one has
for every \(x\) in the interval and every \(J\ge0\)
\begin{equation}
 \sum_{\substack{2<p\le c_{\psi,\eps}(\log\log x)^{1/5}\\
                  p\ {\rm prime}}}
       \quad\sum_{1\le r<\lfloor x^{401}\rfloor}
                \log |M_{p,r}(B_J)|
                  \ \le\ \exp(x^{2\eps/3}).
 \label{eq:log147}
\end{equation}
\end{theorem}
\begin{proof}
Choose \(\eta=\eps/80\) and \(b=401/\eta\). Fix a constant
\(W>2\) large enough for the dyadic interval selection of
Theorem \ref{thm:common16}.
Specifically take \(W>2(2b)^{R_2}+1\), where \(R_2=R_2(\psi)\)
is the bound there for the number of dyadic exceptions.
For a sufficiently large \(Z\), set \(P=\lceil\log\log Z\rceil\)
and choose an integer \(R=\lceil C_\psi P^5\rceil\) using
Lemma~\ref{lem:budget147}. Form
\[
     K=W(Wb)^R,\qquad T=\lceil Z^{401K}\rceil+1 .
\]
The estimate for \(R\) is independent of \(T\), so we may now
use the lemma with this \(T\) and obtain a set \(\mathcal D\).
All the fields in that use of the lemma are the specified
\(B_{a_p}\). They are independent of the auxiliary quadratic
primes.

Start at \(U=K\log Z\). Place the points \((\log r)/401\),
\(r\in\mathcal D\), on the real line. If some point is in
\((U/(Wb),U]\), lower \(U\) to \(U/(Wb)\). The point just
encountered exceeds the new upper endpoint, and so cannot be
encountered again. After at most \(R\) such divisions the band
is empty. Set \(Y=\exp(U/W)\) at this stage. Since
\(U\ge K\log Z/(Wb)^R=W\log Z\), one has \(Y\ge Z\) and
\(Y^W\le Z^K\). Throughout \([Y,Y^W]\), an exceptional \(r\)
with \(r\le x^{401}\) satisfies \(\log r/401\le U\). The empty
band forces \(\log r/401\le U/(Wb)\le(\log x)/b\), or
\(r\le x^\eta\). This conclusion holds for every \(p\le P\)
at once.

Now select the two auxiliary primes at the initial scale \(Y\)
by Theorem \ref{thm:twists}, with the fixed upper power \(W+1\).
The dyadic construction in the proof of Theorem \ref{thm:common16}
takes place inside \([Y,Y^{2(2b)^{R_2}}]\), since that is its
initial range for the dyadic exceptional indices. Its proof
applies at every sufficiently large initial scale. It produces
\([V,V^{3/2}]\subseteq[Y,Y^W]\) with the specified fields,
weights and dyadic bounds. The estimates on ranks, auxiliary
primes and real period levels in Volume~I then apply at the
points of this interval. In particular the auxiliary-prime
selection has used a power fixed before \(Z\), although \(K\)
used for the independent cyclotomic layers grows with \(Z\).

If \(x\) is in this interval, then
\(\log\log x\le\log\log Z+\log K\ll_{\psi,\eps}P^5\).
Choose \(c_{\psi,\eps}\) small enough that every prime in
\eqref{eq:log147} is at most \(P\). All its twists lie in
\([1,T]\). For the twists outside \(\mathcal D\) their total
contribution is at most
\[
         C_\psi x^{401}P^{10}(1+\log(2T)).
\]
Its logarithm is
\(O_{\psi,\eps}(\log x+P^5+\log\log Z)\), which is
\(o(x^{\eps/3})\) uniformly since \(x\ge Z\) and
\(P=\lceil\log\log Z\rceil\).

For an exceptional \(r<x^{401}\) we have \(r\le x^\eta\).
Theorem \ref{thm:effect146} embeds \(M_{p,r}(B_J)\) into
\(M_{p,r}(B_{a_p})\), using the inverse of its isomorphism if
\(J>a_p\). Restriction through the extension of degree at most
two then embeds this latter group into \(\Sel_{p,r}(B_{a_p}(i))\).
The latter field is totally imaginary of degree \(O_\psi(P^3)\)
and has modulus \(O_\psi(P^3)\). Bound its single-prime
logarithmic order by the all-prime estimate \eqref{eq:all}.
The exponent on the right of that estimate is
\[
     O_\psi\bigl((P^3(x^\eta+2))^{10}
                         \log(2P^6(x^\eta+2))\bigr)
                  \ \le\ x^{\eps/3}
\]
eventually, since \(10\eta=\eps/8\). There are at most
\(PR\) pairs to count in this part. Consequently the two
parts together are at most \(\exp(x^{2\eps/3})\) for all
sufficiently large \(Z\), uniformly in \(J\). Letting \(Z\)
tend to infinity proves the assertion.
\end{proof}

\subsection{Testing whole character components on the common interval}\label{nav:II-28}
Keep \(E=E_x\), its degree \(n_E\), modulus \(N\), and
\(t=\lfloor x^{401}\rfloor\) as in Theorem~\ref{thm:res100}.
For each odd \(\ell\nmid Nn_E\), take a finite unramified
coefficient ring \(O/\ZZ_\ell\) containing the character values
of \(\Gal(E/\QQ)\). For a character \(\alpha\) of sign \((-1)^r\)
put \(\vartheta=\kappa_\ell^{-r}\alpha^{-1}\). Denote the
corresponding strict group over \(\QQ\) by \(S(\vartheta)\),
as in Sections \ref{nav:III-02} and \ref{nav:III-14}. It gives the \(\alpha\)-summand of the strict group
over \(E\) after extending scalars. The normalized length is
\(\operatorname{length}_O\); summing it over all characters recovers
the original \(\ZZ_\ell\)-length since \(O/\ZZ_\ell\) is unramified.

\begin{theorem}[A polynomial bound for all components certified by a finite test]
\label{thm:stable102}
Put \(h=1+v_\ell(r)\). For \(\bar\vartheta\ne1\) let \(C_h\) be
the finite class group component of Theorem \ref{thm:ray101}.
Include a triple \((\ell,r,\alpha)\) in the following sum if
\(\ell^{h-1}\) kills \(C_h/\ell^h C_h\). Also include every
triple with \(\bar\vartheta=1\). Then on every indicated field
\[
 \sum_{\substack{\text{included triples}\\1\le r<t}}
       (\log\ell)\operatorname{length}_{O} S(\vartheta)
       \ \ll\ n_E^2 t^2\log(2Nt)\log(2t).
\]
The groups in this estimate are taken with their full exponents.
The bound is \(o(\exp(x^\eps/2))\) uniformly on the common
intervals.
\end{theorem}
\begin{proof}
If \(\bar\vartheta=1\), restrict first to inertia at \(\ell\).
The character \(\alpha\) is unramified there, so \(\ell-1\mid r\).
Then \(\kappa_\ell^{-r}\) is trivial modulo \(\ell\) globally,
and \(\alpha=1\) since its values have order prime to \(\ell\).
The twist \(r\) is even. By Proposition~\ref{prop:realclass100}
its group over \(\QQ\) is zero, since the trivial eigenspace of the
class group modulo \(\ell\) in \(\QQ(\zeta_\ell)\) comes from
\(\operatorname{Cl}(\QQ)=0\). Indeed the norm of an extended
class is its \((\ell-1)\)-st power, and extending the norm gives
the product of all Galois conjugates. Dividing by \(\ell-1\)
on the \(\ell\)-primary group gives the invariant part.

For the other included triples put \(L=\overline\QQ^{\ker\bar\vartheta}\).
Corollary \ref{cor:stopped101}, shows that the entire strict group is
killed by \(\ell^{v_\ell(r)}\). Its number of cyclic summands is
at most \(1+\log|\operatorname{Cl}(L)|/\log\ell\).
For a nonsplit prime this follows directly from the class group
description at level one. For a split prime the ray class sequence
adds at most one generator on the character component, since
the permutation module on the primes over \(\ell\) is the regular
representation. The \(S\)-class group is a quotient of the
ordinary class group.

The field \(L\) has degree at most \(n_E(\ell-1)\) and conductor
dividing \(N\ell\). The conductor--discriminant formula gives
\(\log|\operatorname{disc}L|\le [L:\QQ]\log(N\ell)\).
Every ideal class has an integral representative of norm at most
\(2^{[L:\QQ]}|\operatorname{disc}L|^{1/2}\) by Minkowski's
theorem. There are at most \(Y^2\zeta(2)^{[L:\QQ]}\) integral
ideals of norm at most \(Y\), by bounding the ideal Dirichlet
series prime by prime. It follows that
\(\log|\operatorname{Cl}(L)|\ll n_E\ell\log(2N\ell)\).
Thus the contribution from one triple is at most
\[
   v_\ell(r)\bigl(\log\ell+C n_E\ell\log(2N\ell)\bigr).
\]
Only primes dividing \(r\) give a positive contribution. There
are at most \(n_E\) characters for each prime and each \(r\).
Use \(\ell\le r\) at these primes and
\(\sum_{\ell\mid r}v_\ell(r)\le\log r/\log2\), then sum over
\(r<t\). This proves the stated bound. The field degree, \(t\),
and \(\log(2N)\) are all at most fixed powers of \(x\), which
proves the uniform assertion.
\end{proof}

\begin{proposition}[The rational integer bound in the source]
\leavevmode\label{prop:qbound102}
Let \(r\ge2\) be even. If \(\ell\) is an odd prime with
\(\ell-1\mid 2r\), then \(K_{2r}(\ZZ)\{\ell\}=0\).
For any choice of coefficient fields in the source's rational
integer bound the exact equality
\begin{equation}
 \sup_{P}\ \sum_{\ell\in P}(\log\ell)\,
     \frac{\operatorname{length}_{O_\ell} S_\ell^{(2)}(r;O_\ell)}
          {e(O_\ell/\ZZ_\ell)}
                         =\log|K_{2r}(\ZZ)|                       \label{eq:qbound102}
\end{equation}
holds, where \(P\) runs over finite sets of primes
\(\ell\ge5\) with \(\ell-1\nmid 2r\).
Here \(S_\ell^{(2)}\) uses coefficients
\((\operatorname{Frac}O_\ell/O_\ell)(-r)\), local triviality
at \(\ell\) and at 2, and trivial restriction to inertia at
all other finite primes. Consequently a polynomial bound in \(r\)
for the sums in \eqref{eq:qbound102} is equivalent to one for
\(\log |K_{2r}(\ZZ)|\) at even \(r\).
\end{proposition}
\begin{proof}
By Proposition~\ref{prop:realclass100} the rank modulo \(\ell\)
is the dimension of the \(\omega^{-r}\)-part of the real
cyclotomic class group modulo \(\ell\). The order of that
character divides two when \(\ell-1\mid2r\).
Its eigenspace descends to the field cut out by the character.
Indeed the subgroup fixing that field has order prime to \(\ell\).
Extension of ideals followed by norm multiplies by its order,
and in the reverse order is the sum of its elements on ideal
classes. Thus ideal extension identifies the downstairs
\(\ell\)-primary class group with the invariants upstairs, also
after reduction modulo \(\ell\).

The trivial character gives \(\QQ\). A nontrivial quadratic
character here is even and cuts out \(\QQ(\sqrt{\ell})\), of
discriminant \(\ell\). This follows also from the
conductor--discriminant formula for its single nontrivial
character. In this real quadratic field each ideal class has
an integral ideal of norm at most \(Y=\sqrt{\ell}/2\).
There are at most \(\sum_{n\le Y}\sum_{a\mid n}1\le\lfloor
Y\rfloor^2\) such ideals. Indeed the number at a prime power
\(p^j\) is at most \(j+1\) by prime factorization of quadratic
ideals, and each pair counted in the last sum has two positive
coordinates at most \(Y\). Hence the class number is less than
\(\ell\). Its \(\ell\)-part vanishes. The computed dimension is
zero in both cases, so the finite \(\ell\)-primary \(K\)-group
vanishes.

For odd \(\ell\) the dual of the rational strict group with
\(\QQ_\ell/\ZZ_\ell\) coefficients is all
\(K_{2r}(\ZZ)\{\ell\}\), by the duality argument of
Corollary~\ref{cor:transfer100}. Requiring local triviality at
2 in addition does not change its order. An unramified local
class at 2 is determined by its value at Frobenius, modulo
multiples by \(2^{-r}-1\) on \((\QQ_\ell/\ZZ_\ell)(-r)\).
This is a nonzero scalar in \(\ZZ_\ell\), so multiplication
is surjective. All such local classes vanish. The same statement
holds after finite scalar extension to \(O_\ell\). Concretely
\((\QQ_\ell/\ZZ_\ell)\otimes_{\ZZ_\ell}O_\ell
=\operatorname{Frac}(O_\ell)/O_\ell\). This finite free base
change commutes with the cochains and with the restriction
kernels. On a cyclic summand \(\ZZ_\ell/\ell^a\) the resulting
length is \(a\,e(O_\ell/\ZZ_\ell)\). This accounts for the
denominator in \eqref{eq:qbound102}.

There is no dyadic torsion in \(K_{2r}(\ZZ)\) for positive even
\(r\), by the number-ring computation at two in
\cite[Ch.~VI, Cor.~9.8]{IV-II-I-weibel102}. The prime 3 always satisfies
\(\ell-1\mid2r\). All omitted primes have therefore been
accounted for. There are only finitely many prime divisors of
the group order, so the supremum equals the stated logarithm.
\end{proof}

The rational inequality in the supplied archive is formulated at
\srcloc{Development/AutomorphicArithmetic/AutomorphicVarieties/SmallClassNumbers/ClassNumberThree}{16}.
Its deduction at line 89 of that file uses the comparison of
strict lengths with Frobenius-trace congruences. The single-character
inequality is written at
\srcloc{Development/AutomorphicArithmetic/AutomorphicVarieties/UnitaryThreeGroups/KatoUse}{54}.
Proposition~\ref{prop:qbound102} identifies its rational target.
A uniform upper bound for that target would require the comparison
or an independent bound on the orders. The theta-period height in
Corollary~\ref{cor:alg97} is a further requirement on the same
intervals.

\subsection{All complementary characters for a range of coefficient primes}\label{nav:II-32}
We continue on the field \(E=K_\psi C_{l'}(\sqrt{-q},\sqrt{-d_0})\)
of the period, with \(n=[E:\QQ]\), conductor dividing \(N\),
and \(t=\lfloor x^{401}\rfloor\).
Write \(M_{p,r}=\tfrac12(1+(-1)^r c)\mathcal S_{p,r}(E)\).
This time the bound is for that whole eigenspace on \(E\).
There is no restriction of \(r\) to a prescribed residue class.

\begin{theorem}[A full-field bound away from an exceptional set of twists]
\label{thm:alltw142}
Take any odd prime \(p\nmid n\) with \(p\le x^{199}\). On every
sufficiently large common interval, there is a subset
\(\mathcal E_p\subset\{1,\ldots,t-1\}\) with
\begin{align}
 |\mathcal E_p|
     &\ll n p^2\,\frac{\log(2Np)}{\log p},                 \label{eq:alltw142}\\
 p^{\,\lceil x^{202}\rceil}M_{p,r}&=0
                 \qquad(1\le r<t,\ r\notin\mathcal E_p). \label{eq:exp142}
\end{align}
The implicit constant in \eqref{eq:alltw142} is absolute.
In particular, uniformly over these primes,
\[
       \frac{|\mathcal E_p|}{t-1}
              \ll_{\psi,\eps} x^{-1+33\eps/4}.
\]
\end{theorem}
\begin{proof}
Extend scalars to an unramified \(p\)-adic coefficient ring
containing the values of all characters of
\(\operatorname{Gal}(E/\QQ)\). Character descent as in
Theorem \ref{thm:char}
and the integral idempotents for this prime-to-\(p\) group
split \(M_{p,r}\) into the strict components with
\(\alpha(-1)=(-1)^r\). Theorems~\ref{thm:moment141} and~\ref{thm:inter141}
apply to each component in each class of \(r\bmod(p-1)\).
They include the case \(p\mid N\), with the tame part
\(\omega^i\) at \(p\) separated from \(\alpha'\).
The components with trivial residual character already vanish.

To count exceptions, group the characters \(\alpha\) into
their orbits \(\mathcal O\) under rational Galois conjugacy.
Use \(h=r+1+i\bmod(p-1)\) as the index of the residue class
for each \(\alpha\). Holding \(h\) fixed holds fixed the integer
\(k\in[2,p]\) used to evaluate the annihilator in equation \ref{eq:moment141}. The parity condition becomes \(\alpha'(-1)=(-1)^{h-1}\).
It is invariant throughout each orbit even when its tame
exponent \(i\) changes. For any allowed orbit and \(h\), write
\(V_{\alpha,h}\) for the valuation of that nonzero evaluation.
Then
\begin{equation}
 \sum_{\alpha\in\mathcal O}(V_{\alpha,h}+1)
       \ll |\mathcal O|\left(1+
                    k\,\frac{\log(2Np)}{\log p}\right).
 \label{eq:orbit142}
\end{equation}
Indeed \(V_{\alpha,h}=1+v_p(B_{k,\bar\alpha'})\) by the proof
of Theorem \ref{thm:moment141}. The conductor \(f_0\) of \(\alpha'\)
is fixed on the orbit. The number
\(f_0(k+1)!\,B_{k,\bar\alpha'}\) is an algebraic integer, and
each of its conjugates has modulus at most
\(\exp(Ck\log(2Np))\), as shown there. Regard it in the value
field of \(\alpha\). The embeddings of that field correspond
bijectively to \(\mathcal O\). Repeated restrictions to the
value field of \(\alpha'\) give exactly the multiplicities
in its rational norm. The sum of the valuations of all these
conjugates at a fixed \(p\)-adic embedding is the valuation
of that norm. Its absolute size gives \eqref{eq:orbit142}
since clearing a denominator never decreases valuation.
One may include the already vanishing residual components in
this upper bound using the same evaluation formula.

By Theorem \ref{thm:inter141}, for an individual pair
\((\alpha,h)\) at most \(2(V_{\alpha,h}+1)\) distinct
twists escape its exponent bound. Take the union over all
pairs, including only their stipulated residue classes.
There are at most \(p-1\) choices of \(h\) per orbit and
\(\sum_{\mathcal O}|\mathcal O|=n\). Using \(k\le p\) in
\eqref{eq:orbit142} now proves \eqref{eq:alltw142}.

For an individual component the common upper bound for its
exponent as a power of \(p\) away from that union is
\[
  C\left(1+n p\,\frac{\log(2Np)}{\log p}\right)
                   \left(1+\frac{\log t}{\log p}\right).
\]
This uses the individual estimate of equation \ref{eq:V141}, and
the definition \(L=V+2(V+1)\lfloor\log_p t\rfloor\).
Here \(n\ll_\psi x^{2+8\eps}\),
\(\log(2Np)\ll_{\psi,\eps}x^{\eps/4}\), and \(p\le x^{199}\).
In the bound for the conductor one uses also
\(\log x=O_\eps(x^{\eps/4})\).
Since \(0<\eps<1/52\), the displayed bound is at most
\(x^{202}\) for sufficiently large \(x\), uniformly in the
prime. Taking all the integral character summands and using
faithful scalar extension gives \eqref{eq:exp142}. Finally
\eqref{eq:alltw142}, \(p^2\le x^{398}\), and \(t-1\asymp
x^{401}\) give the asserted proportion. The interval and field
choices are precisely those of Proposition~\ref{prop:short107}.
\end{proof}

\begin{corollary}[The nonexceptional part of the strict cohomology sum]
\leavevmode\label{cor:full142}
On these same intervals,
\[
 \sum_{\substack{2<p\le x^{199}\\p\ \mathrm{prime},\ p\nmid n}}
    \ \sum_{\substack{1\le r<t\\r\notin\mathcal E_p}}
                   \log|M_{p,r}|=o\bigl(\exp(x^\eps/2)\bigr).
\]
At each included pair the remaining summand in
\eqref{eq:tail114} is zero.
\end{corollary}
\begin{proof}
By the elementary-divisor calculation in
Theorem~\ref{thm:deep114}, each logarithm on the left is at
most \(\lceil x^{202}\rceil\log|\mathcal S_{p,r}(E)[p]|\).
Lemma~\ref{lem:rank114} bounds the latter logarithm by
\(Cnp\log(2Np)\). Bounding the sum of the primes by the sum
of all positive integers at most \(x^{199}\), the result is
\(O_{\psi,\eps}(x^{202}nt\,x^{398}\log(2N x^{199}))\).
This is polynomial in \(x\), which proves the claimed estimate
for fixed \(\eps>0\). Moreover \(x^{199}<\exp(x^\eps/12)\)
eventually. The lower bound for \(k_{p,r}\) at these primes in
Theorem~\ref{thm:deep114} therefore exceeds \(\lceil x^{202}\rceil\).
As \(Q_{p,r}\subseteq M_{p,r}\), this gives
\(p^{k_{p,r}}Q_{p,r}=0\), as asserted.
\end{proof}

A bound on the number of exceptions leaves their individual
orders undetermined. Within this prime range, the sum over
\(\mathcal E_p\) and the odd primes dividing \(n\) therefore
still requires an estimate. Primes above \(x^{199}\) also occur
in \eqref{eq:length12}. The height-restricted norm identity and
the compact row types from Theorems~\ref{thm:low136} and
\ref{thm:copies137} apply on exactly the fields treated here.

\subsection{All odd primes in the polynomial range}\label{nav:II-33}
Work on the interval of Proposition~\ref{prop:short107}.
Use its field \(E=K_\psi C_{l'}(\sqrt{-q},\sqrt{-d_0})\) and
retain \(n=[E:\QQ]\), \(N\), \(t=\lfloor x^{401}\rfloor\) and
\(M_{p,r}=\tfrac12(1+(-1)^rc)\mathcal S_{p,r}(E)\).
The next estimate includes the odd prime divisors of \(n\).

\begin{theorem}[Simultaneous control of the complementary groups]
\label{thm:every144}
For each odd prime \(p\le x^{199}\) there is a set
\(\mathcal E_p\subset\{1,\ldots,t-1\}\) with
\[
    |\mathcal E_p|\ll n p^2\,\frac{\log(2Np)}{\log p},
   \qquad p^{\lceil x^{202}\rceil} M_{p,r}=0
                           \quad (1\le r<t,\ r\notin\mathcal E_p).
\]
The implied constant is absolute, and \(x\) is sufficiently
large in the common intervals. In particular the exceptional
proportion is \(O_{\psi,\eps}(x^{-1+33\eps/4})\) uniformly in
these primes. For the smaller range
\(p\le Y=x^{133-3\eps}\) there is a single set
\(\mathcal E\subset\{1,\ldots,t-1\}\) with
\begin{equation}
  \frac{|\mathcal E|}{t-1}\ll_{\psi,\eps}x^{-3\eps/4},
  \qquad
  p^{\lceil x^{202}\rceil}M_{p,r}=0
   \quad\text{for every odd prime }p\le Y,\ r\notin\mathcal E .
  \label{eq:joint144}
\end{equation}
\end{theorem}
\begin{proof}
Write \(\operatorname{Gal}(E/\QQ)=P\times H\) with \(P\) its
\(p\)-Sylow subgroup of order \(d\). Over the unramified
coefficient ring containing the characters of \(H\), the
idempotents for \(H\) are integral because \(p\nmid |H|\).
Complex conjugation lies in \(H\). The components of
\(M_{p,r}\) are therefore exactly those indexed by
\(\alpha\in\widehat H\) with \(\alpha(-1)=(-1)^r\).
Restriction identifies each component with the corresponding
strict group over \(E^H\) by
Lemma \ref{lem:des28}.
This descent keeps the entire action of \(P\) on that component.

For each character and each residue class of \(r\bmod(p-1)\),
Corollary \ref{cor:family143} bounds its exponent outside at most
\(2(V_{\alpha,h}+1)\) twists in that class. Here
\(h=r+1+i\bmod(p-1)\) if the factor of \(\alpha\) at \(p\)
is \(\omega^i\). As in Theorem~\ref{thm:alltw142}, fixing \(h\)
fixes the moment index \(k\le p\), and the parity condition
on the prime-to-\(p\) character is constant on each rational
conjugacy orbit \(\mathcal O\). Taking the union of the
exceptional twists and summing equation \ref{eq:orbit143} over the orbits
and the at most \(p-1\) values of \(h\) proves the asserted
cardinality, since \(\sum_{\mathcal O}d|\mathcal O|=n\).
The estimate also covers characters for which the twisted
coefficient becomes trivial modulo \(p\) on \(G_{E^H}\).

For each character separately, \eqref{eq:orbit143} gives
\(V_{\alpha,h}+1\ll n p\log(2Np)/\log p\).
The two additional scalar factors in the exponent have
total valuation at most \(2(1+v_p(d)+\log_p t)\), by
Lemma \ref{lem:loss143}. By Corollary~\ref{cor:family143}, the power of \(p\)
required for every component outside the union thus has
exponent at most
\[
 C n p\,\frac{\log(2Np)}{\log p}
                  \left(1+\frac{\log t}{\log p}\right)
        +2(1+\log_p n+\log_p t).
\]
Here \(C\) is absolute. On the stated fields we have
\(n\ll_\psi x^{2+8\eps}\) and
\(\log(2Np)\ll_{\psi,\eps}x^{\eps/4}\) for \(p\le x^{199}\).
Also \(\log t=O_\eps(x^{\eps/4})\), so the dominant term is
\(O_{\psi,\eps}(x^{201+17\eps/2})\). Since
\(0<\eps<1/52\), the whole exponent is at most \(x^{202}\)
eventually. Faithful scalar extension yields the bound on
\(M_{p,r}\) itself. Dividing the cardinality bound by
\(t-1\asymp x^{401}\) gives the first exceptional proportion.

For the joint assertion take
\(\mathcal E=\bigcup_{2<p\le Y}\mathcal E_p\), the union
being over primes. Bounding \(\sum_{p\le Y}p^2\) by the sum
of squares of the integers up to \(Y\) gives
\[
       |\mathcal E|\ll_{\psi,\eps}
            x^{2+8\eps+\eps/4}\,Y^3
               =x^{401-3\eps/4}.
\]
Every remaining twist then satisfies the exponent bound for
all the indicated primes simultaneously.
\end{proof}

\begin{corollary}[The remaining prime and twist pairs]\label{cor:sum144}
On these fields and intervals,
\[
 \sum_{\substack{2<p\le x^{199}\\p\ {\rm prime}}}
 \ \sum_{\substack{1\le r<t\\r\notin\mathcal E_p}}
                      \log|M_{p,r}|=o(\exp(x^\eps/2)).
\]
In \eqref{eq:tail114} the sum over odd primes may be restricted
to \(p>x^{199}\) and to \(p\le x^{199}\) with \(r\in\mathcal E_p\).
\end{corollary}
\begin{proof}
At an included pair the elementary-divisor inequality bounds
the logarithmic order by
\(\lceil x^{202}\rceil\log|\mathcal S_{p,r}(E)[p]|\).
Lemma~\ref{lem:rank114} bounds the first-layer logarithm by
\(Cnp\log(2Np)\) for every odd prime. Summing as in
Corollary~\ref{cor:full142} gives a polynomial bound in \(x\),
hence the stated little-oh estimate for every fixed \(\eps>0\).
For \(p\le x^{199}\) the assigned depth \(k_{p,r}\) of
Theorem~\ref{thm:deep114} exceeds \(\lceil x^{202}\rceil\)
eventually. Thus \(p^{k_{p,r}}Q_{p,r}=0\) at each included
pair, which proves the restriction on the sum.
\end{proof}

This reduction leaves the orders at the exceptional twists and
at larger primes to be bounded. The arithmetic normalization
and the untruncated period comparison of Section~\ref{sec:missing}
on \(E\) are also required to deduce the prime-sum estimate
\eqref{eq:window} by this argument.

\subsection{All twists on the quadratic invariants of the period fields}\label{nav:II-34}
Fix \(\psi\) of conductor \(F\), put \(o=[K_\psi:\QQ]\), and
take \(0<\eps<1/52\). For a number field \(L\) in the abelian
family and an odd prime \(p\), write
\[
 M_{p,r}(L)=\tfrac12(1+(-1)^r c)\mathcal S_{p,r}(L).
\]
Here \(c\) is complex conjugation, and \(\mathcal S_{p,r}(L)\)
is the strict degree-one group with coefficients
\(\QQ_p/\ZZ_p(-r)\). A strict class vanishes on inertia away
from \(p\) and locally at each place over \(p\). All logarithms
of finite-group orders below include the full primary exponent.

\begin{theorem}[The quadratic invariants at every prescribed twist]
\label{thm:invar148}
There is a constant \(c_{\psi,\eps}>0\) and there are arbitrarily
distant intervals \([V,V^{3/2}]\) with fixed auxiliary primes
\(q,d_0\) on each interval. For every \(x\) in such an interval
use precisely the period field of Proposition~\ref{prop:short107},
\[
  J=l',\qquad B_J=K_\psi C_J,\qquad
  E=B_J(\sqrt{-q},\sqrt{-d_0}),\qquad H_0=\Gal(E/B_J).
\]
Then
\begin{equation}
 \sum_{\substack{2<p\le c_{\psi,\eps}(\log\log x)^{1/5}\\
                        p\ {\rm prime}}}
       \quad\sum_{1\le r<\lfloor x^{401}\rfloor}
         \log|M_{p,r}(E)^{H_0}|\ \le\ \exp(x^{2\eps/3}).
 \label{eq:invar148}
\end{equation}
On these very fields and intervals, \eqref{eq:low136} holds for
the height-truncated theta-period seeds of
Theorem~\ref{thm:low136}.
\end{theorem}
\begin{proof}
Choose the intervals of
Theorem \ref{thm:log147}.
There the auxiliary prime pair is chosen after the exceptional
twists on the fields \(B_j\) have been located, and the resulting
interval still satisfies the hypotheses of
Proposition~\ref{prop:short107} and Theorem~\ref{thm:low136}.
Neither \(q\) nor \(d_0\) divides \(2F\). Each of the three
quadratic subfields of \(\QQ(\sqrt{-q},\sqrt{-d_0})\) ramifies
at one or both of these primes, whereas \(B_J\) is unramified
at both. Thus \(H_0\) has order four. Averaging over \(H_0\)
is exact on \(p\)-primary modules. Inflation--restriction on
absolute Galois groups identifies degree-one cohomology over
\(B_J\) with the invariants over \(E\). Restriction on local
degree-one groups at \(p\) is injective by corestriction, as
the local degree divides four. Away from \(p\) the coefficient
is unramified and an inertia homomorphism which vanishes
upstairs factors through a group of order dividing four.
Its \(p\)-primary image is zero. The strict local conditions
therefore descend. These identifications commute with \(c\).
Hence each group in \eqref{eq:invar148} is isomorphic to
\(M_{p,r}(B_J)\). Theorem \ref{thm:log147} gives the asserted
sum uniformly in \(J\), in particular at \(J=l'(x)\).
\end{proof}

For the primes in \eqref{eq:invar148} this also shortens
the remaining sum in \eqref{eq:tail114}, including its terms
at the exceptional twists of Theorem~\ref{thm:every144}.
Indeed put \(e_0=|H_0|^{-1}\sum_{h\in H_0} h\). It commutes
with the projector \(\pi_{p,r}\) there since the field is
abelian. The logarithm of the order of
\(p^{k_{p,r}}e_0 Q_{p,r}\) is at most
\(\log|M_{p,r}(E)^{H_0}|\). The direct sum decomposition for
\(e_0\) therefore permits replacing \(Q_{p,r}\) in these terms
by \((1-e_0)Q_{p,r}\) with total error at most
\(\exp(x^{2\eps/3})=o(b^{1/2})\), where \(b=\exp(x^\eps)\).
This leaves the other three characters of \(H_0\) in these
terms.

\subsection{Strict local conditions at these very theta primes}\label{nav:II-35}
Let \(E\) be the period field, \(d=[E:\QQ]\),
\(t=\lfloor x^{401}\rfloor\), and \(R_0\) the squarefree
product of the primes dividing \(2Ng_1g_2T\). For
\(p\) prime, \(r\ge1\), use the actual coefficient
\(A_{p,r}=(\QQ_p/\ZZ_p)(-r)\) of Section~\ref{sec:missing}.

\begin{theorem}[The local cost on the interval of the integral comparison]
\label{thm:both184}
On the same fields and intervals as Theorems~\ref{thm:invar148}
and~\ref{thm:last182}, the two auxiliary prime ideals for
the theta construction can be chosen above distinct rational
primes \(\ell_1,\ell_2\) split completely in \(E\), larger
than \(t+2\) and every prime-power cutoff in the period
count, and with logarithms polynomial in \(x\).
Put \(R=R_0\ell_1\ell_2\) and let \(W\) be all places of
\(E\) over its divisors. In the strict group \(\mathcal
S_{p,r}(E)\), relax the conditions at \(W\) by allowing
all local classes away from \(p\) and unramified classes
above \(p\). Denote the resulting group by
\(\mathcal S^W_{p,r}(E)\). Then
\begin{equation}
  \sum_{1\le r<t}\sum_p
   \log|\mathcal S^W_{p,r}(E)/\mathcal S_{p,r}(E)|
        \ll d t^2\log(2Rdt)\ll_{\psi,\eps} x^C=o(b^{1/2}).
                                                         \label{eq:both184}
\end{equation}
The same bound holds for the increase of the sum of the
orders in logarithmic form on the right of \eqref{eq:tail114}
under this relaxation. The final compact integral norm
\eqref{eq:last182} holds with these choices.
\end{theorem}
\begin{proof}
For each required narrow class use the construction of
Theorem~\ref{thm:common178} with one additional cyclotomic
factor. Take an odd prime \(l\) outside the original rational
primes to avoid and with logarithm at most the logarithm of
one plus their product, by taking a prime divisor of that
integer. Choose \(a\) least positive with \(l^a>Y+1\),
where \(Y\) exceeds the cutoffs and \(t+2\). One may
choose \(Y=O(t+xb)\) with a sufficiently large implied
constant since the numerators of the count are \(O(xb)\)
and all weights are at least one. The extra extension
\(E^+(\zeta_{l^a})/E^+\) is totally ramified at \(l\),
where the prior auxiliary class field is unramified. Thus
its Frobenius can be prescribed as \(-1\bmod l^a\)
independently of the given narrow class. The same
unconditional least-prime theorem gives a degree-one prime
of norm exceeding \(Y\) and with logarithm polynomial in
\(x\), by the abelian conductor bound. Excluding its rational
prime for the next class gives distinct primes. Both are
unramified and split completely in \(E\).

For a finite place \(w\nmid p\), put \(q=Nw\). Unramified
cohomology is zero because the nonzero multiplier
\(q^{-r}-1\) acts surjectively on \(\QQ_p/\ZZ_p\). Inertia
restriction identifies the full local cohomology with the
Frobenius invariants in the maps from the pro-\(p\) tame
inertia \(\ZZ_p(1)\) to \(A_{p,r}\). Its order is therefore
\(p^{v_p(q^{r+1}-1)}\). At \(w\mid p\), the allowable
unramified group is a quotient of the inertia invariants.
The cyclotomic inertia image has index dividing the absolute
ramification index \(e_w\) in \(\ZZ_p^\times\), so these
invariants have logarithmic order at most
\((2+v_p(r)+v_p(e_w))\log p\). At either new rational
prime, when it is the coefficient prime, the completion is
\(\QQ_p\) and \(p-1>r\). A primitive Teichm\"uller unit
then acts nontrivially modulo \(p\) and the invariants vanish.

The quotient of the global groups injects into these
local groups by restriction. Its away-from-\(p\) contributions
sum to at most \(d(t-1)(t+2)\log R/2\) by factoring
\(q^{r+1}-1\). At the residue characteristics the total
for a fixed twist is at most \(d(2\log R+\log r+\log d)\),
since every \(e_w\) divides \(d\). This proves \eqref{eq:both184}.
The original and the relaxed groups are Galois-stable,
and the quotient of their images under the sign and
group-ring operators in \eqref{eq:tail114} is a quotient
of the increase already bounded. The prime choices preserve
the ideals and local properties used in
Theorems~\ref{thm:units180} and~\ref{thm:last182}. Full
proofs are in Section \ref{nav:II-21}.
\end{proof}

Every original strict class already vanishes on restriction
to every place in \(W\). Thus \eqref{eq:both184} measures
the cohomological cost of the finite theta test places.
The orders of the classes in the original strict kernel
on the right of \eqref{eq:tail114} still require a uniform
bound. In the integral period, \eqref{eq:last182} identifies
the last coefficient exactly. Its nonvanishing and height
after the preceding lifts, finite class sums and the
prescribed projection are the other outstanding estimates
for the interval bound.

\subsection{An enlarged range of full orders on the same period field}\label{nav:II-36}
The exponent \(1/5\) from the two auxiliary conductors in
Theorem~\ref{thm:tower104} measures their size at the initial
scale of prime selection. The final interval can start later
while retaining those same primes. This also permits the
depth of the subfield in the real cyclotomic tower to depend
on the twist.

\begin{theorem}[Removing a growing family of complete summands]
\label{thm:region198}
Fix \(0<\delta<1\). On arbitrarily distant admissible intervals
take \(E=K_\psi C_J(\sqrt{-q},\sqrt{-d_0})\) as in
\eqref{eq:tail114}. For each \(1\le r<t\) let \(j(r)\) be the
largest \(0\le j\le J\) with
\[
             [2^j(r+2)]^5\le x^{(1-\delta)\eps}.
\]
Let \(e_r'\) project to the classes invariant over the
subfield \(K_\psi C_{j(r)}(\sqrt{-q},\sqrt{-d_0})\), and put
\(e_r'=0\) if there is no such \(j\).
The contribution of these summands, at every odd prime and
with full primary exponents, is \(o(b^a)\) for every fixed
\(a>0\). Using the cell count of de la Torre Dur\'an, the
same assertion holds with the larger condition
\(2^j(r+2)^2\le x^{(1-\delta)\eps}\).
In particular, for either choice of projector,
\begin{equation}
 \mathcal L_E(x)=\sum_{1\le r<t}\sum_{p\ {\rm odd\ prime}}
       \log\big|p^{k_{p,r}}(1-e_r')Q_{p,r}\big|+o(b^{1/2}).
 \label{eq:region198}
\end{equation}
\end{theorem}
\begin{proof}
Choose a fixed \(A\) such that \(1/(5A)<\delta/2\). At an initial
scale \(Y\) the two auxiliary primes have
\(\log(qd_0)\ll Y^{\eps/5}\log^2(2Y)\). In selecting the common
interval reserve the fixed additional factor \(A\) in the
upper power before choosing these primes. The finitely many
dyadic exceptional indices can then be avoided or placed at
small twists on \([V,V^{3/2}]\) with \(V\ge Y^A\). The earlier
selection over the fields \(K_\psi C_j\) is independent of
the two primes and supplies room for the larger fixed power.
The complete selection is proved in
Lemma \ref{lem:delay197}.
All the conductors in the subfields above now have logarithm
\(O_{\psi,\eps,\delta}(x^{\delta\eps/2})\).

At odd primes the average over the two-power degree extension
to \(E\) identifies the invariant strict group with the strict
group of the chosen subfield. Its degree is at most
\(4\operatorname{ord}(\psi)2^j\). The cell estimate for its
even \(K\)-group bounds the sum of logarithmic strict orders
over all primes. With the displayed choice of \(j(r)\), the sum
over all \(r<t\) is at most
\(\exp(C_{\psi,\eps,\delta}x^{(1-\delta/2)\eps})\), using
the appropriate lattice bound. This is
Theorem \ref{thm:region197}.
Its ratio to \(b^a\) tends to zero for every fixed \(a>0\).
All the projectors in \eqref{eq:tail114} commute with \(e_r'\),
since \(E/\QQ\) is abelian. For each \(p,r\), the logarithmic
order removed by the additional projector is exactly
\(\log|p^{k_{p,r}}e_r'Q_{p,r}|\), bounded by the logarithmic
order of the full invariant strict group. This proves
\eqref{eq:region198} by addition.
\end{proof}

For example, in the second range any fixed \(u,v>0\) with
\(u+2v<1\) are permitted simultaneously for
\(2^j\le x^{u\eps}\) and \(r\le x^{v\eps}\), by taking
\(0<\delta<1-u-2v\) to absorb the constant from \(r+2\).
In the published lattice range the corresponding condition is
\(5(u+v)<1\). Both statements apply to subfields containing
the two auxiliary quadratic fields on the chosen interval.

\begin{corollary}[Zero contribution above the large-prime cutoff]
\label{cor:cut198}
On every sufficiently distant admissible interval put
\(P=\exp(x^\eps/12)\). At \(p>P\), the entire summand of
\(\mathcal S_{p,r}(E)\) descending from \(K_\psi C_j\) vanishes
when
\[
        2^j(r+2)\le c_{\psi,\eps}
            (\log x/\log\log x)^{1/5},\qquad 0\le j\le J .
\]
There is no restriction on the order of the residual cyclotomic
character in this assertion.
\end{corollary}
\begin{proof}
At odd primes restriction through the extension of degree
at most two from \(K_\psi C_j\) to \(K_\psi C_j(i)\) is injective.
This last field has modulus at most \(2^{j+2}F\), independent
of \(q,d_0\). Its bound on the total logarithmic order is
at most \(x^{\eps/2}\) under the indicated hypothesis for a
sufficiently small \(c_{\psi,\eps}\). A single nonzero class at
\(p>P\) costs more than \(x^\eps/12\), which exceeds that bound
eventually. The explicit estimates and the wider range using the
cell count of de la Torre Dur\'an are in
Corollary \ref{cor:cut197}.
\end{proof}

The full period field has \(2^J\ge
x^{2+8\eps}/\operatorname{ord}(\psi)\), and the twist range extends
to \(\lfloor x^{401}\rfloor-1\). In particular the field-degree
bound in Theorem~\ref{thm:region198} does not reach the entire
field even at the first twist. The sum on the right of
\eqref{eq:region198} and the arithmetic lower estimate for the
projected integral theta period at its original finite level
remain to be bounded together on the common intervals.

\subsection{One exceptional set on a whole period interval}\label{nav:II-37}
At a single field the count in Theorem~\ref{thm:every144} sums
separate upper bounds for the different coefficient primes.
For each positive weight \(k\ge2\), the nonzero values
\(-L(1-k,\bar\alpha)\) are algebraic numbers fixed by the
choice of field and character. Factoring their numerator ideals
lets us count the primes using a single integer for each \(k\).
We first give the resulting estimates for the actual complementary
strict groups \(M_{p,r}(E)\) from \eqref{eq:region198}.

\begin{theorem}[Simultaneous bounds on the full period fields]
\label{thm:all202}
Write \(t(x)=\lfloor x^{401}\rfloor\). On arbitrarily distant
common intervals \([V,V^{3/2}]\) in Theorem~\ref{thm:region198},
the following two assertions hold.
\begin{enumerate}
\item At any \(x\) in the interval there is a set
\(\mathcal E(x)\subseteq\{1,\ldots,t(x)-1\}\) of size
\(O_{\psi,\eps}(x^{401-3\eps/2})\). For every \(1\le r<t(x)\) outside
this set and every odd prime \(p\le x^{399/2-5\eps}\),
\[
                   p^{\lceil x^{203}\rceil}M_{p,r}(E_x)=0 .
\]
\item A single set \(\mathcal E_*\subseteq
\{1,\ldots,t(V^{3/2})-1\}\) of size
\(O_{\psi,\eps}(V^{401-3\eps/2})\) works throughout the interval
for the fixed cutoff \(Y_*=V^{199-7\eps}\). For every
\(x\in[V,V^{3/2}]\), every \(1\le r<t(x)\) outside
\(\mathcal E_*\), and every odd prime \(p\le Y_*\),
\[
                    p^{\lceil V^{203}\rceil}M_{p,r}(E_x)=0 .
\]
\end{enumerate}
Both statements include the primes dividing the field degree
or conductor. In each statement the sum of the full logarithmic
orders over the indicated pairs is \(o(\exp(x^\eps/2))\), uniformly
on the intervals.
\end{theorem}

\begin{proof}
We record the finite counting estimate first. Let the abelian field
have degree \(n\), modulus \(N\), and twist cutoff \(t\). For
\(p\nmid nN\), take \(k\in[2,p]\) with \(k\equiv-r\bmod(p-1)\).
Let \(J_k(E)\) be the product, over rational conjugacy orbits of
characters \(\alpha\) of parity \((-1)^k\), of the absolute norms
of the numerator ideals of \(B_{k,\bar\alpha}/k\) in the character
value fields. Each conjugate contributes its positive \(p\)-adic
valuation to \(v_p(J_k(E))\). At each nontrivial residual character
the finite Stickelberger annihilator has a bounded \(p\)-adic
interpolation whose value at \(1-k\) has this character's own
valuation \(v\). Interpolation at \(2v+1\) distinct twists shows
that at most \(2v\) twists in its class escape the exponent bound
\((1+2\lfloor\log_p t\rfloor)v\). The case \(v=0\) gives
vanishing on the whole residue class. Summing over the primes
uses \(\sum_p v_p(J_k(E))\log p=\log J_k(E)\) only once at each
weight. The interpolation, the numerator formula, and the bound
\(\log J_k(E)\ll nk\log(2Nk)\) are proved in
Lemma \ref{lem:shared201} and Theorems \ref{thm:integer201}--\ref{thm:shared201}.

At an odd \(p\mid nN\) use
Corollary \ref{cor:family143}
on the \(p\)-Sylow subgroup and its prime-to-\(p\) complement.
Summing the orbit bound \eqref{eq:orbit143} over at most \(p-1\) residue classes
gives at most \(Cnp^2\log(2Np)/\log p\) exceptional twists.
The exponent outside these twists is bounded as in the proof of
Theorem~\ref{thm:every144}. Consequently for any \(Y\ge3\) the
union \(\mathcal E\) for all odd \(p\le Y\) satisfies
\[
 |\mathcal E|\ll (1+\omega(nN))\,nY^2\log(2NY),
 \qquad
 p^{\lceil L\rceil}M_{p,r}(E)=0
\]
for \(2<p\le Y\) prime and \(1\le r<t\) outside \(\mathcal E\).
Here \(\omega(m)\) counts the distinct prime factors of \(m\) and
\[
                 L=CnY\log(2NY)(1+\log(2t)).
\]
Indeed the unramified sum of weighted counts uses \eqref{eq:count201} in
Volume~V, and each of the at most \(\omega(nN)\) other primes
costs at most \(CnY^2\log(2NY)\). Dividing an individual weighted
count by \(\log p>1\) can only reduce it. The estimates for \(L\)
include the scalar losses at the primes dividing \(n\), bounded
by \(2(1+\log_p n+\log_p t)\) in Theorem~\ref{thm:every144}.

On our fields \(n\ll_\psi x^{2+8\eps}\), and the odd prime factors
of \(nN\) divide the fixed conductor or order of \(\psi\), or are
one of \(q,d_0\). In particular \(\omega(nN)=O_\psi(1)\).
Take the intervals of
Lemma \ref{lem:delay197},
with the fixed initial delay required in Theorem~\ref{thm:region198}.
They give \(\log(2N x^{401})\ll_{\psi,\eps}x^{\eps/2}\);
the same bound holds with \(V\) on the right at the largest
level on \([V,V^{3/2}]\). For the latter assertion the quadratics
are fixed on the interval, and the dyadic modulus is polynomial in
\(V\). More explicitly their log conductor is
\(O_{\psi,\eps}(Y_0^{\eps/5}\log^2(2Y_0))\) at the initial
scale \(Y_0\le V\). Substituting \(Y=x^{399/2-5\eps}\) gives
\[
 nY^2\log(2NY)\ll_{\psi,\eps}
                 x^{2+8\eps+399-10\eps+\eps/2}
                    =x^{401-3\eps/2}.
\]
Since \(\log(2t)\ll_\eps x^{\eps/2}\) and \(\eps<1/52\), one
also has \(L<x^{203}\) eventually. This proves the first assertion.

For the second, let \(E_*\) be the field at the largest dyadic
level used anywhere on the interval, with the same \(K_\psi\)
and the same two quadratic fields as before. Its degree is
\(O_\psi(V^{3+12\eps})\). Every \(E_x\) embeds in \(E_*\) with
relative degree a power of two. Use the finite counting estimate
only once, for \(E_*\), \(t=t(V^{3/2})\), and \(Y=Y_*\).
The exponent for the number of exceptions is now
\[
             3+12\eps+2(199-7\eps)+\eps/2=401-3\eps/2,
\]
and for \(L\) the exponent is at most \(202+6\eps<203\).
Restriction carries strict classes to strict classes and commutes
with complex conjugation. Corestriction after restriction is
multiplication by \([E_*:E_x]\), a unit at every odd \(p\). Hence
the bound on \(M_{p,r}(E_*)\) holds on \(M_{p,r}(E_x)\) for
every \(x\), with exactly the chosen set of twists.

Finally a finite abelian \(p\)-group killed by \(p^a\) has
logarithmic order at most \(a\) times the logarithmic order of
its subgroup killed by \(p\), by its elementary divisors.
Lemma~\ref{lem:rank114} bounds the latter here by
\(Cnp\log(2Np)\). In both ranges \(a,n,p,t\), the number of
primes, and \(\log(2Np)\) are bounded by fixed powers of \(x\).
Their product is \(o(\exp(x^\eps/2))\), uniformly as \(V\)
tends to infinity.
\end{proof}

\begin{corollary}[Divisibility required of a surviving term]
\label{cor:weights202}
Write \(\mathcal D_{p,r}=p^{k_{p,r}}(1-e'_r)Q_{p,r}\) for the group on
the right side of \eqref{eq:region198}. All the pairs covered by
either exponent bound of Theorem~\ref{thm:all202} have
\(\mathcal D_{p,r}=0\). Moreover fix \(A>0\). Uniformly for sufficiently
large \(x\) on the same intervals, if \(p\nmid 2nN\) and
\(\mathcal D_{p,r}\ne0\), then
\[
         p^{1+v_p(r+k)}\mid J_k(E_x)
 \quad\text{for every integer }2\le k\le x^A
                    \text{ with }p-1\mid r+k .
\]
\end{corollary}

\begin{proof}
On the polynomial prime ranges of the theorem,
Theorem~\ref{thm:deep114} gives \(k_{p,r}\ge
\lfloor x^{-420}\allowbreak\exp(x^\eps/3)\rfloor>\lceil x^{203}\rceil\)
eventually. Each group before multiplication by \(p^{k_{p,r}}\)
is a summand of \(M_{p,r}(E_x)\). This proves the first assertion.
For the second, any \(k\) in the stated congruence gives
\(p\le t(x)+x^A+1\), so the same lower bound on \(k_{p,r}\)
applies. Also
\(w=v_p(J_k(E_x))\ll nk\log(2Nk)/\log p\) is polynomially
bounded in \(x\), hence eventually smaller than \(k_{p,r}\).
If \(w<1+v_p(r+k)\), Theorem \ref{thm:integer201}, gives
\(p^wM_{p,r}(E_x)=0\), contradicting \(\mathcal D_{p,r}\ne0\).
\end{proof}

For example, at \(E=\QQ\) the numerator for \(k=2\) is one.
The comparison therefore kills the rational \(p\)-primary
summand at every positive even twist \(r\equiv-2\bmod(p-1)\).
The six rational residue classes and an imaginary quadratic
example are computed in
Corollary \ref{cor:shared201}.
In \eqref{eq:region198} a single surviving group can have large
order. The exceptional-set count gives no bound for its full
order, and if \(p-1>r+x^A\) the divisibility test has no weight
in its stated range. The estimate for the total of those remaining
orders and the arithmetic lower bound for the designated Bessel
pairing in Corollary~\ref{cor:original200} are still needed on
the same interval.

\subsection{Full orders at primes dividing the degree or conductor}\label{nav:II-39}
Keep the fields \(E_x\), the enclosing field \(E_*\), and the integers
\(J_k(E_*)\) of Theorem~\ref{thm:shareddiv230}. There are two reasons
to track the character values integrally. The field conductor can
be divisible by the coefficient prime, and at a divisor of the
degree the character projectors can have denominators. The two
losses are measured differently in the following result.

\begin{theorem}[Full strict orders on the classes detected by a Bernoulli value]
\leavevmode\label{thm:full231}
Fix \(A>0\). On the common intervals of Theorem~\ref{thm:all202}
let \(1\le r<t(x)\), let \(p\) be any odd prime, and put
\(o=[K_\psi:\QQ]\). Say that the pair \((p,r)\) satisfies the
following condition if for some integer \(2\le k\le V^A\) one has
\[
    p-1\mid r+k, \qquad w=v_p(J_k(E_*))<1+v_p(r+k).
\]
Write \(M=M_{p,r}(E_x)\) and \(e_0=|H_0|^{-1}\sum_{h\in H_0}h\),
where \(H_0=\operatorname{Gal}(E_x/K_\psi C_J)\) has order four.
For all sufficiently large \(V\), every such pair satisfies
\[
 \begin{cases}
      p^wM=0, &p\nmid o,\\
      p^{w+v_p(o)}(1-e_0)M=0, &p\mid o.
 \end{cases}
\]
Uniformly for \(x\in[V,V^{3/2}]\), the sum of the full
\(\log|M_{p,r}(E_x)|\) over all pairs satisfying the condition
is \(o(\exp(x^\eps/2))\), counting each pair once.
\end{theorem}
\begin{proof}
The index \([E_x:K_\psi]\) is a power of two, so \(v_p([E_x:\QQ])=
v_p(o)\). In the comparison of Lemma~\ref{lem:fourier231},
use all characters if \(p\nmid o\), and project by \(1-e_0\)
if \(p\mid o\). For such a divisor of the fixed integer \(o\),
we have \(p\ne q,d_0\) for sufficiently large \(V\).
The inertia elements at \(q,d_0\) generate \(H_0\).
Indeed \(K_\psi C_J\) is unramified at both primes, and each
of the two independent quadratic adjunctions ramifies at its
own prime only among this pair. A character receiving a nonzero
image from \(1-e_0\) is nontrivial on one of these two involutions.
Its reduction still takes the value \(-1\) there, and the
cyclotomic character is unramified there. Thus
\(\bar\alpha^{-1}\bar\kappa_p^{-r}\ne1\).
If \(p\nmid o\), a character with trivial residual coefficient
has a zero strict group by Theorem~\ref{thm:all24}.
In this case the character values can be taken in an unramified
extension, since their orders are prime to \(p\).

For every other participating character write \(z=B_{k,\alpha^{-1}}/k\).
Inflating the character along \(\operatorname{Gal}(E_*/\QQ)\to
\operatorname{Gal}(E_x/\QQ)\) preserves its primitive character
and its value \(z\). Equation \eqref{eq:norm201} therefore gives
\(v_p(z)\le w\). Theorem~\ref{thm:scalar231}
also gives \(v_p(z)\ge0\), identifies it with \(v_p(F(X_{1-k}))\),
and supplies the annihilator \(F(X_{r+1})\).
An integral power series has
\(F(u)-F(v)\in(u-v)O\) when \(u,v\in pO\), by subtracting
monomials and passing to the limit. The valuation of
\(X_{r+1}-X_{1-k}\) is \(1+v_p(r+k)\) for \(X_s=(1+p)^s-1\).
It follows that \(v_p(F(X_{r+1}))=v_p(z)\le w\).
Thus \(p^w\) kills this scalar group even in a ramified coefficient
ring, since the scalar divides \(p^w\) there. The integral
comparison proves both exponent assertions.

For the sum, a subgroup killed by \(p^a\) has logarithmic order
at most \(a\log|M[p]|\), by its cyclic decomposition.
Lemma~\ref{lem:rank114} bounds the latter logarithm by
\(C[E_x:\QQ]p\log(2N_*p)\). All primes occurring here satisfy
\(p\le t_*+V^A+1\), and all admissible \(w\) are at most
\(1+\log_3(t_*+V^A+1)\). The number of pairs is at most
\(t_*(t_*+V^A+1)\). These bounds, \([E_x:\QQ]\), and
\(\log(2N_*)\) are polynomial in \(V\) for fixed \(A,\psi,\eps\).
This accounts for all the stated exponent bounds simultaneously.
For the finitely many primes dividing \(o\), the remaining direct
summand is \(e_0M\). Theorem~\ref{thm:invar148} bounds the sum of
its full logarithmic orders over these primes and \emph{all} twists
by \(\exp(x^{2\eps/3})\). The intervals of Lemma~\ref{lem:delay197}
retain precisely the growing odd-prime bound used there.
Both the polynomial and this last bound are
\(o(\exp(x^\eps/2))\), uniformly as claimed.
\end{proof}

\begin{theorem}[A common divisor for all odd coefficient primes]
\label{thm:div231}
Define \(\mathcal B_p\) as in Theorem~\ref{thm:shareddiv230}, now for
each odd prime. For every fixed \(A>0\) and all sufficiently large
\(V\), simultaneously for every \(2\le k\le V^A\),
\[
   \prod_{p\ {\rm odd\ prime}}
    \ \prod_{\substack{r\in\mathcal B_p\\p-1\mid r+k}}
                p^{1+v_p(r+k)}\ \bigm|\ J_k(E_*).
\]
The twists in \(\mathcal B_p\) are counted once each on the entire
interval, even when the corresponding strict groups survive at
several fields.
\end{theorem}
\begin{proof}
Fix \(p,k\) with a factor in the product. For each participating
twist choose a witnessing \(x\), and put \(d=k_{p,r}(x)\).
As in Theorem~\ref{thm:shareddiv230}, \(p\le t_*+V^A+1\) and
\(d\ge\lfloor x^{-420}\allowbreak\exp(x^\eps/3)\rfloor\). A surviving group
gives an element of \(M_{p,r}(E_x)\) of order at least \(p^{d+1}\).
If \(p\mid o\), the invariant summand satisfies
\(\log|e_0 M_{p,r}(E_x)|\le\exp(x^{2\eps/3})<d\log p\)
by Theorem~\ref{thm:invar148}. It is therefore already killed
by \(p^d\). The component of the element in \((1-e_0)M_{p,r}(E_x)\)
still has order at least \(p^{d+1}\).

The first comparison map of Lemma~\ref{lem:fourier231} sends this
element, after extending scalars, to a tuple with some coordinate
of order at least \(p^{d+1-v_p(o)}\). Indeed its composite with
the return map is multiplication by \([E_x:\QQ]\). The map
\(m\mapsto m\otimes1\) preserves orders since \(1\) belongs to
a \(\ZZ_p\)-basis of \(O\). The chosen scalar character has
nontrivial residual coefficient by the proof of
Theorem~\ref{thm:full231}. Regard it as a character of \(E_*\).
The scalar group over \(\QQ\), its integral annihilator, and its
negative special value depend only on this character, \(p\) and
the indicated twist or residue class. Thus witnesses on different
fields can be assigned to the same character.

For each fixed character occurring in this assignment, use its
single series \(F\) from Theorem~\ref{thm:scalar231}. An annihilator
of an element of additive order at least \(p^{d+1-v_p(o)}\)
has valuation strictly greater than \(d-v_p(o)\), including
when \(O\) is ramified. At the nodes \(X_{r+1}\) these lower
bounds exceed \(t_*(1+\log_3(t_*+V^A+1))\), uniformly for
large \(V\). Any two different nodes in this set enlarged by
\(X_{1-k}\) have difference of valuation at most
\(1+\log_3(t_*+V^A+1)\).
Lemma~\ref{lem:samples230} applies also in ramified \(O\) because
division by a monic linear polynomial at a point of \(pO\)
has integral quotient, and its Lagrange interpolation inequality
uses only the normalized valuation. Consequently \(v_p(F(X_{1-k}))\)
is at least the sum of \(1+v_p(r+k)\) over all the twists assigned
to this character. It equals \(v_p(B_{k,\alpha^{-1}}/k)\) by the
unit formula of Theorem~\ref{thm:scalar231}. Summing over the characters
in equation \eqref{eq:norm201} and multiplying over the
finitely many participating primes proves the assertion.
\end{proof}

For a pair failing the strict inequality of Theorem~\ref{thm:full231},
the higher finite differences in Proposition
\ref{prop:digits231}
can still bound the full scalar group. An example is given in its
Corollary~\ref{cor:quartic231}.
A uniform bound for the surviving full orders in \eqref{eq:region198}
and the comparison with the designated integral theta period
remain necessary for Main Theorem~\ref{thm:grh16}.

\subsection{Endpoint bounds on the same fields as the period}\label{nav:II-40}
\begin{theorem}[Full invariant orders up to a logarithmic endpoint]
\label{thm:endpoint232}
Use the polynomial modulus of Theorem~\ref{thm:conductor232}.
For every \(a>0\) there is \(c_a>0\), depending also on
\(\psi,\varepsilon\), for which the following holds.
At each \(1\le r<t(x)\) take the largest \(0\le j\le J\)
satisfying
\[
 [2^j(r+2)]^5\le c_a\,\frac{x^\varepsilon}{\log(2x)}.
\]
Let \(e(r)\) project to the classes coming from
\(L_j=K_\psi C_j(\sqrt{-q},\sqrt{-d_0})\), with \(e(r)=0\)
when the index does not exist. Then, uniformly on the interval,
\[
   \sum_{r<t(x),\ r\ge1}\ \sum_{p\ {\rm odd\ prime}}
         \log|e(r)\mathcal S_{p,r}(E_x)|
                 =o\big(\exp(a x^\varepsilon)\big).
\]
Under the cell estimate in the second part of
Theorem~\ref{thm:cells103}, one may replace the left side of the
index condition by \(2^j(r+2)^2\). Both conclusions include the
full primary orders on both conjugation signs.
\end{theorem}
\begin{proof}
The extension \(E_x/L_j\) has two-power degree. At each odd prime
restriction identifies the strict group of \(L_j\) with the
invariants in \(E_x\), by the global and local descent of
Theorem~\ref{thm:region197}. In either cell bound use
\(d=[L_j:\QQ]\le 4o\,2^j\), \(o=[K_\psi:\QQ]\), and the
modulus \(N_x\le x^D\). Then
\(\log(2N_x d(r+2))\le C_1\log(2x)\), uniformly for all
\(r<x^{401}\). In the first bound the logarithm is multiplied
by at most \(C_2[2^j(r+2)]^5\). In the second it is multiplied
by at most \(C_2 2^j(r+2)^2\). Choose \(c_a\) to make either
exponent at most \(a x^\varepsilon/2\).
Summing the resulting bounds on logarithmic orders costs at
most \(x^{401}\), and \(x^{401}\exp(-a x^\varepsilon/2)\to0\).
\end{proof}

\begin{corollary}[One exceptional set at the enlarged coefficient cutoff]
\label{cor:cut232}
For every fixed \(B>0\) put
\[
 Y_* = \frac{V^{199-6\varepsilon}}
                    {(\log(2V))^{(B+1)/2}} .
\]
There is one set of \(O_{\psi,\varepsilon,B}(V^{401}/(\log(2V))^B)\)
twists for the whole interval of Theorem~\ref{thm:conductor232}.
For every \(x\in[V,V^{3/2}]\), every \(1\le r<t(x)\) outside this
set and every odd prime \(p\le Y_*\), the full complementary group
\(M_{p,r}(E_x)\) is killed by \(p^{\lceil V^{203}\rceil}\).
The sum of its logarithmic orders over these prime--twist pairs
is \(o(\exp(x^\varepsilon/2))\).
\end{corollary}
\begin{proof}
All \(E_x\) embed with two-power index in the largest field
\(E_*\). Its degree satisfies \(n_*\ll_\psi V^{3+12\varepsilon}\),
its modulus \(N_*\le V^{3D/2}\), and \(\omega(n_*N_*)\ll_\psi1\),
since only the two auxiliary primes and fixed prime divisors of
\(2oF\) can occur. Apply the finite counting estimate in the proof
of Theorem~\ref{thm:all202} at this one field with
\(t_*=\lfloor V^{1203/2}\rfloor\). Its bound on the number of
exceptional twists is \(O_\psi(n_*Y_*^2\log(2N_*Y_*))\).
The exponent of \(V\) here is \(3+12\varepsilon+
2(199-6\varepsilon)=401\), and the logarithm contributes at
most one of the \(B+1\) powers in the denominator.
The bound on the exponent of the group is
\(p^{\lceil L\rceil}\) with
\(L\ll n_*Y_*\log(2N_*Y_*)(1+\log(2t_*))\).
Since \(202+6\varepsilon<203\), this gives the asserted exponent.
Restriction to \(E_*\) is injective at odd \(p\), so the same
list works for every \(E_x\). Lemma~\ref{lem:rank114} bounds
\(\log|M_{p,r}(E_x)[p]|\) by \(C n_*p\log(2N_*p)\).
Multiplying by the exponent in powers of \(p\) and summing costs
only a polynomial in \(V\). This proves the order estimate.
\end{proof}

\section{Residue data and exceptional twists}\label{nav:g-II-3}

\subsection{Several twists in the same Bernoulli numerator}\label{nav:II-38}
Corollary~\ref{cor:weights202} tests one surviving twist at a time.
On a whole interval one can charge different twists to the same
numerator, retaining their multiplicities. We first give the
interpolation fact that permits this addition.

\begin{lemma}[Simultaneous divisibility of an interpolated value]
\label{lem:samples230}
Let \(O\) be the integers of a finite unramified extension of
\(\QQ_p\), with \(v(p)=1\), and let \(F(X)\in O[[X]]\).
Take distinct \(a_0,a_1,\ldots,a_m\in pO\). If, for each \(i\ge1\),
\[
 v(F(a_i))\ \ge\ v(a_0-a_i)+
                 \sum_{\substack{1\le j\le m\\j\ne i}}v(a_i-a_j),
\]
then \(v(F(a_0))\ge\sum_{i=1}^m v(a_0-a_i)\).
\end{lemma}
\begin{proof}
For \(a\in pO\), the difference \(F(X)-F(a)\) is divisible
by \(X-a\) in \(O[[X]]\). Indeed for a monomial of degree \(n\)
the quotient is \(\sum_{j=0}^{n-1}X^ja^{n-1-j}\), and the
sum defining each coefficient converges in \(O\). Successive
divisions give \(F=R+H\prod_{i=1}^m(X-a_i)\), with
\(H\in O[[X]]\), \(R\in O[X]\) and \(\deg R<m\).
For \(m=0\) the assertion is integrality. Otherwise
interpolate \(R\) at \(a_1,\ldots,a_m\). At \(a_0\) the term
with coefficient \(R(a_i)=F(a_i)\) has valuation
\[
 v(F(a_i))+
  \sum_{j\ne i,\,1\le j\le m}
       \bigl(v(a_0-a_j)-v(a_i-a_j)\bigr)
 \ \ge\ \sum_{j=1}^m v(a_0-a_j).
\]
The term involving \(H(a_0)\) satisfies the same bound.
\end{proof}

\begin{theorem}[One divisor for all surviving twists on an interval]
\label{thm:shareddiv230}
Use the fields and intervals of Theorem~\ref{thm:all202}. Let
\(E_*\) be the largest field used on \([V,V^{3/2}]\), of degree
\(n_*\) and modulus \(N_*\), and put
\(t_*=t(V^{3/2})=\lfloor V^{1203/2}\rfloor\).
For each odd prime \(p\nmid n_*N_*\), let \(\mathcal B_p\)
consist of the integers \(r\) for which some \(x\) on this
interval has \(1\le r<t(x)\) and \(\mathcal D_{p,r}(E_x)\ne0\),
with \(\mathcal D_{p,r}\) as in Corollary~\ref{cor:weights202}.
Fix \(A>0\). For sufficiently large \(V\), simultaneously for
every integer \(2\le k\le V^A\), one has the integer divisibility
\[
 \prod_{p\nmid 2n_*N_*}
 \ \prod_{\substack{r\in\mathcal B_p\\p-1\mid r+k}}
       p^{\,1+v_p(r+k)}\ \bigm|\ J_k(E_*).
\]
The outer product is over primes. The integer on the right is
the product of the absolute norms of the numerator ideals of
\(B_{k,\bar\alpha}/k\), one per rational orbit of characters
of \(E_*\) of parity \((-1)^k\).
\end{theorem}
\begin{proof}
Fix \(k,p\) with a factor on the left. Then
\(p\le t_*+V^A+1\). At every witnessing \(x\), restriction to
\(E_*\) is injective on the strict group \(M_{p,r}\). The local
conditions and conjugation sign commute with restriction, and
the composite of restriction with corestriction is
multiplication by the power of two \([E_*:E_x]\).
Nonvanishing of the stated image of \(M_{p,r}(E_x)\) implies
that \(M_{p,r}(E_*)\) has an element of order at least
\(p^{k_{p,r}(x)+1}\). Theorem~\ref{thm:deep114} gives
\(k_{p,r}(x)\ge\lfloor x^{-420}\allowbreak\exp(x^\eps/3)\rfloor\),
since this polynomial prime bound is eventually below
\(\exp(x^\eps/12)\) throughout the interval. This lower
bound exceeds \(t_*(1+\log_3(t_*+V^A+1))\), uniformly in
the chosen \(k,p,r,x\).

Extend scalars to an unramified coefficient ring \(O\) containing
all character values of \(\operatorname{Gal}(E_*/\QQ)\).
The idempotents are integral because \(p\nmid n_*\). For each
such \(r\) choose a character \(\alpha\) whose component has
an element of order at least \(p^{k_{p,r}(x)+1}\). Such a
component exists by faithful flatness and the direct sum
decomposition. It has \(\alpha(c)=(-1)^k\) and a nontrivial
residual coefficient character by
Lemma \ref{lem:shared201}.
Group the chosen twists by \(\alpha\).

For one such group take the bounded measure of that lemma with
\(\chi=\bar\alpha\) and \(h=1-k\bmod(p-1)\). Its auxiliary
integer can be fixed for the entire group since the residual
character is fixed. The integral \(A_b(s)\) annihilates the
component at \(s=r+1\), and at \(s=1-k\) it equals a
\(p\)-adic unit times \(B_{k,\bar\alpha}/k\). To apply the
lemma above, put \(X_s=(1+p)^s-1\). There is one
\(F\in O[[X]]\) with \(A_b(s)=F(X_s)\) on \(\ZZ_p\).
Indeed \(\langle a\rangle=(1+p)^{u(a)}\) with
\(u(a)\in\ZZ_p\), since \(p\) is odd. In the integral of
Lemma~\ref{lem:shared201}, expand the factor
\(\langle a\rangle^{-s}\) as
\(\sum_{j\ge0}\binom{-u(a)}j X_s^j\).
Each binomial coefficient is integral by continuity from
nonnegative integers. The remaining character factors are
units. Termwise integration is valid by uniform convergence
on \(pO\), and all coefficients of the resulting series are
integral.

Take \(a_0=X_{1-k}\) and take the other nodes to be \(X_{r+1}\)
for the distinct twists of this group. For distinct integers
\(u,v\) one has \(v_p(X_u-X_v)=1+v_p(u-v)\). In fact
\((1+p)^h-1\) has valuation one for \(p\nmid h\), and taking
the \(p\)-th power increases the valuation by one by binomial
expansion. Inversion handles negative differences. There are
fewer than \(t_*\) nodes besides \(a_0\), and the integer
differences have size at most \(t_*+V^A+1\).
The order of the chosen element forces the annihilator at
each of these nodes to have valuation at least
\(k_{p,r}(x)+1\). The hypotheses of
Lemma~\ref{lem:samples230} therefore hold. It gives
\[
 v_p(B_{k,\bar\alpha}/k)\ \ge       \sum_{r\text{ assigned to }\alpha}(1+v_p(r+k)).
\]
By the factorization of numerator ideals, the sum of the
positive valuations over all embedded characters of the given
parity is exactly \(v_p(J_k(E_*))\). This is equation~\eqref{eq:norm201}. Summing the last inequality and then multiplying
over primes gives the claim. The bound on \(p\) makes the
product finite.
\end{proof}

For example, at \(k=2\) the sum of
\((1+v_p(r+2))\log p\) over the participating prime--twist
pairs on the entire interval is at most \(\log J_2(E_*)\),
which is \(O(n_*\log(2N_*))\) by Theorem~\ref{thm:integer201}.
A pair with \(p-1>r+V^A\) occurs in none of these products.
For pairs that do occur, bounding the full logarithmic orders
in \eqref{eq:region198} still requires control of their
exponents beyond the sampled depth.

\subsection{Frobenius orbits in the exceptional twist count}\label{nav:II-41}
Work on one polynomial-conductor interval of Theorem~\ref{thm:conductor232}.
Write \(E_*\) for its largest field, \(N_*\) for its modulus, and
\(t_*=\lfloor V^{1203/2}\rfloor\). In particular \(\log(2N_*)\ll_{\psi,\eps}
\log(2V)\). For a character \(\alpha\) of \(\operatorname{Gal}(E_*/\QQ)\)
put \(j(\alpha)=v_2(\operatorname{ord}\alpha)\). Write \(J_{k,j}\) for
the portion of the integer \(J_k(E_*)\) in Theorem~\ref{thm:integer201}, obtained by using only rational character
orbits of index \(j(\alpha)=j\). The parity condition
\(\alpha(-1)=(-1)^k\) is retained.

\begin{lemma}[A numerator for a fixed dyadic character order]
\label{lem:orbits233}
For \(k\ge2\) one has
\(\log J_{k,j}\ll_\psi 2^j k\log(2N_* k)\). If \(p\) is an odd
prime away from the field degree and modulus and
\(w(p)=\max(v_2(p-1),v_2(p+1))\), the orbit of such a character under
\(\operatorname{Gal}(\overline{\QQ}_p/\QQ_p)\) has size at least
\(\max(1,2^{j-w(p)})\).
\end{lemma}
\begin{proof}
The characters of the compositum are products of characters of its
factors. Indeed a character of a subgroup of a finite abelian group
extends to the group, by taking a root of its prescribed value when
adjoining each generator. The factors other than the real cyclotomic
tower give at most \(4\operatorname{ord}(\psi)\) choices, all of
bounded order. The tower factor is cyclic of two-power order.
Once \(j\) exceeds the two-adic exponent of the orders of the other
factors, the tower character must have exact order \(2^j\). For
smaller \(j\) its order is bounded in terms of \(\psi\). Thus in
either case at most \(C_\psi 2^j\) characters occur. For each rational
orbit the numerator ideal of \(B_{k,\bar\alpha}/k\) has log norm at
most \([\QQ(\alpha):\QQ]\,Ck\log(2N_*k)\), by the denominator and
complex bounds in the proof of Theorem~\ref{thm:integer201}.
The sum of these degrees counts the characters in the layer. This
proves the first assertion.

The values generate the roots of unity of order \(\operatorname{ord}\alpha\),
prime to \(p\). Frobenius raises them to the \(p\)-th power. If its
orbit has size \(h\), then \(2^j\mid p^h-1\). For odd \(h\) the
two-adic valuation is \(v_2(p-1)\). For even \(h\) it is
\(w(p)+v_2(h)\), by factoring successive differences of squares and
using \(v_2(p^2+1)=1\). Both cases give the asserted lower bound.
\end{proof}

\begin{theorem}[A common exceptional set for almost all coefficient primes up to a square-root cutoff]
\label{thm:cutoff233}
Fix \(B>0\) and \(a>1\). Put \(L=\log(2V)\),
\(W=\lfloor a\log_2 L\rfloor\), where \(\log_2\) denotes logarithm
to base two, and set
\[
  Y_0=V^{199-6\eps}/L^{(B+1)/2},\qquad
  Y=V^{401/2}/L^{(B+a+1)/2}.
\]
There is a single set of \(O_{\psi,\eps,a,B}(V^{401}/L^B)\) twists
for the whole interval. For every \(x\in[V,V^{3/2}]\) and every
\(1\le r<t(x)\) outside this set, the full group \(M_{p,r}(E_x)\)
is killed by \(p^{\lceil V^{203}\rceil}\) at every odd prime
\(p\le Y_0\), and also at every \(Y_0<p\le Y\) with
\(p\not\equiv\pm1\pmod{2^{W+1}}\). The sum of its logarithmic
orders over these pairs is \(o(\exp(x^\eps/2))\), uniformly in \(x\).
At most \(O(Y/L^a+1)\) primes up to \(Y\) fall in the two excluded
residue classes.
\end{theorem}
\begin{proof}
Start with the exceptional set of Corollary~\ref{cor:cut232} for
\(p\le Y_0\). Consider an additional prime \(Y_0<p\le Y\) with
\(w(p)\le W\). It is prime to \([E_*:\QQ]N_*\) for large \(V\).
The primes dividing the degree divide \(2\operatorname{ord}(\psi)\).
Of those dividing the modulus, only \(q,d_0\) vary, and both exceed
\(V^{401}\) by the selection in Theorem~\ref{thm:conductor232}.
Also \(\log p\asymp\log V\) uniformly in this range.

Call \((p,r)\) exceptional here if some \(x\) on the interval has
\(r<t(x)\) and \(M_{p,r}(E_x)\) has an element of additive order
greater than \(p^{\lceil V^{202}\rceil}\). Restriction to \(E_*\)
preserves its order since corestriction composed with restriction is
multiplication by a power of two. Over the unramified coefficient ring
containing the character values, the integral idempotents split the
strict group. Some scalar component has an element of order exceeding
the same threshold. Its residual coefficient character is nontrivial
by Lemma~\ref{lem:shared201}. Every local Frobenius
conjugate has an element of the same order, since conjugating
coefficients sends the cochains, local restrictions and cocycles to
those for the conjugate character. Assign \(r\) to every character of
one such orbit, once for this \((p,r)\), and put \(k\in[2,p]\) for
the unique integer congruent to \(-r\bmod(p-1)\).

Fix \(p,k,\alpha\) in this assignment, with \(j=j(\alpha)\), and set
\(h=v_p(B_{k,\bar\alpha}/k)\), with \(v_p(p)=1\). This is a
nonnegative integer by the integral interpolation formula and
unramifiedness. Every member of the Frobenius orbit contributes the
same \(h\) to \(v_p(J_{k,j})\), by the norm factorization in equation~\eqref{eq:norm201}, restricted to this layer. Hence the lemma
and the bound on \(N_*\) give
\[
  h\le\frac{\log J_{k,j}}{\max(1,2^{j-W})\log p}
    \ll_{\psi,\eps} 2^W Y.
\]
In particular \(h+1=O(V^{201})\) for sufficiently large \(V\).
There can be at most \(h\) distinct twists assigned to this character
at this \(k\). To see this, assume there are \(h+1\) and use their
nodes \(X_{r+1}=(1+p)^{r+1}-1\) and the comparison node \(X_{1-k}\).
Lemma~\ref{lem:shared201}, and the power-series construction
in Theorem~\ref{thm:shareddiv230} give one integral series for this
whole branch. Its value at the comparison node has valuation \(h\).
At every other node its value annihilates a scalar element of the
stated order, hence has valuation greater than \(V^{202}\). The
valuation of a difference of distinct nodes is at most
\(1+\log(t_*+Y+1)/\log p=O(1)\), by the binomial valuation formula.
Lemma~\ref{lem:samples230} applies since the sum of at most \(h+1\)
such valuations is less than \(V^{202}\). It forces the comparison
value to have valuation at least \(h+1\), a contradiction.

For fixed \(p,k,j\), summing over characters therefore counts the
assigned twists with their orbit sizes and gives at most \(v_p(J_{k,j})\).
Each pair assigned to layer \(j\) consumes at least
\(\max(1,2^{j-W})\) of this count. Multiply by \(\log p\), sum over
the primes and use integer factorization. The number of additional
pairs is bounded by
\[
 \frac{C}{\log V}\sum_j\frac{1}{\max(1,2^{j-W})}
           \sum_{2\le k\le Y}\log J_{k,j}
 \ll_{\psi,\eps}(1+\log V)2^W Y^2.
\]
Indeed there are \(O_\psi(1+\log V)\) layers by the degree bound,
and \(\log(2N_*Y)=O_{\psi,\eps}(\log V)\). The definition of \(Y\)
makes the last expression \(O(V^{401}/L^B)\). Adding the twist from
each such pair to the first set proves the uniform exponent bound.
By Lemma~\ref{lem:rank114} the logarithm of the subgroup killed by
\(p\) is \(O([E_*:\QQ]p\log(2N_*p))\). Multiplication by the exponent
in powers of \(p\) bounds the full logarithmic order by a polynomial
in \(V\), as does the number of pairs. This gives the claimed uniform
little-oh. Each omitted residue class contains at most
\(1+Y/2^{W+1}\) integers in the range. Since \(a>1\), these primes
form a vanishing fraction of the primes up to \(Y\), by the elementary
lower bound \(\pi(Y)\gg Y/\log Y\).
\end{proof}

For every pair covered by the theorem the depth in
Theorem~\ref{thm:deep114} exceeds \(\lceil V^{203}\rceil\) for large
\(V\). Its term in \eqref{eq:region198} is therefore zero. An
exceptional twist may still carry a large order.

\subsection{Common coefficient primes for the cohomology and the Gram minors}\label{nav:II-42}
\begin{theorem}[Placing the auxiliary norm primes beyond the coefficient range]
\label{thm:common233}
Fix \(A>0\) and an interval of Theorems~\ref{thm:conductor232}
and~\ref{thm:cutoff233}. For each distinct period field in the interval
the two split representatives \(\mathfrak p,\mathfrak q\) allowed by
Lemma~\ref{lem:good175}, can be chosen above distinct
rational primes exceeding \(\exp(x^A)\) for every \(x\) in the
interval. Their logarithmic norms are still bounded by a fixed power
of \(x\). The single additional prime in Theorem~\ref{thm:norm233}, can be placed above the same range, with
its polynomial logarithmic bound for polynomial test depths.

Use these choices in Corollary~\ref{cor:gram233}. At
every rational coefficient prime \(l\le\exp(x^A)\), every place
\(\mathfrak l\mid l\) of its coefficient field containing \(F=E_x^+\),
and \(v=\mathfrak l\cap F\), the correction for a nonzero minor is
\[
 v_{\mathfrak l}(\det G_{I,J})-v_{\mathfrak l}(\det G'_{I,J})
  =\frac{e(\mathfrak l/v)}2\,v(\mathfrak d_F\mathfrak a_0^{-1})
     \sum_\tau\left(\sum_{c\in I}D_{\tau,c}+\sum_{c\in J}D_{\tau,c}\right).
\]
It is nonnegative and is zero when \(l\nmid 2F_\psi q\), where
\(F_\psi\) is the fixed character conductor. In particular, up to
the cutoff \(Y\) of Theorem~\ref{thm:cutoff233}, every \(l\nmid
2F_\psi\) has identical Gram-minor valuations. On the primes and
twists covered by that theorem the corresponding cohomology
orders have its uniform bound on these same fields.
\end{theorem}
\begin{proof}
We give the prime selection. Let \(K/\QQ\) be one of the fields
\(F,E_x\), let \(d=[K:\QQ]\), and fix a class in a narrow or
ordinary ray class field of \(K\), of finite modulus \(\mathfrak m\).
Take an odd rational prime \(s\) unramified in \(K\) and prime to
the modulus. Every completion of \(K\) at \(s\) is unramified over
\(\QQ_s\), whereas \(\QQ_s(\zeta_{s^b})\) is totally ramified of
degree \(\varphi(s^b)\). Thus \(K(\zeta_{s^b})/K\) has this degree
both globally and on inertia at any such place. Every nontrivial
subextension ramifies there. The ray class field is unramified
there, so the two fields intersect in \(K\). Prescribe the desired
ray class in the first and the identity in the second. For a lower
cutoff \(H\ge2\), take \(b\) least with \(s^b>H\). The abelian
least-prime theorem of Thorner and Zaman gives a prime of absolute
prime norm \(l\) in this class. Its cyclotomic Frobenius sends a
root of unity to its \(l\)-th power, so \(l\equiv1\pmod{s^b}\)
and \(l>H\). Its bound in logarithmic form is
\[
   \log l\ll 1+\log |D_K|+\log N_{K/\QQ}\mathfrak m
                         +d\log(sH)+d\log(2d).
\]
Indeed the largest conductor norm in this abelian compositum is at
most \((N_{K/\QQ}\mathfrak m)s^{bd}\) by the norm calculation in
Theorem~\ref{thm:norm233}. This argument works with
the infinite real places included in the class modulus.

All original rational levels, including every terminal prime on
the interval, have logarithms bounded by fixed powers of \(V\),
by Theorem~\ref{thm:conductor232} and Lemma~\ref{lem:size97}. Start with \(H=\exp(V^{B_0})\), choosing
the fixed \(B_0\) so that \(H\) exceeds these levels and
\(\exp(x^A)\) throughout. Over \(F\) apply the construction with finite modulus
one to the narrow classes of \(\mathfrak d_F\) and \(\mathfrak a_0\),
increasing \(H\) to the first prime before selecting the second.
These classes have trivial character in \(E_x/F\) by Lemma~\ref{lem:good175}, so their representatives are split. The same
choice works whenever this period field occurs on the interval.
For the last prime use the class of the ideal denoted by \(A\) in
Theorem~\ref{thm:norm233}, now over \(E_x\), and first
increase \(H\) to exceed every rational prime in the specified test
support. The number of rational prime divisors of the period
modulus and the prescribed levels is bounded in terms of \(\psi\).
Thus \(s\) can in each case be chosen from a bounded list of primes.
The degree, log discriminant, and the logarithm of every cutoff
used here are polynomially bounded for the indicated depths, by
the displayed estimate successively. The avoidance of all primes
under the modulus also ensures the ray prime is coprime to it.

Finally \((u)=\mathfrak d_F\mathfrak a_0^{-1}/(\mathfrak p\mathfrak q)\).
Its numerator is Galois invariant: the different is invariant, and
\(\mathfrak a_0^2=(d_0)\) uniquely determines the ideal
\(\mathfrak a_0\). All conjugates of the denominator are away
from \(l\) in the asserted range. Consequently
\(v_{\mathfrak l}(\tau u)=e(\mathfrak l/v)\,v(\mathfrak d_F
\mathfrak a_0^{-1})\) for every embedding \(\tau\), using the
Galois group of the totally real field to identify its embeddings.
Substitute this constant into the Gram-minor formula. Integrality
and the support above \(2F_\psi q\) were proved in Theorem~\ref{thm:scale232}. Since \(q>V^{401}>Y\), only the fixed
divisors of \(2F_\psi\) can contribute below \(Y\). Selecting
the three split ideals changes no number field or strict local
condition in the proof of Theorem~\ref{thm:cutoff233}.
\end{proof}

The comparison concerns the minors of the projected vectors. A
lower bound for the designated signed period itself and the orders
at the remaining exceptional twists are still needed for
Main Theorem~\ref{thm:grh16}.

\subsection{A test on the common interval and the prime 37}\label{nav:II-43}
\begin{theorem}[Testing the whole interval at a fixed conductor]
\leavevmode\label{thm:twists235}
Use a polynomial-conductor interval \([V,V^{3/2}]\) of
Theorem~\ref{thm:conductor232}, with largest field \(E_*\),
degree \(n_*\), and modulus \(N_*\). Fix \(B>0\) and set
\(H=\lceil V^{203}\rceil\). For each odd \(p\le V^B\) with
\(p\nmid n_*N_*\), and each primitive character \(\alpha\) of
the Galois group of \(E_*\), form the polynomial of
Theorem~\ref{thm:newton235} on each nontrivial residual branch
of the appropriate parity. One polynomial tests all positive
twists in that branch in every field of the interval on which
\(\alpha\) factors. At \(r\), call its component detected when
\(P(X_{r+1})\not\equiv0\bmod p^{H+v_p(H!)}O\).

If every participating character of \(E_x\) at \((p,r)\) is
detected or has trivial residual coefficient, the whole strict
group \(M_{p,r}(E_x)\) is killed by \(p^{H+v_p(H!)}\). Its
contribution to \eqref{eq:region198} is zero. The sum of the
full logarithmic orders of the groups satisfying this test is
\(o(\exp(x^\eps/2))\) uniformly for \(x\in[V,V^{3/2}]\).
For a pair at which only some components are detected, each
such component separately contributes zero to the remaining
group after multiplication by \(p^{k_{p,r}}\).
\end{theorem}
\begin{proof}
Take \(2\le k\le p\) congruent to \(-r\bmod(p-1)\).
Inflating a character to the largest field leaves its primitive
conductor and all Bernoulli samples unchanged. Its scalar group
is the strict group for the same coefficient character over
\(\QQ\). Theorem~\ref{thm:newton235} gives the asserted
annihilation for each detected component. The trivial residual
case has zero strict group by Lemma~\ref{lem:shared201}.
All idempotents are integral here since \(p\nmid n_*\).
They give a direct sum after extension to the unramified
coefficient ring, and faithful flatness descends the exponent
when all components pass. The two-power projectors and
quotients in \eqref{eq:region198} preserve this decomposition.

For \(p\le V^B\), Theorem~\ref{thm:deep114} gives
\(k_{p,r}\ge H+v_p(H!)\) throughout the interval eventually,
since its lower bounds in both polynomial-prime ranges dominate
every fixed power of \(x\). Thus the asserted terms disappear.
Lemma~\ref{lem:rank114} bounds \(\log|M_{p,r}(E_x)[p]|\) by
\(C n_*p\log(2N_*p)\). The cyclic decomposition bounds the
full logarithmic order by this times \(H+v_p(H!)\) on the
detected pairs. The number of pairs is at most
\(V^B\lfloor V^{1203/2}\rfloor\), and \(n_*,N_*\) are
polynomially bounded in \(V\). The resulting polynomial is
\(o(\exp(x^\eps/2))\) uniformly. The test itself uses
\(H\) weights at most \(pH\) on each branch, each at the
original conductor dividing \(N_*\).
\end{proof}

\begin{proposition}[An irregular branch with one exceptional class at each depth]
\leavevmode\label{prop:example235}
For the rational character at \(p=37\), take \(r\equiv4\bmod36\)
and \(b=2\). There exists \(\xi\in\ZZ_{37}\) such that the
annihilator from Theorem~\ref{thm:scalar231} has valuation
\(1+v_{37}(r-\xi)\) on this branch, with \(v_{37}(0)=+\infty\).
In particular \(37\) kills the strict scalar group unless
\(r\equiv1048\bmod1332\), and \(37^2\) kills it unless
\(r\equiv11704\bmod49284\).
\end{proposition}
\begin{proof}
Put \(a_0=38^{-31}-1\), \(U=(X-a_0)/37\), and
\(G(U)=F(a_0+37U)/37\). The residual character \(\omega^{32}\)
is nontrivial and \(2^{32}\not\equiv1\bmod37\).
Write \(u_s=(38^s-1-a_0)/37\).
For reproducibility the first three samples are displayed
modulo \(37^2=1369\):
\[
\begin{array}{c|rrr}
 k_i&32&68&104\\ \hline
 u_{1-k_i}&0&186&409\\
 F(X_{1-k_i})/37&451&376&745 .
\end{array}
\]
The first row of residues is obtained by exponentiating \(38\)
modulo \(37^3\). For the second, compute \(B_0=1\) and
\(B_j=-(j+1)^{-1}\sum_{i<j}\binom{j+1}{i}B_i\).
For instance \(B_{32}=-7709321041217/510\).
Substituting the three values in
\(2^{1-k_i}(2^{k_i}-1)(1-37^{k_i-1})B_{k_i}/k_i\)
gives \(16687,13912,27565\) modulo \(37^3\), respectively.

All coefficients of \(G\) are integral, because its constant
term is integral by the table and its coefficient of degree
\(j\ge1\) is divisible by \(37^{j-1}\). Modulo \(37\) it is
affine, and its values at \(0,1\) are \(7,6\). Thus its
reduction is \(7-U\). The divided difference
\((G(u)-G(v))/(u-v)\) is a unit for all distinct \(u,v\in
\ZZ_{37}\), since its reduction is \(-1\). This makes \(G\)
a bijection at each residue depth, and hence an isometry onto
\(\ZZ_{37}\). The distance formula in Theorem~\ref{thm:newton235}
likewise makes \(r\mapsto u_{r+1}\) an isometry of \(\ZZ_{37}\)
onto itself. Completeness gives a unique zero \(\xi\) of their
composite. Its distance from any \(r\) gives the claimed valuation.

Here \(u_{r+1}\equiv r+32\bmod37\), so the first exceptional
class is \(r\equiv12\bmod37\), or \(1048\bmod1332\) on the
stated branch. To lift once, \(u_{1049}\equiv673\bmod1369\).
Modulo \(1369\), \(G\) has degree at most two and is determined
by the table, whose three nodes have distinct reductions.
Quadratic interpolation gives \(G(673)\equiv1073=29\cdot37\).
Increasing \(r=1048\) by \(1332t=36\cdot37t\) changes \(u_{r+1}\)
by \(36\cdot37t\bmod1369\), by the binomial expansion,
and changes \(G\) by \(37t\bmod1369\), since its derivative
is \(-1\bmod37\). Thus \(t\equiv8\bmod37\) gives the next
exceptional class. This proves both exponent statements. The
valuation computed here belongs to an annihilator; at the
exceptional twists the order of the group still needs a bound.
\end{proof}

\subsection{One set of Frobenius parameters on the whole interval}\label{nav:II-44}
\begin{corollary}[Common parameters on the support of strict cohomology]
\leavevmode\label{cor:interval238}
Fix one sufficiently distant interval \([Y,Y^{3/2}]\) of
Theorem~\ref{thm:common16}, with its two chosen auxiliary primes.
For all \(x\) on it consider the large coefficient primes
\(\ell>\exp(x^\eps/12)\), \(1\le r<t(x)\), and the complementary
character summands of strict cohomology to which the ordinary
rank-three construction applies. Keep those with nonzero strict group.
For each resulting marked invariant ring \(D_i\) let \(d_i\) be the
dimension of its special fiber. Put \(d_*=\max_i d_i\), taking zero
for the empty list. There are \(d_*\) distinct rational primes splitting
completely in every \(E_x\) and in \(\QQ(i)\), unramified for every
ring in the list, with the following property. For each \(i\), its first
\(d_i\) Frobenius trace differences make \(D_i\) a finite module
over the power series ring in those \(d_i\) parameters over \(O_i\).
The same list can be chosen within any set of these split primes of
positive upper natural density.
\end{corollary}
\begin{proof}
The two auxiliary primes are fixed on this interval. The integer
index of the real cyclotomic layer and the integer \(t(x)\) each take
only finitely many values. For each resulting field \(E\) and each
\(r\), the dual of the odd-primary strict group embeds prime by prime
in the corresponding primary part of the finite group
\(K_{2r}(\mathcal O_E)\), by Theorem~\ref{thm:wild98}.
Thus there are only finitely many \(\ell\) for which a strict summand
can be nonzero on the whole interval taken as a union. For each such
\((E,r,\ell)\) the character list is finite. Fix the coefficient ring
and a marking once per datum, using the same choices whenever the
field and the data recur. This gives a finite list of ordinary rings.
The compositum of the finitely many fields \(E_x\) and \(\QQ(i)\) is
finite and Galois over \(\QQ\). Its completely split primes have
positive natural density by Chebotarev. Deleting the finitely many
ramified and coefficient primes preserves that density. Apply
Theorem~\ref{thm:params238} to these primes, or to the indicated
subset, for all the rings at once. The entries of the list are
auxiliary observations of Frobenius at unramified primes. The local
conditions on strict cohomology and the level defining the theta
period are those for the given \(E_x\) throughout.
\end{proof}

For instance, if on every characteristic-\(\ell\) curve through one
of these rings the traces at its selected primes satisfied monic
polynomial equations over the residue field, their images on the
curve would be constant. Finiteness over the selected parameters
would then force every element on the curve to be algebraic over the
residue field, hence constant. Proposition~\ref{prop:curve71} would
give \(d_i=0\). Establishing such arithmetic control on the original
points, and estimating their integral congruences to the diagonal,
are further steps in the bound for \eqref{eq:region198}.

\subsection{One finite field for the large-prime strict orders}\label{nav:II-54}
Fix a sufficiently distant common interval \([V,V^{3/2}]\)
of Theorem~\ref{thm:conductor232}. Its largest period field is
\(E_*\), of degree \(n_*\) and modulus \(N_*\), and put
\(t_*=\lfloor V^{1203/2}\rfloor\). Let \(\mathcal P\) be the
coefficient primes occurring in a nonzero complementary strict
group at some \(x\) on this interval with
\(\ell>\exp(x^\varepsilon/12)\). Theorem~\ref{thm:wild98}
gives finite support, since the fields form a finite nested
list and \(r<t_*\). At every \(\ell\in\mathcal P\) take the
list of \((r,\alpha)\) from these large-prime data. Repeated
pairs are included once and one representative suffices from
each orbit under residue-field Frobenius. One may include all
the other complementary characters and twists of \(E_*\)
with \(1\le r<t_*\) as well. Write \(d_{\ell,r,\alpha}\) for
the full \(O\)-length of the corresponding strict group,
and \(b_{\ell,\alpha}=[\mathbb F_\ell(\bar\alpha):\mathbb F_\ell]\).
Every \(\ell\) in the list exceeds \(2t_*+3\) and \(N_*\)
when \(V\) is large, since \(N_*\) is polynomial in \(V\).

\begin{theorem}[Simultaneous recovery of the full large-prime lengths]\leavevmode
\label{thm:commonlength267}
For any integers \(h_\ell\ge1\) there is a positive-density set
of primes \(s\), split in the original period fields, for which
\[
 w_{h_\ell}(D_{s,h_\ell}(\alpha,r))
       =\min(h_\ell,d_{\ell,r,\alpha})
       \qquad\text{for every listed triple.}
\]
The primes can also satisfy every splitting and power-residue
condition of Theorem~\ref{thm:detect241}. In the set of primes
split in the common field \(K\) defined below, a subset realizing
the equalities has relative natural density exactly
\[
 \prod_{\ell,r,[\alpha]}(1-\ell^{-b_{\ell,\alpha}})
       \ \ge\ (2V)^{-C_{\psi,\varepsilon} n_*t_*}.
\]
Here the product is over the distinct residual orbits in the
list. In particular one finite field realizes \emph{all}
large-prime orders throughout this interval if each \(h_\ell\)
exceeds the largest length at that prime.
\end{theorem}
\begin{proof}
Take the cyclotomic field \(C_0\) of Theorem~\ref{thm:detect241}
enlarged to contain the roots of unity of order
\(N_*\prod_{\ell\in\mathcal P}\ell^{h_\ell}\). As in that
theorem, adjoin to \(C_0\) the \(\ell\)-th roots of the rational
primes required to be power residues, for every listed \(\ell\),
to form \(K\). It is Galois over \(\QQ\). Complete splitting in
\(K\) implies those residue conditions and all the congruences
needed in Proposition~\ref{prop:moment267}. One can omit these
radicals when only the moments are required.

For each listed orbit choose a primitive basis class \(u\) of
Theorem~\ref{thm:circindex267}. Its reduction is nonzero in
\(H^1(\QQ,k_\alpha(\delta))\), where
\(k_\alpha=\mathbb F_\ell(\bar\alpha)\) and
\(\delta=\bar\alpha\bar\kappa_\ell^{r+1}\).
Indeed reduction of \(H^1(U,T)\) has kernel \(\ell H^1(U,T)\).
The map to Galois cohomology at finite coefficients is
injective in degree one by the low-degree sequence for the
generic-point direct image. One may also perform this
restriction over \(\QQ(\zeta_f)\) and take the exact character
summand. Choose the conjugate bases for the conjugate characters.
Their residual values are related by residue-field Frobenius.

We check that every such residual class remains nonzero on
\(G_K\). Over \(C_0\) its coefficient is trivial. In the
abelian group \(\operatorname{Gal}(C_0/\QQ)\), conjugation at
infinity acts on the coefficient by \(-1\). Evaluating a
cocycle at \(cg=gc\) shows it is a coboundary, since \(2\)
is invertible. Inflation--restriction therefore gives
injectivity on restriction to \(C_0\).
The map sending an automorphism over \(C_0\) to its quotients
on the chosen radicals embeds \(\operatorname{Gal}(K/C_0)\)
into a product of groups of roots of unity. Conjugation acts
on its \(\ell\)-part by \(\bar\kappa_\ell\). On inertia at
\(\ell\) the exponent of \(\delta\) is \(r+1\), strictly
between \(1\) and \(\ell-1\). Thus no equivariant homomorphism
from this radical group can contribute to the kernel of the
second restriction. This proves the assertion even with the
additional cyclotomic powers in \(C_0\).

On \(G_K\) the residual evaluations are additive homomorphisms.
Their joint image for a fixed \(\ell\) is the full product of
the additive groups \(k_\alpha\). In fact the image is stable
under the finite abelian group of the residual character values.
It surjects onto each factor, since a nonzero
\(\mathbb F_\ell\)-subspace stable under all values of \(\delta\)
is stable under the field they generate. This is \(k_\alpha\),
as \(\bar\kappa_\ell\) has values in \(\mathbb F_\ell^\times\).
Different \(r\) give inequivalent modules already on inertia at
\(\ell\), and for fixed \(r\) the inequivalent modules are exactly
the different Frobenius orbits of \(\alpha\). The finite group
algebra in residue characteristic is semisimple and its central
idempotents isolate these simple modules. Primary decomposition
gives the full product across different \(\ell\) as well.
The joint kernel on \(G_K\) is open and normal in \(G_\QQ\).
Among its finite quotient over \(K\), the fraction of vectors
with every coordinate nonzero is the stated product. These
vectors form a conjugacy-stable set over \(\QQ\), so Chebotarev's
theorem gives exactly that relative density among split primes.
Theorems~\ref{thm:circindex267} and~\ref{prop:moment267} now
give the equalities, also for characters with length zero.

For the uniform lower bound, a prime in \(\mathcal P\) contributes
at least \(\log 2\) to the sum of logarithmic strict orders over
the twists of \(E_*\). Indeed restriction from \(E_x\) is
injective because its relative two-power degree is prime to
\(\ell\). The first bound of Theorem~\ref{thm:cells103}, applied
on \(E_*\) and summed up to \(t_*\), gives
\(\#\mathcal P\le Z=\exp(V^C)\) for a fixed \(C\). This follows
since \(n_*,N_*,t_*\) are polynomial in \(V\).
The sum of \(1/\ell\) over listed primes above \(Z\) is at most
one. For primes below \(Z\), put \(\sigma=1+1/\log Z\).
Then \(\sum_{\ell\le Z}1/\ell\le e\log\zeta(\sigma)
\le e\log(1+\log Z)\), by the Euler product and the integral
bound for the Dirichlet series. Thus the sum over the whole
list is \(O_{\psi,\varepsilon}(\log(2V))\).
There are at most \(n_*t_*\) factors at each \(\ell\).
Taking logarithms and using
\(-\log(1-\ell^{-b_{\ell,\alpha}})\le2/\ell\)
proves the bound.
\end{proof}

There is an absolute size consequence for any test prime.
Fix \(x\) in the interval and \(s\equiv1\bmod N_*\). At
every large \(\ell\mid s-1\) take \(h=v_\ell(s-1)\).
For all complementary character orbits and twists of \(E_x\)
whose moments at this precision are nonzero, the sum of their
full logarithmic orders is at most
\[
                 [E_x:\QQ]\,(t(x)-1)\log(s-1).
\]
Indeed each orbit contributes \(b_{\ell,\alpha}d_{\ell,r,\alpha}
\log\ell\), and its length is at most \(h\) by
Proposition~\ref{prop:moment267}. At each twist the sum of the
orbit degrees is at most \([E_x:\QQ]\). Now use
\(\sum_{\ell\mid s-1}v_\ell(s-1)\log\ell=\log(s-1)\).
Choosing precisions above all lengths in the theorem gives
one prime at which every supported large-prime component is
counted. Its nonzero values remain nonzero if one passes to
the precision \(v_\ell(s-1)\), by compatibility of the
cyclotomic class. The density estimate is relative to splitting
in \(K\). The size of a prime detecting the full support at
these precisions is still uncontrolled on the scale required
in \eqref{eq:region198}.

\subsection{Residue valuations over ramified coefficient rings}\label{nav:II-55}
Let \(E_*,N_*,t_*\) be the largest field, modulus and twist
cutoff on an interval of Theorem~\ref{thm:conductor232}.
Take any finite set of odd coefficient primes and positive
precisions \(h_p\). Let \(C\) be an abelian cyclotomic field
containing the roots of order \(N_*\prod_p p^{h_p}\). One
may take \(K=C\). More generally, for any of the listed
\(p>\max(N_*,2t_*+3)\), one may adjoin to \(C\) the
\(p\)-th roots of any finite set of nonzero rationals,
and let \(K\) be the compositum. This permits the power-residue
requirements at the auxiliary primes in
Theorem~\ref{thm:detect241}, after enlarging \(C\) for its
further cyclotomic congruences.

For each prime \(s\) split completely in \(K\), character
\(\alpha\) of \(E_*\) of conductor \(f\), and
\(1\le r<t_*\) with the complementary parity, choose a
generator \(g\) of \(\mathbb F_s^\times\). With \(h=h_p\)
and \(w=g^{(s-1)/(fp^h)}\), set
\[
 D_j(s)=\sum_{\substack{1\le a\le fp^h\\(a,fp)=1}}
       \alpha(a)a^r\operatorname{ind}_g(1-w^a)
         \quad\bmod p^hO_\alpha,\qquad j=(p,r,[\alpha]).
\]
The index is the discrete logarithm modulo \(p^h\).
For \(y\in O_\alpha/p^hO_\alpha\) let \(w_\pi(y)\) be the
largest \(v\le e_\alpha h\) such that \(y\) belongs to
\(\pi^v(O_\alpha/p^hO_\alpha)\). Put \(W_j(s)=w_\pi(D_j(s))\).

\begin{proposition}[The marginal law for every character lattice]\leavevmode
\label{prop:residue269}
Put \(d_j=d_{p,r,\alpha}(N_*)\), and \(m_j=\min(d_j,e_\alpha h)\).
For \(0\le v\le e_\alpha h\), relative to the primes split in \(K\),
\[
                 \Pr(W_j(s)\ge v)
                 =p^{-b_\alpha\max(0,v-m_j)} .
\]
This holds when several twists have the same residual
character and when the character order is divisible by \(p\).
\end{proposition}
\begin{proof}
Take a basis \(u\) of \(H^1(U_{N_*},T)\) from
Theorem~\ref{thm:openindex269}. Its residual class is nonzero since
the kernel of reduction in cohomology is \(\pi H^1(U_{N_*},T)\).
The residual module is unramified outside the open set's
excluded primes, and inflation from the restricted Galois
group to the absolute group is injective in degree one.
Restriction to \(G_C\) keeps the class nonzero. In fact the
abelian quotient \(\operatorname{Gal}(C/\mathbb Q)\) acts on
the residual coefficient, and its central complex conjugation
acts by \(-1\). Evaluating a cocycle on the two expressions
for a product with this central element shows that it is a
coboundary on the quotient.

The group over \(C\) of the radical extension embeds in the
product of the groups of roots of unity by taking quotients
on the chosen roots. If no \(p\)-radicals are present, its
order is prime to \(p\). Otherwise its \(p\)-part has global
conjugation character \(\bar\kappa_p\), since the radicands
are rational. The residual coefficient on inertia at this
large \(p\) is \(\bar\kappa_p^{r+1}\), distinct from
\(\bar\kappa_p\). Thus the kernel of restriction from
\(G_C\) to \(G_K\) cannot contain the restricted class.
Indeed such a class would factor as an equivariant additive
homomorphism from the radical group into the residual
coefficient. This proves nonvanishing on \(G_K\) in both cases.

On \(G_K\) the integral basis cocycle at precision \(p^h\)
is an additive homomorphism into \(O_\alpha/p^hO_\alpha\),
since its coefficient character is trivial there. Its
image is an ideal of this ring. To see this, conjugation
by a global element preserves the image and acts by its
value of \(\alpha\kappa_p^{r+1}\). The image is additive,
so is stable also under every \(p\)-adic integer modulo
\(p^h\). Removing the cyclotomic unit scalar leaves all
values of \(\alpha\), which generate \(O_\alpha\) over
\(\mathbb Z_p\). The image contains a unit by the residual
nonvanishing, hence is the whole ring. Its kernel is open
and normal over \(\mathbb Q\). Divisibility by \(\pi^v\) is
conjugacy-stable, so conditional Chebotarev gives the
fraction \(p^{-b_\alpha v}\) for the basis cocycle.

Finally specialize the actual class \(c\) on the same field.
In the corestriction formula all embeddings at \(s\) send
\(\zeta_{fp^h}\) to the different \(w^a\). The tensor factors
contribute \(\alpha(a)a^r\). Arithmetic Frobenius on a
\(p^h\)-th root of a residue unit \(z\) gives the quotient
\(z^{(s-1)/p^h}\), with coordinate \(\operatorname{ind}_g z\)
in the basis \(w^f\). Thus the moment is the evaluation of
\(c\), up to a basis unit, also when \(p\mid f\).
Multiplication of the primitive class by its coordinate
shifts the valuation by \(d_j\) until it reaches \(e_\alpha h\).
This gives the asserted law.
\end{proof}

\subsection{Frobenius orbits in the local cyclotomic sum}\label{nav:II-56}
Fix the largest modulus \(N_*\) on one interval as in
Theorem~\ref{thm:average269}. Write \(k=r+1\), \(f\mid N_*\)
for the primitive character conductor, \(p\) for an odd coefficient
prime, and \(e=v_\pi(p)\). Take \(v\ge1\) and \(h=\lceil v/e\rceil\).
Let \(s\) be prime with \(s\equiv1\bmod N_*\), \(s\ne p\), and
\(p^h\mid s^k-1\). Put \(n=fp^h\), \(\ell=\operatorname{ord}_n(s)\),
and \(A=(\mathbb Z/n\mathbb Z)^\times\). In \(\mathbb F_{s^\ell}\)
choose a multiplicative generator \(g\), and set
\(w=g^{(s^\ell-1)/n}\).

\begin{proposition}[Specialization without division by the residue degree]\leavevmode
\label{prop:orbits273}
In a choice of coefficient basis, the unramified arithmetic
Frobenius value of \(c\) modulo \(\pi^v\) is
\[
  D_v(s)=\sum_{a\in A/\langle s\rangle}
       \alpha(a)a^r\,\operatorname{ind}_{g}(1-w^a)
                       \quad\bmod\pi^v .
\]
The sum takes one representative in each multiplicative coset.
The discrete logarithm is reduced modulo \(p^h\).
Its vanishing to any order at most \(v\) is independent of the
choices. For fixed \(s,k\) every eligible calculation, including
\(p\mid f\), takes place in subfields of
\(\mathbb F_{s^{\,kN_*}}\).
\end{proposition}
\begin{proof}
At \(s\) the coefficient action of arithmetic Frobenius is
\(\alpha(s)s^k=s^k\), equal to one modulo \(\pi^v\).
The \(s\)-adic factors of the cyclotomic field of order \(n\)
are unramified of degree \(\ell\). They correspond to orbits
of \(a\mapsto sa\) on the primitive roots \(w^a\). Localizing
the global corestriction sums the corestrictions of these
factors. To see this on cocycles, partition the cosets in
the transfer formula by the cycles of Frobenius. Choose
successive powers of Frobenius times one initial representative
around a cycle. All changes of representative cancel except
the return step, which evaluates the conjugated cocycle on
the \(\ell\)-th power of Frobenius. Any power of Frobenius
acting on its value is trivial on the target modulo \(\pi^v\).
This gives one contribution per orbit even if \(p\mid\ell\).

At the representative \(a\) the Kummer unit specializes to
\(1-w^a\). Arithmetic Frobenius to the power \(\ell\) sends a
\(p^h\)-th root of that unit to itself times
\((1-w^a)^{(s^\ell-1)/p^h}\). Its coordinate in the basis
\(w^f=g^{(s^\ell-1)/p^h}\) is the indicated discrete logarithm.
The other Tate factors and the character basis have respective
coordinates \(a^r\) and \(\alpha(a)\), as in
Proposition~\ref{prop:residue269}. Replacing \(a\) by \(sa\)
raises the residue unit to its \(s\)-th power and multiplies
the full term by \(\alpha(s)s^{r+1}=1\bmod\pi^v\).
Changing \(g\) to \(g^j\) with \(j\) prime to \(s^\ell-1\)
permutes the root indices by \(a\mapsto ja\) and scales the
entire sum by \(\alpha(j)^{-1}j^{-k}\), a unit.

The prime-to-\(p\) part of \(f\) divides \(s-1\). If \(p\nmid f\),
then \(\ell\mid k\). Otherwise put \(b_0=v_p(f)\). Since
\(s^k\equiv1\bmod p^h\), the binomial theorem applied
\(b_0\) times gives \(s^{kp^{b_0}}\equiv1\bmod p^{h+b_0}\).
Thus \(\ell\mid k p^{b_0}\mid kN_*\). In the common finite
field one may obtain a generator of the degree-\(\ell\)
subfield by raising a generator to
\((s^{kN_*}-1)/(s^\ell-1)\). The logarithm in the formula is
taken with respect to that subfield generator.
\end{proof}

\subsection{Orders of the detected strict classes on the period interval}\label{nav:II-73}
At each \(x\) on the same original interval, use
\(E=E_x\), \(t=\lfloor x^{401}\rfloor\), and let \(M_x\)
be the integer in Theorem~\ref{thm:simultaneous307} formed
with this field and this cutoff. Consider the triples at
\(\ell>\exp(x^\varepsilon/12)\) of that theorem, with one
term per residue orbit. Write \(d=v_\ell(L(1-r,\alpha^*))\)
and use \(v_j,h_0,b\) from Theorem~\ref{thm:layers308}.
Put
\[
    H(x)=\exp(x^\varepsilon/2)/x^{820},\qquad
    R_{\ell,r,\alpha}=\min\left(v_1,\frac{\log h_0}{b\log\ell}\right).
\]
Let \(\mathcal J_x\) consist of the triples with
\(R_{\ell,r,\alpha}\le x^3H(x)\), and let
\(\mathcal F_x\) consist of those for which \(v_1=0\) or
\(v_{d+1}\le d\).

\begin{theorem}[A bound for orders on the original fields and intervals]\leavevmode
\label{thm:orders308}
Uniformly for all \(x\) of every sufficiently distant original
interval the two estimates
\[
 \begin{split}
 \sum_{\mathcal J_x} b\log\ell\,
       \operatorname{length}_O S^-[\ell^d]
      &\ll_{\psi,\varepsilon} x^{-5}\exp(x^\varepsilon/2),\\
 \sum_{\mathcal F_x} b\log\ell\,\operatorname{length}_O S^-
      &\le \log M_x\ll_{\psi,\varepsilon} x^{812}
 \end{split}
\]
hold for the original complementary strict groups. The first
estimate applies in particular to the full orders of the
subgroups \(S^{\rm ev}\) of
Theorem~\ref{thm:simultaneous307}. Its range includes every
triple with
\(\exp(x^\varepsilon/12)<\ell\le H(x)\).
Every component with \(v_1=0\) vanishes.
\end{theorem}
\begin{proof}
The elementary-divisor decomposition of the finite \(O\)-module
\(S^-\) shows that its \(\ell^d\)-torsion has length at most
\(d\dim_{O/\ell}S^-[\ell]\), also when \(d=0\). Theorem~\ref{thm:layers308}
bounds the dimension by \(R_{\ell,r,\alpha}\).
The exponent bound in Theorem~\ref{thm:simultaneous307}
puts \(S^{\rm ev}\) in this same torsion group. Summing the
length estimate on \(\mathcal J_x\) costs at most
\(x^3H(x)\log M_x\), by the prime valuation formula there
on \(E_x\). Its archimedean estimate gives
\(\log M_x\ll x^{812}\) with the field and conductor bounds
of Theorem~\ref{thm:conductor232}. This proves the first
inequality.

To check the entire asserted prime range, the residual
character cutting out \(K_0\) has conductor dividing
\(f_\alpha\ell\) and order at most
\([E:\QQ](\ell-1)\). The latter follows by multiplying
the orders of the finite character and of the residual
cyclotomic character. The conductor--discriminant bound
and the Minkowski ideal count in Theorem~\ref{thm:evenbudget306}
therefore give
\[
      \log h_0\ll [E:\QQ]\ell\log(2f_\alpha\ell).
\]
For every prime under consideration
\(\log(2f_\alpha\ell)\ll\log\ell\) uniformly, since
\(f_\alpha\mid N_x\) and \(\log N_x\ll x^{\varepsilon/4}\).
Also \([E:\QQ]\ll_\psi x^{2+8\varepsilon}\) with
\(2+8\varepsilon<3\). Hence for large \(x\) one has
\(R_{\ell,r,\alpha}\le x^3\ell\), which includes the
claimed range in \(\mathcal J_x\).

If \(v_1=0\), Theorem~\ref{thm:layers308} gives \(S^-=0\).
In the other case of \(\mathcal F_x\), its stopping test at
\(j=d+1\) bounds the length of the whole group by \(v_{d+1}\le d\).
Summing these lengths costs at most \(\log M_x\) again.
All local conditions in the argument are those of the
original field by character descent. Each bound is pointwise
in \(x\) with uniform constants on the period interval.
\end{proof}

For a contribution from \(S^{\rm ev}\) outside the first estimate,
one needs a positive valuation of \(L(1-r,\alpha^*)\) and
a valuation exceeding \(x^3H(x)\) of the primitive value at
the reflected weight \(s_1\), at the same prime \(\ell>H(x)\).
The remaining full strict orders also include classes above
the controlled torsion levels. A bound for those orders and
the comparison for the arithmetically normalized theta pairing
are still required for the main theorem.

\subsection{Summing full strict orders across weights}\label{nav:II-74}
Use the period field \(E=E_x\) of degree \(n\) and modulus \(N\)
on the interval of Theorem~\ref{thm:conductor232}, and put
\(t=\lfloor x^{401}\rfloor\). For \(3\le p\le Y\), \(p\nmid nN\),
take one \(\alpha\) in each residue Frobenius orbit. Write
\(b_\alpha=[O/p:\mathbb F_p]\). Group the twists of this character
by their residue modulo \(p-1\), retaining the hypotheses
of Theorem~\ref{thm:localvalue309}. A nonempty progression
\(\mathcal T\) consists of all \(1\le r<t\) in the chosen
residue class. Let \(k\in[2,p-1]\) satisfy \(k\equiv-r\bmod(p-1)\),
put \(l=|\mathcal T|\), and define
\[
 v=v_p(L(1-k,\alpha^{-1})),\qquad
 \lambda_r=\operatorname{length}_O S_{\alpha,r}.
\]
The character in the special value is primitive. Here
\(S_{\alpha,r}\) is the full original complementary strict group
in this character component.

\begin{theorem}[Full strict orders outside finitely many twists in each progression]\leavevmode
\label{thm:weightmass309}
There is a subset \(D\subset\mathcal T\) with \(|D|\le v\) and
\[
 \sum_{r\in\mathcal T\setminus D}\lambda_r
       \ \le\ v\bigl(2l+v_p((l-1)!)\bigr).
\]
A removed twist lies within \(p\)-adic distance strictly less
than \(p^{-l}\) of a root in the coordinate specified below.
Using these sets at all the primes and characters above gives
\[
\begin{split}
 \sum_{p,[\alpha],\mathcal T} b_\alpha\log p
      \sum_{r\in\mathcal T\setminus D}\lambda_r
   &\ll n(Y^2+tY)\log(2NY),\\
 \sum_{p,[\alpha],\mathcal T} b_\alpha\log p\,|D|
   &\ll nY^2\log(2NY).
\end{split}
\]
In particular the first sum is \(O_{\psi,\varepsilon}
(x^{-2}\exp(x^\varepsilon/2))\) for
\(Y=x^{-3}\exp(x^\varepsilon/4)\), uniformly on the original
intervals. It includes all twists and characters at
\(\exp(x^\varepsilon/12)<p\le Y\) except the individual
triples removed by these sets.
\end{theorem}
\begin{proof}
Use one smoothing integer \(B\) and series \(F\) from
Theorem~\ref{thm:localvalue309} for the whole progression.
Set \(a=X_{1-k}\) and \(G(U)=F(a+pU)\). Its \(i\)-th coefficient
belongs to \(p^iO\), by the binomial expansion at \(a\in pO\).
In particular the coefficients tend to zero. The constant
coefficient has valuation \(v\). Let \(\mu\) be the minimum
of the coefficient valuations and \(m\) the last index at
which the minimum occurs. Then \(m\le\mu\le v\).
After dividing by \(p^\mu\) and by a unit, the reduction
of the series is a monic polynomial of degree \(m\).

On the closed unit disk this series is a unit times a monic
polynomial \(P\in O[U]\) of that degree. Here is the lifting
argument. Start with a monic lift of the reduction and unit
factor one. Given a product correct modulo \(p^i\), reduce its
error divided by \(p^i\) modulo \(p\). It is a polynomial
since the series is restricted. Divide by the original monic
reduction. Add \(p^i\) times lifts of the quotient to the unit
factor and of the remainder to the polynomial. The product is
now correct modulo \(p^{i+1}\), with the same leading coefficient
and unit reduction. The limits give the factorization. Its
unit has unit value throughout the closed disk. Denote the
roots of \(P\), with multiplicity, by \(\xi_1,\ldots,\xi_m\)
in a finite extension of \(J\). They are integral since \(P\)
is monic.

For \(r\in\mathcal T\) put \(u_r=(X_{1+r}-a)/p\).
The binomial theorem at odd \(p\) gives
\(v_p(u_r-u_s)=v_p(r-s)\) for \(r\ne s\).
For each root choose \(r_i\) maximizing \(v_p(u_{r_i}-\xi_i)\).
Put \(r_i\) in \(D\) if this maximum is greater than \(l\).
At any \(r\ne r_i\) the ultrametric inequality gives
\(v_p(u_r-\xi_i)\le v_p(r-r_i)\).
Index \(\mathcal T\) by \(0,\ldots,l-1\). Its step \(p-1\) is a
unit. Thus the sum of the latter upper bounds is
\(v_p(j!)+v_p((l-1-j)!)\le v_p((l-1)!)\), where \(j\) indexes
\(r_i\). If \(r_i\) is retained its own contribution for this
root is at most \(l\). For every retained twist,
Theorem~\ref{thm:localvalue309} gives
\[
    \lambda_r\le v_p(F(X_{1+r}))
        =\mu+\sum_{i=1}^m v_p(u_r-\xi_i).
\]
A zero at an actual twist has already been removed. Summing
gives \(l\mu+m(l+v_p((l-1)!))\), which proves the first bound
and the assertion about distances.

For fixed \(k\), the sum of \(b_\alpha v\log p\) over the
participating primes and characters is at most \(\log J_k(E)\)
from \eqref{eq:int201}. Indeed the rational numerator norm in
that formula counts the valuations at \(p\) with their residue
degrees. The Bernoulli polynomial estimate in
Theorem~\ref{thm:integer201} gives
\(\log J_k(E)\ll nk\log(2Nk)\). We have \(l\le1+2t/k\) since
\(k\le p\), and \(v_p((l-1)!)\le(l-1)/(p-1)\).
Sum first over \(p,[\alpha]\), then over \(2\le k\le Y\).
The two displayed estimates follow from
\(\sum_{k\le Y}(k+2t)\ll Y^2+tY\) and \(\sum_{k\le Y}k\ll Y^2\).
For the stated choice of \(Y\) we eventually have \(t\le Y\).
The common-field bounds \(n\ll_\psi x^{2+8\varepsilon}\) and
\(\log(2NY)\ll_{\psi,\varepsilon}x^\varepsilon\) give the asserted
order, since \(9\varepsilon<1\) in the standing range. In the
displayed large-prime range \(p\) exceeds \(n,N,t+2\) on
sufficiently distant intervals, so every complementary character
and twist qualifies.
\end{proof}

The estimate for \(|D|\) counts triples with multiplicity
\(b_\alpha\log p\). Bounding the orders of the strict groups at
those triples, and at larger primes, is still necessary for
the uniform cohomology estimate in Section~\ref{sec:missing}.

\subsection{A single list of remaining twists on the period interval}\label{nav:II-75}
Use the nested fields of Theorem~\ref{thm:conductor232} on an
original interval \([V,V^{3/2}]\), and put \(t_*=\lfloor V^{1203/2}\rfloor\).
For a point \(x\) write \(E=E_x\), \(n=[E:\QQ]\), \(N=N_x\),
\(t=\lfloor x^{401}\rfloor\) and \(Y=x^{-3}\exp(x^\varepsilon/4)\).
Let \(\mathcal B\) be the fixed odd prime divisors of the order of
\(\psi\). They are exactly the odd prime divisors of \(n\), since
the rest of the extension has degree a power of two.

Use \(M_{p,r}=\frac12(1+(-1)^rc)\mathcal S_{p,r}(E)\) at odd
\(p\). Write \(W_{p,r}=p^{k_{p,r}}(1-e'_r)Q_{p,r}\subset M_{p,r}\)
for the subgroup already retained in \eqref{eq:region198}. The local
length sum of Section~\ref{sec:missing} differs from
\(\sum_{p,r}\log|W_{p,r}|\) by \(o(\exp(x^\varepsilon/2))\).

\begin{theorem}[Full orders on one list for all fields of an interval]\leavevmode
\label{thm:entireinterval310}
One can select exceptional twists once on \([V,V^{3/2}]\) so that
the following assertions hold at every \(x\). For \(p\le Y\)
outside \(\mathcal B\), selection uses one set
\(D_{p,[\alpha],k}\) for each character orbit and
\(2\le k\le p\) with \(r\equiv-k\bmod(p-1)\) and the required
parity. Globally trivial residual characters have empty sets.
Otherwise, with \(e\) from Theorem~\ref{thm:forced310}, one has
\[
 |D_{p,[\alpha],k}|
   \ \le\ v_p(L(1-k,\alpha^{-1}))-e(1+v_p(k-1)).
\]
At \(p\in\mathcal B\) a single set \(D_p\) suffices for the
whole complementary group. In the first range the sum of the
full logarithmic orders of all unselected components is
\(O_{\psi,\varepsilon}(x^{-2}\exp(x^\varepsilon/2))\).
In the second range the sum of the full logarithmic orders of
the unselected groups is bounded by a polynomial in \(x\).

More precisely define \(W'_{p,r}\) to equal \(W_{p,r}\) for
\(p>Y\). At \(p\le Y\) outside \(\mathcal B\) retain in \(W'_{p,r}\)
only the orbit summands of \(W_{p,r}\) for which
\(r\in D_{p,[\alpha],k}\). At \(p\in\mathcal B\) retain \(W_{p,r}\)
if \(r\in D_p\), and use zero otherwise. Then
\[
       \mathcal L(x)=\sum_{1\le r<t}\sum_{p\ {\rm odd}}
                 \log|W'_{p,r}|+o(\exp(x^\varepsilon/2)).
\]
The additional removals described in Theorem~\ref{thm:orders308}
can be made in this formula as well.
\end{theorem}
\begin{proof}
Make the selection using the largest field \(E_*\) of the
interval, the full twist range \(r<t_*\), and the maximum of
the cutoffs \(Y\) on the interval. For \(p\notin\mathcal B\),
character evaluation decomposes the integral group ring into
unramified integer rings, including at primes ramified in \(E\).
A component of \(O\)-length \(\lambda_r\) therefore contributes
\(b_\alpha\lambda_r\log p\) to the logarithmic order of \(M_{p,r}\),
by Theorem~\ref{thm:char}. Theorem~\ref{thm:criticalall310} kills
the globally trivial residual case. In every other case use the
root selection of Theorem~\ref{thm:weightmass309} with the
integral series \(F^\circ\). Theorem~\ref{thm:forced310} bounds
the full strict length by the valuation of this series wherever
the valuation is finite. Its sample valuation is the displayed
bound on \(|D_{p,[\alpha],k}|\). A zero of the series at an
integer twist is included by the selection rule.

A character occurring earlier on the interval inflates to \(E_*\)
with the same primitive character, coefficient ring, and local
conditions over \(\QQ\). Thus the sets just chosen apply
throughout. At a fixed \(x\) sum only over characters of \(E_x\)
and primes \(p\le Y(x)\). The complete progression in the larger
twist range has at most \(1+2t_*/k\) terms. Bounding its retained
sum also bounds that of its prefix at \(x\). The numerator norm
calculation \eqref{eq:norm201} applies here even if \(p\mid N\):
it counts valuations in the character value fields, which are
unramified at \(p\nmid n\). The corrected valuation is no larger
than the positive part of the primitive valuation. Consequently
the weighted count of the selected triples at \(x\) is at most
\(\sum_{2\le k\le Y}\log J_k(E_x)\). The logarithmic orders
outside the selected triples sum to at most
\[
 C\sum_{2\le k\le Y}(1+2t_*/k)\log J_k(E_x)
       \ll n(Y^2+t_*Y)\log(2NY).
\]
Uniformly on the interval \(t_*\le x^{1203/2}\le Y\) eventually.
The degree and conductor bounds in
Theorem~\ref{thm:weightmass309} now give the asserted estimate.

For each of the finitely many \(p\in\mathcal B\), apply
Theorem~\ref{thm:every144} at \(x_*=V^{3/2}\). Its whole-field
exceptional set \(D_p\) has size \(O_\psi(n_*\log(2N_*))\),
where \(n_*=[E_*:\QQ]\) and \(N_*\) is its modulus. Outside
this set \(p^{\lceil x_*^{202}\rceil}\) kills the whole
complementary group on \(E_*\). Restriction from \(E_x\) is
injective since corestriction composed with restriction is
multiplication by the relative degree, a power of two.
Restriction preserves the local conditions and commutes with
complex conjugation. At \(E_x\), Lemma~\ref{lem:rank114} bounds
the logarithmic order at exponent \(p\) by \(O_\psi(n\log(2Np))\).
Multiplying this by \(\lceil x_*^{202}\rceil\) bounds each full
logarithmic order in question. There are at most \(t\) terms per
prime, and \(x_*\le x^{3/2}\), proving the polynomial bound.

Each part discarded from \(W_{p,r}\) lies in the corresponding
part of \(M_{p,r}\), so the same estimates pay for its removal.
This proves the formula by \eqref{eq:region198}. To see the last
assertion, at \(p>\exp(x^\varepsilon/12)\) discard a remaining
component completely if it belongs to \(\mathcal F_x\) of
Theorem~\ref{thm:orders308}. On the remaining components in
\(\mathcal J_x\), quotient by the intersection with \(S^-[p^d]\)
from that theorem. The total costs are respectively at most
\(\log M_x\) and \(O_{\psi,\varepsilon}(x^{-5}\exp(x^\varepsilon/2))\).
Both are absorbed by the error, independently of overlaps
with the earlier removals.
\end{proof}

The selected twists below \(Y\) and the remaining groups above
\(Y\) enter with their actual orders. A uniform estimate for
their sum and the arithmetic period comparison on this interval
are still required.

\subsection{One auxiliary field for the ordinary congruences}\label{nav:II-45}
\begin{corollary}[Simultaneous removal of auxiliary tame ramification]
\label{cor:common240}
Fix a sufficiently distant period interval \([V,V^{3/2}]\).
For every nonzero large-prime strict summand in
Corollary~\ref{cor:interval238} use its original \(E_x,\ell,r,\alpha\).
There are two rational primes \(s_1,s_2\) which split in every
\(E_x\) and in \(K'=\QQ(i)\), and one totally real abelian
extension \(L/\QQ\) ramified only at these two primes, such
that Theorem~\ref{thm:adjoint240} supplies a cuspidal lift for
each summand. The coefficient primes, all primes of the original
field conductors and the period primes \(g_1,g_2,T\) split
completely in \(L\). On \(K'L\) the corresponding adjoint is
ordinary crystalline over \(\ell\), realizes the given residual
diagonal and has prime-to-\(\ell\) conductor
\(\mathfrak f(\alpha|_{G_{K'L}})^2\), where \(\mathfrak f\)
denotes the conductor ideal. On \(E_xK'L\) it is unramified
away from \(\ell\) and remains cuspidal. All assertions use
the same \(s_1,s_2,L\) throughout the interval.
\end{corollary}
\begin{proof}
Corollary~\ref{cor:interval238} proves that the list of data
with nonzero strict group is finite once the interval is fixed.
In particular the set \(P\) of coefficient primes occurring is
finite. Let \(B\) consist of \(P\), the prime divisors of the
field conductors, and the period primes in the statement.
The last list is finite too. The bounds for \(g_1,g_2,T\)
in Theorem~\ref{thm:prime210} are uniform on this bounded
interval. Choose a modulus \(M\) divisible by the conductors
of \(K'\) and all \(E_x\) and by every required \(8N\ell^a\)
in Lemma~\ref{lem:cusps240}. Form the splitting field of
the cyclotomic polynomial of order \(M\) and the finitely many
polynomials \(X^\ell-p\) for \(\ell\in P,\ p\in B\).
Chebotarev gives two distinct completely split primes outside
all primes used so far. These serve as \(s_1,s_2\).
Each \(p\in B\) is an \(\ell\)-th power modulo each \(s_j\)
for every \(\ell\in P\).

For each pair \(s_j,\ell\), take the unique degree-\(\ell\)
subfield \(L_{j,\ell}\subset\QQ(\zeta_{s_j})\) and a character
\(\xi_{j,\ell}\) cutting it out. Existence and uniqueness follow
from cyclicity of the multiplicative group modulo \(s_j\).
Its kernel consists of the \(\ell\)-th powers in that group.
As \(\ell\) is odd, complex conjugation belongs to the kernel.
The compositum \(L\) of these real fields thus has the indicated
ramification and splitting. For the empty list take \(L=\QQ\).
Apply Theorem~\ref{thm:adjoint240} to each datum with its own
two characters of order \(\ell\).

On \(G_{K'L}\) the characters \(\eta,\xi_{1,\ell},\xi_{2,\ell}\)
are trivial. Thus at every prime away from \(\ell\) the inertia
characters are \(\alpha,1,\alpha^{-1}\), with zero monodromy.
The conductor exponent is twice that of \(\alpha\), since a
character and its inverse have the same ramification groups
as kernels. Restriction to \(G_{E_xK'L}\) kills \(\alpha\).
The ordinary flag and crystalline periods persist under these
restrictions, and cuspidality follows by solvable base change as
in Theorem~\ref{thm:adjoint240}.
\end{proof}

\begin{theorem}[Descent of the original strict group from the same field]
\label{thm:descent240}
For one datum in the corollary, put \(F_x=E_xK'L\) and
\(\Delta=\operatorname{Gal}(F_x/\QQ)\). Let \(O\) be the
unramified coefficient DVR of the character, \(J\) its
fraction field, and \(A=(J/O)(\kappa_\ell^{-r}\alpha^{-1})\).
For a number field \(D\), write \(S_D(A)\) for the classes in
\(H^1(D,A)\) trivial on inertia away from \(\ell\) and on the
whole decomposition groups at \(\ell\). For \(j=1,2\) let
\(I_j\) be inertia at \(s_j\) over \(\QQ\), and \(I'_j\) inertia
over \(F_x\) using the same choice of place. There is an exact sequence
\[
 0\longrightarrow S_{\QQ}(A)
 \longrightarrow S_{F_x}(A)^\Delta
 \longrightarrow
     \bigoplus_{j=1}^2\operatorname{Hom}(I_j/I'_j,A[\ell])
\]
where exactness is asserted at the first two groups. The last
module is isomorphic to \((O/\ell O)^2\). The same sequence
holds with \(A[\ell^n]\) as coefficients, for each \(n\ge1\).
Thus descent to the original strict summand costs at most two
units of \(O\)-length, independently of the degree of \(L\),
the character, the twist and the field on the interval. The
same bound holds for the difference of lengths after multiplying
both groups by any \(\ell^j\), \(j\ge0\).
\end{theorem}
\begin{proof}
Inertia over \(\ell\) in \(G_{F_x}\) has the full residual
cyclotomic image. Indeed \(E_xK'\) is unramified there and
\(\ell\) splits completely in \(L\). Since \(0<r<\ell-1\)
and \(\alpha\) is unramified at \(\ell\), some inertia element
acts on \(A\) by a scalar whose difference from one is a unit
of \(O\). In particular \(A^{G_{F_x}}=0\), also for the finite
coefficient submodules. Inflation--restriction now identifies
\(H^1(\QQ,A)\) with \(H^1(F_x,A)^\Delta\), even though
\(\ell\) may divide \(|\Delta|\). Every invariant strict
class upstairs has a unique global antecedent.

At \(\ell\) the local degree of \(F_x/\QQ\) divides
\([E_xK':\QQ]\), prime to \(\ell\) in this large-prime range.
Local restriction on degree-one cohomology is injective, since
corestriction composed with restriction is multiplication by
that degree. Away from \(\ell,s_1,s_2\) the index of inertia
also divides \([E_xK':\QQ]\), since \(L\) is unramified
there. The same argument gives injectivity on the inertia
cohomology. The antecedent therefore satisfies every original
local condition except possibly those at \(s_1,s_2\).

The action on \(A\) at inertia \(I_j\) is trivial. Restrict
the antecedent to \(I_j\). Its restriction is a homomorphism
vanishing on \(I'_j\), independent of any cocycle representative.
The quotient \(I_j/I'_j\) is cyclic with \(\ell\)-part of
order exactly \(\ell\). Indeed \(E_xK'\) splits at \(s_j\);
the character fields from the other auxiliary prime are
unramified there. At \(s_j\) the fields from this prime are
subfields of \(\QQ(\zeta_{s_j})\) with exactly one factor of
degree \(\ell\) and the remaining factors of distinct prime
degrees. They are totally ramified there. The image of the
homomorphism is therefore in \(A[\ell]\), and the Hom module
for this prime is one copy of \(O/\ell O\).
The two restrictions vanish precisely when the antecedent
is strict. This proves exactness and the length bound.
If the first two groups are \(M\subset B\), multiplication gives
a surjection \(B/M\to\ell^jB/\ell^jM\), proving the assertion
at depth \(j\). In particular, on any collection of character
summands for which \(\ell^j M\ne0\), the sum of the logarithmic
orders upstairs lies between the corresponding sum downstairs
and three times that sum. Indeed a nonzero module has length
at least one, and both logarithmic orders use the same factor
\(\log|O/\ell O|\) per unit of length.
The proof uses the same invariants and the same inertia
quotients for \(A[\ell^n]\). The identification of
\(S_{\QQ}(A)\) with the original character summand over \(E_x\)
is the prime-to-degree descent in Theorem~\ref{thm:char}.
\end{proof}

The comparison thus retains the original strict group, with
its full primary exponents, on a field supporting the stated
cuspidal adjoints. A bound for the orders of the invariant
groups in this sequence is still needed in \eqref{eq:region198}.
The signed integral period comparison concerns the original
field \(E_x\) throughout.

\subsection{One prime detecting all the integral dual classes}\label{nav:II-46}
\begin{theorem}[Simultaneous detection subject to the period splitting conditions]
\label{thm:detect241}
Fix a sufficiently distant interval of Corollary~\ref{cor:common240}
and all its nonzero large-prime strict summands. There is a
positive-density set of rational primes \(s\) satisfying every
individual congruence of Proposition~\ref{prop:level241} and
splitting in every \(E_x\) and \(K'=\QQ(i)\). For every
coefficient prime \(\ell\), every rational prime listed for
splitting in Corollary~\ref{cor:common240} is an \(\ell\)-th
power modulo \(s\). For every datum one can fix an integral
class \(d\) as in Lemma~\ref{lem:dual241} so that its
reduction evaluates nontrivially at Frobenius at every such
\(s\). The reduction character is trivial locally at \(s\),
so nonvanishing is independent of the Frobenius representative.
\end{theorem}
\begin{proof}
There are finitely many data on this interval by
Corollary~\ref{cor:interval238}. Denote the coefficient primes
by \(P\) and the list of primes required to split by \(B\).
Take a cyclotomic field \(C_0\) containing all the original
fields, \(K'\), and the roots of unity for every prescribed
modulus \(8N\ell^a\). Let \(K\) be obtained from \(C_0\) by
adjoining an \(\ell\)-th root of every \(p\in B\) for each
\(\ell\in P\). This is Galois over \(\QQ\).

Fix one residual character \(\delta=\bar\kappa_\ell^{r+1}
\bar\alpha\) and denote its module by \(D=k_\ell(\delta)\).
Restriction \(H^1(\QQ,D)\to H^1(K,D)\) is injective.
Indeed \(D\) is trivial over \(C_0\). For the finite abelian
group \(G=\operatorname{Gal}(C_0/\QQ)\), a central element
acting by a scalar different from one makes \(H^1(G,D)=0\).
Writing a cocycle at \(zg=gz\) expresses it as a coboundary.
The group \(\operatorname{Gal}(K/C_0)\) embeds in a product
of the groups \(\mu_{\ell'}\) for \(\ell'\in P\) by its action
on the chosen radicals. Conjugation on each \(\mu_{\ell'}\)
coordinate is cyclotomic, since the radicands are rational.
Thus its \(\ell\)-part has scalar action \(\bar\kappa_\ell\).
On inertia at \(\ell\), the exponent \(r+1\) lies strictly
between \(1\) and \(\ell-1\) and \(\alpha\) is unramified.
Consequently \(\delta\) is neither trivial nor
\(\bar\kappa_\ell\). There are no \(G\)-equivariant
homomorphisms from \(\operatorname{Gal}(K/C_0)\) to \(D\).
Two applications of inflation--restriction prove injectivity.

Take one character from each orbit under residue-field
Frobenius for each \(\ell\), and choose the class of
Lemma~\ref{lem:dual241}. Conjugate characters use
conjugate classes. Indeed the restriction to inertia at
\(\ell\) determines \(r+1\) in the above range, and
prime-to-\(\ell\) finite characters lift uniquely from their
residues. Use the same class whenever the datum repeats.
On \(G_K\) each reduction is a nonzero additive
homomorphism to \(k_\ell\). Its image is stable under
multiplication by \(\delta(G_\QQ)\), by the cocycle identity
under conjugation. Those values generate the entire minimal
residue field over \(\mathbb F_\ell\). Thus each image is full.

Their joint image is the full product. For a fixed \(\ell\)
it is a submodule of the sum of the \(D\)'s for the
finite abelian group of residual character values. This group
has order prime to \(\ell\). Its group algebra over
\(\mathbb F_\ell\) is a product of finite fields, with
nonconjugate characters in distinct factors. The corresponding
idempotents separate the images, each of which is full.
Different coefficient primes separate by primary decomposition
of the finite additive image. Hence one can choose \(g\in G_K\)
with every evaluation nonzero. The common kernel on \(G_K\)
is open and normal in \(G_\QQ\) by conjugation equivariance.
Chebotarev applied to the conjugacy class of \(g\) in this
finite quotient gives the desired positive-density set.
Its primes split in \(K\), giving all the stated splitting
and power conditions. Conjugation multiplies each evaluation
by a unit. A finite set of ramified or previously used primes
may be discarded. For an empty list ordinary complete splitting
in \(C_0\) suffices.
\end{proof}

\subsection{Interpolation at a fixed conductor}\label{nav:III-35}
Fix an odd prime \(p\), a primitive character \(\chi=\alpha^{-1}\),
and one admissible branch \(2\le k\le p\) of
Theorem~\ref{thm:scalar231}. Thus its residual coefficient
character is nontrivial. Let \(O\) be the integer ring of the
coefficient field, put \(X_s=(1+p)^s-1\), and take the same
smoothing integer \(b\) at all weights in this branch.

\begin{theorem}[Recovery of the strict annihilator from finitely many weights]
\leavevmode\label{thm:newton235}
For \(H\ge1\), set \(k_i=k+i(p-1)\), \(a_i=X_{1-k_i}\)
for \(0\le i<H\). Let \(P\) be the polynomial of degree less
than \(H\) with
\[
 P(a_i)=
 b^{1-k_i}(b^{k_i}-\chi(b)^{-1})(1-\chi(p)p^{k_i-1})
                         B_{k_i,\chi}/k_i .
\]
Then \(P\in O[X]\). If \(F\) is the annihilating series in
Theorem~\ref{thm:scalar231}, then on every admissible twist
\(r\ge1\),
\[
 F(X_{r+1})\equiv P(X_{r+1})
         \pmod {p^{H+v_p(H!)}O}.
\]
In particular a value of \(P\) of valuation \(w<H+v_p(H!)\)
gives \(p^{\lceil w\rceil}S_{\alpha,r}=0\). The same \(H\)
samples determine every coefficient of the restricted series
\(F(a_0+pU)\) modulo \(p^H O\).
\end{theorem}
\begin{proof}
For \(a\in pO\) and \(F\in O[[X]]\), subtract \(F(a)\) and
divide by \(X-a\). The resulting series is integral: the
coefficient at degree \(j\) is
\(\sum_{n>j}[X^n]F\,a^{n-1-j}\), which converges in \(O\).
Apply this division successively at \(a_0,\ldots,a_{H-1}\).
Newton's formula gives \(F=P_0+\prod_i(X-a_i)G\) for
\(P_0\in O[X]\) of degree less than \(H\) and \(G\in O[[X]]\).
The nodes are distinct, and Theorem~\ref{thm:scalar231} gives
exactly the specified values there. Thus \(P=P_0\).

For integers \(s\ne t\),
\(v_p(X_s-X_t)=1+v_p(s-t)\). This follows by factoring out
the unit \((1+p)^t\): a power prime to \(p\) preserves the
valuation of \((1+p)-1\), and each subsequent \(p\)-th power
increases it by one by the binomial expansion, since \(p\)
is odd. Among \(H\) successive terms \(r+k+i(p-1)\), at least
\(\lfloor H/p^j\rfloor\) are multiples of \(p^j\) for each
\(j\ge1\). Consequently the product of \(X_{r+1}-a_i\) has
valuation at least \(H+\sum_{j\ge1}\lfloor H/p^j\rfloor\).
This is the stated precision. Equal residues to this precision
and one valuation below it imply equal valuations. Dividing
\(p^{\lceil w\rceil}\) by the annihilator gives an element of
\(O\), which proves the exponent bound.

Finally every \(a_0+pU-a_i\) has both coefficients in \(pO\).
The substitution of \(a_0+pU\) into \(G\) converges to a
restricted integral series, since both coefficients of the
argument are divisible by \(p\). This proves the
coefficientwise assertion. For explicit numerical evaluation,
the denominator \(\prod_{j\ne i}(a_i-a_j)\) in Lagrange's
formula has valuation \(H-1+v_p(i!(H-1-i)!)\).
The other polynomial factors are integral. It suffices,
for instance, to compute each sample modulo
\(p^{2H+2v_p(H!)}O\) to retain the displayed precision.
\end{proof}

There is also a description by residue classes when \(O\) is
unramified. Write \(F(a_0+pU)=\sum_m d_mU^m\). Here
\(v_p(d_m)\ge m\), by expanding \((a_0+pU)^n\). The constant
coefficient is nonzero by the Bernoulli value and the functional
equation for its primitive character. Put
\(\mu=\min_m v_p(d_m)\), and let \(s\le\mu\) be the degree of
the nonzero reduction of \(p^{-\mu}\sum d_mU^m\).
If \(s=0\), all values on \(\ZZ_p\) have valuation \(\mu\).
For \(s>0\) and \(l\ge1\), twists of valuation
strictly greater than \(\mu+sl\) occupy at most \(s\) residue
classes modulo \(p^l\) inside the fixed branch modulo \(p-1\).

Indeed start with the normalized integral series. At an allowed
residue \(c\in\mathbf F_p\) where its reduction is nonzero all
values are units. At a root of multiplicity \(e\), first translate
by an integer lift of \(c\). The coefficient of degree \(e\) is
a unit, and all lower coefficients are divisible by \(p\).
Now substitute \(pU\). The new minimum valuation is between
one and \(e\), and after division by that power of \(p\) the
degree of the reduction is at most \(e\). At every stage the
sum of the positive degrees over the surviving discs is
therefore at most \(s\). A branch ending by stage \(l\) has
valuation at most \(\mu+sl\). There are at most \(s\) continuing
discs at stage \(l\). The map \(r\mapsto (X_{r+1}-a_0)/p\)
preserves distances on integers by the valuation formula above.
As \(p-1\) is a unit, the twists in the fixed branch run through
all residues at each depth. This proves the claim.

\subsection{Cuspidal congruences with prescribed auxiliary characters}\label{nav:III-36}
Fix an integer \(k\ge2\), a primitive Dirichlet character \(\chi\)
of conductor \(N\) with \(\chi(-1)=(-1)^k\), and a prime
\(\ell>\max(k+1,(8N)^3)\). Fix an embedding of the character
values into \(\overline{\QQ}_\ell\). Put \(h=8N\).
We use the Fourier variable \(Q=\exp(2\pi i z)\).

\begin{lemma}[Two independent cuspidal Eisenstein congruences]
\label{lem:cusps240}
There is an integer \(a\ge1\) such that the following holds for
any two distinct primes \(s_1,s_2\equiv1\bmod h\ell^a\).
At level \(\Gamma(h)\cap\Gamma_0(s_1s_2)\) there are two
independent cusp forms in residue characteristic \(\ell\), of
weight \(k\), with the transformation law of \(\chi\) at level
\(\Gamma_0(Ns_1s_2)\). Each has eigenvalue
\(1+\chi(p)p^{k-1}\) reduced modulo the maximal ideal for every
Hecke operator \(T_p\) with \(p\nmid h s_1s_2\), including \(p=\ell\).
\end{lemma}
\begin{proof}
Except when \((k,\chi)=(2,1)\), take the normalized holomorphic
Eisenstein series \(E_k(1,\chi)\) at level \(N\). Its first
Fourier coefficient is one and its coefficient at \(n\ge1\)
is \(\sum_{d\mid n}\chi(d)d^{k-1}\). Its constant terms at the
cusps of \(\Gamma(h)\) are algebraic by the algebraic
Fourier expansion principle. These standard formulas use the
classical normalization of \(T_p\). \cite{IV-II-V-newforms}
Choose \(a\) so that \(\ell^a\) times every one of these finitely
many constants has positive valuation. For \(s=s_1,s_2\), put
\(F_s(z)=E(z)-E(sz)\). At a cusp represented by the first
column \((u,v)\) of \(\gamma\in\mathrm{SL}_2(\ZZ)\), put
\(d=\gcd(su,v)\), so \(d=1\) or \(s\). The first column of
\(\operatorname{diag}(s,1)\gamma\) is \(d\) times the primitive
column \((su/d,v/d)\). Complete this column to \(\gamma'\)
of determinant one. Then \((\gamma')^{-1}\operatorname{diag}(s,1)\gamma\)
is upper triangular with diagonal \(d,s/d\). Thus the constant
of \(E(sz)|_k\gamma\) is \((d/s)^k\) times that of \(E|_k\gamma'\).
Both primitive columns agree modulo \(h\), since \(s\equiv1\bmod h\).
Their completions modulo \(h\) differ on the right by an upper
unitriangular matrix, which lifts to an integral translation.
The two constants of \(E\) are consequently equal. Multiplication
by \(1-(d/s)^k\) makes the constant vanish after reduction at
every cusp by the choice of \(a\).

For \((2,1)\) use the normalized weight-two Eisenstein series
\(E=-1/24+\sum_{n\ge1}\sigma_1(n)Q^n\) and take
\(F_s(z)=E(z)-sE(sz)\). Its weight-two transformation formula
cancels the nonmodular term, giving a form on \(\Gamma_0(s)\).
The same formula at the cusp zero gives
\(F_s|_2\left(\begin{smallmatrix}0&-1\\1&0\end{smallmatrix}\right)
=E(z)-s^{-1}E(z/s)\). Its two cusp constants are
\((s-1)/24\) and \(-(1-s^{-1})/24\), both zero after reduction.

In each case the expansion at infinity has integral coefficients.
At this prime-to-\(\ell\) level the smooth compact modular curve
with fixed Weil pairing has geometrically integral special fiber.
The expansion at one cusp suffices for integrality. Indeed a
generic-fiber holomorphic section could have a pole only on
the whole special fiber, which its integral expansion excludes.
The vanishing constants at all cusps now make its reduction
cuspidal. We use the good reduction, Tate cusp charts and
Fourier expansion principle for modular curves.
\href{https://virtualmath1.stanford.edu/~conrad/papers/kmpaper.pdf}{Conrad,
\emph{Arithmetic moduli of generalized elliptic curves}}.
The Hecke formula for the divisor sums gives the stated eigenvalues.
Replacing \(z\) by \(sz\) commutes with every indicated \(T_p\).
In weight two with trivial character the same Fourier formula
applies directly to \(F_s\). The two reductions have coefficient
one at \(Q\), while their coefficients at \(Q^{s_1}\) differ by
one up to sign, since \(s_1\equiv1\bmod\ell\). They are independent.
\end{proof}

\begin{theorem}[An ordinary newform of exact auxiliary conductor]
\label{thm:eigen240}
Keep \(s_1,s_2\) as in Lemma~\ref{lem:cusps240}.
For each \(j=1,2\) choose any Dirichlet character \(\xi_j\) of
conductor \(s_j\) and order \(\ell\). There is a normalized
cuspidal newform \(f\) of weight \(k\), exact level \(Ns_1s_2\)
and nebentype \(\chi\xi_1\xi_2\), with a place over \(\ell\) such that
\[
       a_p(f)\equiv 1+\chi(p)p^{k-1}\pmod{\mathfrak m}
                              \qquad(p\nmid h s_1s_2).
\]
In particular \(f\) is ordinary at \(\ell\).
\end{theorem}
\begin{proof}
Write \(s=s_1s_2\), \(\xi=\xi_1\xi_2\).
Work over the integers of a finite extension of \(\QQ_\ell\)
containing the relevant roots of unity and cusp parameters.
On the compact curve \(X\) of level \(\Gamma(h)\cap\Gamma_0(s)\)
the cover which chooses generators of the two cyclic subgroups
is an \'etale torsor for \(Z=(\ZZ/s\ZZ)^\times\).
On the smooth elliptic locus this follows since \(s\) is
invertible and the full level \(h\) removes automorphisms.
At a cusp \((u,v)\) a stabilizing translation of width \(w\)
has lower-left entry \(-v^2w\) and lower-right entry \(1+uvw\).
For \(p\mid s\), either \(p\mid v\) or \(p\mid w\).
Thus its lower-right entry is one modulo every \(p\mid s\).
As \(s\) is squarefree the width stays the same after choosing
both generators. In the tame Tate charts the parameter maps
therefore have degree one and the residue extensions are
\'etale. This proves the assertion on the compactification.
Descent along this torsor gives a line bundle \(L_\xi\) whose
sections have the prescribed diamond character on the cover.
It has degree zero, since its \(\ell\)-th power is trivial.
Its reduction is canonically trivial, since both characters
reduce to one. These statements follow by trivializing the
torsor \'etale locally and using the character for the transition
functions, also when \(\ell\) divides \(|Z|\).

Let \(C\) be the cusp divisor and \(\omega\) the Hodge bundle.
The reduction of integral sections of \(\omega^k(-C)\otimes L_\xi\)
injects into \(H^0(X_{\rm sp},\omega^k(-C))\). Its image has
codimension at most one. To see this, the module of integral
sections is finite free over the coefficient DVR, with rank
the generic-fiber dimension. The injection follows from
multiplication by a uniformizer on the line bundle. The two
Euler characteristics agree by Riemann--Roch, since the fibers
have the same genus and the line bundles have the same degree.
By Kodaira--Spencer, \(\Omega^1=\omega^2(-C)\). Serre duality
identifies special-fiber \(H^1\) with the dual of
\(H^0(X_{\rm sp},\omega^{2-k})\). This has dimension zero
for \(k>2\), because the modular discriminant gives a positive
cusp divisor for \(\omega^{12}\), and one for \(k=2\),
because the fiber is geometrically connected. This proves the
codimension bound.
\href{https://web.math.princeton.edu/~nmk/old/padicpropMFMS.pdf}{Katz,
\emph{\(p\)-adic properties of modular schemes and modular forms}}.
In particular some nonzero vector in the
two-dimensional space of Lemma~\ref{lem:cusps240} lifts.

The change of full \(h\)-level to \(\Gamma_0(N)\) acts through
a subgroup of \(\mathrm{SL}_2(\ZZ/h\ZZ)\) of order less than
\(\ell\). Representatives can be chosen to act trivially modulo
\(s\) by the Chinese remainder theorem and surjectivity of
reduction for \(\mathrm{SL}_2(\ZZ)\). This action commutes with
the auxiliary diamonds.
Average a lift against the inverse of \(\chi\) on this finite
group. The denominator is a unit and the reduction, already
of character \(\chi\), is preserved. We have obtained a nonzero
residual Hecke eigenvector in the reduction of the classical
cusp lattice at level \(Ns\) of character \(\chi\xi\).
The Hecke operators away from \(hs\) preserve this lattice
and its reduction agrees on the chosen vector with the usual
Fourier formula, since \(\xi\) reduces to one. This includes
\(T_\ell\), whose coefficient at \(Q^n\) is
\(a_{\ell n}+\chi\xi(\ell)\ell^{k-1}a_{n/\ell}\),
where the second coefficient is zero if \(\ell\nmid n\).

The commutative algebra generated by these Hecke operators on
the lattice is finite and torsion-free over the DVR. A minimal
prime under the residual eigencharacter has characteristic
zero because a uniformizer is a nonzero divisor. Passing to
the fraction field of this finite domain and choosing a prime
above the eigencharacter gives characteristic-zero eigenvalues
with the desired reductions. They occur on the cusp space
because the Hecke algebra acts faithfully there, so its
localization at this generic factor has a nonzero eigenspace
after scalar extension. Newform decomposition realizes that
system in a primitive form of level dividing \(Ns\).
The character \(\chi\xi\) has conductor \(Ns\), forcing equality.
At \(\ell\) the eigenvalue is one modulo the maximal ideal.
\end{proof}

\section{Relative orders through the norm map}\label{nav:g-II-4}

\subsection{Strict orders on one cyclotomic field per interval and precision}\label{nav:II-77}
Fix the auxiliary primes on an interval of
Section~\ref{sec:missing}. Let \(E_*\) be its largest
field, so \(E_x\subseteq E_*\) with relative degree a
power of two throughout the interval. For an odd \(p\) and
\(h\ge1\) put \(L_h=E_*(\mu_{p^h})\), \(F_h=L_h^+\),
\(T_h=\widetilde C(F_h)\), and
\(\Delta_{x,h}=\Gal(L_h/E_x)\). Let \(B_{x,h,r}\) be the
quotient of \(T_h\) by \(p^h\) and by the relations
\(\delta-\kappa_p(\delta)^{-r}\bmod p^h\) for
\(\delta\in\Delta_{x,h}\). Only the action on \(\mu_{p^h}\)
is used to evaluate this scalar. Write \(S_{x,r}\) for
\(\mathcal S_{p,r}(E_x)^{c=(-1)^r}\).

\begin{theorem}[All twists and all fields of one interval]\leavevmode
\label{thm:loginterval319}
For every \(x\) on the interval one has
\(S_{x,r}[p^h]\simeq\operatorname{Hom}(B_{x,h,r},\ZZ/p^h)\).
The number \(m_{x,h,r}\) of cyclic factors of \(S_{x,r}\)
of order at least \(p^h\) is
\(\dim_{\mathbb F_p}p^{h-1}B_{x,h,r}\). If
\(d_p=[E_*(\mu_p):E_*]\), then for \(t\ge1\) integral
\[
 \sum_{1\le r<t}m_{x,h,r}\le
       \left\lceil{t-1\over d_p}\right\rceil
       \dim_{\mathbb F_p}\bigl(p^{h-1}T_h/p^hT_h\bigr),
 \qquad
 \sum_{1\le r<t}\log_p|S_{x,r}|
             =\sum_{h\ge1}\sum_{1\le r<t}m_{x,h,r}.
\]
There is also a fixed-field formula for full orders. Write
\(b=\log_p\exp T_1\). Whenever \(p^b\mid r\), the full
group \(S_{x,r}\) is dual to
\[
 T_1\Big/\sum_{\delta\in\Delta_{x,1}}
                         (\delta-\omega(\delta)^{-r})T_1 ,
\]
where \(\omega\) is the Teichm\"uller lift of the action on
\(\mu_p\). These assertions include odd primes dividing
\([E_x:\QQ]\).
\end{theorem}
\begin{proof}
Apply Theorem~\ref{thm:logstrict319} first over \(E_*\).
Restriction to \(E_*\) identifies strict classes over
\(E_x\) with the \(\Gal(E_*/E_x)\)-invariants. Indeed at
the odd prime the degree is a unit, so restriction and
averaging give descent in global cohomology and injectivity
on each required local group, including the inertia groups
away from \(p\). This also follows from Lemma~\ref{lem:desc}.
All restriction maps in Theorem~\ref{thm:logstrict319} are
equivariant. Thus descending to \(E_x\) imposes precisely
the equivariance for all lifts to \(\Delta_{x,h}\), with
the indicated action on the coefficients. The quotient
formula follows. For finite abelian \(p\)-groups,
duality into \(\ZZ/p^h\) preserves the elementary divisors
of exponent at most \(p^h\). Multiplication by \(p^{h-1}\)
counts the summands of exact order \(p^h\) at this
precision. Summing over \(h\) counts the length of each
summand in the full strict group. These groups are finite
for the positive twists by Proposition~\ref{prop:finite15}.

The group \(\Gal(L_h/E_*)\) is a subgroup of
\((\ZZ/p^h)^\times\). Its subgroup of order prime to \(p\)
is cyclic of order \(d_p\) and acts through roots of unity
on the coefficient. Its character idempotents are integral
since \(d_p\) is a unit and all \(d_p\)-th roots lie in
\(\ZZ_p\). Consequently \(B_{x,h,r}\) is a quotient of
the single \(\omega^{-r}\)-component of \(T_h/p^h T_h\).
The \(d_p\) residues of \(r\) give all these components
once. Multiplying each quotient by \(p^{h-1}\) gives the
inequality by counting their occurrences up to \(t-1\).

For the fixed-field assertion take \(h=b+1\). The principal
unit part of \(\kappa_p^{-r}\) is one modulo \(p^h\) when
\(p^b\mid r\), by binomial expansion on \(1+p\ZZ_p\).
Thus Theorem~\ref{thm:logstrict319} can use \(L_1\) over
\(E_*\) at this precision. Descend through the same group
of order a power of two. Every resulting quotient has
exponent at most \(p^b\), hence the \(p^{b+1}\)-torsion
has no element of order \(p^{b+1}\). An element of any
greater order would have a multiple of that order. It
follows that this torsion already gives the full group.
Duality can now be taken into \(\QQ_p/\ZZ_p\).
\end{proof}

At fixed \(p,h\) all tests above use the same real field
\(F_h\) for every twist and every \(E_x\) on the interval.
Theorem~\ref{thm:logunits319} bounds its contribution by
ordinary class data and a determinant of local cyclotomic
logarithms. A bound on the sum of the surviving lengths
as \(p,h\) and the interval vary is still needed for
Section~\ref{sec:missing}.

\subsection{A common quotient bound for the residual twists}\label{nav:II-78}
Fix a whole interval in Section~\ref{sec:missing}, an odd
prime \(p\), and integers \(h>b\ge0\), \(T\ge2\). Use the
largest field \(E_*\) on this interval. Put
\(M=\widetilde C((E_*(\mu_{p^h}))^+)/p^h\), as in
Theorem~\ref{thm:loginterval319}, and
\(d=[E_*(\mu_p):E_*]\).
In the infinite cyclotomic extension over \(E_*\), the
principal-unit image is \(1+p^a\ZZ_p\) for an integer
\(a\ge1\). Take a generator \(u\) of this group and a
Galois element \(\gamma\) with cyclotomic value \(u\).
Write \(M_\tau\) for the summand of \(M\) on which the
finite subgroup of order \(d\) acts by the character \(\tau\).
All lengths below are over \(\ZZ_p\).

For each \(\tau\) list in increasing order the weights
\(1\le r<T\) with tame character \(\omega^{-r}=\tau\)
on this subgroup. The list is an arithmetic progression
\(r_i=r_0+di\), \(0\le i<k_\tau\). For a nonempty list put
\[
     w_i=ai+v_p(i!),\qquad D_\tau=w_{k_\tau-1}.
\]
For a field \(E_x\) on the interval let
\(S_{x,r}=\mathcal S_{p,r}(E_x)^{c=(-1)^r}\).

\begin{theorem}[All residual lengths at one finite precision]\leavevmode
\label{thm:jointlog322}
If \(b\ge D_\tau\) for every nonempty list, then
\[
 \sum_{1\le r<T}\max_x
 \len_{\ZZ_p}\big((S_{x,r}/S_{x,r}[p^b])[p^{h-b}]\big)
 \ \le\ \sum_{\tau:k_\tau>0}
                    \len_{\ZZ_p}\big(p^{b-D_\tau}M_\tau\big).
\]
In particular it bounds the same sum for the surviving
direct summands cut out by \(1-J_{p,r}\) in
Section~\ref{sec:missing}. At a point \(x\) of every
sufficiently distant common interval one may use that
section's \(T=\lfloor x^{401}\rfloor\) and its original
cutoff \(b=b_0(x)\) for \(p\le P(x)\), \(b=0\) for \(p>P(x)\).
The exponents \(w_0,\ldots,w_{k_\tau-1}\) are the exact
elementary divisor valuations of polynomial evaluation at
the full list of wild characters, before reduction modulo \(p^h\).
\end{theorem}
\begin{proof}
Because \(p\) is odd, its cyclotomic image is the direct
product of the tame subgroup and the stated principal-unit
group. Indeed the image is closed of finite index in
\(\ZZ_p^\times\); taking limits of powers separates its
prime-to-\(p\) and pro-\(p\) parts. Character averages split
\(M\) over \(\ZZ_p\), since \(d\mid p-1\).
For each weight choose any field \(E_{x_i}\) from the interval.
By Theorem~\ref{thm:loginterval319}, the finite group
\(S_{x_i,r_i}[p^h]\) is dual to a quotient \(B_i\) of
\(M_\tau\). On this quotient \(\gamma\) acts as the scalar
\(z_i=u^{-r_i}\) modulo \(p^h\). The same element acts for
every field, since it fixes \(E_*\). On cyclic factors one
checks that the length in the theorem is \(\len(p^bB_i)\):
a factor of exponent \(e\) contributes
\(\max(0,\min(e,h)-b)\) to either length.

For \(i\ge0\) set \(P_i(Z)=\prod_{l<i}(Z-z_l)\), with
empty product one. For distinct integral indices \(j,l\)
one has \(v_p(z_j-z_l)=a+v_p(j-l)\).
In fact raising a principal unit of depth \(a\) to an
exponent prime to \(p\) leaves its depth unchanged, and
raising to \(p\) increases it by one, by binomial expansion
for odd \(p\). Also \(p\nmid d\). It follows that for
\(j\ge i\) the value \(P_i(z_j)\) has valuation at least
\(w_i\), with equality for \(j=i\): the product of the
\(i\) successive nonzero integers ending at \(j\) is a
multiple of \(i!\), equal to \(i!\) at \(j=i\).
For \(j<i\) the value is zero. Changing the monomial basis
of degrees below \(k_\tau\) to the monic polynomials \(P_i\)
is integral unimodular. In their evaluation matrix, dividing
column \(i\) by \(p^{w_i}\) gives a lower triangular integral
matrix with units on its diagonal. This proves the elementary
divisor assertion, since the \(w_i\) are nondecreasing.

Let \(Y\) be the image of \(p^{b-D_\tau}M_\tau\) in
\(\bigoplus_i B_i\). Filter \(Y\) by requiring its first
\(i\) coordinates to vanish. The quotient at stage \(i\)
is the image of the \(i\)-th coordinate on that subgroup.
For any \(m\in M_\tau\), the image of
\(p^{b-w_i}P_i(\gamma)m\) belongs to this subgroup since
\(w_i\le D_\tau\). Its \(i\)-th coordinate is \(p^b\)
times a unit times the image of \(m\). These coordinates
therefore contain \(p^b B_i\), including when the power is
zero in the quotient. Additivity of length gives the proposed
bound on this list. The fields were chosen independently,
so we may maximize the individual lengths and add over the
lists. The possible \(E_x\) form a finite chain. Passing
to a direct summand of a finite abelian group decreases the
displayed truncated length.

Finally the \(p\)-part of the index of the cyclotomic image
is \(p^{a-1}\), bounded by \([E_*:\QQ]\) by restriction
from the absolute Galois group of \(\QQ\). This field degree
is polynomially bounded in \(V\) on \([V,V^{3/2}]\):
the common biquadratic factors have bounded degree and the
dyadic real layer has degree at most \(2(V^{3/2})^{2+8\eps}\)
in the field construction. Thus
\(D_\tau\le T(1+\log_p[E_*:\QQ]+1/(p-1))\), eventually
less than \(b_0(x)=\lfloor\exp(x^\eps/8)\rfloor\) uniformly
for the small primes. For \(p>P(x)=\exp(x^\eps/12)\),
\(d\ge(p-1)/[E_*:\QQ]>T\) eventually. Those lists contain
at most one term, so their losses are zero.
\end{proof}

The result counts the twists together without a multiplicity
factor in front of the length of a common real class group.
A bound for the right side uniform in the precision, with the
weighted sum over \(p\) required in Section~\ref{sec:missing},
is still needed to control the full residual orders.

\subsection{Full relative strict orders in progressions of twists}\label{nav:II-79}
Fix one of the intervals of Section~\ref{sec:missing} and its
largest field \(E_*\). Fix an odd prime \(p\) and a subfield
\(F_0\subseteq E_*\) of relative degree a power of two. For
\(F_0\subseteq F\subseteq E\subseteq E_*\), write \(S_r(E)\)
for the sign \((-1)^r\) summand of the strict group. Put
\[
 U_r(E/F)=(1-e_F)S_r(E),\qquad
 e_F=|\operatorname{Gal}(E/F)|^{-1}
                         \sum_{\delta\in\operatorname{Gal}(E/F)}\delta .
\]
All the fields here are abelian over \(\QQ\). In particular
complex conjugation acts on their cohomology.

Use \(L_h=E_*(\mu_{p^h})\) and its real logarithmic class
group \(T_h=\widetilde C(L_h^+)\) from
Theorem~\ref{thm:loginterval319}. For a subfield \(A\)
between \(F_0\) and \(E_*\), let \(B_h(A,u)\) be the quotient
of \(T_h/p^h\) by all the relations
\(\delta-\kappa_{p,h}(\delta)^{-u}\) for
\(\delta\in\operatorname{Gal}(L_h/A)\), where \(\kappa_{p,h}\)
is the action on \(\mu_{p^h}\). Define
\[
 D_h(E/F,u)=\ker\bigl(B_h(E,u)\longrightarrow B_h(F,u)\bigr),
 \qquad
 M_h=\operatorname{lcm}\bigl(2,\,
      |\kappa_{p,h}(\operatorname{Gal}(L_h/F_0))|\bigr).
\]
The arrow imposes the extra relations over \(F\).

\begin{theorem}[Full strict orders from one relative quotient]\leavevmode
\label{thm:relativeprogress325}
Let \(u\ge1\) and \(h\ge1\). If
\(p^{h-1}D_h(E_*/F_0,u)=0\), then for every positive
\(r\equiv u\bmod M_h\) and every \(F_0\subseteq F\subseteq
E\subseteq E_*\), the full group \(U_r(E/F)\) has the
elementary divisors of \(D_h(E/F,u)\). For every fixed \(u\)
such an \(h\) exists. For \(h=1\) the test is just the
vanishing of the relative quotient modulo \(p\).
For the whole \(S_r(E)\) the same statements use \(B_h(E,u)\)
in place of the relative quotients and the test at \(E_*\).
\end{theorem}
\begin{proof}
Restriction to \(E_*\) identifies \(S_r(A)\) with the
\(\operatorname{Gal}(E_*/A)\)-invariants in \(S_r(E_*)\).
Indeed averaging is exact in degree a power of two at the odd
prime \(p\). Inflation--restriction gives this assertion in
global degree one. Restriction on the local groups above
\(p\) is injective by corestriction, so the criterion of
local triviality descends as well. The same argument applies
to the unramified conditions outside \(p\). All maps commute
with the sign projection. This also proves the assertion for
a real subfield, where the negative sign part is zero.

The finite-precision identification of
Theorem~\ref{thm:loginterval319} thus identifies \(S_u(A)[p^h]\)
with the characters of \(B_h(A,u)\) into \(\ZZ/p^h\).
One may fix a generator for the coefficient module at this
precision. Under restriction from \(F\) to \(E\) a character
of \(T_h\) is the same character subject to fewer equivariance
conditions. Dualizing restriction gives exactly the quotient
arrow above. Since \(e_F\) is an integral projection,
\(D_h(E/F,u)\) is dual to \(U_u(E/F)[p^h]\). In particular
the test at \(E_*/F_0\) holds precisely when this largest
relative group has no element of order \(p^h\). For a finite
abelian \(p\)-group this means that its full exponent divides
\(p^{h-1}\). By finiteness at the positive twist \(u\) in
Proposition~\ref{prop:finite15}, we can choose \(h\) above
the exponent valuation.

If \(r\equiv u\bmod M_h\), all the defining scalars are
identical at precision \(h\), by Lagrange's theorem on the
image of \(\kappa_{p,h}\), and the sign agrees as well.
Consequently every quotient and every arrow just used is
identical at that precision. The test now bounds the full
exponent at \(r\) over \(E_*/F_0\). Under restriction each
\(U_r(E/F)\) lies in \(U_r(E_*/F_0)\), since averaging first
over \(\operatorname{Gal}(E_*/F)\) and then over the larger
group fixing \(F_0\) gives the latter average. Hence the
same exponent bounds every smaller relative group. Its
subgroup killed by \(p^h\) is the whole group, and the
dual identification gives the claimed elementary divisors.
Omitting the two averages gives the assertion for \(S_r(E)\).
\end{proof}

For the small-twist range one can take \(F_0\) to be the
smallest \(L_x\) on the interval and use \(E=E_x,F=L_x\).
For the large twists at \(p\in\Pi\), use the smallest
\(B(x)\) and \(F=B(x)\). For \(J_{p,r}=0\) use the
whole-group version. At a given \(x\), let \(N_x\) count
the integers in the applicable range \(1\le r<t(x)\)
in a progression tested by the theorem. Put \(b=b_0(x)\)
for \(p\le P(x)\) and \(b=0\) otherwise. Its exact
contribution to the residual length sum of
Proposition~\ref{prop:res23} is
\[
       N_x(\log p)\,\len_{\ZZ_p}\bigl(p^b D_h(E_x/F,u)\bigr),
\]
with \(B_h(E_x,u)\) in the whole-group case. This follows
on each cyclic factor by removing its subgroup killed by
\(p^b\). The relative test can succeed even when the
invariant summand has larger exponent. A bound for the
weighted sum over all progressions and primes on these
intervals is still required.

\subsection{What two class groups determine about every twist}\label{nav:II-80}
Fix the auxiliary primes and let \(E_*\) be the largest CM field on
one interval of Section~\ref{sec:missing}. Fix an odd prime \(p\).
Write \(L_\infty=E_*(\mu_{p^\infty})\), \(F_\infty=L_\infty^+\)
and \(F_0=(E_*(\mu_p))^+\). Let \(F_j/F_0\) be the subextension
of degree \(p^j\) in \(F_\infty\). This is the cyclotomic
\(\ZZ_p\)-extension of \(F_0\). Indeed
\(L_\infty=E_*(\mu_p)\QQ_\infty\), where \(\QQ_\infty\) is
the cyclotomic \(\ZZ_p\)-extension of \(\QQ\). Taking real parts
gives \(F_\infty=F_0\QQ_\infty\), since the quadratic extension
\(E_*(\mu_p)/F_0\) is imaginary. Put \(H=\Gal(F_\infty/F_0)\).
Via restriction, a topological generator \(\gamma\) of \(H\)
has a unique lift to \(\Gal(L_\infty/E_*(\mu_p))\). Write
\(a\in 1+p\ZZ_p\) for its cyclotomic character.

Take the finite prime-to-\(p\) subgroup \(D'\) of a decomposition
group at \(p\) in \(\Gal(F_\infty/\QQ)\). It is the same for
every prime above \(p\), by commutativity. Set
\[
 e=1-\frac1{|D'|}\sum_{g\in D'}g,\qquad
 C_j=e\,\Cl(F_j)\{p\}\quad(j=0,1).
\]
Braces denote the primary subgroup. Let \(\Delta_0\) be the tame
subgroup of \(\Gal(L_\infty/E_*)\). The cyclotomic character
identifies it with a subgroup of the \((p-1)\)-st roots of unity
in \(\ZZ_p^\times\). For a character \(\tau\) of \(\Delta_0\)
with these values, use the usual integral character idempotent
\(e_\tau\), acting on the class groups through restriction to the
real field. Define
\[
 C_{j,\tau}=e_\tau C_j,\qquad
 K_\tau=\ker\!\left(\mathrm N:C_{1,\tau}\longrightarrow C_{0,\tau}\right),
 \qquad b_\tau=\dim_{\mathbf F_p}(K_\tau/pK_\tau).
\]

For \(E=E_x\) on the interval, let \(A_r(E)\) be the Pontryagin
dual of the complementary strict group of Theorem~\ref{thm:loginterval319},
in its \emph{untwisted class group coordinates}. Thus the relation
for \(\delta\in\Gal(L_\infty/E)\) in these coordinates is action by
\(\kappa_p(\delta)^{-r}\). Put \(B_r(E)=eA_r(E)\),
\(\tau=\kappa_p^{-r}|_{\Delta_0}\) and \(s_r=v_p(a^{-r}-1)\).

\begin{theorem}[Two ordinary class groups and all positive twists]\leavevmode
\label{thm:twolayers327}
For every \(E_x\) on this interval and every positive \(r\),
the norm defining \(K_\tau\) is onto, \(B_r(E_x)/p^{s_r}B_r(E_x)\)
is a quotient of \(C_{0,\tau}\), and at most \(b_\tau\) cyclic
factors of \(B_r(E_x)\) have order greater than \(p^{s_r}\).
If \(K_\tau=0\), the full group in question is
\[
 B_r(E_x)\simeq
 C_{0,\tau}\big/
 \overline{\sum_{\delta\in\Gal(L_\infty/E_x)}
      (\delta-\kappa_p(\delta)^{-r})C_{0,\tau}} .
\]
The finite group on the right uses the action on \(F_0\).
The same assertions hold after imposing any further integral
character idempotent commuting with the field action, applied
to all the modules in the statements.
\end{theorem}
\begin{proof}
Let \(\mathcal N\) be the maximal abelian pro-\(p\) extension of
\(F_\infty\) unramified away from \(p\) and split at every prime
over \(p\). Write \(X=\Gal(\mathcal N/F_\infty)\). Conjugation acts
on \(X\) through \(\Gal(F_\infty/\QQ)\), since the local conditions
are preserved by that group and \(X\) is abelian. For each \(j\)
there is a canonical isomorphism
\[
 X/(\gamma^{p^j}-1)X \simeq \widetilde C(F_j).
\]
To see this, the Galois group of \(\mathcal N/F_j\) is an
extension of the procyclic group \(H^{p^j}\) by \(X\). It is
pro-\(p\). A lift of a topological generator defines a continuous
section because its closed cyclic subgroup maps isomorphically
onto \(H^{p^j}\). The commutator subgroup after closure is
\((\gamma^{p^j}-1)X\), which is already compact. The maximal
abelian subextension allows at \(p\) only subextensions of the
local cyclotomic tower, and is unramified elsewhere. Conversely
any finite abelian \(p\)-extension of \(F_j\) with these
conditions becomes split at \(p\) and unramified elsewhere
after adjoining \(F_\infty\). Its compositum with \(F_\infty\)
is abelian over \(F_\infty\), so it lies in \(\mathcal N\).
Taking the kernel of the map to \(H^{p^j}\) gives the stated
isomorphism by the definition of the degree-zero logarithmic
class group. Inclusion of Galois groups induces the quotient
map on \(X\) when \(j\) decreases. Under class field reciprocity
this inclusion on abelianizations is the norm on ideles.

At each finite level \(D'\) fixes every prime over \(p\). Its
action on each local cyclotomic degree group is trivial because
the cyclotomic Galois group is abelian. In the exact sequence of
Theorem~\ref{thm:logunits319}, \(e\) therefore kills the entire
local kernel and the quotient of the global degree group.
Further, inverting \(p\) in the ordinary ideal class group kills
precisely the subgroup generated by classes of primes above
\(p\), which \(e\) also annihilates. Projection by an integral
idempotent is exact. These identifications give, compatibly
with norms, for \(M=e_\tau eX\),
\[
 C_{0,\tau}=M/(\gamma-1)M,\qquad
 C_{1,\tau}=M/(\gamma^p-1)M .
\]
In particular the norm is onto. Put \(U=(\gamma-1)M\) and
\(N_\gamma=1+\gamma+\cdots+\gamma^{p-1}\).
The norm kernel is canonically \(U/N_\gamma U\).

Theorem~\ref{thm:loginterval319} identifies \(B_r(E_x)\)
with the quotient of \(M\) by all the displayed character
relations for \(\Gal(L_\infty/E_x)\). In detail, modulo \(p^h\)
the kernel of the action on \(E_*(\mu_{p^h})\) has trivial
twisting scalar. Its relation on \(X\) gives precisely the
logarithmic group of the corresponding finite real layer by
the preceding isomorphism. The remaining relations are those
of that theorem. Projecting by \(e_\tau\) retains the whole
\(e\)-part of the quotient because \(\Delta_0\) already acts
there by \(\tau\). Passing over all \(h\) recovers the compact
quotient with the closed relation subgroup. Its finite
quotients of \(p\)-power exponent separate its points, and
the strict group is finite by Proposition~\ref{prop:finite15}.

Under \(M\to B=B_r(E_x)\), the image of \(U\) is
\((a^{-r}-1)B=p^{s_r}B\). The image of \(N_\gamma U\) is
\(p^{s_r+1}B\), since \(1+a^{-r}+\cdots+a^{-(p-1)r}\) has
valuation one. Indeed \(a^{-r}=1+z\), \(z\in p\ZZ_p\),
and the sum is \(p+\binom p2 z+\cdots+z^{p-1}\),
congruent to \(p\bmod p^2\) for odd \(p\).
It follows that \(B/p^{s_r}B\) is a quotient of \(C_{0,\tau}\),
and \(p^{s_r}B/p^{s_r+1}B\) is a quotient of \(K_\tau/pK_\tau\).
The latter dimension counts the indicated cyclic factors exactly.

If \(K_\tau=0\), then \(U=N_\gamma U\). On any finite continuous
\(H\)-stable quotient of \(U\), \(N_\gamma\) acts nilpotently.
Modulo \(p\) it acts as \((\gamma-1)^{p-1}\), and the action
of \(\gamma\) has \(p\)-power order. A power of \(N_\gamma\)
therefore maps the finite quotient into its multiple by \(p\),
and iteration kills it. Surjectivity forces the quotient to
vanish. These quotients separate points: intersect a separating
open subgroup with all its \(H\)-translates. The orbit of the
subgroup is finite by continuity and compactness (an open
subgroup of \(H\) preserves it, by applying compactness to
the subgroup of \(U\) and then using inverses).
Thus \(U=0\), so \(M=C_{0,\tau}\) with trivial action of \(H\).
The full quotient formula follows. Commuting integral projectors
preserve all the quotient maps and identities above.
\end{proof}

\subsection{Real logarithmic class groups at arbitrary precision}\label{nav:III-74}
Let \(p\) be odd. For a number field \(D\), let \(X(D)\)
be the maximal abelian pro-\(p\) Galois quotient unramified
outside \(p\) in which the decomposition map at each \(w\mid p\)
factors through the local cyclotomic \(\ZZ_p\)-extension.
Let \(\Gamma_D\) be the Galois group of the cyclotomic
\(\ZZ_p\)-extension of \(D\). Its character is the principal
unit part \(\langle\kappa_p\rangle\) of the cyclotomic
character. The quotient \(X(D)\to\Gamma_D\) is surjective.
Write \(\widetilde C(D)\) for its kernel, the logarithmic
class group of degree zero at \(p\). It is finite for
fields abelian over \(\QQ\), by the abelian case of the
Gross--Kuz'min theorem. For the identification with
logarithmic classes and this finiteness result see
\href{https://www.math.u-bordeaux.fr/~jjaulent/Articles/bnflog.pdf}
{K.~Belabas and J.-F.~Jaulent, \emph{The logarithmic class
group package in PARI/GP}, Section~1}. The connection of
these classes with twists of wild kernels is developed by
\href{https://arxiv.org/abs/0801.0919}{J.-F.~Jaulent and
A.~Michel, \emph{Approche logarithmique des noyaux
\'etales sauvages des corps de nombres}}.

Fix an abelian field \(E/\QQ\), \(r\ge1\), \(h\ge1\),
and \(R_h=\ZZ/p^h\). Let \(L\) be a finite cyclotomic
extension of \(E\) inside \(E(\mu_{p^\infty})\), Galois
over \(\QQ\), containing \(\mu_p\), such that
\(\kappa_p^{-r}\bmod p^h\) is trivial on \(G_L\).
One may always take \(E(\mu_{p^h})\). Put \(F=L^+\),
\(\Delta=\Gal(L/E)\), and let \(\theta:\Delta\to R_h^\times\)
be the induced character \(\kappa_p^{-r}\bmod p^h\).
Use compatible cyclotomic coordinates on the coefficients.

\begin{theorem}[The complementary strict group on a real field]\leavevmode
\label{thm:logstrict319}
Let \(\Delta\) act on \(\widetilde C(F)\) by restriction
of its automorphisms to \(F\). Form the finite module
\[
 B_{h,r}=
 {\widetilde C(F)\over
   p^h\widetilde C(F)+
     \sum_{\delta\in\Delta}(\delta-\theta(\delta))\widetilde C(F)} .
\]
In this expression scalar residues can be lifted arbitrarily
to \(\ZZ_p\) since the denominator contains \(p^h\widetilde C(F)\).
With \(c\) complex conjugation there is an isomorphism
\[
    \bigl(\mathcal S_{p,r}(E)^{c=(-1)^r}\bigr)[p^h]
             \simeq\operatorname{Hom}(B_{h,r},R_h).
\]
The formula allows \(p\) to divide \(|\Delta|\) and \([E:\QQ]\).
\end{theorem}
\begin{proof}
Let \(G\) be the Galois group over \(E\) unramified outside
\(p\), allowing all archimedean places, and put
\(N=\ker(\kappa_p|G)\), \(J=G/N\). The image \(J\) is a
closed subgroup of \(\ZZ_p^\times\) with an open procyclic
pro-\(p\) part. For \(D_h=R_h(\kappa_p^{-r})\), restriction
induces
\[
 H^1(G,D_h)/\operatorname{inf}H^1(J,D_h)
                      \simeq\operatorname{Hom}(N,D_h)^J .
\]
Indeed \(N\) acts trivially and \(H^2(J,D_h)=0\) in
inflation--restriction. For the latter fact a copy of
\(\ZZ_p\) has the length-one completed group-ring resolution
given by multiplication by \(\gamma-1\), followed by
augmentation. Identifying its completed group ring with
\(\ZZ_p[[T]]\) sends \(\gamma-1\) to \(T\), proving exactness.
Passing to \(J\) takes invariants under a finite group of
order prime to \(p\), an exact operation by averaging.

Let \(\mathfrak X\) be the maximal abelian pro-\(p\) quotient
of \(N\) after killing the images of the local cyclotomic
kernels at all places over \(p\) and their conjugates in
\(G\). Lemma~\ref{lem:period30} now identifies the full
strict \(p^h\)-torsion with
\(\operatorname{Hom}_J(\mathfrak X,D_h)\). In detail the
finite class must restrict to zero on each such kernel,
where the coefficient action is trivial. The quotient by
global cyclotomic classes in that lemma is precisely the
kernel of restriction to \(N\) above. Outside \(p\) the
unramified condition is already imposed in \(G\).

Over \(L\) the same \(N\) and \(\mathfrak X\) occur.
In fact \(L/E\) lies in the cyclotomic tower and is
unramified away from \(p\). The maximal extensions
unramified outside \(p\) over the two fields coincide,
as is seen by taking the compositum of the conjugates over
\(E\) of each finite extension over \(L\). At a place
over \(p\) the absolute local cyclotomic kernel already
fixes \(L_w\). Ranging over all places of \(L\) gives
the same collection of kernels up to conjugation.
Write \(J_L=\kappa_p(G_L)\). Since \(D_h\) is trivial
over \(L\), the \(J_L\)-invariant homomorphisms out of
\(\mathfrak X\) identify naturally with
\(\operatorname{Hom}(\widetilde C(L),R_h)\) in the
untwisted coordinates. To see this, apply the same
restriction calculation over \(L\) to the trivial finite
module. Its numerator consists of characters on \(X(L)\),
as each local character factors through the cyclotomic
image at \(p\). Its denominator consists of the pullbacks
of characters on \(\Gamma_L\). Restriction from \(X(L)\)
to its degree-zero kernel realizes their quotient, since
\(X(L)\to\Gamma_L\simeq\ZZ_p\) splits as a map of
abelian pro-\(p\) groups by lifting a topological generator.
The restriction map is canonical and is equivariant under
field automorphisms. Requiring equivariance under all of
\(J\) imposes \(f(\delta z)=\theta(\delta) f(z)\) for
\(\delta\in J/J_L=\Delta\). This argument takes no average
over \(\Delta\).

On the coefficient, complex conjugation acts by \((-1)^r\).
Thus in the desired eigenspace the corresponding untwisted
character satisfies \(f(cz)=f(z)\). The untwisted invariant
characters of \(\widetilde C(L)\) are exactly those on
\(\widetilde C(F)\). Here are the local details. Restriction
identifies \(R_h\)-valued global characters over \(F\)
with the invariant characters over \(L\), since \(L/F\)
has order two. At each place away from \(p\) its inertia
subgroup after restriction has index at most two in the
original inertia subgroup. A character zero on the former
is zero on both. At a place over \(p\), intersect the
kernel of \(\langle\kappa_p\rangle\) with the local
absolute group of \(L\). This intersection again has
index at most two in that kernel, proving descent of the
local condition in both directions. The pro-\(p\) global
cyclotomic images for \(L,F\) agree by the same index
argument. Taking invariants is exact at the odd prime
\(p\), so the identification remains valid on quotienting
by the global cyclotomic characters. All these maps commute
with \(\Delta\). Factoring a homomorphism to \(R_h\) with
the stated equivariance through \(B_{h,r}\) now gives the
required isomorphism.
\end{proof}

\subsection{An interval bound outside the factors detected by the norm kernel}\label{nav:III-76}
Use the fields and the class groups of Theorem~\ref{thm:twolayers327}
on \([V,V^{3/2}]\), with the same pair of auxiliary primes that
occurs in the periods. For each odd prime \(p\) write \(B_{p,r}\)
for the summand \(B_r(E_x)\) in that theorem, and put
\[
 t=\lfloor x^{401}\rfloor,\qquad
 Y=x^{-3}\exp(x^\varepsilon/4),\qquad n_*=[E_*:\QQ].
\]
Write the elementary divisor exponents of \(B_{p,r}\) in decreasing
order as \(v_1(p,r),v_2(p,r),\ldots\), padded by zeros. For each
\(p,r\) use \(b_{p,\tau}=\dim_{\mathbf F_p}K_\tau/pK_\tau\)
and \(s_{p,r}=v_p(a^{-r}-1)\) from that theorem. Define the
nonnegative integer
\[
 D_{p,r}=\sum_i\min\{v_i(p,r),s_{p,r}\}.
\]

\begin{theorem}[A uniform bound for the part determined by the first layer]\leavevmode
\label{thm:normsum327}
At every sufficiently large \(x\) on the common intervals,
\[
 \sum_{\substack{2<p\le Y\\p\ {\rm prime}}}
  (\log p)\sum_{1\le r<t} D_{p,r}
  \ll_{\psi,\varepsilon} x^{-2}\exp(x^\varepsilon/2).
\]
The difference between the left side and the corresponding
sum of the full logarithmic orders of \(B_{p,r}\), taken as
full sum minus displayed sum, is exactly
\[
 \sum_{2<p\le Y}(\log p)\sum_{1\le r<t}
       \sum_{i=1}^{b_{p,\tau}}(v_i(p,r)-s_{p,r})_+ ,
       \qquad \tau=\kappa_p^{-r}|_{\Delta_0}.
\]
Here the outer sum is over primes and \(z_+=\max(z,0)\).
Consequently entire summands satisfying the two-layer
norm isomorphism test enter the first bound simultaneously
at all their positive twists below \(t\).
\end{theorem}
\begin{proof}
For a fixed prime, \(D_{p,r}\log p\) is the logarithm
of the order of \(B_{p,r}/p^{s_{p,r}}B_{p,r}\).
Theorem~\ref{thm:twolayers327} bounds it by
\(\log |C_{0,\tau}|\). The integral character idempotents
for \(\Delta_0\) are orthogonal and sum to one, so summing
\(\log |C_{0,\tau}|\) once over all characters gives
\(\log |e\,\Cl(F_0)\{p\}|\). As \(r\) varies, every character
occurs at most \(1+t/d_p\) times, with
\[
 d_p=|\Delta_0|=[E_*(\mu_p):E_*]\ge(p-1)/n_*.
\]
This degree bound follows by intersecting \(E_*\) with
\(\QQ(\mu_p)\). The count therefore uses a single ordinary
class number at each prime, even though the field in the
period varies along the interval.

Write \(N_*\) for a cyclotomic modulus for \(E_*\) as in
Theorem~\ref{thm:conductor232}. Since \(F_0\subseteq
\QQ(\zeta_{pN_*})\), its degree \(d\) is at most \(n_*p\)
and the conductor--discriminant formula gives
\(|\operatorname{disc}(F_0)|\le(pN_*)^d\).
Minkowski's theorem supplies an integral ideal of norm at most
\(A=(4/\pi)^d|\operatorname{disc}(F_0)|^{1/2}\) in every
ideal class. The number of such ideals is at most
\(A^2\zeta_{F_0}(2)\); this follows by bounding each indicator
by \(A^2/(\mathrm N\mathfrak a)^2\). The Euler product gives
\(\zeta_{F_0}(2)\le\zeta(2)^d\) because the number of prime
ideals over a rational prime is at most \(d\).
In particular
\[
 \log|\Cl(F_0)|\ll n_* p\log(2pN_*).
\]
Summing over the integers from 3 to \(Y\) bounds the desired
sum by \(O((n_*Y^2+n_*^2tY)\log(2N_*Y))\), using
\(p/(p-1)\le 2\). On the specified interval,
\(n_*\ll_\psi V^{3+12\varepsilon}\le C_\psi x^{3+12\varepsilon}\).
This follows from the dyadic degree at its upper endpoint and
the fixed degree of the other factors. The modulus consists
of the same fixed character conductor, two fixed auxiliary
primes, and a power of two of polynomial size in \(V\).
Thus \(\log(2N_*Y)\ll_{\psi,\varepsilon}x^\varepsilon\),
by the conductor bounds of Section~\ref{sec:missing}.
The first summand is
\(O_{\psi,\varepsilon}(x^{-3+13\varepsilon}
  \exp(x^\varepsilon/2))\).
The second is bounded by a fixed power of \(x\) times
\(\exp(x^\varepsilon/4)\). Both satisfy the asserted bound
since \(0<\varepsilon<1/52\).

Finally the same two-layer theorem gives
\(v_i(p,r)\le s_{p,r}\) for \(i>b_{p,\tau}\).
Subtracting the truncated length on each cyclic factor proves
the exact expression for the excess. If \(K_\tau=0\) the
inner sum is empty.
\end{proof}

The estimate covers the summand away from the trivial action
of the prime-to-\(p\) decomposition group. It leaves the excess
orders in the formula, the invariant summand, and primes above
\(Y\) to the other arithmetic estimates of Section~\ref{sec:missing}.

\subsection{The strict group at each cyclotomic norm level}\label{nav:specialization329}
Fix a power interval and use its largest field \(E_*\), as in
Section~\ref{sec:missing}. Fix an odd prime \(p\) and write
\(L_\infty=E_*(\mu_{p^\infty})\), \(F_\infty=L_\infty^+\),
\(F_0=E_*(\mu_p)^+\). Let \(F_j/F_0\) be the layer of degree
\(p^j\) in \(F_\infty\). Write \(T_j=\widetilde C(F_j)\) for
the degree-zero logarithmic class group. Here the subscript counts
layers above \(F_0\). For \(E=E_x\), put \(G_x=\Gal(L_\infty/E)\)
and let \(A_{x,r}\) be the Pontryagin dual of
\(\mathcal S_{p,r}(E)^{c=(-1)^r}\), in the untwisted class group
coordinates of Theorem~\ref{thm:loginterval319}. Thus the character
on these coordinates is \(\eta_r(\delta)=\kappa_p(\delta)^{-r}\).
For a compact module \(T\) with the indicated action set
\[
 Q_{x,r}(T)=T\big/
 \overline{\sum_{\delta\in G_x}(\delta-\eta_r(\delta))T}.
\]
On \(T_j\) the field action is the one on \(F_j\). The scalar uses
the indicated lift to \(L_\infty\).

\begin{theorem}[The exact strict quotient at a fixed norm level]\leavevmode
\label{thm:strictdigits329}
Let \(a\) be the cyclotomic value of a topological generator \(\gamma\)
of \(\Gal(L_\infty/E_*(\mu_p))\). Write \(s_r=v_p(a^{-r}-1)\).
For every positive \(r\), every \(j\ge0\), and every field \(E_x\)
on the interval there is an isomorphism of finite abelian groups
\begin{equation}\label{eq:strictdigits329}
 Q_{x,r}(T_j)\simeq A_{x,r}/p^{s_r+j}A_{x,r}.
\end{equation}
In particular the base-\(p\) logarithm of
\( |Q_{x,r}(T_{j+1})|/|Q_{x,r}(T_j)|\) counts the cyclic factors
of \(A_{x,r}\) of order greater than \(p^{s_r+j}\).
If the two orders agree, \(Q_{x,r}(T_j)\) already gives the full
order and all the elementary divisors of \(A_{x,r}\).
This applies to every tame character, including the invariant part
for the prime-to-\(p\) decomposition group.
\end{theorem}

\begin{proof}
The intersection of \(E_*(\mu_p)\) with the field of all \(p\)-power
roots of unity has degree \(p^u\) over \(\QQ(\mu_p)\) for some
\(u\ge0\). Indeed the cyclotomic group over \(\QQ(\mu_p)\) is
\(1+p\ZZ_p\simeq\ZZ_p\). Its subgroup of index \(p^u\) is
\(1+p^{u+1}\ZZ_p\). Thus \(v_p(a-1)=u+1\), and \(p^u\)
divides \([E_*(\mu_p):\QQ(\mu_p)]\le [E_*:\QQ]\).
In the tower one consequently has
\[
 E_*(\mu_{p^{u+1+i}})^+=F_i\qquad(i\ge0).
\]
The extension before taking real parts has degree \(p^i\) over
\(E_*(\mu_p)\). Taking real parts preserves that degree because
both fields are CM. If \(z\in p\ZZ_p\) has valuation \(d\), the
binomial expansion shows that \((1+z)^p-1\) has valuation
\(d+1\), and that raising \(1+z\) to an exponent prime to \(p\)
preserves \(d\). Inverting a unit preserves the valuation of its
difference from one. It follows that
\[
 s_r=u+1+v_p(r),\qquad
 v_p(a^{-rp^j}-1)=s_r+j.
\]

For use with the norm, let \(X\) be the Galois group of the maximal
abelian pro-\(p\) extension of \(F_\infty\) unramified away from
\(p\) and split at the places over \(p\). The conjugation action is
well defined on this abelian group. Recall the class field argument
in Theorem~\ref{thm:twolayers327}. Over \(F_i\) the group upstairs
is an extension of \(\langle\gamma^{p^i}\rangle\simeq\ZZ_p\)
by \(X\). A lift of its topological generator supplies a section.
The closed commutator subgroup is therefore
\((\gamma^{p^i}-1)X\). Abelianizing, then taking the kernel of the
map to the cyclotomic group, gives \(T_i=X/(\gamma^{p^i}-1)X\).
In fact a finite abelian extension with the indicated local
cyclotomic conditions splits at the places over \(p\) after
adjoining \(F_\infty\), and the converse follows by restriction
at each completion. Under reciprocity, inclusion from the higher
field gives the norm. Hence for \(k\ge j\) the norm \(T_k\to T_j\)
is onto with kernel \((\gamma^{p^j}-1)T_k\).

Set \(h=s_r+j\) and \(k=j+v_p(r)\). Then
\(E_*(\mu_{p^h})^+=F_k\). The finite-coefficient description in
Theorem~\ref{thm:loginterval319}, dualized using
\(1\mapsto p^{-h}\bmod\ZZ_p\) on \(\ZZ/p^h\), identifies
\(A_{x,r}/p^h A_{x,r}\) with \(Q_{x,r}(T_k)/p^h Q_{x,r}(T_k)\).
All lifts of a given automorphism of \(E_*(\mu_{p^h})\) give
the same scalar modulo \(p^h\). On this quotient the operator
\(\gamma^{p^j}-1\) acts by \(a^{-rp^j}-1\), which is zero modulo
\(p^h\). Thus passage by norm to \(T_j\) imposes precisely relations
already satisfied, giving \(Q_{x,r}(T_j)/p^h Q_{x,r}(T_j)\).
On \(T_j\) itself \(\gamma^{p^j}\) acts trivially. Its character
relation kills \((a^{-rp^j}-1)Q_{x,r}(T_j)=p^hQ_{x,r}(T_j)\).
This proves \eqref{eq:strictdigits329} with the asserted precision.

For a cyclic group of order \(p^v\), the quotient modulo \(p^h\)
has order \(p^{\min(v,h)}\). Its order increases by a factor of
\(p\) at the next precision exactly when \(v>h\). Multiplication
over the cyclic factors proves both final assertions, since
\(A_{x,r}\) is finite by Proposition~\ref{prop:finite15}.
\end{proof}


\subsection{The total contribution of the primes required to split}\label{nav:frobenius330}
Fix a field \(E=E_x\) on the interval in Section~\ref{sec:missing}
and write \(n=[E:\QQ]\). For each odd prime \(p\) form
\(F_\infty=(E_*(\mu_{p^\infty}))^+\) using the same largest
field \(E_*\) on this interval. Let \(M_p\) be the Galois group
of the maximal everywhere unramified abelian pro-\(p\) extension
of \(F_\infty\). Let \(X_p\) be its quotient obtained by imposing
complete splitting at the primes above \(p\).
Use \(Q_{x,r}\) and \(A_{p,x,r}\) as in
Theorem~\ref{thm:strictdigits329}. Put \(U_{p,x,r}=Q_{x,r}(M_p)\).
The following estimate compares the full orders at all primes
using only the degree of the original field and the range of weights.

\begin{theorem}[The cost of the splitting condition at all odd primes]\leavevmode
\label{thm:decomploss330}
For every positive \(r\) there is a natural surjection of finite
groups \(U_{p,x,r}\longrightarrow A_{p,x,r}\). Write
\(e_p\) for the ramification index of \(p\) in \(E/\QQ\).
This map is an isomorphism whenever \(p-1>e_pr\).
For every integer \(t\ge2\),
\begin{equation}\label{eq:decomploss330}
 0\le \sum_{p>2}\sum_{1\le r<t}
       \log\frac{|U_{p,x,r}|}{|A_{p,x,r}|}
     \le 6 n(t-1)\log\bigl(2n t\bigr).
\end{equation}
Here \(p\) runs over primes. At each pair \((p,r)\), one may
also apply multiplication by any power \(p^{k_{p,r}}\) with
\(k_{p,r}\ge0\) to both groups and use the orders of their
images in the same inequality.
\end{theorem}

\begin{proof}
First identify \(Q_{x,r}(X_p)\). In the notation of
Theorem~\ref{thm:strictdigits329} the logarithmic class group
at level \(j\) is \(X_p/(\gamma^{p^j}-1)X_p\).
On its twisted quotient the additional relation coming from
\(\gamma^{p^j}-1\) is multiplication by \(a^{-rp^j}-1\).
It generates the ideal \(p^{s_r+j}\ZZ_p\). The finite quotients
of \(Q_{x,r}(X_p)\) modulo \(p^{s_r+j}\) are therefore the
quotients of \(A_{p,x,r}\) at exactly those precisions,
by \eqref{eq:strictdigits329}. The identifications use restriction
of characters and norm between layers, so they commute with
reduction. A compact abelian pro-\(p\) group is the inverse limit
of these quotients. Indeed every finite quotient has bounded
exponent, and compactness gives surjectivity to the inverse limit.
This identifies \(Q_{x,r}(X_p)\) with \(A_{p,x,r}\), which is finite
by Proposition~\ref{prop:finite15}.

The primes above \(p\) in \(F_\infty\) form a finite set \(\mathcal P\),
and each has finite residue field. In fact the underlying real
cyclotomic tower of \(\QQ\) has one such prime with residue
field \(\mathbb F_p\), and \(F_\infty\) has degree at most
\([E_*:\QQ]\) over this tower. At \(w\in\mathcal P\) the local
unramified pro-\(p\) group is generated by arithmetic Frobenius
on the residue field. Denote its image in \(M_p\) by \(\phi_w\).
This element is independent of the chosen prolongation since
\(M_p\) is abelian. A field automorphism sends \(\phi_w\) to
\(\phi_{\delta w}\), because its action on residues carries the
power map defining one arithmetic Frobenius to the other.
Consequently the kernel \(D_p\) of \(M_p\to X_p\) is the image
of the permutation module \(\ZZ_p[\mathcal P]\) under
\([w]\mapsto\phi_w\). Its image is compact, hence closed.
Quotienting by the character relations is right exact. Lifting
the relations on \(X_p=M_p/D_p\) gives the quotient by \(D_p\)
and the same relations on \(M_p\). Thus the kernel of
\(U_{p,x,r}\to A_{p,x,r}\) is a quotient of
\(Q_{x,r}(\ZZ_p[\mathcal P])\).

Write \(G_x=\Gal(E_*(\mu_{p^\infty})/E)\). The primes of
\(\mathcal P\) in a single \(G_x\)-orbit contribute at most
one cyclic summand to the latter quotient. Indeed
\([\delta w]=\kappa_p(\delta)^{-r}[w]\) there. Every value
\(1-\kappa_p(\delta)^{-r}\) for a stabilizer element kills
this generator. To bound its order, extend \(w\) to the CM
tower and let \(v\) be its restriction to \(E\). The image
of the local absolute Galois group of \(E_v\) fixes this
prolongation and hence \(w\). Its cyclotomic image is
\[
 H_v=\Gal\bigl(\QQ_p(\mu_{p^\infty})/
                 (E_v\cap\QQ_p(\mu_{p^\infty}))\bigr)
       \ \subseteq\ \ZZ_p^\times .
\]
The index is the degree of the intersection. This finite
intersection is totally ramified over \(\QQ_p\), so its degree
divides \(e_p\). Write \(H_v=\mu_d\times(1+p^h\ZZ_p)\), where
\(d\mid p-1\) and \(h\ge1\). This decomposition follows by
projecting a closed subgroup of
\(\mu_{p-1}\times(1+p\ZZ_p)\) to its prime-to-\(p\) and
pro-\(p\) factors, as in Theorem~\ref{thm:alltrunc329}.
In particular
\[
        \frac{p-1}{d}\,p^{h-1}\ \mid\ e_p .
\]
The quantities \(d,h\) are the same above each \(v\), since
\(E/\QQ\) is Galois. If \(d\nmid r\), some element of \(\mu_d\)
gives a unit among the annihilators, so the generator vanishes.
If \(d\mid r\), a generator of \(1+p^h\ZZ_p\) gives an
annihilator of valuation \(h+v_p(r)\) by the binomial formula
at an odd prime. These bounds show the kernel is finite,
and they prove the asserted isomorphism whenever \(p-1>e_pr\).

The number of orbits is the number of primes of \(E^+\) above
\(p\). Indeed restriction identifies \(G_x\) with
\(\Gal(F_\infty/E^+)\), since \(EF_\infty=E_*(\mu_{p^\infty})\)
and \(E\cap F_\infty=E^+\). The number of primes of \(E^+\)
is at most the number of primes of \(E\), which is at most
\(n/e_p\). Put \(m=\lfloor(t-1)/d\rfloor\). The sum of \(v_p(r)\)
over the positive multiples \(r=di<t\) equals
\(\sum_{j\ge1}\lfloor m/p^j\rfloor\le m/(p-1)\).
The total logarithmic kernel order at \(p\) is therefore at most
\[
 \frac{n}{e_p}\frac{t-1}{d}
          \left(h+\frac1{p-1}\right)\log p
 \ \le\ n(t-1)\frac{\log p}{p-1}
          \left(1+\frac1{p-1}\right).
\]
For the last inequality use \(e_pd\ge (p-1)p^{h-1}\) and
\(h\le p^{h-1}\). All the kernels vanish if \(p>1+n(t-1)\).
For \(M=1+n(t-1)\) the elementary bound
\(\sum_{p\le M}(\log p)/p\le2\log M\) follows from
\(\lfloor M/p\rfloor\ge M/(2p)\) and
\(\sum_{p\le M}\lfloor M/p\rfloor\log p\le \log(M!)\).
For \(p\ge3\), \(p/(p-1)\le3/2\) and \(1+1/(p-1)\le3/2\).
These estimates imply \eqref{eq:decomploss330}. Finally
\(p^k U_{p,x,r}\to p^k A_{p,x,r}\) is onto, and its kernel
is a subgroup of the original kernel. This proves the estimate
at all the stated precisions.
\end{proof}

On the intervals of Section~\ref{sec:missing}, \(n\le
[E_*:\QQ]\ll_\psi x^{3+12\varepsilon}\) and
\(t=\lfloor x^{401}\rfloor\). Thus even the full correction
is \(O_{\psi,\varepsilon}(x^{404+12\varepsilon}\log(2x))\).
This also bounds the change in the excess above the precisions
of Theorem~\ref{thm:alltrunc329}. It leaves the sizes of the
ordinary unramified specializations themselves to be estimated.


\subsection{Strict classes killed by one integer throughout an interval}\label{sec:integerlayers333}

Fix an interval \([V,V^{3/2}]\) of Theorem~\ref{thm:conductor232}
and \(0<\varepsilon<1/52\). At \(x\) write \(E=E_x\), \(g=[E:\QQ]\),
\(t=\lfloor x^{401}\rfloor\), and let \(N\) be the cyclotomic modulus
used for the theta periods. For an odd prime \(p\) put
\[
  M_{p,r}(E)=\tfrac12(1+(-1)^rc)\mathcal S_{p,r}(E),
  \qquad P=\exp(x^\varepsilon/12).
\]
Here \(c\) is complex conjugation on the field. If \(B\) is a positive
integer, \(M[B]\) denotes the kernel of multiplication by \(B\) on \(M\).
The prime powers occurring in \(B\) can have very different exponents.
We estimate their contributions using the same integer at every weight.

\begin{theorem}[Order of the strict subgroup killed by a common integer]
\label{thm:integerlayers333}
On every sufficiently distant such interval, for \(Y\ge P\) and \(B\ge1\),
\[
 \begin{split}
 \sum_{\substack{p\le Y\\p\ {\rm odd\ prime}}}\ \sum_{1\le r<t}
        \log|M_{p,r}(E)[B]|
 \ \ll\ 
 g\left\{tP\left(1+\frac{\log(2N)}{\log 2}\right)
        +Y\left(1+\frac{\log(2N)}{\log P}\right)\right\}\log B .
 \end{split}
\]
The implicit constant is absolute. Take \(B=M_*\), the single positive
integer formed from the primitive critical values in
Theorem~\ref{thm:simultaneous307} on the largest field of the interval.
Then, uniformly in \(x\) on the entire interval, the left side with
\[
             Y=Y(x)=x^{-1240}\exp(x^\varepsilon/2)
\]
is \(O_{\psi,\varepsilon}(x^{-15}\exp(x^\varepsilon/2))\).
This also holds with \(B=M_*B_*\), where \(B_*\) is the positive
integer for the two critical groups in Corollary~\ref{cor:criticaledges333}.
With \(B=M_x\) formed on the individual field, the sum up to
\(x^{-820}\exp(x^\varepsilon/2)\) is
\(O_{\psi,\varepsilon}(x^{-5}\exp(x^\varepsilon/2))\).
\end{theorem}
\begin{proof}
Write \(v=v_p(B)\). In a cyclic group of order \(p^a\), the subgroup
killed by \(B\) has order \(p^{\min(a,v)}\). Summing over cyclic
factors gives \(\log|M[B]|\le v\log|M[p]|\) for each of the finite
groups in the statement, also when \(v=0\).
For \(p\le P\), Lemma~\ref{lem:rank} gives
\(\log|M_{p,r}(E)[p]|\ll gp\log(2Np)\) separately at each \(r\).
For \(p>P\), Lemma~\ref{lem:rank114} gives the same bound for the
\emph{sum} over \(r\). Its simultaneous estimate applies because
\(p>t+1\) and \(p\nmid gN\) on these intervals. The different weights
then use distinct character components in the first layer of the real
cyclotomic class group. Now
\[
  v p\log(2Np)=(v\log p)\,p\left(1+\frac{\log(2N)}{\log p}\right).
\]
In the first range the last two factors have product at most
\(P(1+\log(2N)/\log 2)\), and in the second at most
\(Y(1+\log(2N)/\log P)\). Sum the weights \(v_p(B)\log p\), whose
sum over all primes is \(\log B\). This proves the first assertion.

For the second, the field bounds give
\(g\ll_\psi x^{2+8\varepsilon}\), \(\log(2N)\ll_{\psi,\varepsilon}
x^{\varepsilon/4}\), and \(t\le x^{401}\). At the upper endpoint
write \(g_*,N_*,t_*\) for these parameters. Theorem~\ref{thm:simultaneous307}
gives
\[
   \log M_* \ll g_*t_*^2\log(2N_*t_*)
       \ll_{\psi,\varepsilon} V^{1220}\le C_{\psi,\varepsilon}x^{1220}.
\]
Indeed \(g_*\ll V^{3+12\varepsilon}\),
\(t_*^2\le V^{1203}\), and \(\log(2N_*t_*)\ll V^{3\varepsilon/8}+\log V\).
For large \(V\) we have \(Y\ge P\) and \(\log(2N)/\log P\le1\).
The first term in the asserted bound is at most a constant times
\(x^{1625}\exp(x^\varepsilon/12)\), and the second at most a
constant times \(x^{1223}Y\). Both have the required size.
Corollary~\ref{cor:criticaledges333} gives the same polynomial bound
for \(\log B_*\) directly from Lemma~\ref{lem:int13}. Finally
\(\log M_x\ll x^{812}\) by the same primitive-value estimate on
\(E_x\). Substitution into the first assertion at the larger cutoff
gives the last bound, using \(g\ll x^3\).
\end{proof}

\begin{corollary}[Full orders of the classes detected by even orthogonal points]
\label{cor:evenorders333}
At every \(x\) on the same interval, consider all \(P<p\le Y(x)\),
all \(1\le r<t\), and all complementary character orbits \([\alpha]\).
Give each orbit its minimal unramified coefficient ring \(O_\alpha\)
of degree \(b_\alpha\). The sums \(S^{\rm ev}_{p,r,\alpha}\) of
the classes in Theorem~\ref{thm:simultaneous307} satisfy
\[
       \sum_{p,r,[\alpha]} b_\alpha\log p\,
           \len_{O_\alpha} S^{\rm ev}_{p,r,\alpha}
         \ll_{\psi,\varepsilon} x^{-15}\exp(x^\varepsilon/2).
\]
In the full complementary order at \(x\), removing
\(M_{p,r}(E)[M_*]\) at every odd \(p\le Y(x)\) costs at most this
same upper bound, with a possible change of the implicit constant.
\end{corollary}
\begin{proof}
Theorem~\ref{thm:simultaneous307} kills the entire detected subgroup
by \(p^{v_p(L(1-r,\alpha^*))}\). The valuation identity there shows
that this prime power divides \(M_*\), also for characters inflated
from \(E_x\). Since \(p>g\), the ring
\(\ZZ_p[\operatorname{Gal}(E/\QQ)]\) splits into the unramified
integer rings indexed by the residue Frobenius orbits. Restriction
identifies the strict character group over \(\QQ\) with its factor
in \(M_{p,r}(E)\). Indeed averaging over the field group gives
descent in degree one and commutes with all the local restrictions,
as in Lemma~\ref{lem:des28}. A simple factor of an \(O_\alpha\)-module
has \(p^{b_\alpha}\) elements. Thus the left side is the sum of
the logarithmic orders of subgroups of distinct factors of
\(M_{p,r}(E)[M_*]\). The theorem bounds it.
Theorem~\ref{thm:orders308} gives a further componentwise range for
the detected subgroups. The common-integer estimate here includes
the \(M_*\)-torsion in every character component, even when its
individual primitive value is a \(p\)-adic unit.
For the final assertion use the exact sequence with subgroup
\(M_{p,r}(E)[M_*]\) and quotient by that subgroup, separately
at each prime and weight. The terms left over are the orders
of those quotients and the full orders above \(Y(x)\).
\end{proof}


\subsection{Simultaneous vanishing at the marked coefficient primes}\label{sec:markedstrict337}
Use the polynomial-conductor intervals of Theorem~\ref{thm:conductor232}.
Fix one such interval \([V,V^{3/2}]\), with its single pair of
quadratic primes. The fields \(E_x\) are nested there, since only the
index in the real cyclotomic tower varies. Denote their largest
member by \(E_*\), and choose its cyclotomic modulus \(N_*\le V^{3D/2}\)
from that theorem. All subfields considered below lie in this
one cyclotomic field. We use the full strict group
\(\Sel_{p,r}\) of Section~\ref{nav:II-13}.

\begin{theorem}[Selecting the period primes with vanishing descended groups]\label{thm:markedstrict337}
Fix \(0<9\theta_g<\eps\) and
\(0<9\theta_T<1/2+15\eps/16\). On every sufficiently distant
interval just specified, and at every \(x\) on it, both auxiliary
primes \(g_1,g_2\) of Theorem~\ref{thm:levels88} and the common
terminal prime \(T\) of Theorem~\ref{thm:common111} may be chosen
with the following additional property. For \(p=g_1,g_2\) use
\(\theta=\theta_g\), and for \(p=T\) use \(\theta=\theta_T\). Then
\[
 \Sel_{p,r}(L)=0
 \quad\text{for every totally imaginary } L\subseteq E_*
 \quad\text{and }r\ge1\text{ with }[L:\mathbb Q](r+2)\le V^\theta .
\]
For \(L\subseteq E_x\) the full \(\operatorname{Gal}(E_x/L)\)-invariant
subgroup of \(\Sel_{p,r}(E_x)\) is also zero. The terminal choices
still number \(Y^{1-o(1)}\) for the same selected family of words
and for the prescribed admissible factors \(\Gamma_c\) fixed before
selecting \(T\).
\end{theorem}
\begin{proof}
For any fixed \(\theta>0\), form an exceptional set of rational primes
for the whole interval. Include each \(p>3^{4V^\theta}\) for which
one of the pairs \((L,r)\) in the displayed range has
\(\Sel_{p,r}(L)\ne0\). For one pair put \(d_L=[L:\mathbb Q]\),
\(n=4(r+2)\). We have \(3^{nd_L}\le3^{4V^\theta}\).
Lemma~\ref{lem:top}, with \(M=N_*\), bounds the sum of the logarithmic
orders of all its strict groups above this threshold by
\(\exp(C(nd_L)^9\log(2N_*nd_L))\). Each nontrivial group contributes
at least \(\log2\). Thus the number of excluded primes for this pair
is at most that bound divided by \(\log2\).

Write \(G=\operatorname{Gal}(E_*/\mathbb Q)\).
Its order is at most \(N_*\). Any subgroup can be generated by at
most \(\lceil\log_2|G|\rceil\) elements, since adjoining an element
outside a generated subgroup at least doubles its size.
Counting lists of generators, padded with identities, gives at most
\(\exp(O_{\psi,\eps}((\log V)^2))\) subgroups. There are at most
\(V^\theta\) possible \(r\). Adding the bounds over all pairs gives
the following upper bound for the size of the single exceptional set:
\[
 \exp\!\left(O_{\psi,\eps,\theta}
       (V^{9\theta}\log(2V)+(\log V)^2)\right).
\]
Use this once for each of the two fixed exponents.

For the auxiliary primes the proof of Theorem~\ref{thm:levels88}
provides \(U^{1-o(1)}\) candidates in \((U,2U]\), where
\(\log U\asymp_\eps x^\eps\). They exceed \(3^{4V^{\theta_g}}\)
for large \(V\), and the exceptional set has size \(U^{o(1)}\).
After the first prime is chosen, the additional exclusions for the
second number \(O(r_0^3(1+U^{3/4}))\), with \(r_0\) the row rank,
bounded by a fixed power of \(x\). This leaves \(U^{1-o(1)}\) choices
with the required disjointness. For these primes carry out the word
selection in Theorem~\ref{thm:common111} with the stated factors.
It gives \(Y^{1-o(1)}\) eligible primes in one window for the
selected words, with
\[
 \log Y\gg m x^\eps\gg V^{1/2+15\eps/16}.
\]
These candidates exceed \(3^{4V^{\theta_T}}\), and the second
exceptional set has size \(Y^{o(1)}\). Removing it proves the
simultaneous assertions. The fields and the quadratic primes were
chosen before either exceptional set was formed.

Finally \([E_*:\mathbb Q]\le V^{3D/2}<p\) for all three chosen
primes once \(V\) is large. On global \(H^1\) with coefficients
\(\mathbb Q_p/\mathbb Z_p(-r)\), restriction to \(E_x\) identifies
the group over \(L\) with the invariants. Indeed the two outside
terms of inflation--restriction in degrees one and two vanish for
a quotient of order prime to \(p\), by averaging. The same argument
gives injection upon restriction on local \(H^1\) at \(p\).
At every other place it gives injection for the inertia groups,
since the index of the smaller inertia group divides the local
extension degree. Consequently a class whose restriction is
unramified outside \(p\) and locally zero above \(p\) already has
both properties over \(L\). Applying the global identification to
these kernels proves the assertion about invariants.
\end{proof}

For instance one may take \(\theta_g=\eps/10\) and \(\theta_T=1/18\).
For the actual period field this includes every layer
\(K_\psi C_j(\sqrt{-q},\sqrt{-d_0})\subseteq E_x\) for which
\(4\operatorname{ord}(\psi)2^j(r+2)\le V^\theta\), since the
three constituent degrees multiply to at most the coefficient on
the left. This gives zero contribution at the designated coefficient
primes in either conjugation eigenspace of these descended parts.
In particular it applies at the same \(T\) as the period transport
in Theorem~\ref{thm:graphtransport337}, with no truncation of the
cohomology exponent. The relative summands of larger degree, the
other coefficient primes and the full range \(r<x^{401}\) still
enter the length required in Section~\ref{sec:missing}.


\subsection{Vanishing of the critical sign at every required weight}\label{sec:criticalmarks339}
Use \(E_*\), \(V\) and the two exceptional sets of
Theorem~\ref{thm:markedstrict337}. Put
\(t_*=\lfloor V^{(3/2)401}\rfloor\). For an odd prime \(p\)
write \(e_n=(1+(-1)^n c)/2\), where \(c\) denotes complex
conjugation in the field group.

\begin{theorem}[The full critical sign at the three period primes]\label{thm:criticalmarks339}
The choices of \(g_1,g_2,T\) in Theorem~\ref{thm:markedstrict337}
can simultaneously satisfy, at each of these three coefficient
primes and for every \(1\le r<t_*\),
\[
                 e_{r+1}\Sel_{p,r}(L)=0
       \quad\text{for every CM subfield }L\subseteq E_* .
\]
At the same primes all the complementary strict subgroups detected
by the even orthogonal specializations of
Theorem~\ref{thm:simultaneous307} are zero throughout its range.
These choices retain the same number of terminal candidates in
the window for the selected words.
\end{theorem}
\begin{proof}
Use the single larger CM field \(E'_* = E_*(i)\) of
Corollary~\ref{cor:criticaledges333}, and the positive integer
\[
  B_*=\prod_{n=2}^{2t_*+1}|D_n(E'_*)|
\]
defined there. For every odd \(p\) prime to \(B_*\), the strict
critical-sign group over \(E'_*\) vanishes for \(n=r+1\) in the
displayed range. Indeed in the proof of Theorem~\ref{thm:critical13}
the length of its dual for each fixed \(n,p\) is bounded by
a finite multiple of \(v_p(D_n(E'_*))\), using the injection of
Lemma~\ref{lem:fit13}. For the complementary detected subgroups
Theorem~\ref{thm:simultaneous307} supplies the annihilating
integer \(M_*\). Multiplication by \(M_*\) is an automorphism of
any finite \(p\)-group if \(p\nmid M_*\), so each such subgroup
then vanishes as well.

The two integer bounds in Theorem~\ref{thm:integerlayers333}
give \(\log(B_*M_*)\ll_{\psi,\eps}V^{1220}\). The number of their
distinct prime divisors is at most this logarithm divided by
\(\log2\). Add these primes to both exceptional sets in
Theorem~\ref{thm:markedstrict337} before either auxiliary prime
is selected. A fixed polynomial in \(V\) is \(U^{o(1)}\) and
\(Y^{o(1)}\) in the notation of that theorem. Its proof therefore
gives the same counts with the extra exclusions.
Finally each chosen \(p\) exceeds \([E'_*:\mathbb Q]\) for large
\(V\). Restriction from any indicated \(L\) to \(E'_*\) is injective
on strict groups since corestriction followed after restriction
multiplies by the prime-to-\(p\) degree. It preserves complex
conjugation, as these fields are abelian over \(\mathbb Q\).
This proves the first assertion on every \(L\), including all
of the actual \(E_x\).
\end{proof}

Thus the coefficient primes used by the exchange formulas in
Theorem~\ref{thm:markedglobal339} can have zero critical-sign length
over the full period fields up to the entire weight cutoff.
On the complementary sign Theorem~\ref{thm:markedstrict337}
also removes the indicated descended components at their full
exponent. The full sum over complementary relative classes at
all coefficient primes remains in Proposition~\ref{prop:res23}.


\subsection{The full complementary order outside the class norm exceptions}\label{sec:fullnorm341}
Fix the largest field $E_*$ on $[V,V^{3/2}]$ as in
Section~\ref{sec:missing}. For every odd prime $p$ use its real
cyclotomic tower $F_i$, split class module $X_p$, generator $\gamma$
and tame group $\Delta_0$ from Theorem~\ref{thm:strictdigits329}.
Put $n_*=[E_*:\QQ]$, $j=\lfloor\log_p n_*\rfloor$ and
$B_i=\widetilde C(F_i)$. In this subsection use the whole of $B_i$,
including the invariant part for the decomposition group. Let
$g_i$ be the number of primes over $p$ in $F_i$. Set
\[
 K_j=\ker(\mathrm N:B_{j+1}\longrightarrow B_j),\qquad
 C_i=\mathrm{Cl}(\mathcal O_{F_i}[1/p])\{p\}.
\]

\begin{proposition}[The local degree contribution to a class norm kernel]\leavevmode\label{prop:localnorm341}
Every prime above $p$ in $F_{j+1}$ is fixed by
$\mathrm{Gal}(F_{j+1}/F_j)$. The map $K_j\to C_{j+1}$ has image
in the kernel of the ordinary norm to $C_j$. Its kernel is killed
by $p$ and has dimension at most $g_{j+1}-1$ over $\mathbf F_p$.
In particular
\[
 \log|K_j|\le\log|\ker(C_{j+1}\to C_j)|+(g_{j+1}-1)\log p .
\]
The same kernel assertion holds on each integral tame character
summand. The dimensions of these kernels sum to at most $g_{j+1}-1$.
\end{proposition}
\begin{proof}
At every finite layer $F_i$ the number of primes above $p$ is at
most $n_*$, by total ramification in the real cyclotomic field used
in the proof of Theorem~\ref{thm:alltrunc329}. The restriction maps
between the finite prime sets are onto. Their sizes are therefore
nondecreasing and bounded, so their inverse limit has at most $n_*$
elements. Every orbit of the procyclic pro-$p$ group generated by
$\gamma$ has size a power of $p$ at most $n_*$. All its orbit sizes
divide $p^j$. Hence $\gamma^{p^j}$ fixes every such prime, including
after restriction to $F_{j+1}$.

At that level the kernel of $B_{j+1}\to C_{j+1}$ is a quotient
of the degree-zero local lattice of rank $g_{j+1}-1$ in
Theorem~\ref{thm:logunits319}. The cyclic group of the indicated
extension acts trivially on that lattice. It fixes each prime and
its action on the local cyclotomic degree group is conjugation in
an abelian group. For a degree-zero logarithmic class $z$, extension
of its norm equals the sum of its conjugates. Indeed on an idele
the local product for the norm, embedded back into the higher
field, is the product of the conjugate ideles. Principal ideles,
local norm kernels and the degree maps respect these operations,
so the identity passes to the class group. On the local kernel this
composition is multiplication by $p$. An element in that kernel
with zero norm is consequently killed by $p$. The $p$-torsion of
a finite quotient of a free $\ZZ_p$-module of rank $g_{j+1}-1$ has
dimension at most this rank, by its cyclic decomposition. The map
to the ordinary class group commutes with norm by the same idele
calculation. Projection onto the distinct $\Delta_0$ characters is
integral and exact, giving the last assertion.
\end{proof}

Write $d=v_p(\kappa_p(\gamma)-1)$, $s_r=d+v_p(r)$ and take
$t_*=2+\lceil V^{602}\rceil$ and $H=d+\lfloor\log_p t_*\rfloor$.
As before $A_{p,x,r}$ is the dual of the complementary strict
cohomology group on $E_x$. For each character $\tau$ of $\Delta_0$
put $k_{p,\tau}=\log_p|e_\tau K_j|$.

\begin{theorem}[A bound from a finite class norm in every tame character]\leavevmode\label{thm:fullnorm341}
There are sets $D_{p,\tau}$ in the progressions
$\kappa_p^{-r}|_{\Delta_0}=\tau$, $1\le r<t_*$, with
$|D_{p,\tau}|\le k_{p,\tau}$. For every $x$ on the interval and
all unselected weights in their respective progressions,
\[
             \log_p|p^{s_r+j} A_{p,x,r}|\le H k_{p,\tau} .
\]
This uses the full order above the specified precision, through every
exponent. The bound of Proposition~\ref{prop:localnorm341} applies
to the sum of $k_{p,\tau}\log p$ over all characters, even when
$p$ divides the conductor or the degree of $E_*$.
\end{theorem}
\begin{proof}
Put $M=e_\tau X_p$ and $U=(\gamma^{p^j}-1)M$. These are finitely
generated over $\ZZ_p$ by Theorem~\ref{thm:unramdet330}. With
$N(Y)=\Phi_{p^{j+1}}(Y)$, the cyclotomic polynomial, the group
$U/N(\gamma)U$ equals $e_\tau K_j$. Indeed
$B_i=X_p/(\gamma^{p^i}-1)X_p$, the norm is the quotient map, and
$Y^{p^{j+1}}-1=(Y^{p^j}-1)N(Y)$. For $b_r=\kappa_p(\gamma)^{-r}$,
the map $M\to A_{p,x,r}$ is onto. Its quotient by $\gamma-b_r$
is the full group on $E_*$ in this progression before imposing
the additional field relations from $E_x$. The image of $U$ is
$(b_r^{p^j}-1)A_{p,x,r}=p^{s_r+j}A_{p,x,r}$. On the rational span
of $M$ and its stable subspace spanned by $U$, $\gamma-b_r$ is
injective by finiteness of the strict group on $E_*$.

For a primitive $p^{j+1}$-th root $\zeta$ put
$e'=p^j(p-1)$ and $O=\ZZ_p[\zeta]$. Via evaluation of $Y$ at
$\gamma$, $(U\otimes O)/(\gamma-\zeta)(U\otimes O)=U/N(\gamma)U$.
The extension has ramification degree $e'$ and residue field
$\mathbf F_p$. Apply Lemma~\ref{lem:simuldet340} to $U$ at $\zeta$
and all the $b_r$ of the progression. For the resulting monic
polynomial $P$, $v_p(P(\zeta))=k_{p,\tau}/e'$. Every $b_r$ has
$v_p(b_r-\zeta)=1/e'$. Roots $\alpha$ with
$v_p(\alpha-\zeta)<1/e'$ contribute the same sum $S$ to the
valuation at every $b_r$. There are at most $k_{p,\tau}-e'S$ other
roots, counted with multiplicities. For each of them choose one
closest $b_r$ when the progression is nonempty. Away from that
selection its valuation at any other $b_r$ is at most $H$, since
two distinct sample points have difference of valuation
$d+v_p(r-r')\le H$. Hence $v_p(P(b_r))\le S+H(k_{p,\tau}-e'S)
\le Hk_{p,\tau}$ outside the union of the selected weights. This
bounds the length of $U/(\gamma-b_r)U$ and hence of its indicated
image on every $E_x$.
\end{proof}

\begin{corollary}[The full strict orders outside the norm selections]\leavevmode\label{cor:allsample341}
Use the polynomial-conductor intervals, $t=\lfloor x^{401}\rfloor$,
$Y=x^{-3}\exp(x^\varepsilon/4)$, and $Z=x^{-3}\exp(x^\varepsilon/6)$.
For each odd $p\le Z$ use $D_{p,\tau}$ above. For $Z<p\le Y$ use
the first selection $D^0_\tau$ in Theorem~\ref{thm:normsample340}.
Denote the applicable selection by $D(p,\tau;x)$. Then, for each fixed $0<\varepsilon<1/52$, uniformly at sufficiently
large $x$ on these intervals,
\[
 \sum_{\substack{2<p\le Y\\p\ {\rm prime}}}
 \sum_{\substack{1\le r<t\\ r\notin D(p,\kappa_p^{-r}|_{\Delta_0};x)}}
           \log|A_{p,x,r}|
        \ll_{\psi,\varepsilon}x^{-2}\exp(x^\varepsilon/2).
\]
All selections of weights use $E_*$ and the interval cutoff $t_*$.
Only the numerical ranges $Z,Y$ in which they are applied vary with
$x$. They may be intersected with the earlier exceptional component
selections of Theorem~\ref{thm:entireinterval310} when retaining
terms in the residual length sum.
\end{corollary}
\begin{proof}
For $p\le Z$, one has $p^{j+1}\le n_*Z\le Y$ eventually, since
$n_*$ is bounded by a fixed power of $x$. The sum of the orders
in logarithmic scale through precision $s_r+j$ is therefore bounded
by Theorem~\ref{thm:alltrunc329}, which includes the invariant part.
Beyond this precision use Theorem~\ref{thm:fullnorm341}. Summing its
bound over the weights costs at most $1+tn_* /(p-1)$ repetitions per
tame character. For $p\le n_*$, the fields $F_{j+1}$ have degrees
at most $n_*p^{j+2}\le n_*^2p^2$ and cyclotomic moduli at most
$N_*n_*^2p^2$. Here $p^{d-1}\le n_*$ as in
Theorem~\ref{thm:strictdigits329}. The conductor and class number
bound of Theorem~\ref{thm:alltrunc329}, together with
Proposition~\ref{prop:localnorm341}, bounds the entire cost in
this range by a fixed power of $x$, including $H$ and the number
of weights. For $n_*<p\le Z$ one has $j=0$. The logarithmic size
of $K_0$ is $O(n_*p^2\log(2N_*n_*p^2)+n_*\log p)$. Multiplying
by $H(1+tn_* /(p-1))$ and summing costs
$O_{\psi,\varepsilon}(n_* x^\varepsilon\log(2x)(Z^3+tn_*Z^2))$.
This is within the bound, by the calculation in
Corollary~\ref{cor:normbudget340}.

For $Z<p\le Y$, eventually $p>t_*+2$ and $p>N_*$. Indeed both
$t_*$ and $N_*$ are bounded by fixed powers of $V$, while $Z$
grows faster than each such power uniformly for $x\ge V$. Now
$E_*$ is unramified at $p$ and disjoint from the $p$-power
cyclotomic field, since its modulus is prime to $p$. The latter
extension is totally ramified at $p$, also after unramified local
base change. Hence the full tame group of order $p-1$ acts by
inertia at $p$, also on restriction to the real tower. Its action
on the real tower is faithful because the only possible nontrivial kernel element,
complex conjugation, acts nontrivially on the CM field $E_*$. Thus the image
of $\Delta_0$ is contained in the prime-to-$p$ decomposition
subgroup $D'$. Every $1\le r<t$ gives a nontrivial $\tau$, since
$r<p-1$. Averaging over $D'$ kills this summand. The group denoted
$eA_{p,x,r}$ in Corollary~\ref{cor:normbudget340} is the entire
$A_{p,x,r}$ at these weights. Its first selection and bound finish
the sum. For the final assertion, terms discarded by either theorem
have nonnegative logarithmic orders, and the union of the two
discarded sums costs at most the sum of their bounds.
\end{proof}

The estimate leaves the selected full orders and the primes above
$Y$. Proposition~\ref{prop:normexample340} shows why class group
orders and maps between the underlying norm groups alone do not
bound a selected order, even with just one exceptional weight.


\subsection{The local conditions at the conductor}\label{sec:levelstrict345}

Fix one of the intervals $[V,V^{3/2}]$ of
Theorem~\ref{thm:conductor232}. Write $E=E_x$, $n=[E:\QQ]$,
$t=\lfloor x^{401}\rfloor$, and let $N_*$ be the cyclotomic modulus
of the largest period field $E_*$ on the interval. We use this same
modulus even for characters coming from a smaller $E_x$.

Fix an odd prime $\ell$, $1\le r<t$, and a character $\alpha$ of
$G=\Gal(E/\QQ)$ with values in $\overline{\QQ}_\ell$. Set
$O_\alpha=\ZZ_\ell[\alpha(G)]$, write $J_\alpha$ for its fraction
field, $e_\alpha$ for its ramification index, and put
$D=(J_\alpha/O_\alpha)(\alpha^{-1}\kappa_\ell^{-r})$.
The character group of a finite abelian extension is viewed as a
group of characters of the absolute Galois group by inflation.
Let $S(\alpha,r)$ consist of the classes in $H^1(\QQ,D)$
unramified away from $\ell$ and zero at $\ell$. In
$S^N(\alpha,r)$ require in addition that the class be zero on the
decomposition group at every $q\mid N$, $q\ne\ell$.
Write $\lambda^N_{\ell,r,\alpha}=\len_{O_\alpha}S^N(\alpha,r)$.
These are the local conditions on the scalar group underlying the
rank-three argument. The distinction matters even for divisible
coefficients when the finite character has ramified $\ell$-power image.

\begin{lemma}[The correction at one finite prime]\leavevmode\label{lem:locallevel345}
For $q\ne\ell$, let $I_q$ denote inertia, and choose a lift $F_q$ of
arithmetic Frobenius. The possible local values of an unramified
class belong to a group of length zero unless $\alpha(I_q)$ has
nontrivial $\ell$-power order and no prime-to-$\ell$ part. If this
image has order $\ell^a$, $a\ge1$, the ambient unramified local
group has $O_\alpha$-length
\[
 \min\!\left\{\frac{e_\alpha}{\varphi(\ell^a)},\
     v_{\pi_\alpha}\big(\alpha(F_q)^{-1}q^{-r}-1\big)\right\}.
\]
Here $v_{\pi_\alpha}$ assigns value one to a uniformizer.
The quotient $S(\alpha,r)/S^N(\alpha,r)$ injects into the direct
sum of these local groups for $q\mid N$, $q\ne\ell$.
\end{lemma}
\begin{proof}
Inflation--restriction identifies the unramified subgroup with
$D^{I_q}/(F_q-1)D^{I_q}$. Indeed the procyclic residue Galois
group acts on $D^{I_q}$, and a cocycle is specified by its value
on arithmetic Frobenius modulo the values of coboundaries. Every
value defines a continuous cocycle because it lies in a finite stable
submodule, and a multiple of the order of the action times the
exponent of that submodule makes the cyclic norm of the value zero.
The cyclotomic character is unramified at $q$.
If inertia has trivial image then $D^{I_q}=D$ and
the scalar $\alpha(F_q)^{-1}q^{-r}-1$ is nonzero. In fact
$q^r>1$ and $\alpha(F_q)$ is a root of unity. A nonzero scalar
acts surjectively on $J_\alpha/O_\alpha$. If inertia has
nontrivial prime-to-$\ell$ image, some inertia eigenvalue differs
from one by a unit and $D^{I_q}=0$.

In the remaining case a generator of the inertia image is a
primitive $\ell^a$-th root $\zeta$. Its difference from one has
valuation $e_\alpha/\varphi(\ell^a)$. This follows by evaluating
the cyclotomic polynomial at one. The product of the
$1-\zeta^j$ for $j$ prime to $\ell$ modulo $\ell^a$ equals $\ell$,
and each has the same valuation. Thus $D^{I_q}=D[\zeta-1]$ is
cyclic of this length. Its quotient by $F_q-1$ has the displayed
length by the elementary divisor theorem. Changing the lift of
Frobenius preserves its action on $D^{I_q}$. Finally localization
at the indicated primes on $S(\alpha,r)$ has kernel exactly
$S^N(\alpha,r)$.
\end{proof}

\begin{theorem}[Full character lengths on the period fields]\leavevmode\label{thm:levelstrict345}
Let $M_{\ell,r}(E)=\mathcal S_{\ell,r}(E)^{c=(-1)^r}$.
For every $x$ on these intervals put
\[
 Z_{\ell,r}(E)=
   \sum_{\substack{\alpha\in\widehat G(\overline{\QQ}_\ell)\\
                       \alpha(c)=(-1)^r}}
       \frac{\lambda^{N_*}_{\ell,r,\alpha}}{e_\alpha}.
\]
All the character groups here have finite length. For $\ell\nmid n$
one has $\len_{\ZZ_\ell}M_{\ell,r}(E)=Z_{\ell,r}(E)$. More generally,
uniformly on the interval,
\[
 \sum_{\ell>2}(\log\ell)\sum_{1\le r<t}
 \left|\len_{\ZZ_\ell}M_{\ell,r}(E)-Z_{\ell,r}(E)\right|
       \ll_\psi t\,n^2\log(2N_*).
\]
The sum over geometric characters in $Z_{\ell,r}$ can equivalently
be taken over the $\QQ_\ell$-orbits, with coefficient
$[O_\alpha/\pi_\alpha:\mathbf F_\ell]$ multiplying each
$\lambda^{N_*}_{\ell,r,\alpha}$.
\end{theorem}
\begin{proof}
Choose the integer ring $O$ of a finite extension of $\QQ_\ell$
containing all the character values, with uniformizer $\pi$, residue
field $k$ and ramification index $e$. A finite free extension of
coefficients commutes with the global and local cochains and the
kernels under consideration. This follows in a basis of the
extension, including for the divisible coefficients by writing
$(J_\alpha/O_\alpha)\otimes_{O_\alpha}O=\operatorname{Frac}(O)/O$.
After this extension put
\[
 A=M_{\ell,r}(E)\otimes_{\ZZ_\ell}O,\qquad
 B=\bigoplus_{\alpha(c)=(-1)^r} S(\alpha,r)\otimes_{O_\alpha}O .
\]
By Lemma~\ref{lem:fourier231} the character evaluation and its
inverse up to scalar give maps $A\to B\to A$ whose composites
in both orders equal multiplication by $n$. The local assertion
in that lemma uses the inertia cosets on restriction of the
induced representation, so applies also at primes ramified in $E$.
When $\ell\nmid n$ these maps are isomorphisms after dividing
one by $n$. No inertia image of any $\alpha$ then has nontrivial
$\ell$-power order. Lemma~\ref{lem:locallevel345} identifies
$S(\alpha,r)$ with $S^{N_*}(\alpha,r)$, giving the asserted
equality upon dividing the $O$-lengths by $e$.

It remains to bound the kernels when $m=v_\ell(n)>0$.
The extension $E/K_\psi$ has degree a power of two, so
$\ell^m\mid\operatorname{ord}(\psi)$. In particular both $\ell$
and $m$ are bounded in terms of $\psi$. The kernels of the two maps
are contained in $A[\ell^m]$ and $B[\ell^m]$, respectively.
The group $A$ is finite by Proposition~\ref{prop:finite15}.
The Kummer estimate in the proof of Theorem~\ref{thm:trunc15} gives
$\dim_{\mathbf F_\ell}M_{\ell,r}(E)[\ell]\ll_\psi n\log(2N_*)$.

Here is a bound for the other kernel that also proves its finiteness.
Put $\Sigma=\{q:q\mid\ell N_*\}$, and write $G_{\QQ,\Sigma}$ for
the Galois group of the maximal extension unramified at all other
finite primes. Outside $\Sigma$ the coefficient is unramified, and an
unramified cocycle is zero on inertia. Thus the scalar groups lie in
the cohomology of $G_{\QQ,\Sigma}$. By the coefficient sequence
for multiplication by $\pi$, each scalar $\pi$-torsion group over
$O$ has dimension at most that of
$H^1(G_{\QQ,\Sigma},k(\bar\alpha^{-1}\bar\kappa_\ell^{-r}))$.
Let $P$ be the $\ell$-Sylow of $G$ and put $F'=E^P(\mu_\ell)$.
Its degree divides $(\ell-1)n/|P|$ and is prime to $\ell$.
Restriction to $F'$ is injective by corestriction, and the residual
coefficient becomes trivial there. In a basis of $k$ the resulting
dimension over $k$ is bounded by the $\mathbf F_\ell$-dimension
of $H^1(G_{F',\Sigma},\mu_\ell)$, since $\mu_\ell\subset F'$.
Kummer theory bounds this by the dimension of the $\Sigma$-units
modulo $\ell$ and the $\ell$-rank of the $\Sigma$-class group.
The unit rank and number of inverted places are
$O_\psi(n\log(2N_*))$. The field has cyclotomic modulus dividing
$N_*\ell$ and degree $O_\psi(n)$. The conductor--discriminant
formula and the Minkowski ideal count used in
Theorem~\ref{thm:trunc15} give the same bound for the logarithm
of its class number. This proves
$\dim_k B[\pi]\ll_\psi n^2\log(2N_*)$, since there are at most
$n$ characters.

For any $\pi$-primary module of finite socle dimension $d$, the
subgroup killed by $\pi^j$ has length at most $jd$, since multiplication
by $\pi^{i-1}$ embeds each successive quotient of $\pi^i$-torsion
by $\pi^{i-1}$-torsion in the socle. Hence $B[\ell^m]$ is finite.
The map to $A$ now gives finiteness of $B$. Using the two maps in
opposite directions bounds $\left|\len_O A-\len_O B\right|/e$
by $m\max(\dim_k A[\pi],\dim_k B[\pi])$. For $A$ its socle
dimension is the number of $\ZZ_\ell$-cyclic factors of
$M_{\ell,r}(E)$, by base change of its elementary divisors.
This proves an $O_\psi(n^2\log(2N_*))$ bound at each such prime
and twist. Passing to $S^{N_*}$ costs at most one unit of normalized
length per character and prime dividing $N_*$ by
Lemma~\ref{lem:locallevel345}. Summing over $r$ and the fixed
prime divisors of $\operatorname{ord}(\psi)$ gives the claim.
Finally an orbit has $e_\alpha[O_\alpha/\pi_\alpha:\mathbf F_\ell]$
members, and conjugation identifies their lengths.
\end{proof}

In the source the divisible character module and the modulus used for
its local conditions are specified in
\srcloc{Development/AutomorphicArithmetic/AutomorphicVarieties/CharacterLevelBounds/ShaFamMono}{30}.
Theorem~\ref{thm:levelstrict345} bounds the change of local conditions
and of integral character lattices. A bound for the scalar lengths
themselves is still needed.



\subsection{Ordinary class groups at every precision of the strict group}
\label{sec:ordinarydigits514}

Fix the fields \(E=E_x\subseteq E_*\) on the interval of
Section~\ref{sec:missing}. Write \(n=[E:\QQ]\) and \(n_*=[E_*:\QQ]\).
Fix an odd prime \(p\). Use the real cyclotomic tower
\(F_0\subset F_1\subset\cdots\subset F_\infty\) of
Theorem~\ref{thm:strictdigits329}, with \([F_j:F_0]=p^j\).
In the full CM tower \(L_\infty=E_*(\mu_{p^\infty})\), the group
over \(E_*(\mu_p)\) restricts isomorphically to
\(\Gamma=\Gal(F_\infty/F_0)\). Choose a topological generator
\(\gamma\), put \(a=\kappa_p(\gamma)\), and put
\(s_r=v_p(a^{-r}-1)\) for each positive integer \(r\).
The ordinary group used below is
\[
 C_j=\Cl(\mathcal O_{F_j}[1/p])\{p\}.
\]
It is obtained from the ideal class group of \(F_j\) by killing
the classes of the primes above \(p\) and taking the \(p\)-primary
part. Write \(T_j=\widetilde C(F_j)\) for the degree-zero
logarithmic class group. For \(G_x=\Gal(L_\infty/E)\), define
\[
 Q_{r,m}(Y)=Y\Big/\overline{p^mY+
         \sum_{\delta\in G_x}(\delta-\kappa_p(\delta)^{-r})Y}
         \qquad(m\ge1).
\]
The action on each class group is induced by restriction to
its field. Let \(A_{p,x,r}\) be the Pontryagin dual of
\(\mathcal S_{p,r}(E)^{c=(-1)^r}\) in the untwisted class
group coordinates of Theorem~\ref{thm:strictdigits329}.

\begin{lemma}[Local degrees under norm after the primes cease splitting]
\leavevmode\label{lem:locnorm514}
There is an integer \(\sigma_p\ge0\), with \(p^{\sigma_p}\le n_*\),
such that restriction gives a bijection on the primes above
\(p\) from \(F_\infty\) to \(F_j\) for every \(j\ge\sigma_p\).
Let \(g_p\le n_*\) be their number in these layers. For each such
\(j\) there is a natural surjection \(T_j\to C_j\) whose kernel
\(K_j\) is a quotient of
\[
 P^0=\{(z_w)_{w\mid p}\in\ZZ_p^{g_p}:\ \sum_w z_w=0\}.
\]
The field action on this lattice permutes the places. For each
integer \(i\ge0\) one has \(N_{j+i,j}K_{j+i}=p^iK_j\).
\end{lemma}
\begin{proof}
Write \(u=v_p(a-1)-1\). As in the proof of
Theorem~\ref{thm:strictdigits329},
\(F_j=(E_*(\mu_{p^{u+1+j}}))^+\). Its degree over
\(\QQ(\mu_{p^{u+1+j}})^+\) is at most \(n_*\), by the compositum
degree bound and the fact that both CM fields have index two over
their real subfields. The real cyclotomic field has a unique prime
over \(p\). Thus the number of places above \(p\) in each layer is
at most \(n_*\). These finite sets map onto one another, so their
inverse limit has at most \(n_*\) elements. The decomposition
subgroup of a place in \(\Gamma\) consequently has finite index
\(p^{\sigma_p}\le n_*\). It is the same subgroup for all places,
since \(F_\infty/\QQ\) is abelian and transitive on the places.
For \(j\ge\sigma_p\), the group \(\Gamma_j=\Gal(F_\infty/F_j)\)
lies in this stabilizer. Each place of \(F_j\) then has a unique
extension to \(F_\infty\) and its local degree group is all of
\(\Gamma_j\). This is the degree group for the local cyclotomic
extension because \(F_\infty\) is the compositum of \(F_j\) with the
cyclotomic \(\ZZ_p\)-extension of \(\QQ\).

By local reciprocity a local factor at \(w\mid p\) maps onto this
degree group. In the global class group defined by the locally
cyclotomic condition, killing all these local factors gives \(C_j\)
by ideal class field theory. The global cyclotomic degree of a
tuple of the local factors is the sum of its local degrees. Since
one local factor already maps onto \(\Gamma_j\), any lift of a
class in \(C_j\) can be corrected to have degree zero. The resulting
surjection from \(T_j\) has as kernel the image of the tuples of
degree sum zero, exactly as in Theorem~\ref{thm:logunits319}.
Use \(\gamma^{p^j}\) as generator in each coordinate. Conjugation
acts trivially on the cyclotomic group and permutes the places,
which proves the asserted lattice description. By functoriality
of local reciprocity the norm on local degrees is the inclusion
\(\Gamma_{j+i}\subseteq\Gamma_j\). In these coordinates each entry
is multiplied by \(p^i\), at the same place. The image of all
degree-zero tuples from above is therefore exactly \(p^iP^0\)
below. Taking images in the logarithmic class group proves the
norm assertion, including in the presence of relations from units.
\end{proof}

Let \(v\mid p\) in \(E\). Write the cyclotomic image of its local
absolute Galois group as \(\mu_d\times(1+p^h\ZZ_p)\), with
\(d\mid p-1\) and \(h\ge1\). This open subgroup splits into its
prime-to-\(p\) and pro-\(p\) factors by taking the corresponding
projections in its finite cyclic quotients. Its index
\((p-1)p^{h-1}/d\) divides the ramification index \(e_p(E/\QQ)\),
since the intersection of \(E_v\) with the local cyclotomic tower
is totally ramified. The quantities are independent of the place
by Galois conjugacy. Define
\[
 b=b_{p,x,r,m}=
 \begin{cases}
  0 & d\nmid r,\\
  \min\{m,h+v_p(r)\} & d\mid r .
 \end{cases}
\]

\begin{theorem}[Recovery from ordinary ideal classes at a specified layer]
\leavevmode\label{thm:ordinarydigits514}
Fix \(r,m\ge1\), put \(A=A_{p,x,r}/p^m A_{p,x,r}\), and choose
\(j\ge\max\{\sigma_p,m-s_r,0\}\). There is a subgroup \(B\subseteq A\)
generated by at most \(g_p-1\) elements and killed by \(p^b\) such
that, compatibly with norm, for every \(i\ge0\)
\[
        Q_{r,m}(C_{j+i})\simeq A/p^i B.
\]
In particular the layer \(j+b\) recovers \(A\) exactly. Equality
of the orders of the ordinary quotients at two consecutive layers
\(j+i,j+i+1\) already certifies recovery of \(A\) at precision
\(m\). At fixed \(p,m\) the single layer
\(F_{m+\lfloor\log_p n_*\rfloor}\) suffices for every
\(E_x\) in the interval and every positive weight.
\end{theorem}
\begin{proof}
Theorem~\ref{thm:strictdigits329} identifies \(Q_{r,m}(T_{j+i})\)
with \(A\) because \(s_r+j+i\ge m\). These identifications use
restriction of characters and are compatible with norm, which
induces the identity on \(A\). The quotient functor \(Q_{r,m}\)
is right exact. Taking the quotient by \(K_{j+i}\) before imposing
the scalar and precision relations has the same effect as imposing
all these relations in one quotient. Its kernel on the surjection
to the ordinary group is consequently the image of \(K_{j+i}\).
Denote the image of \(K_j\) in \(A\) by \(B\). The lemma and the
norm identification show that the image of \(K_{j+i}\) is exactly
\(p^i B\). The lattice in the lemma gives the claimed generator bound.

Choose a prolongation of \(v\) to \(L_\infty\). The local Galois
group at \(v\) restricts into \(G_x\) and fixes that place. Its
elements fix every place over \(p\) in \(F_\infty\), since the
Galois group over \(\QQ\) is abelian. They therefore act trivially
on \(P^0\). In the twisted quotient each of the scalars
\(1-\kappa_p(\delta)^{-r}\) for these elements annihilates \(B\).
If \(d\nmid r\), a tame generator gives a unit. Otherwise a generator
of \(1+p^h\ZZ_p\) gives a scalar of valuation \(h+v_p(r)\) by
binomial expansion at the odd prime \(p\). The precision also kills
\(p^m B\). Thus \(p^bB=0\). If consecutive ordinary orders agree,
then \(p^iB=p^{i+1}B\), and iteration in the finite group gives \(p^iB=0\).

Take \(j=\max\{\sigma_p,m-s_r,0\}\). When \(j=\sigma_p\), one has
\(j+b\le\lfloor\log_p n_*\rfloor+m\). In the other case use
\(s_r=1+u+v_p(r)\), \(p^{h-1}\le e_p\le n_*\), and
\(b\le h+v_p(r)\) to obtain the same inequality. Every larger
layer gives the same quotient, proving the common-field assertion.
\end{proof}

\begin{corollary}[Total error over all odd primes, independent of precision]
\leavevmode\label{cor:allprimes514}
For \(t\ge2\) integral, choose for each \(p>2\) prime and \(1\le r<t\)
any \(m\ge1\), any \(j\) as in the theorem and any integer \(k\ge0\).
With \(A=A_{p,x,r}/p^m A_{p,x,r}\) at each pair,
\[
 0\le\sum_{p>2}\sum_{1\le r<t}
     \log\frac{|p^k A|}{|p^k Q_{r,m}(C_j)|}
 \le 6 n_* n(t-1)\log(2nt).
\]
At \(t=\lfloor x^{401}\rfloor\) the right side is polynomial in \(x\)
uniformly on the fields of Section~\ref{sec:missing}.
\end{corollary}
\begin{proof}
The restricted surjection on the images of \(p^k\) has kernel
contained in \(B\), whose order is at most \(p^{(g_p-1)b}\).
A positive discrepancy requires \(d\mid r\), so \(p-1\le e_p r\)
and \(p\le 1+n(t-1)\). At this prime put \(M=\lfloor(t-1)/d\rfloor\).
The sum of \(v_p(r)\) over these weights is
\(v_p(M!)\le M/(p-1)\), because \(d\) is prime to \(p\).
Using \(g_p-1\le n_*\), \(e_p\le n\) and \(h\le p^{h-1}\), the
logarithmic contribution is at most
\[
 n_*n(t-1)\frac{\log p}{p-1}\left(1+\frac1{p-1}\right).
\]
For integral \(X\ge2\), the inequality
\(\sum_{p\le X}\lfloor X/p\rfloor\log p\le\log(X!)\)
gives \(\sum_{p\le X}(\log p)/p\le2\log X\).
Apply it with \(X=1+n(t-1)\), using \(p/(p-1)\le3/2\) and
\(1+1/(p-1)\le3/2\). This proves the displayed bound. The degree
estimate from Theorem~\ref{thm:conductor232} gives the last assertion.
\end{proof}

The ordinary quotient at the common layer accounts for all classes
in the complementary sign at the prescribed precision. The remaining
cohomological estimate in Section~\ref{sec:missing} asks for an upper
bound on the orders of the surviving quotients themselves as the
precisions and weights vary.


\subsection{The exceptional weights at small coefficient primes}
\label{sec:finiteweights515}

We now locate the weights where the full complementary group at
a small prime can exceed a polynomial bound. The set of weights
will be determined by ordinary ideal classes in finitely many
fields attached to the largest field \(E_*\) of the interval. The
same set will work on every \(E_x\). Keep \(F_j,C_j,T_j,\gamma,a\)
and \(\sigma=\sigma_p\) from Section~\ref{sec:ordinarydigits514}
at an odd prime \(p\). Let \(X\) be the abelian pro-\(p\) class
module over \(F_\infty\) unramified away from \(p\) and split at
the primes over \(p\). Write \(T=\gamma-1\) for the Iwasawa
variable. The class field description in
Theorem~\ref{thm:strictdigits329} gives
\[
       T_j=X/\big((1+T)^{p^j}-1\big)X ,
\]
with norm induced by the quotient map. Indeed abelianizing over
\(F_j\) imposes the commutators with \(\gamma^{p^j}\), and the
kernel of the cyclotomic degree in that abelianization is the
displayed quotient.

The group of \(L_\infty/E_*\) is a subgroup of
\(\ZZ_p^\times\) under the cyclotomic character. It is the product
of its finite prime-to-\(p\) part \(\Delta\) and the closed
subgroup generated by \(\gamma\). The order of \(\Delta\)
divides \(p-1\). Thus for each character \(\tau\) of \(\Delta\)
the idempotent \(e_\tau=|\Delta|^{-1}\sum_{\delta\in\Delta}
\tau(\delta)^{-1}\delta\) is integral over \(\ZZ_p\). These
idempotents split every module in the following calculation.
Take \(M_\tau=e_\tau X\) and \(P(Y)=\Phi_{p^{\sigma+1}}(Y)\).
Define nonnegative integers from the ordinary group at a finite level by
\[
 D_\tau=e_\tau C_{\sigma+2}/P(\gamma)e_\tau C_{\sigma+2},
       \qquad h_\tau=\log_p|D_\tau|.
\]

\begin{lemma}[An ordinary class group at a primitive cyclotomic character]
\leavevmode\label{lem:finiteprimitive515}
There is an isomorphism \(M_\tau/P(\gamma)M_\tau\simeq D_\tau\).
At the preceding level one has
\[
       \sum_\tau h_\tau\log p
       \le \log |C_{\sigma+1}|+(g_p-1)\log p .
\]
These statements apply also at primes dividing the conductor or
degree of \(E_*\).
\end{lemma}
\begin{proof}
The polynomial \(P(Y)\) divides \(Y^{p^j}-1\) whenever
\(j\ge\sigma+1\). Hence \(X/P(\gamma)X\) identifies with
\(T_j/P(\gamma)T_j\) at all these levels, compatibly with norm.
On the kernel \(K_{\sigma+1}\) of \(T_{\sigma+1}\to C_{\sigma+1}\),
the element \(\gamma^{p^\sigma}\) acts trivially, by the local
lattice of Lemma~\ref{lem:locnorm514}. Consequently
\(P(\gamma)=\sum_{i=0}^{p-1}\gamma^{i p^\sigma}\) acts there as
multiplication by \(p\). The image of this kernel after imposing
\(P(\gamma)=0\) therefore has order at most \(p^{g_p-1}\).
The surjection to the corresponding quotient of \(C_{\sigma+1}\)
gives the order bound for \(X/P(\gamma)X\).

The norm from level \(\sigma+2\) sends its local kernel onto
\(pK_{\sigma+1}\) by the same lemma. The latter maps to zero
after imposing \(P(\gamma)=0\), and norm on the specialized
logarithmic groups is an isomorphism. The upper local kernel
therefore has zero image. Right exactness identifies the upper
ordinary quotient with \(X/P(\gamma)X\). Projecting onto
\(e_\tau\) gives the isomorphism. The orders of the summands
multiply to the order of the whole group, proving the inequality.
\end{proof}

For an integer \(R\ge2\) put \(w=v_p(a-1)\),
\(W=w+\lfloor\log_p R\rfloor\), and
\[
 I_\tau=\{1\le r<R:\ \kappa_p^{-r}|_\Delta=\tau\},
 \qquad H_\tau=W h_\tau,\qquad m_p=1+\max_\tau H_\tau .
\]
Compute the ordinary quotients on the right below at
\(j_p=m_p+\lfloor\log_p n_*\rfloor\), with \(G_x\) replaced by
\(\Gal(L_\infty/E_*)\) in \(Q_{r,m_p}\). Define
\[
 \mathcal E_{p,\tau}=
   \{ r\in I_\tau:\ \log_p|Q_{r,m_p}(C_{j_p})|>H_\tau\}.
\]
Thus testing membership uses the action on ordinary ideal classes
at a specified finite level and scalar congruences modulo \(p^{m_p}\).

\begin{theorem}[A finite exceptional set for every tame character]
\leavevmode\label{thm:finiteweights515}
One has \(|\mathcal E_{p,\tau}|\le h_\tau\).
For every field \(E_x\) of the interval and every
\(r\in I_\tau\setminus\mathcal E_{p,\tau}\), the full order satisfies
\(\log_p|A_{p,x,r}|\le H_\tau\).
At \(E_*\) the exceptional set consists exactly of the weights
whose full length exceeds \(H_\tau\).
\end{theorem}
\begin{proof}
For any finite abelian \(p\)-group \(B\) and integers \(m>H\ge0\),
one has \(\log_p|B|>H\) if and only if \(\log_p|B/p^mB|>H\).
Indeed any cyclic factor cut short by reduction still contributes
\(m>H\), and if none is cut short the orders agree. Apply
Theorem~\ref{thm:ordinarydigits514} with \(E=E_*\) to prove the
exact identification of the exceptional set. In the progression
\(I_\tau\), put \(t_r=a^{-r}-1\). Imposing the field and scalar
relations on \(X\) first over \(E_*\) gives
\(A_{p,*,r}=M_\tau/(T-t_r)M_\tau\), as in the proof of
Theorem~\ref{thm:decomploss330}. Indeed the tame relations
project onto \(e_\tau\), and equivariance for the topological
generator gives equivariance for its closed subgroup by continuity.
The corresponding group on \(E_x\) imposes all of these relations
and possibly further ones. It is a quotient of \(A_{p,*,r}\).
It remains to count the weights of large length at \(E_*\).

Choose a primitive \(p^{\sigma+1}\)-th root of unity \(\zeta\).
The ring \(O=\ZZ_p[\zeta]\) has uniformizer \(\pi=\zeta-1\) and
ramification degree \(e=p^\sigma(p-1)\), by the Eisenstein
polynomial \(P(1+T)\). Evaluation at \(T=\pi\) maps
\(\Lambda=\ZZ_p[[T]]\) onto \(O\) with kernel \((P(1+T))\),
by division by this monic distinguished polynomial. The lemma
gives \(M_\tau\otimes_\Lambda O=D_\tau\). Its \(O\)-length is
\(h_\tau\), since the residue field has order \(p\).

We give the presentation argument to keep track of the integral
length. The reduction of \(M_\tau\) modulo \((p,T)\) is finite
dimensional, as it is a quotient of \(D_\tau\). Lifts of a basis
generate \(M_\tau\) over \(\Lambda\). In fact the submodule they
generate is closed by compactness of a finite free \(\Lambda\)-module.
The quotient \(Y\) then satisfies \(Y=(p,T)Y\). Its finite
\(\Gamma\)-stable quotients separate points. To see this, an
open subgroup contains a smaller open subgroup stable under
\(\Gamma\), obtained by intersecting its translates. There are
only finitely many translates because continuity and compactness
of the open subgroup make its stabilizer open. On any finite
stable quotient of \(Y\), the action of \(\gamma\) has \(p\)-power
order. Modulo \(p\), some power of \(T=\gamma-1\) is zero, by
the binomial formula in characteristic \(p\). Since the finite
quotient is also killed by a power of \(p\), some power of
\((p,T)\) kills it. Thus all these quotients vanish, proving
generation.

Take a surjection \(\Lambda^q\to M_\tau\) from the chosen
generators. If \(q=0\), all the groups in question are zero.
Otherwise its relations map at \(\pi\) onto a full-rank sublattice
of \(O^q\) with cokernel \(D_\tau\). This follows by right
exactness and surjectivity of \(\Lambda\to O\). Lift a basis of
the sublattice to relation columns over \(\Lambda\), and let
\(f(T)\) be their determinant. For the valuation normalized by
\(v(p)=1\), Smith reduction gives \(v(f(\pi))=h_\tau/e\).
Evaluation at any \(t_r\) gives relation columns on the full
group \(A_{p,*,r}\). Therefore its length is at most \(v(f(t_r))\)
whenever this determinant is nonzero.

Factor \(f=p^\mu U\prod_i(T-\beta_i)\) by Weierstrass
preparation over \(\ZZ_p\), allowing a finite extension for
the roots. The unit power series \(U\) has valuation zero at
all these arguments, and the roots, with multiplicities, have
positive valuation. Since \(v(\pi)=1/e<1\), a root with
\(v(\beta_i)\ge1\) contributes exactly \(1/e\) at \(\pi\).
There are at most \(h_\tau\) such roots. For each, in a nonempty
progression choose one \(r_i\) maximizing \(v(t_{r_i}-\beta_i)\),
allowing an infinite maximum. At every other weight,
\[
 v(t_r-\beta_i)\le v(t_r-t_{r_i})
       =w+v_p(r-r_i)\le W .
\]
The equality uses the binomial formula for a principal unit at
an odd prime. For a root with \(v(\beta_i)<1\), its contribution
at every \(t_r\) is \(v(\beta_i)\le e\,v(\pi-\beta_i)\).
Indeed if \(v(\beta_i)<1/e\) the two valuations coincide, and
otherwise \(e\,v(\pi-\beta_i)\ge1\).
Away from the chosen weights every individual contribution to
\(v(f(t_r))\), including \(\mu\), is consequently at most
\(eW\) times its contribution at \(\pi\). In particular the
determinant is nonzero there and \(v(f(t_r))\le W h_\tau\).
Every weight in \(\mathcal E_{p,\tau}\) must be among the at
most \(h_\tau\) chosen weights.
\end{proof}

\begin{corollary}[Simultaneous vanishing of the residual groups at small primes]
\leavevmode\label{cor:smallresidual515}
Take \(0<\varepsilon<1/52\) and a polynomial-conductor interval
\([V,V^{3/2}]\) of Theorem~\ref{thm:conductor232}. Set
\(R=\lfloor V^{1203/2}\rfloor\) and \(Y=V^{100}\).
There is a set \(\mathcal E_V\subset[1,R)\cap\ZZ\) determined by
the finite ordinary groups above, of size
\[
 |\mathcal E_V|\ll_{\psi,\varepsilon}
            V^{306+24\varepsilon}\log(2V).
\]
On every sufficiently distant such interval, for all
\(x\in[V,V^{3/2}]\), all odd primes \(p\le Y\), and all
\(1\le r<\lfloor x^{401}\rfloor\) outside \(\mathcal E_V\),
the residual group \(R_{p,r}\) of Section~\ref{sec:missing}
is zero. The full complementary groups in this range have total
weighted length bounded by a fixed power of \(V\).
\end{corollary}
\begin{proof}
Let \(N_*\) be the cyclotomic modulus of \(E_*\). The real field
\(F_{\sigma+1}\) has degree at most \(n_*^2p^2\) and admits the
modulus \(N_*p^{u+\sigma+2}\le N_* n_*^2 p^2\).
Indeed \(p^u,p^\sigma\le n_*\), and its degree over \(F_0\)
is \(p^{\sigma+1}\), with \([F_0:\QQ]\le n_*(p-1)\).
For a real abelian field of degree \(d\) and modulus \(N\), the
conductor--discriminant formula gives \(|\mathrm{disc}|\le N^d\).
Every ideal class contains an integral ideal of norm at most
\(\sqrt{|\mathrm{disc}|}\) by Minkowski. The number of these
ideals is at most \(|\mathrm{disc}|\zeta_F(2)\le(2N)^d\),
using the Euler factors. Lemma~\ref{lem:finiteprimitive515}
therefore gives
\[
   \sum_\tau h_\tau\log p
      \ll n_*^2p^2\log(2N_*n_*p).
\]
Take \(\mathcal E_V\) to be the union of the finite sets
\(\mathcal E_{p,\tau}\) for \(2<p\le Y\). Summing the count in
the theorem over primes using \(\sum_{p\le Y}p^2\le Y^3\) gives
\(O(n_*^2Y^3\log(2N_*n_*Y))\). Now
\(n_*\ll_\psi V^{3+12\varepsilon}\) and
\(\log(2N_*)\ll_{\psi,\varepsilon}\log(2V)\), which give the
displayed cardinality. In particular it is \(o(V^{401})\),
even at the smallest weight cutoff of this interval.

Here \(W\le1+\log_p n_*+\log_p R\). Thus every \(H_\tau\)
in the stated prime range is bounded by a fixed power of \(V\).
For large \(V\), \(Y<P(x)=\exp(x^\varepsilon/12)\) and every
\(H_\tau<b_0(x)=\lfloor\exp(x^\varepsilon/8)\rfloor\), uniformly
for \(x\ge V\). At a weight outside the union, the full
complementary group has length at most \(H_\tau\) and hence is
killed by \(p^{b_0(x)}\). This kills its summand \(U_{p,r}\) in
the definition of \(R_{p,r}\), so the residual quotient is zero.
For the final assertion, there are at most \(YR\) prime and
weight pairs, each with log order bounded by
\((\max_{p,\tau}H_\tau)\log Y\). Their product bound is polynomial.
\end{proof}

Membership in \(\mathcal E_V\) can be tested without locating
any root of an Iwasawa power series. The calculation bounds the
number of exceptional weights. Their full orders, including on the
smaller fields \(E_x\), still enter the sum in Section~\ref{sec:missing}.


\subsection{A fixed level for the rational strict groups}
\label{sec:rationalfixed538}

Write \(S_{p,r}=\mathcal S_{p,r}(\QQ)\). Its coefficients are
\(A_{p,r}=(\QQ_p/\ZZ_p)(-r)\). A class is unramified away
from \(p\) and has zero restriction to \(G_{\QQ_p}\).
If \(v\ne p\) is a finite prime, the unramified local group is
\(A_{p,r}/(\operatorname{Frob}_v-1)A_{p,r}\). This is zero for
\(r\ge1\), since multiplication by the nonzero scalar
\(v^{-r}-1\) on \(\QQ_p/\ZZ_p\) is surjective. Requiring zero
localization at any additional such primes therefore gives
the same strict group.

\begin{lemma}[The automorphic labels for the trivial character]
\leavevmode\label{lem:rationallevel538}
Put \(K=\QQ(\sqrt{-3})\) and \(u=7\), and fix one place of \(K\)
above \(u\). For every even \(r\ge2\) form the multiset of
Theorem~\ref{thm:integer} with this \(K,u\), level \(N=2\), and
the unique character modulo two. Set \(d_r=1+u^r+u^{-r}\).
For every prime \(p\ge11\) with \(p-1\nmid2r\), the one-character
comparison of Theorem~\ref{thm:adjacentweights427} for \(S_{p,r}\)
has the following form. There is a submultiset \(\mathcal T_{p,r}\)
of the labelled trace values at this fixed level, with their
coefficient embeddings, such that
\[
 \operatorname{length}_{\ZZ_p} S_{p,r}
       \ \le\ 2\sum_{b\in\mathcal T_{p,r}}v_p(b-d_r),
       \qquad v_p(p)=1 .
\]
\end{lemma}
\begin{proof}
Use the initial character of conductor one in
Theorem~\ref{thm:conductor232}. In
Proposition~\ref{prop:quadratic345} take \(s=3,u_1=7,u_2=13\).
The last two primes are \(1\bmod3\), so split in \(K\).
The Minkowski bound for \(K\) is \(2\sqrt3/\pi<2\), giving
\(|\mathrm{Cl}(K)|=1\). Fix \(p,r\) as in the lemma. Choose
an interval far enough that \(r<t(V)\) and its two moving
ramified primes exceed \(p\). Its upper field has degree a
power of two and its modulus is supported at two and at those
two moving primes. Consequently
\(p\nmid su_1N_*[E_*:\QQ]|\mathrm{Cl}(K)|\), as required in
Theorem~\ref{thm:adjacentweights427}. The trivial tame character
has a singleton residue orbit, coefficient ring \(\ZZ_p\) and
the complementary sign at this even weight. The strict scalar
group in the comparison is \(S_{p,r}\) by the local calculation
above. In particular, this group is independent of the chosen
interval.

For this one character, the proof of
Theorem~\ref{thm:adjacentweights427} applies the intersection
inequality of Proposition~\ref{prop:localmult399} before it sums
over characters or weights. With the multiplicities of
Lemma~\ref{lem:markmultiplicity425}, this bounds the strict
length by twice the sum of \(v_p(b-d_r)\) over nonsplit
geometric trace points. Equivalently, each of these points
occurs once in the determinant of multiplication by the good
Frobenius difference on the torsion-free augmentation ideal,
as in Theorems~\ref{thm:strictlarge425} and~\ref{thm:originalorder424}.
The hypothesis \(p-1\nmid2r\) is exactly what makes the
three residual inertia characters distinct. If the strict
group is zero we can take the empty submultiset.

It remains to verify the level at which we count the points.
At an irreducible point Proposition~\ref{prop:alloriginalpoints424}
gives a conjugate-self-dual cuspidal representation over \(K\).
At each finite place away from \(p\) its Galois representation
has finite inertia of order prime to \(p\), by the tame condition
in Theorem~\ref{thm:adjacentweights427}. The three residual
characters on this inertia are trivial, since the tame character
is trivial and the cyclotomic character is unramified there.
For every inertia element, all its eigenvalues thus reduce to one.
They have order prime to \(p\), and reduction is injective on
roots of unity of such order. The element is diagonalizable
in characteristic zero since it has finite order. It follows
that the entire inertia group acts trivially. The Tate twist
used in descent is unramified at these places. Thus the associated
Weil--Deligne parameters have trivial inertia and zero monodromy.
At places above \(p\), crystallinity from
Proposition~\ref{prop:filteredalgebra427} gives the same assertion
about the Weil--Deligne parameters. The local--global compatibility,
weight calculation and spherical level criterion used in
Theorem~\ref{thm:auxdescent361} now give hyperspecial fixed
vectors at every finite place and coefficient weights of absolute
value at most \(4(r+2)\). In particular the representation has
vectors of level \(U_K(2)\).

At each remaining nonsplit point
Proposition~\ref{prop:planelabels379} gives a primitive cuspidal
form of weight \(2r+1\) and level equal to the conductor of
the quadratic character of \(K\), namely three. Its normalized
trace in the arithmetic convention is \(1+u^{-r}a_u(f)\).
Both this trace and the inverse convention are among the plane
labels of Theorem~\ref{thm:traces} at \(N=2\), since \(3\mid 2|D_K|\).
For both types of points the full embedded automorphic label with
Frobenius convention determines at most one trace point by
Proposition~\ref{prop:alloriginalpoints424}. Distinct labels
which have the same numerical trace remain separate factors
in the multiset. This proves the asserted inclusion with
multiplicity and hence the inequality.
\end{proof}

\begin{theorem}[A polynomial logarithmic bound for the rational groups]
\leavevmode\label{thm:rationalpolynomial538}
There are absolute constants \(C,A>0\) and, for each even
\(r\ge2\), a positive integer \(H_r\) such that
\[
       |K_{2r}(\ZZ)|\mid H_r^2,\qquad
       \log H_r\le C(r+2)^A .
\]
Consequently the normalized sum of strict lengths in
\eqref{eq:qbound102} is at most \(2C(r+2)^A\), uniformly
in the coefficient rings and in the finite set of primes.
\end{theorem}
\begin{proof}
Let \(H_r\) be the absolute value of the integer \(I_{u,H,r}\)
of Theorem~\ref{thm:integer} for the data of
Lemma~\ref{lem:rationallevel538}, where \(H=(\ZZ/2\ZZ)^\times\).
The integer is the product, over the full labelled list with
conjugates, of \(u^{18h+1}(b-d_r)\), with \(h=4(r+2)\).
Each factor is an algebraic integer and Galois permutes the
factors. None vanishes since \(|b|\le3<d_r\) under complex
embeddings. Theorem~\ref{thm:integer} bounds its logarithm
as asserted because \(u,K,N\) are fixed. For every \(p\ge11\)
with \(p-1\nmid2r\), all factors are integral at \(p\) and the
power of \(u\) is a unit. The lemma gives
\[
       \operatorname{length}_{\ZZ_p}S_{p,r}\le2v_p(H_r).
\]

By Proposition~\ref{prop:qbound102}, \(S_{p,r}\) is dual
to the \(p\)-part of \(K_{2r}(\ZZ)\) for odd \(p\), this part
is zero if \(p-1\mid2r\), and there is no factor at two.
For completeness the primes five and seven also contribute
zero at every even \(r\). Their maximal real cyclotomic fields
have degrees two and three and respective discriminants
\(5\) and \(7^2\) by the conductor--discriminant formula.
Their Minkowski bounds are \(\sqrt5/2\) and \(6\cdot7/27\),
both less than two. Thus both real class groups are trivial,
and Proposition~\ref{prop:realclass100} makes the corresponding
quotient of \(K_{2r}(\ZZ)\) modulo \(p\) zero. The group is finite,
so its \(p\)-part is zero. The prime three always satisfies
\(p-1\mid2r\). Multiplication of the local divisibilities gives
the asserted divisibility of ordinary integers.
Equation~\eqref{eq:qbound102} then gives the coefficient statement.
\end{proof}

\begin{corollary}[Rational classes on the period fields]
\leavevmode\label{cor:rationalperiod538}
Let the original character \(\psi\) be arbitrary and let
\(E_x\) be its fields on \([V,V^{3/2}]\). For each odd \(p\)
and each even \(1\le r<t=\lfloor x^{401}\rfloor\), let
\(D_{p,r}(E_x)\) be the image of \(S_{p,r}\) under restriction.
Then, uniformly on these intervals,
\[
 \sum_{p>2}\ \sum_{\substack{1\le r<t\\2\mid r}}
      \log |D_{p,r}(E_x)|\ \ll\ (t+2)^{A+1}.
\]
For every sufficiently distant interval,
\(R_{p,r}^{\operatorname{Gal}(E_x/\QQ)}=0\) for every odd
\(p\le P(x)\) and every \(1\le r<t\), with \(R\) as in
Proposition~\ref{prop:res23}.
\end{corollary}
\begin{proof}
Restriction preserves vanishing at the places over \(p\) and
unramifiedness elsewhere. Each image has order at most that
of its source, so summing the preceding theorem over the weights
proves the first assertion. Put \(G=\operatorname{Gal}(E_x/\QQ)\).
For odd \(r\) complex conjugation acts as minus one on the
complementary summand, and two is invertible on its \(p\)-group.
There are therefore no invariants. For even \(r\) with
\(p\nmid|G|\), the norm divided by \(|G|\) is an integral
projection onto invariants and commutes with \(J_{p,r}\).
On strict groups, restriction identifies \(S_{p,r}\) with
the invariants over \(E_x\). Indeed the two restriction and
corestriction composites are \(|G|\) and the sum of \(G\).
Corestriction preserves strictness, since its local restrictions
are sums of local corestrictions and a strict class is zero at
every finite place. The latter assertion follows by the same
unramified calculation as above with the local residue cardinality.
Thus the invariant part of \(U_{p,r}\) has length at most
\(2C(r+2)^A/\log p<b_0\) for large \(V\). In taking
\(U_{p,r}/U_{p,r}[p^{b_0}]\), the integral averaging projection
is exact and its entire image is killed by the truncation.

Finally, every odd prime dividing \(|G|\) divides the order of
\(\psi\) and belongs to the fixed set \(\Pi\) in
Proposition~\ref{prop:res23}. At each weight for such a \(p\),
\(J_{p,r}\) is the average over a subgroup of two-power order.
This average is the identity on every \(G\)-invariant element
of the quotient as well, whereas it is zero on
\(U_{p,r}=(1-J_{p,r})M_{p,r}\) and on all of its quotients
by \(G\)-stable subgroups. This proves the assertion at these
primes.
\end{proof}

The remaining complementary classes can transform under nontrivial
characters of the period field. The preceding argument treats the
trivial character and supplies a single integer at each weight on
every interval, independently of its moving ramified primes.


\subsection{Odd quadratic and quartic characters on the common fields}
\label{sec:oddquartic539}

Let \(S(\alpha,r)\) be the strict group on \(\QQ\) with coefficients
\((J/O)(\alpha^{-1}\kappa_p^{-r})\), where \(\kappa_p\) is the
\(p\)-adic cyclotomic character and \(O\) is the ring generated
over \(\ZZ_p\) by the values of the finite character \(\alpha\)
in its fraction field \(J\). We use zero localization at \(p\)
and zero restriction to inertia elsewhere. The extra conductor
conditions in Theorem~\ref{thm:levelstrict345} define a subgroup
of this group. Write \(\eta_p\) for the unique quadratic character
of conductor \(p\). At \(r\equiv(p-1)/2\bmod(p-1)\) the reduction
of \(\kappa_p^{-r}\) is \(\eta_p\).

\begin{proposition}[Two conductor bounds at the half residue]
\leavevmode\label{prop:quarticconductor539}
Let \(\alpha\) have primitive conductor \(a\), let \(p\) be
an odd prime with \(p\nmid a\), and suppose
\(\alpha(-1)=(-1)^r\) and \(r\equiv(p-1)/2\bmod(p-1)\).
\begin{enumerate}
\item If \(\alpha\) is quadratic, then, given \(0<\delta<4\),
one has \(S(\alpha,r)=0\) whenever \(a\le(4-\delta)p\) and
\(p\) is sufficiently large in terms of \(\delta\) alone.
\item Suppose \(\alpha\) is odd of order four. Write
\(a_2=\operatorname{cond}(\alpha^2)\) and
\(Y=\max\{1,3a\sqrt{a_2}\,p/32\}\).
If \(Y(1+\log Y)^3<p^2\), then \(S(\alpha,r)=0\).
\end{enumerate}
For each fixed odd quadratic or quartic character these
conclusions cover every sufficiently large prime at which an
odd complementary weight lies in the half residue.
\end{proposition}
\begin{proof}
Let \(\lambda=\alpha^{-1}\eta_p\) and let \(L\) be the field
cut out by \(\lambda\). The parity condition makes \(L\) totally
real. The two conductor factors are coprime, and inertia at \(p\)
has image of order two. The residual character \(\bar\lambda\)
has the same image order \(m\) as \(\lambda\) since \(m=2\)
or \(4\) is prime to \(p\). If the strict group is nonzero,
Lemma~\ref{lem:collisionclass457} gives
\[
              p^{\operatorname{ord}_m(p)}\mid h_L ,
\]
where \(h_L\) is the ordinary class number. To recall the
divisibility, a strict class of order \(p\) lifts to a residual
class. The coefficient invariants vanish globally and at \(p\)
because some inertia element acts by a scalar different from
one modulo \(p\). Thus the residual class is still zero locally
at \(p\). At the other finite places the inertia action on the
coefficient group has order prime to \(p\). Its invariants
are either the divisible group or zero, so passage to its
\(p\)-torsion preserves the unramified condition. Restriction
to \(G_L\) is injective by averaging over the quotient of order
\(m\). The resulting homomorphism factors through the ordinary
class group by reciprocity, since it is unramified everywhere
and \(p\) is odd. Equivariance makes its nonzero image a vector
space over the field generated by a primitive \(m\)-th root
in characteristic \(p\), of degree \(\operatorname{ord}_m(p)\).

In the quadratic case the conductor--discriminant formula
gives \(D=\operatorname{disc}L=ap\). We record the uniform
estimate
\[
       h_L\le
       \left(\frac12+
        O\!\left(\frac{\log\log(3D)}{\log(3D)}\right)\right)\sqrt D .
\]
Indeed, if \(\epsilon>1\) is a fundamental unit, then
\(\epsilon-\epsilon'\) is a nonzero integral multiple of
\(\sqrt D\). Since \(|\epsilon'|=\epsilon^{-1}\), it follows
that \(\epsilon+\epsilon^{-1}\ge\sqrt D\), whence
\(2\log\epsilon\ge\log D-O(1/D)\) as \(D\to\infty\).
For the primitive character \(\chi_D\), its sum on an interval
is \(O(\sqrt D\log(2D))\). This follows by the additive Fourier
expansion modulo \(D\), since the primitive Gauss sums have
absolute value \(\sqrt D\) at unit frequencies and vanish at
the other frequencies, and the geometric sum at frequency
\(1\le j<D\) has absolute value at most \(D/(2\min(j,D-j))\).
Partial summation at \(y=\lceil\sqrt D\log^2(2D)\rceil\) bounds
the tail of \(L(1,\chi_D)\) by \(O(1/\log(2D))\).
The initial sum is at most \(1+\log y\). The quadratic
class number formula \(2h_L\log\epsilon=\sqrt D L(1,\chi_D)\)
now gives the estimate. Under the displayed condition on \(a\),
\(h_L/p\le(1/2+o(1))\sqrt{4-\delta}<1\) uniformly as
\(p\to\infty\). This contradicts the divisibility.

In the quartic case \(r\) is odd, hence \(p\equiv3\bmod4\) and
\(\operatorname{ord}_4(p)=2\). The three nontrivial characters
of \(L\) have conductors \(ap,a_2,ap\), respectively. Thus
\(\operatorname{disc}L=a^2a_2p^2\). Minkowski's bound for an
integral ideal in each class of this real quartic field is
\((4!/4^4)\sqrt{\operatorname{disc}L}\le Y\). The number of
ideals of norm \(k\) is at most the number of ordered
factorizations of \(k\) into four positive factors, by comparison
of Euler factors with those of \(\zeta(s)^4\). Sum first over
three factors, bounding the number of choices of the last by
\(Y\) divided by their product. This gives
\(h_L\le Y(1+\log Y)^3\), contradicting \(p^2\mid h_L\) under
the stated condition. At fixed \(a\), this upper bound is
\(O_a(p(1+\log p)^3)=o(p^2)\).
\end{proof}

Let \(C_\infty=\bigcup_j C_j\) be the real cyclotomic
\(\ZZ_2\)-extension of \(\QQ\), and let \(D\) be the subfield of
\(K_\psi C_\infty\) cut out by all its characters of order
dividing four. This is a fixed finite abelian field of two-power
degree. Indeed its ambient Galois group has an open procyclic
subgroup isomorphic to \(\ZZ_2\), whose fourth powers form
another open subgroup. Every character in question kills those
fourth powers, and the resulting joint character quotient has
exponent dividing four. In particular \(D\subset B_J=K_\psi C_J\)
for all sufficiently large \(J\).

\begin{theorem}[All odd weights on the fixed subfield of exponent four]
\leavevmode\label{thm:oddquarticfield539}
Take the actual fields and intervals of
Theorem~\ref{thm:conductor232}, and set \(t=\lfloor x^{401}\rfloor\).
There is a constant \(A\), depending only on the original data,
such that uniformly there
\[
 \sum_{p>2}\ \sum_{\substack{1\le r<t\\ r\ {\rm odd}}}
   \log\big|\mathcal S_{p,r}(D)^{c=-1}\big|
        \ \ll_{\psi,\varepsilon}\ V^A+\exp(x^{2\varepsilon/3}).
\]
For every such odd weight and every odd \(p\le P(x)\), one has
\[
                 R_{p,r}(E_x)^{\operatorname{Gal}(E_x/D)}=0
\]
on every sufficiently distant interval.
\end{theorem}
\begin{proof}
There are finitely many characters of \(D\). Its odd characters
have order two or four. Choose a fixed \(p_0\) exceeding their
conductors, all primes in the finite set of
Lemma~\ref{lem:fullinteger517}, the prime divisors of
\([K_\psi:\QQ]\), and the thresholds in
Proposition~\ref{prop:quarticconductor539} for these finitely
many characters. At each \(p>p_0\) and each odd half-residue
weight the full minus summand on \(D\) vanishes. To see the
character comparison at this fixed field, extend coefficients
to contain fourth roots of unity. The degree of \(D\) is a
unit at \(p\), so the character projections on the cochains
are integral. Restriction identifies each scalar strict group
with the corresponding character summand over \(D\) by inflation
and restriction with a prime-to-\(p\) quotient. This holds on
local decomposition groups at \(p\) and on inertia elsewhere
as their quotients also have order prime to \(p\). Equivalently
one can use Shapiro's lemma with the regular representation
and then project to a character. Each odd character summand is
therefore zero by the proposition. The scalar extension used here
is faithfully flat.

The zero residue modulo \(p-1\) contains no odd weights.
Thus every other possible pair with \(p>p_0\) satisfies
\(p-1\nmid2r\). For all sufficiently distant intervals
\(D\subset E_*\), the actual upper endpoint field. The full
orders at all these separated pairs on \(D\) have product
dividing \(J(V)\) by Lemma~\ref{lem:fullinteger517}, because
\(t\le\lfloor V^{1203/2}\rfloor\). The logarithmic sum for
\(p>p_0\) is at most \(\log J(V)\ll V^A\).

For \(2<p\le p_0\) restrict from \(D\) to the actual \(B_J\)
inside \(E_x\). Its kernel is killed by \([B_J:D]\), whose
\(p\)-adic valuation is at most \(s_p=v_p([K_\psi:\QQ])\).
For fixed \(p,D\), the order of
\(\mathcal S_{p,r}(D)[p]\) is bounded independently of \(r\).
Indeed it is bounded by the order of the global degree-one
group with coefficients \(\mathbb F_p(-r)\) unramified away
from \(p\), using the multiplication-by-\(p\) sequence.
Restriction to \(D(\mu_p)\) is injective since the relative degree
divides \(p-1\). On this field the coefficient module is isomorphic
to \(\mu_p\). Kummer theory puts the unramified-away-from-\(p\)
group between the group of \(p\)-units modulo \(p\)-th powers
and the \(p\)-torsion of the ideal class group with primes over
\(p\) inverted. Both are finite and independent of \(r\).
Successive multiplication by \(p\) now bounds the logarithm
of the kernel on restriction to \(B_J\) by a constant depending
only on \(p,D,s_p\). Summing these losses costs \(O_\psi(t)\).
The image lies in the corresponding complementary group on \(B_J\).
For sufficiently large \(V\) the growing prime range of
Theorem~\ref{thm:log147} contains every \(p\le p_0\).
Its bound on the same interval gives total logarithmic order
of all these images at most \(\exp(x^{2\varepsilon/3})\).
This proves the first assertion.

For the second let \(H=\operatorname{Gal}(E_x/D)\). When
\(p\nmid |H|\), restriction and corestriction identify the
\(H\)-invariants in the strict group with the strict group on
\(D\), as in Corollary~\ref{cor:rationalperiod538}. Its minus
part has length bounded by the preceding display divided by
\(\log p\), including the implied constant. This is less than
\(b_0=\lfloor\exp(x^\varepsilon/8)\rfloor\) uniformly for large
\(V\), since \(V\le x\) and \(x^{2\varepsilon/3}=o(x^\varepsilon)\).
The averaging operator over \(H\) commutes with all the field
projections and passes exactly to invariants in their quotients.
The truncation in Proposition~\ref{prop:res23} thus kills the
whole invariant image. If \(p\mid |H|\), then \(p\) belongs
to the fixed set \(\Pi\) there. The field of its averaging
operator \(J_{p,r}\) contains \(D\), since both candidate
fields contain any fixed subfield of \(K_\psi C_\infty\) for
sufficiently large \(V\). The subgroup used in this average
is contained in \(H\). Its average fixes every \(H\)-invariant
class and acts as zero on the entire residual quotient by
definition. This also proves vanishing.
\end{proof}

\begin{corollary}[Every complementary weight over a fixed imaginary quadratic field]
\leavevmode\label{cor:fixedimaginary539}
Let \(F\subset K_\psi C_\infty\) be any imaginary quadratic
field. On the same intervals, for some fixed \(A'\),
\[
 \sum_{p>2}\sum_{1\le r<t}
   \log\big|\mathcal S_{p,r}(F)^{c=(-1)^r}\big|
           \ll_{\psi,\varepsilon} V^{A'}+\exp(x^{2\varepsilon/3}).
\]
The sum divided by \(\exp(x^\varepsilon/2)\) tends uniformly to
zero. At every odd \(p\le P(x)\) and every weight \(1\le r<t\),
the invariants of the actual \(R_{p,r}(E_x)\) under
\(\operatorname{Gal}(E_x/F)\) are zero for large \(V\).
\end{corollary}
\begin{proof}
For odd weights \(F\subset D\) and restriction to \(D\) is
injective since \([D:F]\) is a power of two. Invariants under
\(\operatorname{Gal}(E_x/F)\) are in particular invariants
under \(\operatorname{Gal}(E_x/D)\). Apply the theorem.
At an even weight \(c\) generates \(\operatorname{Gal}(F/\QQ)\),
so the plus group on \(F\) is the rational strict group by
prime-to-\(p\) descent. Sum Theorem~\ref{thm:rationalpolynomial538}
over these weights and use \(t\le V^{1203/2}\). An invariant
residual class under \(\operatorname{Gal}(E_x/F)\) is already
fixed by \(c\) at this weight. Together these elements generate
\(\operatorname{Gal}(E_x/\QQ)\). Corollary~\ref{cor:rationalperiod538}
therefore gives zero. Finally both terms of the displayed
upper bound are \(o(\exp(x^\varepsilon/2))\) since
\(x\ge V\to\infty\).
\end{proof}

The bound for \(D\) concerns all its odd characters and all odd
weights on the original intervals. The quadratic fields added
at the two moving ramified primes need not be contained in \(D\).


\subsection{Normalized Frobenius values on the original period fields}
\label{sec:normalizedstrict600}

We now evaluate the full groups \(\mathcal S_{\ell,r}(E_x)\)
of Section~\ref{sec:coho12}. Fix a common algebraic closure
of \(\QQ\). First let \(E/\QQ\) be a finite abelian field
unramified at the odd prime \(\ell\), with
\(\ell\nmid[E:\QQ]\), and let \(r\ge1\). Write
\(M=(\QQ_\ell/\ZZ_\ell)(-r)\), \(\chi=\kappa_\ell^{-r}\),
and \(S=\mathcal S_{\ell,r}(E)\). We inflate cocycles
to \(G_E\). Choose a decomposition group \(D_v\subset G_E\)
at one place over \(\ell\). For a finite abelian
\(\ell\)-group \(U\), put
\(U^\vee=\operatorname{Hom}(U,\QQ_\ell/\ZZ_\ell)\),
using the additive coefficient group.

\begin{lemma}[A normalization at the coefficient prime]
\leavevmode\label{lem:normalizedstrict600}
Every \(c\in S\) has precisely one representing cocycle
\(y_c\) which is zero on \(D_v\). The correspondence is
additive, and a class killed by \(\ell^a\) has normalized
values in \(M[\ell^a]\). With \(H=G_{\QQ^{\rm ab}}\)
there is a surjection
\[
    H\longrightarrow S^\vee,\qquad
           n\longmapsto\bigl(c\longmapsto y_c(n)\bigr).
\]
\end{lemma}
\begin{proof}
By strictness, subtracting a global coboundary makes a
representative zero on \(D_v\). Locally \(E_v/\QQ_\ell\)
is unramified, whereas the \(\ell\)-power cyclotomic tower
of \(\QQ_\ell\) is totally ramified. Thus the cyclotomic
image of inertia is all of \(\ZZ_\ell^\times\), also
over \(E_v\). The invariants of \(D_v\) in \(M\) are
therefore the same as the global invariants. A global
coboundary vanishing on \(D_v\) is identically zero.
This proves uniqueness; addition and multiplication by
\(\ell^a\) preserve the normalization.

The group \(H\) is contained in \(G_E\) and acts
trivially on \(M\), since all roots of unity belong to
\(\QQ^{\rm ab}\). Suppose \(y_c\) vanishes on \(H\).
It descends to a cocycle on the abelian quotient
\(B=\operatorname{Gal}(\QQ^{\rm ab}/E)\). On this
quotient, for \(z\) in the image of inertia at \(v\)
and \(g\in B\), the equation \(zg=gz\) and \(y_c(z)=0\)
give \((\chi(z)-1)y_c(g)=0\). The inertia image under
\(\chi\) is the full global image. Thus all values are
globally invariant, and the descended cocycle is a
homomorphism. Its image is finite, because it is a
continuous map from a compact group to a discrete group.
It vanishes on every inertia group away from \(\ell\)
by the original unramified condition. All inertia images
in \(B\) at places above \(\ell\) agree, since
\(\QQ^{\rm ab}/\QQ\) is abelian and \(E/\QQ\) is Galois.
The homomorphism is unramified there as well.

A nonzero image would give an extension \(E'/E\) of
degree \(\ell\), unramified at all finite places, with
\(E'/\QQ\) abelian. The \(\ell\)-primary subgroup of
\(\operatorname{Gal}(E'/\QQ)\) would have order \(\ell\).
Each finite inertia group injects into
\(\operatorname{Gal}(E/\QQ)\) and consequently projects
trivially onto that primary subgroup. The degree-\(\ell\)
field fixed by the complementary subgroup would be
unramified at every finite prime of \(\QQ\), contradicting
Minkowski's discriminant bound. We conclude that
\(y_c=0\). The group \(S\) is finite by the argument of
Proposition~\ref{prop:finite15} for odd \(\ell\).
The evaluation map on \(H\) is a homomorphism, and the
argument just given says its image has zero annihilator
in \(S\). A proper subgroup of the finite group \(S^\vee\)
has nonzero annihilator by finite abelian duality.
This proves surjectivity.
\end{proof}

\begin{theorem}[Full strict groups at one list of rational primes]
\leavevmode\label{thm:normalizedprimes600}
Take a nonempty finite list of triples \((E_i,\ell_i,r_i)\)
as above, writing \(S_i=\mathcal S_{\ell_i,r_i}(E_i)\).
Let \(C/\QQ\) be any finite abelian extension containing
all the \(E_i\), and prescribe an element
\(\sigma_0\in\operatorname{Gal}(C/\QQ)\) fixing them
all. Put \(b=\#\{i\}\), \(d_i=\dim_{\mathbb F_{\ell_i}}
S_i[\ell_i]\), and \(d=\max_i d_i\). There exist
\[
               h=d+\lceil\log_2(2b)\rceil
\]
distinct good rational primes with this prescribed
Frobenius in \(C\). At one compatible choice of places
above them let \(s_1,\ldots,s_h\) be their Frobenius
elements. For every \(i\) the normalized values give
an injection
\[
       S_i\ \longrightarrow\
        (\QQ_{\ell_i}/\ZZ_{\ell_i})^h,\qquad
             c\longmapsto (y_c(s_j))_{1\le j\le h}.
\]
In particular the exact additive order of each class
is the largest of the orders of its displayed values.
One can avoid any finite set of primes.
\end{theorem}
\begin{proof}
Lift \(\sigma_0\) to \(\sigma\in G_\QQ\). It lies in
every \(G_{E_i}\). For Haar-uniform \(n\in H\), the
functional \(c\mapsto y_c(n)\) is uniform on \(S_i^\vee\)
by Lemma~\ref{lem:normalizedstrict600}. The functional
at \(\sigma n\) is obtained by multiplication by the
unit \(\chi_i(\sigma)\) and addition of the fixed
functional at \(\sigma\), by the cocycle identity.
It too is uniform. Take \(h\) independent such elements
of the coset. For a fixed \(i\), their reductions modulo
\(\ell_i S_i^\vee\) are independent uniform vectors in
a vector space of dimension \(d_i\). If these vectors
fail to span, they lie in a hyperplane. For \(d_i>0\)
there are \((\ell_i^{d_i}-1)/(\ell_i-1)\) hyperplanes;
the probability that all vectors lie in any particular
one is \(\ell_i^{-h}\). For \(d_i=0\) failure has
probability zero. Summing the resulting bounds over
the indices gives at most \(b\,2^{d-h}<1\).
This uses no independence between different groups.
Fix a tuple for which every reduction spans.
The integral functionals generate \(S_i^\vee\):
the quotient by their span is a finite \(\ell_i\)-group
equal to its multiple by \(\ell_i\), and is therefore
zero. Finite duality now gives the asserted injections
at the chosen group elements.

For each \(i\) take the exponent \(\ell_i^{a_i}\)
of \(S_i\), omitting trivial groups. All normalized
representatives have values killed by \(\ell_i^{a_i}\).
The finite semidirect product recording the character
modulo \(\ell_i^{a_i}\) and the values of the cocycles
on a generating set of \(S_i\) gives a finite quotient
of \(G_{E_i}\). This follows directly from the cocycle
law. Take the compositum of the normal closures of its
fixed fields over \(\QQ\) for all \(i\), together with
\(C\). For each chosen coset element, Chebotarev gives
infinitely many rational primes with that conjugacy
class in the finite Galois group. A choice of place
realizes the specified element itself. All the elements
restrict to \(\sigma_0\) in \(C\), and every normalized
evaluation is preserved. Excluding previous choices
and all bad primes proves the assertion. Orders are
preserved by an injection, and the order of a tuple of
elements of \(\ell_i\)-power order is their largest order.
\end{proof}

\begin{corollary}[The weighted lengths on the common interval]
\leavevmode\label{cor:originalinterval600}
Use the fields in Section~\ref{sec:sourcelevels563} on
\([V,V^{3/2}]\), with \(N^\natural\) the common
conductor modulus and \(n_*=[E_*:\QQ]\) at the largest
layer. Fix \(A_0>0\). For all the distinct triples
\[
 E=E_x,\qquad 1\le r<t(x),\qquad
 2<\ell\le V^{A_0},\qquad \ell\nmid n_*N^\natural ,
\]
one list of \(V^{O(1)}\) rational primes in a prescribed
compatible finite abelian Frobenius class suffices in
Theorem~\ref{thm:normalizedprimes600}. In particular they
can all split in every period and starting field, with
any given cyclotomic depth at each coefficient prime.
The number of primes has the stated bound independently
of those depths. For each triple write \(s_j\) for the
resulting Frobenius at the chosen place, and put
\[
 u_{\ell,r,E,j}=
    \log_\ell\#\{y_c(s_j):c\in\mathcal S_{\ell,r}(E)\}.
\]
For every subset of these triples the sum of
\((\log\ell)\operatorname{length}_{\ZZ_\ell}
\mathcal S_{\ell,r}(E)\) over the subset is at most the
sum of \((\log\ell)\sum_{j=1}^h u_{\ell,r,E,j}\)
over that same subset. Also
\(\ell^{\max_j u_{\ell,r,E,j}}\) is exactly the
exponent of the corresponding strict group.
\end{corollary}
\begin{proof}
All degrees divide \(n_*\), and each conductor divides
\(N^\natural\). Thus each triple satisfies the lemma.
The number of different fields on the interval is at
most \(1+\log_2 n_*\), since consecutive distinct fields
in this nested tower differ by degree at least two.
The number of twists is at most \(V^{1203/2}\), and
the number of coefficient primes is at most \(V^{A_0}\).
The degree \(n_*\) and \(\log(2N^\natural)\) are
polynomially bounded on the intervals of
Theorem~\ref{thm:conductor232}. The Kummer estimate in
the proof of Theorem~\ref{thm:trunc15} gives uniformly
\[
 d_i\log\ell_i\ \ll\
        n_*\ell_i\log(2N^\natural\ell_i).
\]
It bounds the dimension of the \(\ell_i\)-torsion via
the units and the class group of \(E_i(\mu_{\ell_i})\).
Hence both \(d\) and \(h\) meet the claimed bound.
For the splitting assertion take \(C\) to contain all
the period fields, the finite abelian starting fields
of Corollary~\ref{cor:polarsplit593}, and the indicated
roots of unity, and take \(\sigma_0=1\).

Each coordinate image is a finite subgroup of
\(\QQ_\ell/\ZZ_\ell\), hence cyclic of order
\(\ell^{u_{\ell,r,E,j}}\). The group embeds into the
product of those cyclic images. Taking logarithms of
their orders and summing gives the stated inequality.
The exponent of the subgroup is at most the largest
coordinate order and at least the exponent of each
coordinate image, proving equality.
\end{proof}

The value \(y_c(s_j)\) uses the normalization at the
chosen place over \(\ell\). Ordinary localization at a
good split rational prime \(q\) would take its value in
\(M/(q^{-r}-1)M\), which is zero: multiplication by the
nonzero \(\ell\)-adic scalar \(q^{-r}-1\) is surjective
on \(\QQ_\ell/\ZZ_\ell\). The theorem uses finite
quotients determined by the actual strict groups. It
gives no upper bound for the orders in the coordinate
images or the sizes of the rational primes chosen.
An estimate for those orders, or an integral formula
relating them to the projected periods, is still needed
for the remaining lengths in Section~\ref{sec:missing}.


\subsection{The genus field contribution to the full strict length}
\label{sec:genusstrict601}

Let \(E/\QQ\) be finite abelian. Write \(T_E\) for its
genus field, the maximal extension of \(E\) abelian over
\(\QQ\) and unramified at all finite places over \(E\).
For every prime \(\ell\), including two, and \(r\ge1\), set
\(M=(\QQ_\ell/\ZZ_\ell)(-r)\), \(H=G_{\QQ^{\rm ab}}\), and
\[
 S=\mathcal S_{\ell,r}(E),\qquad
 K_{\ell,r}(E)=
 \ker\bigl(S\longrightarrow\operatorname{Hom}_{\rm cont}(H,M)\bigr).
\]
The restriction makes sense on classes since \(H\) acts
trivially on \(M\). We compute this kernel without a
restriction on the ramification of \(E\) at \(\ell\).

\begin{lemma}[The genus field on the varying interval]
\leavevmode\label{lem:genusfield601}
Let \(\psi=\prod_{p\mid F}\psi_p\) be the decomposition
of the original primitive character into its prime-power
conductor factors. Put
\[
   T_\psi=\prod_{p\mid F} K_{\psi_p},\qquad
   g_\psi=\frac{\prod_{p\mid F}\operatorname{ord}(\psi_p)}
                   {\operatorname{ord}(\psi)} .
\]
Here a product of fields is their compositum, and the
empty product gives \(T_\psi=\QQ\) and \(g_\psi=1\).
For the period fields of Theorem~\ref{thm:conductor232},
\[
        T_{E_x}=E_xT_\psi,\qquad [T_{E_x}:E_x]\mid g_\psi .
\]
\end{lemma}
\begin{proof}
We give the Dirichlet character description of the genus
field using finite primes throughout.
\href{https://doi.org/10.1002/mana.19530090604}{Leopoldt (1953)}.
For an abelian field \(E\) inside \(\QQ(\mu_N)\), let
\(X\) be its group of Dirichlet characters. For each
\(p\mid N\), let \(X_p\) consist of their components at
\(p\), viewed as characters of the single \(p\)-power
factor. In the cyclotomic Galois group the inertia group
at \(p\) is precisely this local factor: the \(p\)-power
cyclotomic extension is totally ramified there and all
the other factors are unramified. Its image in the
Galois group of \(E\) has character group \(X_p\), by
restriction of characters of a finite abelian group.
Thus \(e_p(E/\QQ)=|X_p|\).

Let \(T\) have character group generated by all the \(X_p\).
It contains \(E\), since each character is the product of
its components. Distinct factors are independent, so
\([T:\QQ]=\prod_p |X_p|\), and its inertia index at
each \(p\) is still \(|X_p|\). It follows that \(T/E\)
is unramified at every finite place. Conversely, if
\(B/E\) is any finite extension unramified at finite
places with \(B/\QQ\) abelian, then all its finite
inertia groups have orders \(e_p(E/\QQ)\). They generate
\(\operatorname{Gal}(B/\QQ)\), since their fixed field
is unramified at all finite primes of \(\QQ\) and hence
is \(\QQ\) by Minkowski's bound. The group is abelian,
so \([B:\QQ]\le\prod_p e_p(E/\QQ)\). Applying the same
bound to \(BT\), which is still unramified over \(E\),
gives \(B\subset T\). This proves \(T=T_E\).

The characters of \(E_x\) are generated by \(\psi\)
and the characters of its three factors
\(C_J,\QQ(\sqrt{-q}),\QQ(\sqrt{-d_0})\).
Each of the latter characters is supported at a single
prime, respectively \(2,q,d_0\). For the quadratic
fields this follows from \(q,d_0\equiv3\bmod4\), so
their discriminants are \(-q,-d_0\). Taking all
local components therefore adjoins exactly the
\(\psi_p\) to this group. This gives \(T_{E_x}=E_xT_\psi\).
The single-prime components of \(\psi\) are independent
in the product of cyclotomic factors, so
\([T_\psi:K_\psi]=g_\psi\). Since \(K_\psi\subset E_x\),
the degree of its base change divides \(g_\psi\).
\end{proof}

\begin{theorem}[The strict classes vanishing on the absolute commutator subgroup]
\leavevmode\label{thm:genuskernel601}
Fix a place \(v\mid\ell\) of \(E\), with absolute
decomposition group \(D_v\). Let \(U=M^{D_v}\) and set
\[
   C_{E,\ell}=\operatorname{Gal}(T_E/E)/
            \langle\operatorname{Frob}_v(T_E/E)\rangle .
\]
The action on \(U\) factors through \(C_{E,\ell}\).
Inflation induces an isomorphism
\[
     H^1(C_{E,\ell},U)\ \simeq\ K_{\ell,r}(E),\qquad
     \operatorname{length}_{\ZZ_\ell}K_{\ell,r}(E)
                    \le v_\ell(|C_{E,\ell}|).
\]
In particular for each \(x\) on the period interval,
\[
 \sum_{\ell}(\log\ell)\sum_{1\le r<t(x)}
          \operatorname{length}_{\ZZ_\ell}K_{\ell,r}(E_x)
          \ \le\ (t(x)-1)\log g_\psi .
\]
\end{theorem}
\begin{proof}
Set \(B=\operatorname{Gal}(\QQ^{\rm ab}/E)\).
In \(B\) divide out by the closed subgroup \(J\)
generated by the images of \(D_v\) and of inertia
at all finite primes away from \(\ell\). Inertia at all
places over \(\ell\) has the same image, since \(E/\QQ\)
is Galois and the larger extension is abelian over \(\QQ\).
Dividing out first by all finite inertia gives
\(\operatorname{Gal}(T_E/E)\), by the definition just
proved. The image of \(D_v\) there is generated by its
unramified Frobenius. Hence \(B/J=C_{E,\ell}\).
Inertia away from \(\ell\) acts trivially on \(M\),
so \(M^J=U\). This is a finite cyclic \(\ell\)-group:
the local cyclotomic image contains an open subgroup
of \(\ZZ_\ell^\times\), and a unit of infinite order
there has a nontrivial \(r\)-th power. Its difference
from one kills only a finite subgroup of \(\QQ_\ell/\ZZ_\ell\).

Take \(c\in K_{\ell,r}(E)\). A representative \(y\)
vanishes on \(H\), since all coboundaries do and its
restriction is zero. It therefore descends to \(B\).
Strictness permits subtraction of a coboundary so that
\(y=0\) on \(D_v\). At the remaining finite primes
away from \(\ell\) it is zero on inertia, as the
original cocycle factors through the quotient unramified
there and inertia acts trivially on \(M\). Thus the
adjusted cocycle vanishes on \(J\). Its values belong
to \(M^J\), by applying the cocycle identity to
\(jg=gj\). It descends to a cocycle of \(C_{E,\ell}\)
with coefficients \(U\). Conversely, inflating any
such cocycle gives an unramified class away from \(\ell\)
and zero restriction at every place above \(\ell\),
since the images of their decomposition groups in \(B\)
are the same. Inflation is injective: if an inflated
cocycle is the coboundary of \(m\in M\), then its
vanishing on \(J\) forces \(m\in U\). This proves the
isomorphism.

For any finite group \(C\) acting on a finite cyclic
\(\ell\)-group \(U\), restriction to an \(\ell\)-Sylow
subgroup \(P\) is injective in first cohomology, by
corestriction. We prove
\(\operatorname{length}H^1(P,U)\le\log_\ell |P|\)
by induction. For nontrivial \(P\) take a central
subgroup \(Z\) of order \(\ell\). The kernel of
restriction to \(Z\) is the inflation of
\(H^1(P/Z,U^Z)\), bounded by induction since \(U^Z\)
is cyclic. Evaluation at a generator of \(Z\) makes
\(H^1(Z,U)\) a subquotient of \(U\); restriction to
the trivial subgroup followed by corestriction shows
it is killed by \(\ell\). Its length is at most one.
The trivial group starts the induction. This proves
the stated bound and, by Lemma~\ref{lem:genusfield601},
the inequality for each \(r\) with \(v_\ell(g_\psi)\)
on the right. Summing uses
\(\sum_\ell v_\ell(g_\psi)\log\ell=\log g_\psi\).
\end{proof}

\begin{corollary}[Common split primes for every strict group up to the genus contribution]
\leavevmode\label{cor:allprimesstrict601}
Fix one of the intervals and a finite list of its triples
\((E_x,\ell,r)\), now allowing every coefficient prime.
At one list of good rational primes split in any prescribed
finite abelian field containing the \(E_x\), the normalized
evaluations give homomorphisms on each \(S\) with common
kernel exactly \(K_{\ell,r}(E_x)\). Normalization requires
the representing cocycle to vanish on the chosen \(D_v\).
The values at the selected primes are independent of
the choice of such a representative. The number of primes
may be taken as \(d+\lceil\log_2(2b)\rceil\), where
\(b\ge1\) is the number of triples and \(d\) is the largest
minimal number of generators of any \(S/K_{\ell,r}(E_x)\)
in the list. For \(\ell\le V^{A_0}\) it is polynomial
in \(V\) for each fixed \(A_0\); unrestricted prescribed
finite cyclotomic depths are allowed. A finite list can
also include every nonzero strict group on the interval,
without a size restriction on \(\ell\).
\end{corollary}
\begin{proof}
The restriction pairing on \(H\) has annihilator \(K\)
by definition. Therefore \(H\to(S/K)^\vee\) is onto,
as in Lemma~\ref{lem:normalizedstrict600}. The hyperplane
count in Theorem~\ref{thm:normalizedprimes600} selects
the stated number of group elements of \(H\) giving
generation for all the quotient duals together.
For each \(c\in S\) choose a representing cocycle with
value zero on \(D_v\). Any two choices differ by the
coboundary of an element of the finite group \(U=M^{D_v}\)
from the theorem. Add to the common splitting field
enough \(\ell\)-power roots of unity to make all the
characters at the chosen primes act trivially on these
groups \(U\), as well as the roots requested in the
statement. At such a Frobenius the difference of the two
representatives is exactly zero.

Each selected representative has finite image by compactness.
At a finite coefficient precision containing that image,
its values and character are recorded by a finite
semidirect product. There are only finitely many classes
in each group by Proposition~\ref{prop:finite15}. Apply
Chebotarev to the joint normal closures and the finite
abelian splitting field, matching each selected element
of \(H\) with a rational Frobenius and a choice of place.
All evaluations agree with their values on \(H\). Thus
they are additive and their joint kernel is precisely
\(K\). The same rank estimate from
Theorem~\ref{thm:trunc15} as in
Corollary~\ref{cor:originalinterval600} applies also
to the ramified primes and to two; a quotient of a
finite abelian \(\ell\)-group has at most as many
generators as the group itself. Finally there are only
finitely many fields and weights on each interval.
Proposition~\ref{prop:finite15} gives finite prime support
for each pair. This proves the last assertion.
\end{proof}

For these evaluations let \(I_{\ell,r}(E_x)\) denote
the image of the whole group in its product of coordinate
groups. Its order has the exact relation
\[
 \log |S|=\log |I_{\ell,r}(E_x)|
                         +\log |K_{\ell,r}(E_x)|.
\]
The total second term on an interval, summed at each
\(x\) over all primes and permitted weights, is at most
\((t(x)-1)\log g_\psi\), a polynomial in \(x\).
For a character of prime-power conductor, or for the
trivial character, \(g_\psi=1\) and the common tests
recover the full groups even at the ramified primes.
Bounding the orders of the images \(I_{\ell,r}(E_x)\)
on the scale required by Section~\ref{sec:missing}
and comparing the signed integral periods there are
the remaining arithmetic estimates.


\subsection{Local unit tests at arbitrary precision on the period fields}
\label{sec:idealprecision608}

Fix the fields \(E=E_x\subseteq E_*\) of Section~\ref{sec:missing}
on one of its intervals. Write \(\kappa=\kappa_p\) for the
cyclotomic character at an odd prime \(p\). For a CM field \(F\)
let \(D_p(F)=\mathcal S_{p,-1}(F)^{c=-1}\). Its coefficients
are \(\mu_{p^\infty}\). Theorem~\ref{thm:compare30} uses the
group \(D_p(E(\mu_p))\) through precision \(1+v_p(r+1)\).
To go beyond that precision we need descent through extensions
of \(p\)-power degree. We carry this out on the fields
\[
 B_k=E_*(\mu_{p^k})\quad(k\ge1),\qquad
 B_\infty=\bigcup_{k\ge1}B_k .
\]
They depend only on \(p\) and the chosen interval.

\begin{lemma}[Descent of strict classes in a cyclotomic tower]\leavevmode
\label{lem:cyclostrict608}
Let \(F\) be a number field, \(F_\infty=F(\mu_{p^\infty})\),
and \(A=(\QQ_p/\ZZ_p)(s)\) with \(s\ne0\). Write \(S(F,A)\)
for the classes in \(H^1(F,A)\) trivial on inertia away from
\(p\) and on the full local groups at \(p\).
If \(F\subseteq K\subset F_\infty\) is finite over \(F\), then
restriction identifies \(S(F,A)\) with
\(S(K,A)^{\Gal(K/F)}\).
\end{lemma}
\begin{proof}
Put \(\Gamma=\Gal(F_\infty/F)\). Its image in \(\ZZ_p^\times\)
is open and is the product of a group of order dividing \(p-1\)
and a procyclic copy of \(\ZZ_p\). If \(g\) generates the latter,
\(\kappa(g)^s-1\) is nonzero and multiplication by it on \(A\)
is surjective. The two-term resolution for the procyclic group
over its completed group ring, with differential \(g-1\), gives
zero in degrees one and two. Averaging over the finite factor
therefore gives \(H^1(\Gamma,A)=H^2(\Gamma,A)=0\).
The cyclotomic image at each \(w\mid p\) is also open, so its
first cohomology vanishes by the same proof.

Since \(G_{F_\infty}\) acts trivially on \(A\), inflation and
restriction identify \(H^1(F,A)\) with the \(\Gamma\)-invariants
in \(\operatorname{Hom}_{\rm cont}(G_{F_\infty},A)\).
At a place above \(p\), local injectivity says exactly that the
homomorphism must vanish on the corresponding local subgroup
of \(G_{F_\infty}\). Away from \(p\), inertia already lies in
\(G_{F_\infty}\) and must also be killed. These are the same
subgroups when one starts at \(K\): the infinite field is the
same, the places above \(p\) are all taken, and \(K/F\) is
unramified away from \(p\). Passing from the invariants for
\(\Gal(F_\infty/K)\) to the invariants for \(\Gamma\) proves
the lemma.
\end{proof}

\begin{theorem}[The exact ideal-class subgroup at every precision]\leavevmode
\label{thm:idealprecision608}
For \(r,m\ge1\) put \(k=\max\{1,m-v_p(r+1)\}\) and
\(M_{p,r}(E)=\mathcal S_{p,r}(E)^{c=(-1)^r}\). Then
\(\kappa^{r+1}\bmod p^m\) defines a character \(\xi\) of
\(\Gal(B_k/\QQ)\). With a choice of compatible Tate-module bases
one has
\[
 M_{p,r}(E)[p^m]\simeq
 \{z\in D_p(B_k)[p^m]:\sigma z=\xi(\sigma)z
                               \text{ for all }\sigma\in\Gal(B_k/E)\}.
\]
To calculate \(D_p(B_k)[p^m]\), take \(F=B_k\) in the following
map and compute its kernel:
\[
 \Cl(F)[p^m]^-\ \longrightarrow\
 \left(\prod_{w\mid p}
  \frac{\mathcal O_{F_w}^\times}
       {\mu(F_w)_p\,(\mathcal O_{F_w}^\times)^{p^m}}\right)^- .
\]
Represent the ideal class by \(\mathfrak a\) prime to \(p\),
write \(\mathfrak a^{p^m}=(a)\), and take the minus projection
of the tuple of \(a\) at \(w\mid p\). For fixed \(p,m\) one may
use \(B_m\) for every pair \(E_x,r\) on the interval.
\end{theorem}
\begin{proof}
The extension \(E_*/E\) has two-power degree. Restriction
identifies \(p\)-primary cohomology over \(E\) with the
invariants over \(E_*\) by averaging. At a place above \(p\)
restriction is injective by local corestriction. The same
holds on inertia elsewhere, since its index divides this
two-power degree. Thus this identification respects strictness,
even when a place ramifies. Apply Lemma~\ref{lem:cyclostrict608}
to \(B_k/E_*\). By transitivity of invariants we may now impose
the full \(\Gal(B_k/E)\)-action on the strict classes over \(B_k\).

For \(u\in1+p^k\ZZ_p\), binomial expansion gives
\(v_p(u^{r+1}-1)\ge k+v_p(r+1)\). This proves the assertion
about \(\xi\). On the subgroup \(G_{B_k}\), the actions
\(\kappa^{-r}\) and \(\kappa\) agree modulo \(p^m\).
Apply the proof of the lemma to both over \(B_k\).
At this precision the homomorphisms on \(G_{B_\infty}\) take
values in the same cyclic group, subject to the same local
vanishing and equivariance conditions. On conjugating such a
homomorphism by an element over \(\QQ\), the two actions on its
argument agree. The ratio of the actions on its value, from
twist \(-r\) to twist \(+1\), is \(\kappa^{r+1}\bmod p^m\).
Thus the former invariants give exactly the indicated
\(\xi\)-eigenvectors. The original sign of \(c\) becomes
\((-1)^r(-1)^{r+1}=-1\). The argument applies to every larger
\(k\) as well, giving the common layer \(B_m\).

For the ideal formula use Proposition~\ref{prop:kummer30}.
Here are the details that allow the local conditions to be
tested on any generators. Weak approximation supplies a
representative prime to \(p\). Changing the representative
multiplies \(a\) by a \(p^m\)-th power and a global unit.
For a global unit \(u\), every conjugate of the algebraic
integer \(u/c(u)\) has absolute value one. It is a root of
unity by Kronecker's theorem. Since \(2\) is invertible on
the target and a prime-to-\(p\) root is a \(p^m\)-th power,
the minus part of the global unit disappears there.
The displayed map is therefore well defined and additive.
If its value is zero, \(e^-[a]\), with \(e^-=(1-c)/2\)
on the Kummer group modulo powers, has image in
\(\mu(F_w)_p F_w^{\times p^m}/F_w^{\times p^m}\) at every
\(w\mid p\). It has divisible valuations away from \(p\).
These are precisely the conditions of Proposition~\ref{prop:kummer30},
and its ideal class is \(e^-[\mathfrak a]=[\mathfrak a]\).
The converse follows by taking a generator supplied by that
proposition. This proves the formula at all finite places.
\end{proof}

The local power test is finite. In \(F_w\), normalize the
valuation by \(v_w(\mathfrak p_w)=1\) and put \(e=v_w(p)\).
For \(b=1+\lfloor e/(p-1)\rfloor\), the logarithm and exponential
identify \(1+\mathfrak p_w^b\) with \(\mathfrak p_w^b\).
Hence \(1+\mathfrak p_w^{b+me}=(1+\mathfrak p_w^b)^{p^m}\).
One can test the images of local roots and \(p^m\)-th powers in
the units modulo \(\mathfrak p_w^{b+me}\).

\begin{proposition}[Equality of two local kernels controls all higher layers]
\leavevmode\label{prop:idealkernels608}
Let \(h\ge1\) be such that the cyclotomic image of
\(G_{E_*(\mu_p)}\) is \(1+p^h\ZZ_p\), and fix \(k\ge h\).
If \(|D_p(B_{k+1})|=|D_p(B_k)|\), then restriction identifies
\(D_p(B_k)\) with \(D_p(B_i)\) for every \(i\ge k\).
For all \(E_x,r\), the full group \(M_{p,r}(E_x)\) is then
isomorphic to the eigenspace in Theorem~\ref{thm:idealprecision608}
in \(D_p(B_k)[p^{k+v_p(r+1)}]\) at precision \(k+v_p(r+1)\).
\end{proposition}
\begin{proof}
Lemma~\ref{lem:cyclostrict608} at twist \(+1\) gives inclusions
of the finite \(D_p(B_i)\) identifying each earlier group with
the corresponding fixed subgroup. Put \(Z=\bigcup_iD_p(B_i)\)
under these maps. Choose \(\gamma\) with cyclotomic value
\(1+p^k\) generating \(\Gal(B_\infty/B_k)\).
Then \(Z^\gamma=D_p(B_k)\) and \(Z^{\gamma^p}=D_p(B_{k+1})\).
By inclusion and equality of orders these fixed subgroups agree.
Let \(T=\gamma-1\), \(N=1+\gamma+\cdots+\gamma^{p-1}\).
If \(NTz=0\), then \(\gamma^p z=z\), hence \(Tz=0\).
Consequently \(N\) acts injectively on \(TZ\).

Each \(z\) lies in a finite stable \(p\)-group on which \(\gamma\)
has \(p\)-power order, say dividing \(p^a\). On its reduction
modulo \(p\), \(T^{p^a}=0\) and \(N=T^{p-1}\). A power of \(N\)
therefore maps the finite group into its multiple by \(p\).
Iteration gives zero. Injectivity on \(TZ\) now forces \(Tz=0\)
for every \(z\). This proves stabilization. Each element of a
strict group has finite \(p\)-power order. Apply the theorem at
its precision using a sufficiently high layer. Its image in \(Z\)
is fixed by \(\gamma\), while \(\gamma\) acts on it by
\((1+p^k)^{r+1}\). The difference from one has valuation
\(k+v_p(r+1)\), again by binomial expansion. This bounds the
exponent and proves the stated full comparison.
\end{proof}

\begin{proposition}[The quadratic summand of discriminant \(-1208\)]
\leavevmode\label{prop:quadraticorders608}
Let \(F=\QQ(\sqrt{-302})\). For every positive odd \(r\),
\(M_{3,r}(F)\) is cyclic of order three or nine. It has order
three when \(r\equiv5\bmod6\). Among the two progressions
\(r\equiv1,3\bmod6\), at most one can have order nine. Its
order is constant on each progression. For the original
quadratic character of discriminant \(-1208\) these are the
full orders of the \(\Gal(E_x/F)\)-invariant summands on every
one of its period fields.
\end{proposition}
\begin{proof}
The class number of \(F\) is \(12\). For verification, a reduced
primitive positive binary quadratic form \((a,b,c')\) of
discriminant \(-1208\) has \(a\le\sqrt{1208/3}\),
\(|b|\le a\le c'\), and \(b\ge0\) on either boundary. Checking
\(4a\mid b^2+1208\) gives the pairs
\[
 (1,0),(2,0),(3,\pm2),(6,\pm4),(9,\pm4),
                 (13,\pm12),(17,\pm4).
\]
The form--ideal correspondence gives the asserted number.
Proposition~\ref{prop:example30} constructs a nonzero class
in \(D_3(F)[3]\) from the generator \(73+9\sqrt{-302}\).
Thus \(D_3(F)\) has order three. Since \(r+1\) is even,
Theorem~\ref{thm:compare30} gives
\(|M_{3,r}(F)[3]|=3\), and it gives the whole group when
\(3\mid r+1\).

For the other progressions we need an upper bound at higher
precision. Let \(\chi\) be the quadratic character of \(F\)
and \(f=1208\). Since \(3\nmid2f\) and the minus character of
\(\Gal(F/\QQ)\) is \(\chi\), Lemma~\ref{lem:shared201} applies
to the whole group under consideration. The required Bernoulli
values are
\[
\begin{array}{c|rrrrrr}
 k&3&5&9&11&15&17\\ \hline
 (B_{k,\chi}/k)\bmod27&9&18&9&18&9&18 .
\end{array}
\]
Here is an integer and rational calculation of the table. Set
\(S_j=\sum_{a=1}^f\chi(a)a^j\). Skip even \(a\) and multiples
of \(151\). For the remaining \(a\), give the sign \(+1\) to
\(a\bmod8\) equal to \(1,7\) and the sign \(-1\) to \(3,5\).
Multiply by \(+1\) or \(-1\) according as \(a^{75}\bmod151\)
is \(1\) or \(-1\). This is \(\chi\) since
\(-1208=8(-151)\). Starting with \(b_0=1\), the two finite
formulas used are
\[
 b_i=-\frac1{i+1}\sum_{j=0}^{i-1}\binom{i+1}{j}b_j,\qquad
 \frac{B_{k,\chi}}k=\frac1{kf}
       \sum_{i=0}^k\binom{k}{i}b_i f^i S_{k-i}.
\]
The latter is the defining sum of the Bernoulli polynomial
at \(a/f\). For example it gives \(839844\) and
\(-364139453484\) at the first two indices. All six entries
have valuation two. Every odd \(r\) with \(3\nmid r+1\) has
\(-r\bmod18\) in the six columns. For that \(k\),
\(2<1+v_3(r+k)\), so the cited lemma gives \(9M_{3,r}(F)=0\).
The order-three subgroup at precision one now implies cyclicity.
The finite coefficient comparison of Lemma~\ref{lem:period30}
at precision two shows that each order depends only on \(r\bmod6\).

To distinguish how many progressions can attain the upper bound,
apply Theorem~\ref{thm:idealprecision608} with \(E=E_*=F\)
and with the common field \(F(\mu_9)\) at precision two.
The proof of that theorem applies to this constant field as well.
Here \(F\cap\QQ(\mu_9)=\QQ\), since the latter field has unique
quadratic subfield \(\QQ(\sqrt{-3})\). Choose the element
\(\gamma\) over \(F\) with cyclotomic residue \(4\bmod9\).
The involution over \(F\) with residue \(-1\) fixes the desired
classes since \(r+1\) is even. In the subgroup fixed by this
involution, the elements also fixed by \(\gamma\) form \(D_3(F)\)
by Lemma~\ref{lem:cyclostrict608}. All the elements of order
three belonging to the eigenspaces in the theorem are in this
line, as \(4^{r+1}\equiv1\bmod3\).

If both remaining progressions had order nine, choose classes
\(y_1,y_2\) of order nine in the two eigenspaces, rescaled so
that \(3y_1=3y_2=v\ne0\). On the two progressions,
\(4^{r+1}\equiv1+3(r+1)\bmod9\), with \(r+1\) equal to
\(1\) and \(2\) modulo three in some order. Thus
\((\gamma-1)(y_1+y_2)=(1+2)v=0\), whereas
\(3(y_1+y_2)=2v\ne0\). This would give an element of order
nine in \(D_3(F)\), a contradiction. Finally \(E_x/F\) has
two-power degree for this original character. Restriction,
corestriction and averaging, including at the local subgroups
as in the theorem, identify \(M_{3,r}(F)\) with the invariant
summand on \(E_x\).
\end{proof}

In particular the weighted contribution of this summand at three
over odd \(r<t\) is bounded by \(\lceil(t-1)/2\rceil\log9\) on
the original intervals. For general fields the theorem gives
the actual subgroup at each depth. A uniform upper bound for the
sizes of the groups persisting to greater depths, and the upper
comparison for the integral projected periods in
Section~\ref{sec:missing}, are still required.


\subsection{The exact order of the quadratic summand at three}
\label{sec:quadraticexact609}

The upper bound in Proposition~\ref{prop:quadraticorders608} leaves
open the possibility of order nine. To decide it we use the
ordinary class groups in a real quadratic tower. This also gives
a bound for a collection of quadratic characters at all their
weights on the original intervals.

Let \(\alpha\) be a quadratic character with
\(\alpha(-1)=(-1)^r\), where \(r\ge1\). Write \(\kappa\)
for the \(3\)-adic cyclotomic character, \(\omega\) for its
quadratic Teichm\"uller part, and \(\langle\kappa\rangle=\kappa/\omega\).
The character \(\beta=\alpha\omega^r\) is even. When \(\beta\ne1\)
write \(R\) for its real quadratic field and let
\[
  Q_n=\QQ(\zeta_{3^{n+1}})^+,\qquad R_n=RQ_n,\qquad
  R_\infty=\bigcup_{n\ge0}R_n .
\]
Let \(S(\alpha,r)\) denote the cohomology classes over \(\QQ\)
with coefficient \(A=(\QQ_3/\ZZ_3)(\alpha^{-1}\kappa^{-r})\)
whose restrictions to inertia away from three and to the local
group at three are zero.

\begin{lemma}[A bound from the ordinary class numbers of the real layers]
\leavevmode\label{lem:realorders609}
Write \(h_{3,n}=|\Cl(R_n)\{3\}|\). Every finite subgroup of
\(S(\alpha,r)\) has order at most \(h_{3,n}\) for some \(n\).
In particular, if \(h_{3,n}\le H\) at every layer, then
\(|S(\alpha,r)|\le H\) for every \(r\) giving the same \(\beta\).
\end{lemma}
\begin{proof}
The coefficient character is
\(\theta=\beta^{-1}\langle\kappa\rangle^{-r}\).
It is trivial on \(G_{R_\infty}\). The group
\(\Gal(R_\infty/\QQ)\) is the product of a group of order two
and \(\ZZ_3\), since \(R\) is quadratic and each \(Q_n\) has
degree a power of three. A generator of the infinite factor
acts on \(A\) by a scalar \(u\ne1\). Since \((u-1)A=A\), every
continuous cocycle of that factor is a coboundary, as is seen
by matching its value at the generator. Averaging over the
finite factor therefore gives zero first cohomology of the
quotient. Restriction now injects the global classes into
\(\operatorname{Hom}_{\rm cont}(G_{R_\infty},A)\).

Choose cocycles representing generators of a finite subgroup
of \(S(\alpha,r)\). Each cocycle \(z\) has finite image by
compactness. On one decomposition group at three it has the
form \(z(g)=(\theta(g)-1)b\) for some \(b\in A\). On a
conjugate decomposition group the cobounding element can be
taken to be \(\theta(s)b-z(s)\), where \(s\) conjugates the
first group to the second. All these elements and all the
cocycle images are killed by a common power of three. Take \(n\)
large enough that \(\theta(G_{R_n})\) fixes that coefficient
torsion. Each restriction to \(G_{R_n}\) is then a homomorphism
and vanishes on every decomposition group above three. At a
prime away from three it vanishes on inertia: its original
restriction was a coboundary there, and on inertia over \(R_n\)
both factors of \(\theta\) are trivial.

By global class field theory the homomorphisms so obtained
factor through the \(3\)-part of the ideal class group of
\(\mathcal O_{R_n}[1/3]\). This group is the quotient of
\(\Cl(R_n)\{3\}\) by the \(3\)-primary images of the primes
over three. Infinite places affect only the prime two.
After restriction to \(G_{R_\infty}\), all our generators belong
to the image of the character group of this one finite quotient.
Its character group into \(\QQ_3/\ZZ_3\) has the same order
as the quotient. The initial restriction was injective, proving
the finite-subgroup assertion. Each global class has finite
order by the finite-image argument, so any finite list generates
a finite subgroup. The uniform bound follows.
\end{proof}

\begin{theorem}[Order three at every positive odd weight]\leavevmode
\label{thm:exactthree609}
Let \(F=\QQ(\sqrt{-302})\). For every positive odd integer \(r\),
\[
                  M_{3,r}(F)\simeq\ZZ/3\ZZ .
\]
For the period fields of the character of discriminant \(-1208\),
this is the whole \(\Gal(E_x/F)\)-invariant summand of
\(M_{3,r}(E_x)\). Its contribution at the odd weights \(1\le r<t\)
to the sum of logarithmic orders is exactly
\(\lfloor t/2\rfloor\log 3\).
\end{theorem}
\begin{proof}
Take \(\alpha\) to be the character of \(F\).
At every odd weight, \(\beta=\alpha\omega\) cuts out
\[
                R=\QQ(\sqrt{906}),\qquad
                \operatorname{disc}(R)=3624 .
\]
Indeed \(\omega\) cuts out \(\QQ(\sqrt{-3})\), and \(906\)
is squarefree and is \(2\bmod4\). Restriction from \(\QQ\)
to \(F\) identifies \(S(\alpha,r)\) with \(M_{3,r}(F)\).
To see this, averaging over the quadratic extension identifies
global cohomology with the invariants of its twisted action.
Removing the twist over \(F\) gives the \(\alpha\)-eigenspace,
which is the minus part for odd \(r\). Restriction at three
and restriction on inertia away from three are injective on
the corresponding downstairs local cohomology by corestriction,
since their indices divide two. Thus the identification
preserves both directions of the strict local conditions.

We specify the external class-number input. For a positive
quadratic discriminant \(f\equiv0\bmod3\), Mercuri, Paoluzi
and Schoof identify the order of the \(3\)-class group in layer
\(n\) with the order of
\[
  C(f)/\bigl((1+T)^{3^n}-1\bigr)C(f),
      \qquad C(f)=\ZZ_3[[T]]/J .
\]
This uses equality of the orders of the class group component
and the unit quotient by cyclotomic units, followed by finite
duality and their formula for \(C_n\) in Section~2.
Their discriminant table records \([3624,0,[T,3]]\).
Here \(0\) denotes the stabilization level and the last entry
gives generators of \(J\). This is the level-zero case in the
first column of their Table~3.1. Thus \(C(3624)\) has order
three and \(T\) acts by zero. The displayed quotient has order
three for every \(n\), including \(n=0\).
~\cite{IV-II-mps609,IV-II-mpsdata609}

Lemma~\ref{lem:realorders609} gives \(|M_{3,r}(F)|\le3\).
For the lower bound, Proposition~\ref{prop:example30} constructs
a nonzero class from \(73+9\sqrt{-302}\) at precision one and
transfers it to every positive odd weight. This proves equality.
Restriction from \(F\) to \(E_x\) gives the invariant summand
by Proposition~\ref{prop:quadraticorders608}. There are
\(\lfloor t/2\rfloor\) positive odd integers less than \(t\).
\end{proof}

\begin{corollary}[Quadratic contributions with bounded real discriminant]
\leavevmode\label{cor:quadraticrange609}
For each \(x\) in the original interval and each \(1\le r<t(x)\),
consider the quadratic characters of \(E_x/\QQ\) with the
prescribed parity for which \(\beta\ne1\) and
\(\operatorname{disc}(R)<100000\). For all these pairs together
one has
\[
                       \sum_{r,\alpha}\log |S(\alpha,r)|\ll t(x)
\]
with an absolute implied constant. The same bound applies
to the sum of logarithmic orders of their restriction images
on \(E_x\). When \(3\nmid[E_x:\QQ]\), these images give the
whole indicated character summands.
\end{corollary}
\begin{proof}
Mercuri, Paoluzi and Schoof, Theorem~1.1, prove boundedness
of the \(3\)-class numbers in the cyclotomic tower of every
real quadratic field in this discriminant range.
~\cite{IV-II-mps609}
Let \(H\) be the maximum of the finitely many resulting bounds.
For each \(r\), multiplication by \(\omega^r\) is injective on
quadratic characters, so each real quadratic field occurs at
most once. The lemma bounds each order by \(H\). There are
fewer than \(100000\) choices of discriminant, giving an upper
bound \(100000\,t(x)\log H\). Restriction can only decrease
the order. When the degree of \(E_x\) is prime to three,
restriction followed by averaging identifies global and local
conditions with invariants just as for \(F\) in the theorem.
Removing the coefficient twist gives exactly the corresponding
character summand.
\end{proof}

This bound controls the stated part of the sum on the original
fields. The remaining character components and the integral
upper comparison for the projected theta periods in
Section~\ref{sec:missing} still require estimates there.


\subsection{The exact change in strict cohomology at split primes}
\label{sec:matchedstrict617}

Fix one of the intervals of Section~\ref{sec:missing}. Write \(E_*\)
for its largest period field, \(N_*\) for a cyclotomic modulus for
\(E_*\), \(n_*=[E_*:\QQ]\), and \(t_*\) for the largest twist cutoff.
Let \(\mathcal P\) be any finite set of the large coefficient primes.
We use \(\ell\nmid 4N_*\), \(\ell>n_*t_*\) and \(\ell>2t_*+3\)
for \(\ell\in\mathcal P\). These conditions hold in the large-prime
range of Theorem~\ref{thm:ring39} throughout sufficiently distant
intervals by the polynomial bounds for \(N_*,n_*,t_*\).
For each \(\ell\) work over the integers \(\OO\) of a finite
unramified extension \(J/\QQ_\ell\) containing all the values of
characters of \(\Gal(E_*/\QQ)\).

For \(1\le r<t_*\) and \(\alpha(c)=(-1)^r\), put
\[
  \theta=\alpha^{-1}\kappa_\ell^{-r},\qquad
  \lambda=\theta^{-1}\kappa_\ell,\qquad
  M_m=(\OO/\ell^m)(\theta),\qquad D_m=(\OO/\ell^m)(\lambda).
\]
Use arithmetic Frobenius. Denote by \(S_m\) the classes of
\(H^1(\QQ,M_m)\) which vanish at \(\ell\) and are unramified at all
other finite primes. For a finite set \(Q\) of good primes away from
\(\ell\), let \(S_m^Q\) allow all local classes at \(Q\) and keep
the other conditions. Write \(S\) and \(S^Q\) for the same groups
with coefficients \((J/\OO)(\theta)\). Theorem~\ref{thm:exact241}
proves exact preservation for one auxiliary prime on the nonzero
summands. The issue for several primes is that global reciprocity
pairs the sum of their local contributions with each dual class.

\begin{lemma}[A primitive odd class]\label{lem:oddprimitive617}
There is an integral class \(w\in H^1(\QQ,\OO(\lambda))\) unramified
at every finite prime different from \(\ell\) whose reduction modulo
\(\ell\) has nonzero restriction to \(H=G_{\QQ^{\rm ab}}\).
\end{lemma}
\begin{proof}
Take \(S_0\) to consist of infinity and the primes dividing
\(4N_*\ell\), and write \(G_{\QQ,S_0}\) for the Galois group of the
maximal extension unramified outside \(S_0\). The integral argument
of Lemma~\ref{lem:dual241} applies also when the original strict
group vanishes. Indeed its global Euler characteristic formula for
\(V=J(\lambda)\) gives
\[
    \dim_J H^1(G_{\QQ,S_0},V)-\dim_J H^2(G_{\QQ,S_0},V)=1 .
\]
Indeed the global invariants vanish and \(c\) acts by \(-1\), so
the minus space at the real place has dimension one. The integral
degree-one group for the lattice \(T=\OO(\lambda)\) is finitely
generated by the usual finiteness theorem for Galois cohomology
with restricted ramification. It spans the rational group, as a
continuous cocycle on a compact group has bounded image. It is
torsion-free since \(H^0(G_{\QQ,S_0},V/T)=0\), again by the action
of \(c\). A basis member \(w\) has nonzero reduction by the long
exact sequence for multiplication by \(\ell\).~\cite{IV-II-II-duality13}

At \(p\ne\ell\), pass locally to the field cut out by \(\alpha\);
its degree is prime to \(\ell\). There inertia acts trivially on
\(T\), so the inertia restriction is a homomorphism. Its only possible
contribution is from the pro-\(\ell\) tame quotient \(\ZZ_\ell(1)\).
If the residue field has order \(q\), Frobenius acts on that quotient
by \(q\) and on \(T\) by \(q^{r+1}\). The nonzero scalar \(q^r-1\)
therefore kills its image in a torsion-free module. Restriction on
inertia cohomology across this finite extension is injective by
inflation--restriction and averaging. Thus \(w\) is unramified
over \(\QQ_p\).

On \(H\) the coefficients are trivial, so restriction of a cocycle
is a homomorphism independent of the representative. If the reduction
of \(w\) vanished there it would descend to the abelian quotient.
For a cocycle \(v\) on that quotient, the equality \(cg=gc\) gives
\((\lambda(c)-1)v(g)=(\lambda(g)-1)v(c)\) modulo \(\ell\).
The first multiplier is \(-2\), a unit. This makes \(v\) a
coboundary, contradicting the choice of \(w\).
\end{proof}

\begin{theorem}[Split primes and the change of local conditions]
\label{thm:matchingprimes617}
Assign any integers \(s_\ell\ge1\) to the primes of \(\mathcal P\).
There is an infinite set of rational primes split completely in
\(\QQ(\mu_{4N_*})\) with \(v_\ell(q-1)=s_\ell\) for every
\(\ell\in\mathcal P\). For any \(k\ge1\) distinct members forming
\(Q\), any of the indicated \((\ell,r,\alpha)\), and every \(m\ge1\),
\[
               S_m^Q/S_m
      \ \simeq\ \big(\OO/\ell^{\min(m,s_\ell)}\big)^{k-1}.
      \tag{\(\ast\)}
\]
Under local Tate duality its image is the kernel of addition in \(k\)
copies of the dual of \(\OO/\ell^{\min(m,s_\ell)}\).
\end{theorem}
\begin{proof}
Choose \(\sigma\in G_\QQ\) acting trivially on
\(\QQ(\mu_{4N_*})\) and by \(1+\ell^{s_\ell}\) on
\(\mu_{\ell^{s_\ell+1}}\) for each \(\ell\). Such an element exists by
the Chinese remainder theorem and the cyclotomic Galois isomorphism.
For each character and twist choose a cocycle representing the class
in Lemma~\ref{lem:oddprimitive617}. We can choose \(h\in H\) so
that every one of these cocycles takes a nonzero value modulo
\(\ell\) on \(\gamma=\sigma h\).
In fact the joint image of \(H\) under their residual homomorphisms
is a finite abelian group. Its parts at different primes form a
direct product. On the part at \(\ell\), cancellation of one value
\(\bar w(\sigma)+\bar\lambda(\sigma)\bar w(h)\) excludes at most a
fraction \(1/\ell\), because the homomorphism has nonzero image.
There are at most \(n_*(t_*-1)<\ell\) pairs at that prime, counting
each character of \(E_*\) just once per twist. Their excluded sets
cannot cover the image. Independent choices in its primary parts
give the required \(h\).

Form a finite Galois extension of \(\QQ\) as follows. Include the
cyclotomic fields just used and record every \(D_j\)-valued cocycle
on \(G_{\QQ,S_0}\) for \(1\le j\le s_\ell\), for every indicated
pair at every \(\ell\). There are finitely many because
\(H^1(G_{\QQ,S_0},D_j)\) and the group of coboundaries are finite.
A cocycle \(v\) is recorded by the homomorphism
\(g\mapsto(v(g),\lambda(g)\bmod\ell^j)\) into a finite affine group.
Take the compositum of the corresponding fixed fields of the kernels.
Chebotarev's theorem supplies infinitely many primes with the
conjugacy class of \(\gamma\) in this extension, after omitting all
the primes in the sets \(S_0\). Choose at each of them a place
giving the actual image of \(\gamma\) as Frobenius. These primes
have the required splitting and exact depths.

At each \(q\in Q\), evaluation at the chosen Frobenius identifies the
unramified local group for \(D_m\) with
\[
             U=\OO/(\ell^m,q^{r+1}-1)
               =\OO/\ell^{\min(m,s_\ell)} .
\]
Here \(\alpha(q)=1\), and the binomial expansion of \(q^{r+1}\)
gives \(v_\ell(q^{r+1}-1)=s_\ell\) since \(\ell\nmid r+1\).
The image of the old dual Selmer group in \(U^k\) is exactly
the diagonal. To see the inclusion, this dual group has the local
conditions orthogonal to those of \(S_m\). At every prime outside
\(S_0\) its classes are unramified, since the two unramified
conditions on unramified dual modules are exact annihilators.
Its classes thus come from \(G_{\QQ,S_0}\). For \(m\le s_\ell\)
all their Frobenius values agree by construction. For \(m>s_\ell\)
reduce the cocycles to level \(s_\ell\), which induces an
isomorphism on these local unramified quotients. This proves
the inclusion at every level.

The chosen integral \(w\) reduces to a member of the old dual
Selmer group at every level. At \(\ell\) the primal condition is
zero. At all other finite places both \(w\) and the primal classes
are unramified, and their cup product vanishes by inflation from the
procyclic unramified quotient of cohomological \(\ell\)-dimension
one. The real place contributes zero for odd \(\ell\). The common
value of \(w\) in \(U\) is a unit by our choice of \(h\).
Its \(\OO\)-multiples supply the entire diagonal.

For clarity, the change-of-conditions sequence from finite
Poitou--Tate duality identifies the image of
\[
 S_m^Q\longrightarrow
 \bigoplus_{q\in Q}H^1(\QQ_q,M_m)/H^1_{\rm unr}(\QQ_q,M_m)
\]
with the exact annihilator of these restrictions of the old dual
Selmer group. Indeed the obstruction to lifting a tuple of new
local classes, with the old conditions at the other places, is its
sum of local pairings against that group. This is the middle
localization sequence of global duality after quotienting by the
old local subgroups. Its kernel on global classes is \(S_m\).
Locally the displayed summands are the duals of \(U\). The trace
pairing on the unramified coefficient ring is perfect, so we can
equally use \(\OO\)-linear duals with values in
\(\ell^{-m}\OO/\OO\). The annihilator of the diagonal consists of
tuples whose sum is zero. Its first \(k-1\) entries are arbitrary
and determine the last. This proves \((\ast)\).~\cite{IV-II-II-duality13}
\end{proof}

\begin{corollary}[Elementary divisors on the original fields]
\label{cor:matchedlength617}
For the same \(Q\) there is a noncanonical decomposition restricting
to the given inclusion of \(S\),
\[
                   S^Q\simeq S\oplus(\OO/\ell^{s_\ell})^{k-1}.
\]
It holds on each \(\alpha\)-component of the original groups over
every \(E_x\) through which \(\alpha\) factors. Each prime of \(Q\)
alone leaves the full strict group on this component unchanged.
For \(k\) primes, removing the subgroup killed by \(\ell^a\) from
each of the two groups increases the resulting length by exactly
\((k-1)\max(s_\ell-a,0)\), for every integer \(a\ge0\).
\end{corollary}
\begin{proof}
Let \(A=(J/\OO)(\theta)\) and identify \(M_m\) with \(A[\ell^m]\)
by division by \(\ell^m\). Inertia at \(\ell\) acts through a
nontrivial residual cyclotomic power, so \(A^{G_\ell}=A^{G_\QQ}=0\).
For \(p\ne\ell\), \(A^{I_p}\) is divisible. Indeed the inertia
character has order prime to \(\ell\), giving either \(A\) or zero
as invariants. The long exact sequences for multiplication by
\(\ell^m\), globally, on \(G_\ell\), and on \(I_p\), now give
\(S_m=S[\ell^m]\) and \(S_m^Q=S^Q[\ell^m]\).
The resulting subquotients \(S_m^Q/S_m\) inject into \(S^Q/S\).
Their union is the whole quotient, since a continuous cocycle with
discrete coefficients on a compact group has finite image. By
\((\ast)\) their orders are constant for \(m\ge s_\ell\). A basis
of the quotient of the asserted type thus lifts already to
\(S^Q[\ell^{s_\ell}]\). Sending its generators to these lifts
gives an \(\OO\)-linear splitting.

Since \([E_x:\QQ]\) is prime to \(\ell\), restriction and averaging
identify cohomology on \(\QQ\) for the twist by \(\alpha^{-1}\)
with the \(\alpha\)-part of cohomology on \(E_x\) for the untwisted
module \((J/\OO)(-r)\). This identification also identifies the
local conditions. On decomposition groups above \(\ell\) and on
inertia away from \(\ell\), restriction has zero kernel because the
finite indices are prime to \(\ell\). The relaxed condition allows
all classes on both sides at every place over \(Q\). Finally on each
new cyclic summand the quotient by its \(\ell^a\)-torsion has length
\(\max(s_\ell-a,0)\). Additivity in the direct sum proves the last
assertion.
\end{proof}

In particular the increase in the sum of retained lengths in
\eqref{eq:minors33}, when the same retained triples are compared,
is obtained by multiplying these increments by \(\log\ell\) and
summing them. For \(k=2\) and \(r=1\), the same two primes may also
be used in Theorem~\ref{thm:ststrict616}. The construction prescribes
their Frobenius evaluations and leaves their sizes unrestricted.
The original summand \(S\) still enters the absolute length in
Section~\ref{sec:missing}.


\subsection{Division of strict classes by allowing tame ramification}
\label{sec:dividestrict618}

Keep \(E_*,N_*,t_*,\mathcal P,\theta,\lambda,M_m,D_m\) of
Section~\ref{sec:matchedstrict617}. Write \(O=\ZZ_\ell[\alpha]\)
for the minimal unramified coefficient ring, with fraction field
\(J\) and residue field \(k_\alpha\). We make the local choices over
this ring. Extension to the common unramified coefficient ring
commutes with all the cohomology groups and localization kernels
below, since it is finite free. Use \(S,S_m,S^Q,S_m^Q\) with the
same meaning as in that section. The set \(\mathcal P\) can contain
the full support of the nonzero large-prime summands on the given
interval by Corollary~\ref{cor:interval238}.

\begin{lemma}[Finite layers of the odd dual group]\label{lem:oddlayers618}
Let \(C\) be the group of classes with coefficients
\((J/O)(\lambda)\) unramified away from \(\ell\), with the full
local condition at \(\ell\). The Selmer group dual to the local
conditions defining \(S_m\), with coefficients \(D_m\), equals
\(C_m=C[\ell^m]\) under division by \(\ell^m\). Its length is
\[
             \operatorname{length}_O C_m
               =m+\operatorname{length}_O S[\ell^m].
\]
The divisible subgroup of \(C\) is one copy of \(J/O\).
The finite quotient by that subgroup has the elementary divisors
of \(S\). The compatible reductions \(w_m\) of a primitive
class \(w\) as in Lemma~\ref{lem:oddprimitive617} generate the
\(\ell^m\)-torsion of the divisible subgroup.
\end{lemma}
\begin{proof}
At a finite place \(p\ne\ell\) where the characters are unramified,
local Tate duality makes the unramified subgroups exact annihilators.
Where they are ramified, both inertia characters have prime-to-\(\ell\)
order and are nontrivial after reduction. Both groups of local
invariants vanish, so the entire local degree-one groups vanish by
local duality and the local Euler characteristic. At \(\ell\) the
annihilator of the zero condition is the full group. Global
invariants of \((J/O)(\lambda)\) vanish by oddness, and its inertia
invariants away from \(\ell\) are either zero or divisible.
The exact sequence for multiplication by \(\ell^m\) now identifies
the finite dual Selmer group with \(C[\ell^m]\).

Apply the finite Selmer order formula to \(M_m,D_m\), as in the
proof of Theorem~\ref{thm:positive306}. The two global invariant
groups are zero. For \(p\ne\ell\) finite the order of the unramified
condition equals the order of the local invariants. At \(\ell\)
both the condition on \(M_m\) and its invariants are zero, the
latter by cyclotomic inertia. The real local condition is zero
and the real invariants are all of \(M_m\), since \(\theta(c)=1\).
Thus \(|S_m|/|C_m|=1/|O/\ell^m|\). Corollary~\ref{cor:matchedlength617}
identifies \(S_m\) with \(S[\ell^m]\).~\cite{IV-II-II-duality13}

The Pontryagin dual of \(C\) is a finitely generated \(O\)-module:
its quotient by \(\ell\) is dual to the finite group \(C[\ell]\),
and lifts of a basis generate by compact Nakayama. Equivalently,
the closed span of the lifts has a compact pro-\(\ell\) quotient
equal to its multiple by \(\ell\), and such a quotient has no
nonzero finite quotient. The group \(S\) is finite by
Proposition~\ref{prop:finite15} and prime-to-degree descent.
Apply the structure theorem over \(O\) and take successive
differences of the displayed lengths. For large \(m\) the
difference is one, giving the divisible rank of \(C\). At each
\(m\) the difference in excess of one counts the elementary
divisors of \(S\) of exponent at least \(m\). Finally \(w_m\)
defines a class of order \(\ell^m\), since
\(\ell^{m-1}w_m=w_1\ne0\) in \(C\). Its compatible lifts at all
larger exponents put it in the divisible subgroup. There it
generates the indicated layer.
\end{proof}

\begin{proposition}[Prescribing the residual dual evaluations]
\label{prop:oddrows618}
For each \((\ell,r)\) take one \(\alpha\) from each orbit under
residue-field Frobenius. Choose any positive depths \(s_\ell\).
On every \(C_1\) one may prescribe an arbitrary
\(k_\alpha\)-linear functional. Infinitely many good rational primes
split in \(\QQ(\mu_{4N_*})\), have \(v_\ell(q-1)=s_\ell\)
simultaneously, and realize all these functionals as evaluation
of unramified classes at Frobenius at a choice of place. For a
conjugate character the evaluations are given by conjugation.
\end{proposition}
\begin{proof}
Let \(H=G_{\QQ^{\rm ab}}\). On \(H\) each residual cocycle is a
homomorphism independent of its representative. Restriction of
\(C_1\) is injective by the argument with the central element \(c\)
in Lemma~\ref{lem:oddprimitive617}. Evaluation therefore maps \(H\)
to \(C_1^*=\operatorname{Hom}_{k_\alpha}(C_1,k_\alpha)\) with image
whose annihilator in \(C_1\) is zero. We verify that the joint
image on all the factors is the full product.

At a fixed \(\ell\) this additive image is stable under conjugation
by \(G_\QQ\). Conjugating an element of \(H\) multiplies its
functional on each factor by the scalar \(\bar\lambda(g)\).
The finite abelian group acting in this way is a quotient of
\(\Delta=\Gal(E_*/\QQ)\times\mathbb F_\ell^\times\), with every
element of \(\Delta\) available. Indeed the cyclotomic conductors
\(N_*\) and \(\ell\) are coprime. The characters of \(\Delta\)
on the respective factors are \((a,u)\mapsto\bar\alpha(a)u^{r+1}\).
Their orbits are distinct: restricting to the second factor
determines \(r+1<\ell-1\), and the prime-to-\(\ell\) values of
\(\alpha\) lift uniquely from their residues. The group algebra
\(\mathbb F_\ell[\Delta]\) is a product of fields with one factor
for each such orbit of characters of \(\Delta\). This follows by
diagonalizing the finite abelian group over
\(\overline{\mathbb F}_\ell\) and adding the character idempotents
in each Frobenius orbit, using \(\ell\nmid|\Delta|\).
The orbit idempotents consequently separate all the projections
of the image. In each projection the acting field is \(k_\alpha\),
which is generated by the character values. An image with zero
annihilator that is a \(k_\alpha\)-subspace of \(C_1^*\) is all
of \(C_1^*\). Distinct \(\ell\)'s are separated by the primary
decomposition of the finite additive image.

Take \(\sigma\) with the cyclotomic actions specified in the proof
of Theorem~\ref{thm:matchingprimes617}. It acts trivially on each
residual coefficient, so evaluation there is a well-defined linear
functional on classes. Multiplying \(\sigma\) by an element of \(H\)
adds an arbitrary joint functional by the preceding paragraph.
Record the residual cocycles and the prescribed cyclotomic actions
in the finite extension used for Chebotarev, now using only level
one for the cocycles. Choosing a place with the specified Frobenius
proves the assertion.
\end{proof}

\begin{theorem}[Division by a prescribed power on the same interval]
\label{thm:tamedivision618}
Put \(d_{\ell,r,\alpha}=\dim_{k_\alpha}S[\ell]\) and choose an integer
\(h\ge1+\max_{\ell,r,\alpha}d_{\ell,r,\alpha}\) for a nonempty
list of coefficient primes. For any depths \(s_\ell\ge1\) there
are \(h\) common split rational primes \(Q\) of these exact depths
with the following properties on every indicated original component.
Writing \(s=s_\ell\) and \(P=S^Q\), one has
\[
             S=\ell^sP,\qquad
             P/S\simeq(O/\ell^s)^{h-1}.
\]
If \(b_1,\ldots,b_{h-1}\) are the exponents of \(S\), padded by
zeros, the elementary divisors of \(P\) have exponents
\(s+b_1,\ldots,s+b_{h-1}\). In particular every strict class of
order \(\ell^a\), \(a\ge1\), is \(\ell^s\) times a class of order
\(\ell^{a+s}\) ramified at most at these additional primes in the
sense of the original local conditions. Only residual cohomology
is used to select the Frobenius evaluations at the \(h\) primes.
\end{theorem}
\begin{proof}
For each representative in Proposition~\ref{prop:oddrows618}, choose
a basis of \(C_1\) beginning with \(w_1\). The dimension is \(1+d\),
where \(d=d_{\ell,r,\alpha}\), by Lemma~\ref{lem:oddlayers618}.
Choose \(h\) rows with first entry one that span the dual. For
example take the first coordinate functional, its sum with each
other coordinate functional in turn, and repeat the first row if
needed. Use Proposition~\ref{prop:oddrows618} for each simultaneous
row, excluding the previously chosen primes. The residual
localization \(C_1\to k_\alpha^h\) is then injective and the value
of \(w\) at each chosen Frobenius is a unit. Finite Poitou--Tate
duality as in Theorem~\ref{thm:matchingprimes617} gives
\[
 \dim_{k_\alpha}P[\ell]
   =d+\big(h-\dim_{k_\alpha}C_1\big)=h-1.
\]
This uses \(S_1^Q=P[\ell]\), as proved for finite and divisible
coefficients in Corollary~\ref{cor:matchedlength617} independently
of the particular good primes.

Next consider a level \(m\ge s\). At each chosen prime the
unramified dual group is \(O/\ell^s\), and reduction from level
\(m\) to level \(s\) is an isomorphism on this local group. On the
global groups \(C_m\to C_s\), the same reduction corresponds to
multiplication by \(\ell^{m-s}\) on \(C[\ell^m]\). For all
sufficiently large \(m\) it kills the finite quotient of \(C\)
and its image is exactly the \(\ell^s\)-torsion of the
divisible line. Here the statement about the image follows also
by splitting the divisible line off as an injective \(O\)-module.
Lemma~\ref{lem:oddlayers618} identifies that layer with
\((O/\ell^s)w_s\). The joint unramified local image is therefore
a single vector with \(h\) unit entries and all of its multiples.
Its annihilator in the \(h\) local ramification quotients is
isomorphic to \((O/\ell^s)^{h-1}\), by solving for the last
coordinate. Finite global duality identifies \(S_m^Q/S_m\) with
this annihilator. These subquotients inject into \(P/S\) and
exhaust it, so their eventual common order and isomorphism type
give \(P/S\simeq(O/\ell^s)^{h-1}\). In particular \(P\) is finite.

The equality \(\dim_{k_\alpha}P[\ell]=h-1\) says that \(P\)
requires \(h-1\) generators over the discrete valuation ring \(O\).
Thus \(P/\ell^sP\) has length at most \(s(h-1)\). It surjects
onto \(P/S\), which has exactly that length, since the latter is
killed by \(\ell^s\). This proves \(S=\ell^sP\), and all \(h-1\)
cyclic summands of \(P\) have length at least \(s\). Multiplying
them by \(\ell^s\) gives exactly the stated relation of elementary
divisors. Any lift of a nonzero class of order \(\ell^a\) under
this multiplication has order \(\ell^{a+s}\).
Prime-to-degree restriction and character projection carry the
whole calculation to each \(E_x\), by the descent in
Corollary~\ref{cor:matchedlength617}.
\end{proof}

\begin{corollary}[The layers contributing to the length sum]
\label{cor:tamelayers618}
Choose \(Q\) as in Theorem~\ref{thm:tamedivision618}, and a disjoint
set \(Q'\) of the same cardinality and depths as in
Theorem~\ref{thm:matchingprimes617}. For every \(m\ge1\),
\[
 \operatorname{length}_O S_m^{Q'}-\operatorname{length}_O S_m^{Q}
    =\operatorname{length}_O S[\ell^m]
       -\operatorname{length}_O S[\ell^{\max(m-s_\ell,0)}].
\]
In particular for \(s_\ell=1\) the difference counts exactly the
cyclic summands of the original group of length at least \(m\).
\end{corollary}
\begin{proof}
Put \(s=s_\ell\). Corollary~\ref{cor:matchedlength617} gives
\(\operatorname{length}_O S_m^{Q'}=(h-1)\min(m,s)
+\operatorname{length}_O S[\ell^m]\). For the other set, sum
\(\min(m,b_i+s)=\min(m,s)+\min((m-s)_+,b_i)\) over the elementary
divisors in Theorem~\ref{thm:tamedivision618}. Subtraction proves
the formula. Both sets can avoid any further finite set of primes.
\end{proof}

For \(s_\ell=1\), summing these differences for \(m>a\) gives
the exact length after removal of the part killed by \(\ell^a\)
in \eqref{eq:minors33}. All but finitely many terms vanish for
each triple. The local construction gives no bound on the last
nonzero \(m\), which is still needed for the absolute estimate.


\subsection{Measurement of the total residual length by ordinary class groups}
\label{sec:totalresidual649}

We now include the primes dividing the field degree or the
conductor, and \(p=3\), in the measurement of
Proposition~\ref{prop:res23}. The important degree in the
finite averaging argument is the relative degree from the
chosen field to the largest field in its interval. These
relative degrees are powers of two even when \(p\) divides
\([E_*:\QQ]\).

Fix \(x\in[V,V^{3/2}]\), \(E=E_x\), and let \(R=t_*\) and
\(n_*=[E_*:\QQ]\). Fix an odd \(p\le 2R+1\). As in
Theorem~\ref{thm:strictdigits329}, set
\[
 L_\infty=E_*(\mu_{p^\infty}),\qquad F_0=(E_*(\mu_p))^+,
 \qquad F_\infty=L_\infty^+ .
\]
Write \(F_j\) for the relative degree \(p^j\) layer and
\(C_j=\Cl(\mathcal O_{F_j}[1/p])\{p\}\).
Let \(\Gamma=\operatorname{Gal}(L_\infty/E_*(\mu_p))\),
choose its topological generator \(\gamma\), and write
\(a=\kappa_p(\gamma)\), \(u=v_p(a-1)-1\),
\(l=\lfloor\log_p n_*\rfloor\). We have \(p^u\le n_*\).
Write \(\mathcal T\) for the subgroup of elements of finite
order prime to \(p\) in \(\operatorname{Gal}(L_\infty/\QQ)\).
It is finite since it injects into the finite quotient by
\(\Gamma\simeq\ZZ_p\).

For \(D=E\), or for either field used for the average
\(J_{p,r}\), put
\[
 T_D=\mathcal T\cap\operatorname{Gal}(L_\infty/D),\qquad
 d_{D,r}=\frac{1}{|T_D|}\sum_{\delta\in T_D}
                              \kappa_p(\delta)^r\delta .
\]
These are idempotents in \(\ZZ_p[\mathcal T]\).
Set \(e_r=d_{E,r}\) if the original average is zero, and
\(e_r=d_{E,r}-d_{D,r}\) if it is the average for \(E/D\).
Here \(D=L_x\) or \(B(x)\) according to
Section~\ref{sec:missing}.

\begin{lemma}[The original averages in the cyclotomic class coordinates]
\leavevmode\label{lem:averages649}
Put \(j=m+l\) for any integer \(m\ge1\).
The finite dual, through precision \(m\), of the complementary
strict group over \(E\) after applying \(1-J_{p,r}\) has the
same elementary divisors as
\[
             e_r C_j/(p^m,\gamma-a^{-r})e_r C_j.
\]
Each \(e_r\) is an integral idempotent. The assertion includes
the primes at which the full Galois group of \(E/\QQ\) has
nontrivial \(p\)-Sylow subgroup.
\end{lemma}
\begin{proof}
For every \(D\) above, \(E_*/D\) has degree a power of two.
Thus \(E_*(\mu_p)/D\) has degree prime to the odd prime
\(p\). The abelian group \(G_D=\operatorname{Gal}(L_\infty/D)\)
is the direct product of \(\Gamma\) and \(T_D\). To see this,
let \(h\) be the order of its quotient by \(\Gamma\).
For \(g\in G_D\) there is a unique \(z\in\Gamma\) with
\(z^h=g^h\), since the \(h\)-th power acts bijectively
on \(\ZZ_p\). Then \(gz^{-1}\) has order dividing \(h\).
This also gives \(T_D/T_E=\operatorname{Gal}(E/D)\) when
\(D\subset E\).

Theorem~\ref{thm:ordinarydigits514} identifies the full
complementary dual at precision \(m\) with the quotient of
\(C_j/p^m C_j\) by the relations
\(\delta=\kappa_p(\delta)^{-r}\) for \(\delta\in G_E\).
On its finite prime-to-\(p\) factor these relations project
onto \(d_{E,r}\). On its procyclic factor the single relation
\(\gamma=a^{-r}\) suffices, by continuity. This construction
is compatible with the field action: on equivariant
homomorphisms to the twisted coefficient the action of a
lift \(s\) sends \(f(z)\) to the function
\(z\mapsto\kappa_p(s)^{-r} f(s^{-1}z)\). On the dual
untwisted class coordinates it therefore acts by
\(\kappa_p(s)^r s\). This calculation commutes with
restriction and norm in the finite-layer identification.
Lifts differing by \(G_E\) induce the same action.

On the \(d_{E,r}\)-summand the average of these operators over
\(T_D/T_E\) is exactly \(d_{D,r}\), as is seen by summing
first over \(T_E\). Also \(d_{D,r}d_{E,r}=d_{D,r}\).
Averaging and taking its complement commute with the
precision and cyclotomic quotients, since they are direct
summand projections. Dualizing a finite-group average gives
the average of the inverse actions, which is the same
projector. This proves the stated identification.
\end{proof}

Decompose the finite abelian prime-to-\(p\) algebra
\(\ZZ_p[\mathcal T]\) by its Frobenius orbit idempotents
\(e_\nu\). On each factor \(e_r\) is zero or one. Put
\[
 \mathcal W_\nu=\{1\le r<t(x):e_\nu e_r=e_\nu\},\qquad
 P_x(Y)=\sum_\nu e_\nu\prod_{r\in\mathcal W_\nu}(Y-a^{-r}).
\]
Each empty product is one. For this \(x,p\) define
\[
 Z_m=C_{m+l}/(p^m,P_x(\gamma))C_{m+l},\qquad
 h_p=\max_{\nu,\ r\in\mathcal W_\nu}
       \sum_{r'\in\mathcal W_\nu\setminus\{r\}}
                        (1+u+v_p(r-r')) ,
\]
with zero for the empty maximum.

\begin{theorem}[A finite measurement for the entire residual local sum]
\leavevmode\label{thm:totalresidual649}
For \(X>0\) put \(b=b_0(x)\), \(m_p=b+\lfloor X/\log p\rfloor+1\)
and form the two natural logarithmic sums
\[
 A_-=\sum_{2<p\le2R+1}\log|p^{b+h_p}Z_{m_p}|,\qquad
 A_+=\sum_{2<p\le2R+1}\log|p^{b-h_p}Z_{m_p}|
\]
over primes. For all sufficiently large \(V\), \(b\ge h_p\)
for every prime in these sums. There is a fixed
\(C=C_{\psi,\varepsilon}\) such that independently of \(x,X\)
\[
 \begin{gathered}
        0\le A_+-A_-\ll_{\psi,\varepsilon} V^C,\\
 A_++\log J(V)\le X\ \Longrightarrow\ \mathcal R(x)\le X,
 \qquad
 \mathcal R(x)\le X\ \Longrightarrow\ A_-\le X .
\end{gathered}
\]
In particular these quantities account for the full range of
odd coefficient primes in the residual sum of Section~\ref{sec:missing}.
\end{theorem}
\begin{proof}
Let \(Y_m\) be the direct sum over weights of the finite
groups in Lemma~\ref{lem:averages649}. On each orbit \(\nu\)
apply the interpolation maps of Theorem~\ref{thm:alltails648}
to the distinct values \(a^{-r}\) for \(r\in\mathcal W_\nu\).
Their difference valuations are \(1+u+v_p(r-r')\) by the
binomial theorem at an odd prime. Thus the two composites
are \(p^{h_p}\), and
\[
                  |p^{b+h_p}Z_m|\le |p^bY_m|
                                      \le |p^{b-h_p}Z_m|.
\]
Here \(h_p\le R(1+l+\log_p R)\), which is polynomial in
\(V\) and less than \(b\) uniformly. The difference between
the two outer logarithms to base \(p\) is at most
\(2h_p\dim_{\mathbb F_p}Z_m/pZ_m\).
As \(P_x\) reduces to a power of \(Y-1\) of degree at most
\(R\) on each factor, the filtration by \(\gamma-1\) bounds
this dimension by \(R\dim_{\mathbb F_p}
C_{m+l}/(p,\gamma-1)C_{m+l}\).

We bound the last dimension without assuming total
ramification at the bottom. In the cyclotomic
\(\ZZ_p\)-extension of \(\QQ\), inertia at \(p\) surjects
onto the whole Galois group. Passing to inertia of a
completion of \(F_0\) reduces its image by an index dividing
the ramification index of that completion: the local
inertia subgroup has precisely that index in the original
inertia group. The \(p\)-part of this index is at most
\(n_*\) since \([F_0:\QQ]\) divides \(n_*(p-1)\).
Under the cyclotomic identification, the subgroup
\(\Gamma^{p^l}\) is therefore in this inertia image.
Indeed relative to the generator \(1+p\) of the whole
cyclotomic direction it has index \(p^{u+l}\).
It follows that every prime over \(p\) in \(F_l\) is totally
ramified in \(F_{m+l}/F_l\).

The dimension of \(C_{m+l}/(p,\gamma-1)C_{m+l}\) is at
most that of the quotient by \(p,\gamma^{p^l}-1\).
Over \(F_l\) the split class field argument in
Theorem~\ref{thm:alltails648} applies to this latter
quotient with \(\gamma^{p^l}\) as generator. It gives
the bound \(v_p(h_{F_l})+g_l-1\), where \(g_l\) is the
number of primes of \(F_l\) above \(p\). Both quantities
are polynomially bounded in \(V\), uniformly for
\(p\le2R+1\). Indeed
\([F_l:\QQ]\le n_*(p-1)p^l\), and by
Theorem~\ref{thm:strictdigits329} its conductor divides
\(N_*p^{u+1+l}\). Here \(p^l,p^u\le n_*\).
Apply the conductor-discriminant and Minkowski bounds of
Corollary~\ref{cor:periodfieldweights644}. Summing the
width with factors \(\log p\) over at most \(2R+1\) primes
proves the uniform polynomial estimate.

All primes in the finite test are below \(P(x)\) for large
\(V\). Hence \(p^bY_{m_p}\) measures exactly the residual
tail there with each elementary divisor capped at \(m_p\).
A divisor longer than this cap would already contribute more
than \(X\) to the capped total, by the choice of \(m_p\).
If \(A_++\log J(V)\le X\), the capped total is therefore
the true total through \(2R+1\). For larger primes the
separated-weight bound of Lemma~\ref{lem:fullinteger517}
applies at all the weights and costs at most \(\log J(V)\),
including primes above \(P(x)\). Its fixed excluded primes
are all below \(2R+1\) eventually. This proves the first
implication; the capped total is always at most the true
one, which proves the second.
\end{proof}

For example one can take \(X=\tfrac14\exp(x^\varepsilon/2)\),
the sufficient budget in Proposition~\ref{prop:res23}.
The polynomial width is smaller than any fixed positive
fraction of this budget eventually, uniformly on the
interval. An arithmetic upper bound for \(A_+\) at this
scale is still required. The width estimate bounds the
difference of the two measured sums and is compatible
with both sums exceeding the budget. The signed period
comparison must also hold for the same original data.


\subsection{Which prime ideals can detect the residual strict classes?}
\label{sec:primeideals650}

Fix an interval of Theorem~\ref{thm:conductor232} and a point
\(x\) in it. Write \(E=E_x\), \(n=[E:\QQ]\), \(t=t(x)\),
\(R=t_*\) and \(n_*=[E_*:\QQ]\). Use the fields, the action
of \(\gamma\), and the idempotents of
Lemma~\ref{lem:averages649}. For each odd \(p\le2R+1\)
choose an arbitrary precision \(m_p\ge1\), set
\(j=m_p+\lfloor\log_p n_*\rfloor\), and write
\[
 Q_{p,r}=e_r C_j/(p^{m_p},\gamma-a^{-r})e_r C_j
                   \qquad(1\le r<t).
\]
There is a surjection from the ordinary ideal class group of
\(F_j\) to \(Q_{p,r}\). It kills the primes above \(p\), takes
the \(p\)-primary projection, and then applies the displayed
projection and quotient. For a rational prime \(s\), let
\(H_{p,r}(s)\) be the subgroup generated by the images of the
prime ideals above \(s\). For a finite set \(S\) of rational
primes put \(H_{p,r}(S)=\sum_{s\in S}H_{p,r}(s)\).
The ramification index, residue degree and number of places
of \(E\) over \(s\) are denoted by \(e_s,f_s,g_s\). These
satisfy \(e_s f_s g_s=n\), since \(E/\QQ\) is Galois.

\begin{lemma}[An integer annihilating the projected prime classes]
\leavevmode\label{lem:primeann650}
For \(s\ne p\), the group \(H_{p,r}(s)\) has at most \(g_s\)
generators and is annihilated by \(s^{f_s r}-1\).
For \(s=p\) the group is zero. These assertions include
primes \(s\) ramified in \(E\) and primes \(p\) dividing
\([E_*:\QQ]\).
\end{lemma}
\begin{proof}
Let \(G_E=\operatorname{Gal}(L_\infty/E)\). Its action on
\(Q_{p,r}\) is scalar multiplication by \(\kappa_p^{-r}\).
Indeed Lemma~\ref{lem:averages649} imposes this relation on
the finite subgroup \(T_E\) and on a generator of the
procyclic factor. It persists under the field average because
the entire Galois group is abelian.

Fix a prime of \(F_j\) above \(s\ne p\) and prolong it to an
algebraic closure containing \(L_\infty\). Write \(v\) for
the place it induces on \(E\). An arithmetic Frobenius lift
at \(v\) preserves the chosen prime of \(F_j\). On all
\(p\)-power roots of unity its action is raising to
\(s^{f_s}\). In fact reduction is injective on roots of unity
of order prime to \(s\), and arithmetic Frobenius on the
residue field raises to the cardinality of the residue field
of \(v\). This also proves independence of the lift in the
presence of inertia. Thus on the image of this ideal class
the same element acts both as the identity and as
\(s^{-f_s r}\). Multiplication by the \(p\)-adic unit
\(s^{f_s r}\) gives the asserted annihilator. The cyclotomic
Frobenius formula also appears in
\srcloc{ReciprocityTheory/GaloisValuations/Global}{364}.

The image of \(G_E\) on \(F_j\) is
\(\operatorname{Gal}(F_j/(E\cap F_j))\), by Galois
correspondence. Its orbits on the primes above \(s\) thus
correspond to the primes of \(E\cap F_j\) above \(s\).
Each extends to a prime of \(E\), so there are at most
\(g_s\) orbits. The class map is equivariant, and scalar
multiplication on a finite abelian \(p\)-group preserves
the cyclic subgroup generated by each element. One prime
per orbit therefore suffices. The case \(s=p\) follows from
the definition of \(C_j\).
\end{proof}

\begin{theorem}[The total contribution of a set of rational primes]
\leavevmode\label{thm:primemass650}
Let \(S\) be finite. The full natural logarithmic order of its
images, summed over all measured coefficient primes and weights,
satisfies
\[
 \sum_{2<p\le2R+1}\ \sum_{1\le r<t}\log |H_{p,r}(S)|
 \ \le\ \frac{n\,t(t-1)}{2}\sum_{s\in S}\frac{\log s}{e_s}.
\]
For any integer \(b\ge0\), the same bound holds for the
decrease of \(\sum_{p,r}\log|p^bQ_{p,r}|\) when the primes
over \(S\) are inverted in the ordinary class groups.
The bounds are uniform in the precisions \(m_p\).
\end{theorem}
\begin{proof}
For \(s\ne p\), the preceding lemma bounds
\(\log |H_{p,r}(s)|\) by
\(g_s v_p(s^{f_s r}-1)\log p\). For fixed \(s,r\) this is
the valuation of a single positive integer, with a multiplier
independent of \(p\). Summing over \(p\) gives at most
\[
 g_s\log(s^{f_s r}-1)\ \le\ (n/e_s)\,r\log s.
\]
The formula includes the integer one, whose logarithm is zero.
The product of the groups for the individual \(s\) surjects
onto their sum. Addition over \(S\) and then over \(r\)
proves the first assertion.

Inverting all the primes over \(S\) quotients the ideal class
group by their classes. This subgroup is stable under every
field automorphism. Primary projection, the idempotent and the
scalar quotient therefore give \(Q_{p,r}/H_{p,r}(S)\).
For \(Q=Q_{p,r}\) and \(H=H_{p,r}(S)\) the kernel of
\(p^b Q\longrightarrow p^b(Q/H)\) is \(H\cap p^b Q\).
Its logarithmic order is bounded by \(\log |H|\), proving
the second assertion.
\end{proof}

For instance, on our interval one may invert all rational
primes dividing \(pN_*\) in the class group for \(p\). The
primes above \(p\) have already been removed, so the additional
set can be chosen as \(S=\{s:s\mid N_*\}\), independent of
\(p\) and \(x\). Its cost in the actual capped residual sum
is at most \(n_*R^2\log N_*/2\). This is polynomial in \(V\)
and includes the two moving quadratic primes. The same
conclusion holds for any set with
\(\sum_{s\in S}\log s\le V^A\) for a fixed \(A\).

\begin{corollary}[Cyclotomic splitting forced by a surviving prime class]
\leavevmode\label{cor:detectingsplit650}
Put \(b=b_0(x)\) and
\(k_p=\max\{0,b-\lfloor\log_p(n_*R)\rfloor\}\).
Map the class group to the tail \(p^bQ_{p,r}\) by multiplying
the original class map by \(p^b\). For \(s\ne p\) a nonzero
image of a prime above \(s\) requires \(s\) to split
completely in the degree-\(p^{k_p}\) subfield of
\(\QQ(\mu_{p^{k_p+1}})\). Its size must also satisfy
\[
                f_s r\log s>(b+1)\log p.
\]
In particular every prime with \(\log s\le
b_0(V)\log 3/(n_*R)\) has zero image in all these tail
maps, simultaneously throughout the interval.
\end{corollary}
\begin{proof}
If \(z\) is the projected class, \(p^b z\ne0\) forces its
order to be at least \(p^{b+1}\). Hence
\(p^{b+1}\mid s^{f_s r}-1\). This proves the strict size
inequality. Let \(d\) be the order of \(s\bmod p\). It
divides \(f_s r\) and \(p-1\). For an odd \(p\), binomial
expansion at \(1\) gives
\[
 v_p(s^{f_s r}-1)
    =v_p(s^{p-1}-1)+v_p(f_s r).
\]
Indeed raising \(s^d\) to an exponent prime to \(p\)
preserves the valuation of its difference from one, and
each subsequent \(p\)-th power increases it by one.
Since \(f_s r\le n_*R\), the valuation on the right coming
from \(s^{p-1}-1\) is at least \(1+k_p\) when \(k_p>0\).
Modulo \(p^{k_p+1}\), the residue of \(s\) then lies in the
subgroup of order \(p-1\). This is exactly the Frobenius
criterion for complete splitting in the indicated subfield.
When \(k_p=0\) that field is \(\QQ\). The size inequality,
\(p\ge3\), and \(b_0(x)\ge b_0(V)\) give the last assertion.
\end{proof}

This splitting restriction concerns the rational primes
which can represent nonzero elements under the tail map.
Even a set of primes of very small density can represent
every class of a finite quotient. An absolute bound for the
order of the quotient in Theorem~\ref{thm:totalresidual649}
is still needed, as is the signed period comparison
\eqref{eq:gap} on the same fields.


\subsection{Recovering strict torsion from ray characters}
\label{sec:raylength656}

Fix an interval of Theorem~\ref{thm:conductor232} and \(x\) in
that interval. Put \(E=E_x\) and use its common modulus \(N_*\).
Let \(p\) be odd and \(1\le r<t(x)\). Write
\(\operatorname{Gal}(E/\QQ)=P\times\Delta\), where \(P\) is its
\(p\)-Sylow subgroup. Take a finite extension \(J/\QQ_p\) containing
the values of all characters of this group. Let \(O,\pi,k,e\)
be its ring of integers, a uniformizer, its residue field and
\(v_\pi(p)\), respectively. Write \(A=J/O\), \(R_h=O/\pi^hO\).
For characters \(\alpha\) of \(\Delta\) and \(\eta\) of \(P\),
inflated to the absolute Galois group, put
\[
        \theta=\kappa_p^{-r}\alpha^{-1}\eta^{-1},
        \qquad \Sigma=\{v:v\text{ is a prime dividing }pN_*\}.
\]
Let \(S(\theta)\subset H^1(\QQ,A(\theta))\) be the classes
trivial on the decomposition groups at \(\Sigma\) and on inertia
elsewhere. These are the strong scalar local conditions of
Section~\ref{sec:levelstrict345}. Their source definition is
\srcloc{Development/AutomorphicArithmetic/AutomorphicVarieties/CharacterLevelBounds/ShaFamMono}{30}.

Write \(H\) for the fixed field of the kernel of
\(\chi=\theta\bmod\pi\), and \(\Delta_H=\operatorname{Gal}(H/\QQ)\).
This group embeds in \(k^\times\). Denote by \(\widetilde\chi\)
the lift of \(\chi\) to roots of unity of order prime to \(p\).
Normalize all local valuations of \(\pi\)-adic character values
by \(v_\pi(\pi)=1\), with \(v_\pi(0)=+\infty\), and define
\[
 c_H=\min_{g\in G_H}v_\pi(\theta(g)-1),\qquad
 c_w=\min_{g\in G_{H_w}}v_\pi(\theta(g)-1)\quad(w\mid\Sigma).
\]
Here \(G_H,G_{H_w}\) denote absolute Galois groups. We have
\(1\le c_H\le c_w<\infty\). At \(w\mid p\) the cyclotomic
image contains an open subgroup of \(1+p\ZZ_p\), even after
killing the finite characters. Elsewhere the value of the
cyclotomic character at Frobenius is the residue cardinality,
whose \(r\)-th power is greater than one and is not a root of unity.

For \(1\le h\le c_H\) put \(u=\lceil h/e\rceil\). For \(w\mid p\)
write \(e_w=v_w(p)\) with \(v_w(H_w^\times)=\ZZ\). Form the ray modulus
\[
 \mathfrak m_h=
 \prod_{w\mid p}\mathfrak p_w^{\,1+\lfloor e_w/(p-1)\rfloor+u e_w}
 \prod_{\substack{w\mid\Sigma\\w\nmid p}}\mathfrak p_w .
\]
At \(w\mid\Sigma\) define the following character with values in
the additive group of \(R_h\) by the formula
\[
         \ell_w(g)=\frac{\theta(g)-1}{\pi^{c_w}}\pmod{\pi^h}
                          \quad(g\in G_{H_w}).
\]
Indeed the discarded cross term in the product rule has valuation
at least \(c_w\ge h\) after division. One value of \(\ell_w\)
is a unit, so \(\ell_w\) spans a free \(R_h\)-line.
Write \(\mathrm{Cl}_{\mathfrak m_h}(H)\) for the ray class group.
Let \(L_h\) be the following explicitly defined finite \(O\)-module.
Take the homomorphisms
\[
 \varphi:\mathrm{Cl}_{\mathfrak m_h}(H)\longrightarrow (R_h,+),
 \qquad \varphi(\delta z)=\widetilde\chi(\delta)\varphi(z)
                         \quad(\delta\in\Delta_H).
\]
Retain those whose pullback by reciprocity at each \(w\mid\Sigma\)
is in \(R_h\ell_w\). These are finite linear conditions. Under local reciprocity one
can evaluate them on a uniformizer and generators of the unit
group modulo the specified principal units. Both sides factor
through these quotients.

\begin{theorem}[The finite torsion of a strict character group]
\leavevmode\label{thm:raylength656}
For \(1\le h\le c_H\) there are isomorphisms of \(O\)-modules
\[
 S(\theta)[\pi^h]\ \simeq\
 \begin{cases}
 L_h,&\chi\ne1,\\
 L_h/R_h\ell_H,&\chi=1 .
 \end{cases}
\]
In the second line \(H=\QQ\) and
\(\ell_H(g)=(\theta(g)-1)/\pi^{c_H}\bmod\pi^h\).
The residual field is the same for all \(\eta\) and all weights
congruent to \(r\) modulo \(p-1\) at fixed \(\alpha\).
\end{theorem}
\begin{proof}
Identify \(A[\pi^h]\) with \(R_h\) by
\(\pi^{-h}z+O\mapsto z\bmod\pi^h\). The action modulo \(\pi^h\)
is trivial on \(G_H\). It factors through \(\Delta_H\) on which
it is \(\widetilde\chi\bmod\pi^h\), since the kernel of reduction
on units of \(R_h\) is a \(p\)-group. Inflation--restriction
identifies the global finite-coefficient cohomology with the
equivariant continuous homomorphisms on \(G_H\) displayed above,
initially without a conductor condition. The outer cohomology
groups of \(\Delta_H\) in degrees one and two vanish by averaging
over its order, a unit in \(O\). At each rational prime restriction
to \(H_w\) likewise detects the zero class for the divisible
coefficients, by restriction followed by corestriction.

The coefficient sequence for multiplication by \(\pi^h\) on
\(A(\theta)\) maps finite-coefficient cohomology onto the subgroup
killed by \(\pi^h\). Its kernel is the image under the connecting
map of the invariants modulo \(\pi^h\) times the invariants. At
\(H_w\) the invariants are \(\pi^{-c_w}O/O\). Lifting their
generator to \(\pi^{-c_w-h}+O\) in the coefficient sequence
gives exactly \(\ell_w\). Thus the image of a finite class in
divisible cohomology vanishes locally precisely when its
restriction lies on this line. At primes outside \(\Sigma\)
both the character and \(H/\QQ\) are unramified. On inertia all
of \(A\) is invariant and divisible, so the connecting map there
vanishes. The finite class must itself vanish on inertia.
At a real place the positive cohomology is zero since \(p\) is odd.

Global class field theory now describes the finite characters
unramified outside \(\Sigma\). Away from \(p\) the higher
units are pro-residue-characteristic and have no nonzero image
in \(R_h\). Above \(p\) put \(s_w=1+\lfloor e_w/(p-1)\rfloor\).
The local logarithm and exponential are inverse between
\(1+\mathfrak p_w^{s_w}\) and \(\mathfrak p_w^{s_w}\).
Raising to \(p^u\) therefore maps this unit group onto
\(1+\mathfrak p_w^{s_w+u e_w}\). Every additive character
with values in \(R_h\) kills those powers. These facts give
exactly all the required characters from the stated ray group.
They apply to the local characters \(\ell_w\) as well.

Finally, if \(\chi\ne1\), some \(\theta(g)-1\) is a unit
and the global invariants vanish. If \(\chi=1\), the global
connecting map gives the line \(R_h\ell_H\), by the same
calculation as at \(w\). Its restriction is
\(\pi^{c_w-c_H}\ell_w\), hence it belongs to \(L_h\).
This gives the asserted quotient. The character of \(P\)
reduces to one and \(\bar\kappa_p\) has order \(p-1\),
which proves the assertion about \(H\).
\end{proof}

\begin{corollary}[Two consecutive precisions]
\leavevmode\label{cor:raystable656}
Let \(d_0=0\) and let \(d_h\) be the \(O\)-length of the
displayed quotient at precision \(h\). If \(d_h=d_{h-1}\) for
some \(1\le h\le c_H\), then \(\pi^{h-1}S(\theta)=0\)
and \(S(\theta)\) has length \(d_h\).
Writing \(n_*=[E_*:\QQ]\), the ray groups used for these tests
satisfy, for \(n_H=[H:\QQ]\le (p-1)n_*\),
\[
 \log|\mathrm{Cl}_{\mathfrak m_h}(H)|
 \ll n_H\{\log(2pN_*n_H)+(2+\lceil h/e\rceil)\log p\}.
\]
\end{corollary}
\begin{proof}
The two successive torsion subgroups are nested and finite,
so equality of lengths gives equality of subgroups. Every
cohomology class is killed by a power of \(\pi\), since a
continuous cocycle on a compact group has finite image in \(A\).
Any element of order \(\pi^j\) with \(j\ge h\) would give an
element of order exactly \(\pi^h\) by multiplication. This proves
the first assertion. The field \(H\) is contained in
\(E^P(\mu_p)\) and has conductor dividing \(pN_*\). Its absolute
discriminant is at most \((pN_*)^{n_H}\) by the
conductor--discriminant formula. Minkowski's bound gives an ideal
of norm at most \(2^{n_H}|\operatorname{disc}H|^{1/2}\) in
each ordinary class. The number of ideals of norm at most \(y\)
is at most \(y^2\zeta_H(2)\le y^2\zeta(2)^{n_H}\).
The ray group adds a factor at most the norm of its modulus.
Away from \(p\) this is at most \(N_*^{n_H}\). Above \(p\)
it is at most \(p^{(2+u)n_H}\), since
\(1+\lfloor e_w/(p-1)\rfloor\le2e_w\) and
\(\sum_{w\mid p}e_w f_w=n_H\). These inequalities prove the bound.
\end{proof}

\begin{theorem}[Descent on the integral character block]
\leavevmode\label{thm:integralblock656}
Write \(|P|=p^a\). Let \(M_\alpha\) be the \(\alpha\)-eigenspace
for \(\Delta\) in \(\mathcal S_{p,r}(E)\otimes_{\ZZ_p}O\).
For \(b\ge a\),
\[
       |p^b M_\alpha|
         \ \le\ \prod_{\eta\in\widehat P}
                  |p^{b-a}S(\kappa_p^{-r}\alpha^{-1}\eta^{-1})|.
\]
In particular \(p^bM_\alpha=0\) if every \(\eta\) satisfies
the consecutive-precision test at some \(h_\eta\) with
\(h_\eta-1\le e(b-a)\).
\end{theorem}
\begin{proof}
Over \(E\) the coefficient is \(A(\kappa_p^{-r})\).
At a finite place \(w\nmid p\) its unramified local group
is \(A/((Nw)^{-r}-1)A=0\), since the nonzero scalar acts
surjectively. Thus the strict kernel there also satisfies the
strong condition at all places above \(\Sigma\). Extending
coefficients commutes with these kernels since \(O\) is finite
free over \(\ZZ_p\). Twist the descent by \(\alpha^{-1}\).
Restriction to \(E\) from \(E^\Delta\) identifies the strict
kernel with \(M_\alpha\). Indeed averaging gives descent on global
cohomology and injectivity on local cohomology, including on
inertia. Shapiro's lemma from \(E^\Delta\) to \(\QQ\) now
identifies this group with the strong kernel on
\(O[P]\otimes_O A(\kappa_p^{-r}\alpha^{-1})\). On a decomposition
group the induced module splits over its cosets into induced
local modules. The same holds on inertia, so the asserted kernel
is preserved under Shapiro.

Character evaluation \(O[P]\to\bigoplus_\eta O(\eta^{-1})\)
has an integral reverse map given by the inverse Fourier matrix
multiplied by \(|P|\). Orthogonality makes both composites
multiplication by \(p^a\). Tensoring these identities with the
twisted coefficient and taking the strong kernels preserves the
identities. In particular \(p^b M_\alpha\) is contained in the
image under the reverse map of \(p^{b-a}\bigoplus_\eta S(\theta)\).
The order bound and the vanishing assertion follow.
\end{proof}

These tests apply at the cutoff \(b=b_0(x)\) of
Section~\ref{sec:missing}, also to any further direct summand
of the indicated block when \(b_0(x)\ge a\). Their validity requires no change of
period field or interval. An upper bound for the sum of the
blocks for which no test has yet given stabilization is still
needed for the full local estimate.


\subsection{Ordinary class groups before cyclic descent}
\label{sec:beforecyclic661}

Fix \(x\) on an interval of Theorem~\ref{thm:conductor232},
write \(E=E_x\), and let \(p\) be odd. Write
\(\operatorname{Gal}(E/\QQ)=P\times\Delta\), with \(P\) its
\(p\)-Sylow subgroup. Set \(K=E^\Delta\). This is the cyclic
\(p\)-power subfield of the fixed character field \(K_\psi\),
including \(K=\QQ\). Indeed \(E/K_\psi\) has degree a power
of two. In particular \(K\) is totally real and its ramified
primes form a set independent of \(x\).

Let \(O\) be an unramified extension of \(\ZZ_p\) containing
the values of the characters of \(\Delta\), let \(k=O/pO\),
and put \(A=O[1/p]/O\). For \(1\le r<t(x)\) let \(\alpha\)
be a character of \(\Delta\) with \(\alpha(c)=(-1)^r\), where
\(c\) is complex conjugation. Write
\[
 M=\big(\mathcal S_{p,r}(E)\otimes_{\ZZ_p}O\big)_\alpha,\qquad
 \theta=\kappa_p^{-r}\alpha^{-1},\qquad
 \chi=\bar\kappa_p^{-r}\bar\alpha^{-1}.
\]
We regard \(\alpha,\theta,\chi\) as characters of the absolute
Galois group by inflation where necessary. The residual field
\(H=\overline{\QQ}^{\,\ker\chi}\) is totally real, since
\(\chi(c)=1\), and \(\Delta_H=\operatorname{Gal}(H/\QQ)\)
has order prime to \(p\). Thus \(H\) and \(K\) are disjoint.
Write \(\widetilde\chi\) for the Teichm\"uller lift of \(\chi\)
and \(R_h=O/p^hO\). For \(L=H\) or \(KH\) set
\[
 D_h(L)=
 \{\varphi:\mathrm{Cl}(L)\longrightarrow (R_h,+)
      \mid \varphi(\delta z)=\widetilde\chi(\delta)\varphi(z)
                     \text{ for every }\delta\in\Delta_H\}.
\]
The maps are homomorphisms of groups. On \(KH\) the action of
\(\Delta_H\) fixes \(K\). Use \(D_\infty(L)\) for the same
definition with target \(A\). Its exponent and order are finite.
For the local hypotheses below, \(D_v\) denotes an absolute
decomposition group at the rational prime \(v\).

\begin{theorem}[Torsion on the entire \(p\)-Sylow block]
\leavevmode\label{thm:wholeblock661}
Suppose \(\chi(D_p)\ne1\). For \(1\le h\le1+v_p(r)\) one has
\[
                        M[p^h]\simeq D_h(KH)
\]
as \(O\)-modules with \(P\)-action. The action on the character
group on the right is by precomposition with the inverse action
on the class group. If \(p^d\) is the exponent of
\(D_\infty(KH)\), with \(d=0\) for the zero module, then
\(v_p(r)\ge d\) implies \(p^d M=0\) and
\(M\simeq D_\infty(KH)\).
\end{theorem}
\begin{proof}
First descend through \(E/K\). Since its degree \(|\Delta|\)
is a unit in \(O\), global restriction identifies the group
on the indicated eigenspace with cohomology over \(K\) twisted
by \(\alpha^{-1}\). The same averaging gives injectivity on
local restrictions. More precisely the resulting group \(S_K\)
consists of classes in \(H^1(K,A(\theta))\) which vanish at the
places over \(pN_*\) and on inertia elsewhere. Over \(E\) these
conditions are the original ones. At a place \(u\nmid p\) the
unramified local subgroup on \(E\) is
\(A/((Nu)^{-r}-1)A=0\). This proves the claim about the kernels
by applying restriction to each decomposition or inertia group.
The coefficient extension commutes with kernels since \(O\) is
finite free over \(\ZZ_p\). The descent is equivariant for \(P\)
with the usual cohomological action. In particular it retains
the integral \(p\)-Sylow action throughout.

Let \(w\mid p\) in \(K\). The subgroup \(G_{K_w}\subset D_p\)
has \(p\)-power index, while \(\chi\) has order prime to \(p\).
Hence its restriction has the same nontrivial image. Some local
value of \(\theta-1\) is a unit, and \(A(\theta)\) has no
invariants at \(w\) or globally over \(K\). The sequence for
multiplication by \(p^h\) thus identifies \(S_K[p^h]\) with
finite-coefficient classes that vanish at every \(w\mid p\)
and are unramified at every \(w\nmid p\). To verify the latter
assertion, the inertia image of \(\theta\) away from \(p\) has
order prime to \(p\). If it is nontrivial there are no inertia
invariants, in either coefficient module. If it is trivial
the divisible module of inertia invariants is \(A\). In both
cases the finite-to-divisible map on inertia cohomology is
injective. An unramified finite class maps to zero on the
whole local group in divisible cohomology, since the latter's
unramified subgroup is zero. In the unramified-coefficient case
the scalar on Frobenius is a root of unity times \((Nw)^{-r}\)
and differs from one, which proves this last statement.

For \(h\le1+v_p(r)\) the finite action is
\(\widetilde\chi\bmod p^h\). Indeed \(\alpha\) takes
prime-to-\(p\) roots of unity as values, and for
\(u\in1+p\ZZ_p\) one has \(u^r\equiv1\bmod p^{1+v_p(r)}\)
by binomial expansion. Restrict from \(K\) to \(KH\).
Inflation--restriction and averaging identify the finite classes
with \(\widetilde\chi\)-equivariant additive characters there,
as \([KH:K]\) is prime to \(p\). Unramifiedness is detected by
restriction on inertia since the local inertia index divides
this degree. Local vanishing at \(p\) is detected in the same
way on decomposition groups. Class field theory gives the
ordinary ideal class group for everywhere-unramified characters,
real places imposing no further condition at odd \(p\).
Moreover each prime ideal over \(p\) has trivial image under
every equivariant character. Its stabilizer in \(\Delta_H\)
has nontrivial \(\chi\)-image by the local calculation at \(w\).
One element of the stabilizer acts on the target by a root whose
difference from one is a unit. Thus the character vanishes on
the ideal class, exactly the additional splitting condition at
this place. This gives the isomorphism.

For the \(P\)-action one may use representatives in \(G_H\),
since \(KH/H\) maps isomorphically to \(K/\QQ\). They act
trivially on \(R_h(\theta)\) in this range, so conjugation of
ideal classes gives the stated dual action. Finally
\(D_h(KH)\simeq D_\infty(KH)[p^h]\) via
\(z\bmod p^h\mapsto z/p^h+O\). Under the exponent hypothesis
the orders of \(M[p^d]\) and \(M[p^{d+1}]\) agree, taking
\(M[1]=0\). Every cohomology class is killed by a power of
\(p\) since a continuous cocycle on a compact group has finite
image in \(A\). An element of greater exponent would yield
one of order exactly \(p^{d+1}\). This proves the assertion
about the whole module.
\end{proof}

\begin{theorem}[A vanishing test on the prime-to-\(p\) field]
\leavevmode\label{thm:cyclicdescent661}
Assume in addition that \(\chi(D_v)\ne1\) at every finite prime
ramified in \(K/\QQ\). With \(s=1+v_p(r)\) one has
\[
                  (M[p^h])^P\simeq D_h(H)\quad(1\le h\le s),
       \qquad       M=0\ \Longleftrightarrow\ D_1(H)=0.
\]
The number of indecomposable summands of the \(k[P]\)-module
\(M[p]\) is \(\dim_k D_1(H)\). In particular
\(\dim_k M[p]\le |P|\dim_k D_1(H)\).
\end{theorem}
\begin{proof}
Use the finite coefficient \(R_h(\theta)\) of the preceding
proof. There are no global invariants over \(K\), so
inflation--restriction identifies its first cohomology over
\(\QQ\) with the \(P\)-invariants over \(K\). At primes
outside \(p\) and the ramification of \(K\), unramifiedness
descends because the inertia groups are the same. At a ramified
prime \(v\ne p\) the unramified finite subgroup over \(K_w\)
is zero. Indeed a nontrivial residual inertia value kills the
invariants. If inertia acts trivially, the nontrivial residual
local character has Frobenius value different from one and the
unramified quotient is zero. The same alternatives hold over
\(\QQ_v\). The character image stays nontrivial over \(K_w\)
by the \(p\)-power index argument. Restriction on local first
cohomology is injective here, and at \(p\), since local
coefficient invariants over \(K_w\) vanish. Thus the finite
local conditions descend exactly. The first theorem's finite
class field calculation over \(\QQ\) now gives \(D_h(H)\).

The finite space \(M[p]\) is a module over
\(k[P]\simeq k[T]/(T^{|P|})\), with \(T=\sigma-1\) for a
generator of the cyclic group. On each nilpotent Jordan block
the kernel of \(T\) is one-dimensional. This proves the count,
the dimension bound and the equivalence at precision one.
Since \(M\) is \(p\)-primary torsion, its vanishing is detected
by \(M[p]\).
\end{proof}

\begin{corollary}[The test throughout the deep dyadic layers]
\leavevmode\label{cor:dyadiccyclic661}
Write \(\alpha=\beta\nu\), with \(\beta\) from \(C_J\) of order
\(2^j\) and \(\nu\) of order dividing \(\operatorname{lcm}(2,o)\),
where \(o=\operatorname{ord}\psi\). Put
\[
 \mathcal B_p=\{p\}\cup\{v:v\text{ ramifies in }K/\QQ\},
 \qquad
 j_p=\max_{v\in\mathcal B_p}\big(v_2(v^2-1)-3\big)
                          +v_2(\operatorname{lcm}(2,o,p-1)).
\]
Both theorems apply for \(j>j_p\) at every weight of the
specified sign, for every \(x\) on the original intervals and
with either choice of each moving quadratic character.
If \(p-1\mid 2r\), the fields \(KH\) for these tests have degree
at most \(2[E:\QQ]\) and conductor dividing \(pN_*\).
\end{corollary}
\begin{proof}
The characters of \(E\) are generated by the original character,
the dyadic characters and the two quadratic characters. This
gives the decomposition with the asserted order bound. All
primes in \(\mathcal B_p\) are odd since a quotient of
\((\ZZ/2^a\ZZ)^\times\) has no odd-order character.
The Frobenius of such a prime \(v\) acts on the real dyadic
tower by the class of \(v\) modulo signs. Represent that class
by the choice of \(v\) or \(-v\) congruent to one modulo four.
The valuation of its difference from one is \(v_2(v^2-1)-1\).
The group \(1+4\ZZ_2\) is generated topologically by \(5\),
and \(v_2(5^{2^a}-1)=a+2\) by repeated squaring. In the
quotient of order \(2^j\), the Frobenius therefore has order
\(2^{\max(0,j-(v_2(v^2-1)-3))}\). The dyadic character is
faithful on that quotient and unramified at \(v\). In the
residual character all the other factors on a Frobenius lift
have order dividing \(\operatorname{lcm}(2,o,p-1)\).
For \(j>j_p\) they cannot cancel its two-part. This proves the
local hypotheses, including for \(v=p\). At \(p-1\mid2r\)
the cyclotomic character modulo \(p\) to exponent \(r\) has
order at most two. The compositum \(KH\) then lies in the
compositum of \(E\) with the quadratic subfield of
\(\QQ(\mu_p)\), giving both asserted bounds.
\end{proof}

For the cubic character of conductor \(3367\) in
Section~\ref{sec:cyclicray657}, one has \(p=3\),
\(\mathcal B_p=\{3,7,13,37\}\), and \(j_p=2\). Thus every
dyadic character of order at least eight satisfies these tests
at every positive weight of the prescribed sign. For an even
weight and \(\nu=1\), the residual field is \(C_j\). The
class number of \(C_j\) is prime to three for every \(j\) by
Fukuda and Komatsu, Corollary~1.1 of their paper on Weber's
class number problem~\cite{IV-II-weberFK524}. Theorem~\ref{thm:cyclicdescent661}
therefore kills the entire cubic Sylow block above this dyadic
character, even though the invariant block of the original
cubic field in Theorem~\ref{thm:strictcyclic657} is nonzero.
For a moving quadratic twist the class group in the vanishing
test is the one for its own residual character. Controlling
the groups that survive these tests in all weights remains
necessary in Section~\ref{sec:missing}.

\clearpage
\section{References}\label{nav:bibliography-IV}
\begingroup\def\section*#1{}\begin{thebibliography}{99}\small
\bibitem{IV-II-II-duality13} J. S. Milne, \emph{Arithmetic Duality Theorems}, 2nd ed., 2006, Chapters I and II. \href{https://www.jmilne.org/math/Books/ADTnot.pdf}{Book}.
\bibitem{IV-II-II-iwasawa15} K. Iwasawa, \emph{On \(\ZZ_\ell\)-extensions of algebraic number fields}, Ann. of Math. \textbf{98} (1973), 246--326.
\bibitem{IV-II-III-weibel} C. A. Weibel, \emph{The \(K\)-book}, AMS, 2013. \href{https://sites.math.rutgers.edu/~weibel/Kbook/Kbook.VI.pdf}{Chapter VI, Section 8}.
\bibitem{IV-II-III-kahn} B. Kahn, \emph{The Quillen--Lichtenbaum conjecture at the prime 2}. \href{https://webusers.imj-prg.fr/~bruno.kahn/preprints/QL2.pdf}{Theorem 1}.
\bibitem{IV-II-III-beh} E. Bayer-Fluckiger, V. Emery and J. Houriet, \emph{Hermitian lattices and bounds in \(K\)-theory of algebraic integers}, Doc. Math., Extra Volume Merkurjev (2015), 71--83. \href{https://arxiv.org/abs/1406.5953}{Paper}, especially Section 3.1 and Corollary 5.2.
\bibitem{IV-II-III-mathew} A. Mathew, \emph{Torsion exponents in stable homotopy and the Hurewicz homomorphism}, Algebr. Geom. Topol. \textbf{16} (2016), 1025--1041. \href{https://msp.org/agt/2016/16-2/agt-v16-n2-p12-s.pdf}{Theorem 1.3}.
\bibitem{IV-II-III-fw} B. Ferrero and L. C. Washington, \emph{The Iwasawa invariant \(\mu_p\) vanishes for abelian number fields}, Ann. of Math. \textbf{109} (1979), 377--395. \href{https://doi.org/10.2307/1971116}{Paper}.
\bibitem{IV-II-III-wash27} L. C. Washington, \emph{The non-\(p\)-part of the class number in a cyclotomic \(\ZZ_p\)-extension}, Invent. Math. \textbf{49} (1978), 87--97. \href{https://www.researchgate.net/profile/Lawrence-Washington/publication/226757118_The_non-p-part_of_the_class_number_in_a_cyclotomic_Zp-extension/links/0deec51929879ecc4a000000/The-non-p-part-of-the-class-number-in-a-cyclotomic-Zp-extension.pdf}{Theorem, p.~87}.
\bibitem{IV-II-III-torre103} A. de la Torre Dur\'an, \emph{Perfect forms and torsion bounds in algebraic \(K\)-theory over totally real and CM fields}, preprint, November 2026. \href{https://arxiv.org/abs/2611.06300}{Theorem 3.9 and Remark 4.2}.
\bibitem{IV-II-I-weibel102} C. A. Weibel, \emph{The \(K\)-book}, Graduate Studies in Mathematics 145, AMS, 2013. \href{https://sites.math.rutgers.edu/~weibel/Kbook/Kbook.VI.pdf}{Chapter VI, Corollary 9.8}.
\bibitem{IV-II-V-newforms} F. Diamond and J. Im, \emph{Modular forms and modular curves}, in \emph{Seminar on Fermat's Last Theorem}, CMS Conf. Proc. \textbf{17} (1995), 39--133. \href{https://wstein.org/edu/Fall2003/252/references/diamond-im/Diamond-Im-Modular_forms_and_modular_curves.pdf}{Authors' text}.
\bibitem{IV-II-mps609} P. Mercuri, M. Paoluzi and R. Schoof, \emph{Greenberg's conjecture for real quadratic number fields}, Journal of Experimental Mathematics \textbf{1} (2025), 207--217. Theorem~1.1, Section~2 and Table~3.1. \href{https://reneschoof.github.io/MPSfinal.pdf}{Published paper}.
\bibitem{IV-II-mpsdata609} P. Mercuri, M. Paoluzi and R. Schoof, data accompanying \emph{Greenberg's conjecture for real quadratic number fields}, file \emph{Ideals with discriminant 0 mod 3.txt}, entry for \(3624\). \href{https://github.com/mercuri-pietro/Iwasawa-modules}{Authors' data repository}.
\bibitem{IV-II-weberFK524} T. Fukuda and K. Komatsu, \emph{Weber's class number problem in the cyclotomic $\mathbb Z_2$-extension of $\mathbb Q$, II}, J. Th\'eor. Nombres Bordeaux \textbf{22} (2010), 359--368, Corollary 1.1. \href{https://www.numdam.org/article/JTNB_2010__22_2_359_0.pdf}{Published paper}.
\end{thebibliography}\endgroup
\end{document}
